% concatenated from CTY_2026_JASA--SA.tex, tables/empapp_main_itt_bwitt.tex, tables/empapp_main_fs_bwitt.tex, tables/empapp_main_fuzzy_bwitt.tex, tables/empapp_covbal_itt_bwitt.tex, tables/empapp_main_itt0_bwitt.tex, tables/empapp_main_percent_bwitt.tex, tables/empapp_main_itt_bwmain.tex, tables/empapp_main_fs_bwmain.tex, tables/empapp_main_fuzzy_bwmain.tex, tables/empapp_covbal_itt_bwmain.tex, tables/empapp_main_itt0_bwmain.tex, tables/empapp_main_percent_bwmain.tex, tables/simuls_itt_s1_bwitt.tex, tables/simuls_itt_s2_bwitt.tex, tables/simuls_itt_s3_bwitt.tex, tables/simuls_fs_s1_bwitt.tex, tables/simuls_fs_s2_bwitt.tex, tables/simuls_fs_s3_bwitt.tex, tables/simuls_fuzzy_s1_bwitt.tex, tables/simuls_fuzzy_s2_bwitt.tex, tables/simuls_fuzzy_s3_bwitt.tex, tables/simuls_itt_s1_bwmain.tex, tables/simuls_itt_s2_bwmain.tex, tables/simuls_itt_s3_bwmain.tex, tables/simuls_fs_s1_bwmain.tex, tables/simuls_fs_s2_bwmain.tex, tables/simuls_fs_s3_bwmain.tex, tables/simuls_fuzzy_s1_bwmain.tex, tables/simuls_fuzzy_s2_bwmain.tex, tables/simuls_fuzzy_s3_bwmain.tex
\documentclass[10pt]{article}

\usepackage[top=1in,bottom=1in,left=1in,right=1in]{geometry}
\usepackage[onehalfspacing]{setspace}
\usepackage{amssymb, amsmath, amsthm}
\usepackage{graphicx, hyperref, enumitem, booktabs, subcaption}
\usepackage[sans]{dsfont}
\usepackage[scr=rsfs,cal=boondox]{mathalpha}
\hypersetup{pdfborder = {0 0 0},colorlinks=true,allcolors=blue}
\usepackage{natbib}


\newtheorem{thm}{Theorem}
\newtheorem{lem}{Lemma}
\newtheorem{assumption}{Assumption}
\newtheorem{remark}{Remark}

% Labelling for SA
\renewcommand{\thesection}{SA-\arabic{section}}
\renewcommand{\theequation}{SA-\arabic{equation}}
\renewcommand{\thefigure}{SA-\arabic{figure}}
\renewcommand{\thetable}{SA-\arabic{table}}
\renewcommand{\thelem}{SA-\arabic{lem}}
\renewcommand{\thethm}{SA-\arabic{thm}}
\renewcommand{\theassumption}{SA-\arabic{assumption}}
\renewcommand{\theremark}{SA-\arabic{remark}}

\allowdisplaybreaks[1]

\setlength{\parindent}{1em}

\renewcommand*{\arraystretch}{.6}


% fixing ToC for SA
\usepackage{tocloft}
\addtolength{\cftsecnumwidth}{20pt}
\addtolength{\cftsubsecnumwidth}{20pt}
\addtolength{\cftsubsubsecnumwidth}{20pt}
\addtolength{\cftparanumwidth}{20pt}
%%%%%%%%%%%%%%%%%%%


% New operators
\DeclareMathOperator*{\argmin}{arg\,min}
\DeclareMathOperator*{\diag}{diag}
\DeclareMathOperator{\Tstat}{T}
\DeclareMathOperator{\CI}{I}
\DeclareMathOperator{\setbd}{bd}
\DeclareMathOperator{\setint}{int}
\DeclareMathOperator{\supp}{supp}
\DeclareMathOperator{\diam}{diam}


% % % % % % % % % % % % % % Some Notations % % % % % % % % % % % % % %
\renewcommand{\P}{\mathbb{P}}
\newcommand{\E}{\mathbb{E}}
\newcommand{\Var}{\mathbb{V}}
\newcommand{\Cov}{\mathbb{C}\mathrm{ov}}
\newcommand{\Indicator}{\mathds{1}}
\newcommand{\En}{\mathbb{E}_n}
\newcommand{\leb}{\mathfrak{m}}
\newcommand{\Haus}{\mathfrak{H}}

\newcommand{\reals}{\mathbb{R}}


\newcommand{\dif}{\mathop{}\!\mathrm{d}}

\newcommand{\norm}[1]{\left\lVert#1\right\rVert}
\newcommand{\infnorm}[1]{\lVert #1 \rVert_{\infty}}
\newcommand{\lessp}{\lesssim_{\mathbb{P}}}

\newcommand{\pnorm}[1]{\lVert #1 \rVert_{\mathbb{P},2}}

\newcommand{\fr}{\mathfrak{r}}

% % % % % % % % % % % % % % Bold letters % % % % % % % % % % % % % %
\newcommand{\ba}{\mathbf{a}}
\newcommand{\bA}{\mathbf{A}}
\newcommand{\bb}{\mathbf{b}}
\newcommand{\bB}{\mathbf{B}}
\newcommand{\bc}{\mathbf{c}}
\newcommand{\be}{\mathbf{e}}
\newcommand{\bh}{\mathbf{h}}
\newcommand{\bH}{\mathbf{H}}
\newcommand{\bI}{\mathbf{I}}
\newcommand{\bM}{\mathbf{M}}
\newcommand{\bQ}{\mathbf{Q}}
\newcommand{\br}{\mathbf{r}}
\newcommand{\bS}{\mathbf{S}}
\newcommand{\bu}{\mathbf{u}}
\newcommand{\bv}{\mathbf{v}}
\newcommand{\bV}{\mathbf{V}}
\newcommand{\bx}{\mathbf{x}}
\newcommand{\bX}{\mathbf{X}}
\newcommand{\by}{\mathbf{y}}
\newcommand{\bz}{\mathbf{z}}
\newcommand{\bZ}{\mathbf{Z}}

\newcommand{\bbeta}{\boldsymbol{\beta}}
\newcommand{\bnu}{\boldsymbol{\nu}}
\newcommand{\bomega}{\boldsymbol{\omega}}
\newcommand{\bgamma}{\boldsymbol{\gamma}}
\newcommand{\bGamma}{\boldsymbol{\Gamma}}
\newcommand{\bTheta}{\boldsymbol{\Theta}}
\newcommand{\bSigma}{\boldsymbol{\Sigma}}
\newcommand{\bUpsilon}{\boldsymbol{\Upsilon}}
\newcommand{\bPi}{\boldsymbol{\Pi}}

\newcommand{\hGammap}{\widehat{\boldsymbol{\Gamma}}_{1, \bx}}
\newcommand{\hGamman}{\widehat{\boldsymbol{\Gamma}}_{0, \bx}}
\newcommand{\hGamma}{\widehat{\boldsymbol{\Gamma}}_{t, \bx}}
\newcommand{\Gammap}{\boldsymbol{\Gamma}_{1, \bx}}
\newcommand{\Gamman}{\boldsymbol{\Gamma}_{0, \bx}}
\newcommand{\Gammax}{\boldsymbol{\Gamma}_{t, \bx}}


% % % % % % % % % % % % % % Script letters % % % % % % % % % % % % % %
\newcommand{\A}{\mathcal{A}}
\newcommand{\B}{\mathcal{B}}
\newcommand{\C}{\mathcal{C}}
\newcommand{\F}{\mathcal{F}}
\newcommand{\G}{\mathcal{G}}
\newcommand{\K}{\mathscr{K}}
\newcommand{\M}{\mathcal{M}}
\newcommand{\Q}{\mathcal{Q}}
\newcommand{\R}{\mathcal{R}}
\newcommand{\V}{\mathcal{V}}
\newcommand{\X}{\mathcal{X}}
\newcommand{\Y}{\mathcal{Y}}
\newcommand{\q}{\mathcal{q}}

\newcommand{\calP}{\mathcal{P}}


\newcommand{\bbQ}{\mathbb{Q}}

\newcommand{\ttc}{\mathtt{c}}
\newcommand{\ttd}{\mathtt{d}}
\newcommand{\ttE}{\mathtt{E}}
\newcommand{\ttF}{\mathtt{F}}
\newcommand{\ttk}{\mathtt{k}}
\newcommand{\ttK}{\mathtt{K}}
\newcommand{\ttL}{\mathtt{L}}
\newcommand{\ttM}{\mathtt{M}}
\newcommand{\ttN}{\mathtt{N}}
\newcommand{\ttv}{\mathtt{v}}


\newcommand{\gn}{\mathbb{G}_n}
\newcommand{\ttTV}{\mathtt{TV}}
\newcommand{\px}{\mathbf{x}'}

\begin{document}
% % % % % % % % % % % % % % Title Page % % % % % % % % % % % % % %
\title{Estimation and Inference in Boundary Discontinuity Designs: Location-Based Methods\\ Supplemental Appendix\bigskip}
\author{Matias D. Cattaneo\thanks{Department of Operations Research and Financial Engineering, Princeton University.} \and
	    Rocio Titiunik\thanks{Department of Politics, Princeton University.} \and
	    Ruiqi (Rae) Yu\thanks{Department of Operations Research and Financial Engineering, Princeton University.}
	    }
\maketitle


\begin{abstract}
    This supplemental appendix presents more general theoretical results encompassing those reported in the paper, their theoretical proofs, and other technical results. It also provides a discussion of implementation details, and additional numerical results.
\end{abstract}

\textit{Keywords}: regression discontinuity, treatment effects estimation, causal inference.


\thispagestyle{empty}
\clearpage

\onehalfspacing
\setcounter{page}{1}
\pagestyle{plain}

\setcounter{tocdepth}{2}
\setcounter{secnumdepth}{4}
\tableofcontents

\clearpage


%%%%%%%%%%%%%%%%%%%%%%%%%%%%%%%%%%%%%%%%%%%%%%%%%%%%%%%%%%%%%%%%%%%%%%
\section{Introduction}
%%%%%%%%%%%%%%%%%%%%%%%%%%%%%%%%%%%%%%%%%%%%%%%%%%%%%%%%%%%%%%%%%%%%%%

This supplemental appendix considers four generalizations of the setup in the paper.
\begin{enumerate}
    \item The score $\bX_i$ is $d$-dimensional with $d\geq2$ and support $\mathcal{X}\subseteq\mathbb{R}^d$.
    \item The boundary region $\B$ is a substantively more general low-dimensional manifold with ``effective dimension'' $d-1$, and hence not necessarily piecewise linear with finitely many line segments.
    \item Derivative estimation and inference for all parameters of interest.
    \item Different bandwidths $\bh = (h_1, \ldots, h_d)^{\top}$ are allowed for each dimension of the score variable.
\end{enumerate}
In addition, this supplement presents several other technical and methodological results not covered in the main text due to space limitations.

The paper presents results for the special case $d=2$, $\B$ piecewise linear with finitely many line segments, $h=h_1=h_2$, and $\bnu=\mathbf{0}$. The assumptions in the main paper are more specialized primitive conditions; after imposing these restrictions on the supplemental results, the mapping between the paper and this supplemental appendix is as follows.
\begin{itemize}%[leftmargin=*]
    \item Theorem \ref{sa-thm: Convergence Rates: Sharp BATEC} $\implies$ Theorem 1.
    \item Theorem \ref{sa-thm: MSE Expansions: Sharp BATEC} $\implies$ Theorem 2.
    \item Theorems \ref{sa-thm: Confidence Intervals: Sharp BATEC} and \ref{sa-thm: Confidence Bands: Sharp BATEC -- Suprema} $\implies$ Theorem 3.
    \item Theorem \ref{sa-thm: MSE Expansion: Sharp WBATE} $\implies$ Theorem 4.
    \item Theorem \ref{sa-thm: Asymptotic Normality: Sharp WBATE} $\implies$ Theorem 5.
    \item Theorem \ref{sa-thm: Confidence Intervals: Sharp LBATE} $\implies$ Theorem 6.
    \item Theorem \ref{sa-thm: Convergence Rates: Fuzzy BATEC} $\implies$ Theorem 7.
    \item Theorem \ref{sa-thm: MSE Expansions: Fuzzy BATEC} $\implies$ Theorem 8.
    \item Theorems \ref{sa-thm: Confidence Intervals: Fuzzy BATEC} and \ref{sa-thm: Confidence Bands: Fuzzy BATEC -- suprema} $\implies$ Theorem 9.
    \item Theorem \ref{sa-thm: MSE Expansion: Fuzzy WBATE} $\implies$ Theorem 10.
    \item Theorem \ref{sa-thm: Asymptotic Normality: Fuzzy WBATE} $\implies$ Theorem 11.
    \item Theorem \ref{sa-thm: Confidence Intervals: Fuzzy LBATE} $\implies$ Theorem 12.
\end{itemize}

The generalizations considered in this supplemental appendix are conceptually straightforward, but technically more involved. Following the structure of the paper, we first consider sharp designs. Section \ref{sa-sec: Imperfect Compliance} generalizes the results to fuzzy designs. Section \ref{sa-sec: Gaussian Strong Approximation} presents a novel Gaussian strong approximation result that may be of independent interest.

\subsection{Notation and Definitions}

For textbook references on empirical process, see \cite{van-der-Vaart-Wellner_1996_Book}, \cite{dudley2014uniform}, and \cite{Gine-Nickl_2016_Book}. For textbook references on geometric measure theory, see \cite{simon1984lectures}, \cite{federer2014geometric}, and \cite{folland2002advanced}.

\begin{enumerate}[label=(\roman*)]
    \item \textit{Multi-index Notations}. For a multi-index $\bnu = (\nu_1, \ldots, \nu_d) \in \mathbb{N}_0^d$, let $|\bnu| = \sum_{l = 1}^d \nu_l$, $\bnu! = \nu_1!\cdots \nu_d!$, and $\bu^{\bnu} = \prod_{l=1}^d u_l^{\nu_l}$ for $\bu=(u_1,\ldots,u_d)\in\mathbb{R}^d$. For any integer $p\geq0$, let $\br_p(\bu) = (\bu^{\boldsymbol{\alpha}}: |\boldsymbol{\alpha}|\leq p)^\top \in \mathbb{R}^{\mathfrak{p}_p}$ collect all monomials of total degree at most $p$ in a fixed order, where $\mathfrak{p}_p = \frac{(d + p)!}{d! p !}$. Define $\be_{\bnu}$ to be the $\mathfrak{p}_p$-dimensional unit vector such that $\be_{\bnu}^{\top} \br_p(\bu) = \bu^{\bnu}$ for all $\bu \in \reals^d$ and $|\bnu|\leq p$. In addition, let $f^{(\bnu)}(\bx) = \partial^{\nu_1}\cdots\partial^{\nu_d} f(\bx)$.

    \item \textit{Norms}. For a vector $\ba \in \reals^k$, $\norm{\ba} = (\sum_{i = 1}^k \ba_i^2)^{1/2}$, $\infnorm{\ba} = \max_{1 \leq i \leq k}|\ba_i|$. For a matrix $\bA \in \reals^{m \times n}$, $\norm{\bA}_p = \sup_{\norm{\bx}_p = 1} \norm{\bA\bx}_p$, $p \in \mathbb{N} \cup \{\infty\}$, $\lambda_{\min}(\bA)$ denotes its minimum eigenvalue, and $|\bA|=|\det(\bA)|$ denotes the absolute value of its determinant. For a function $f$ on a metric space $(S, d)$, $\infnorm{f} = \sup_{s \in S} |f(s)|$. For a probability measure $Q$ on $(\mathcal{S}, \mathscr{S})$ and $p \geq 1$, define $\norm{f}_{Q,p} = (\int_{\mathcal{S}} |f|^p \dif Q)^{1/p}$, and $Q(f) = \int f \dif Q$. For a set $E \subseteq \reals^d$, denote by $\mathfrak{m}(E)$ the Lebesgue measure of $E$.

    \item \textit{Empirical Process}. We use standard empirical process notations: $\En[g(\bv_i)] = \frac{1}{n} \sum_{i = 1}^n g(\bv_i)$ and $\gn[g(\bv_i)] = \frac{1}{\sqrt{n}} \sum_{i = 1}^n (g(\bv_i) - \E[g(\bv_i)])$. Let $(\mathcal{S},d)$ be a semi-metric space. The covering number $N(\mathcal{S}, d, \varepsilon)$ is the minimal number of balls $B_s(\varepsilon) =\{t: d(t,s) < \varepsilon\}$ needed to cover $\mathcal{S}$. A \emph{$\P$-Brownian bridge} is a mean-zero Gaussian random function $\mathbb{G}_{\P}(f), f \in L_2(\X, \P)$ with covariance $\E[\mathbb{G}_{\P}(f)\mathbb{G}_{\P}(g)] = \P(fg) - \P(f)\P(g)$, for $f,g \in L_2(\X,\P)$. A class $\mathcal{F} \subseteq L_2(\X, \P)$ is \emph{$\P$-pregaussian} if there is a version of $\P$-Brownian bridge $\mathbb{G}_{\P}$ such that $\mathbb{G}_{\P} \in C(\mathcal{F}; \rho_{\P})$ almost surely, where $\rho_{\P}$ is the semi-metric on $L_2(\X,\P)$ is defined by $\rho_{\P}(f, g) = (\|f - g\|_{\P,2}^2 - (\int f \dif\P - \int g \dif\P)^2)^{1/2}$, for $f, g \in L_2(\X,\P)$.

    \item \textit{Geometric Measure Theory}. For a set $A \subseteq \X$, the \emph{De Giorgi perimeter of $A$ related to $\X$} is $\mathcal{L}(A) = \ttTV_{\{\Indicator_{A}\},\X}$. For $d \in \mathbb{N}$ and $0 \leq m \leq d$, the $m$-dimensional Hausdorff (outer) measure is $\Haus^m(A) = \lim_{\delta \downarrow 0}\Haus^m_{\delta}(A)$, $A \subseteq \reals^d$, where for each $\delta > 0$, $\Haus^m_{\delta}(A)$ is defined by taking $\Haus^m_{\delta}(\emptyset) = 0$, and for any non-empty $A \subseteq \reals^d$, $\Haus^m_{\delta}(A) = \frac{\pi^{m/2}}{\Gamma(m/2+1)} \inf \sum_{j = 1}^{\infty} (\diam(C_j)/2)^m$, and the infimum is taken over all countable collections $C_1, C_2, \cdots$ of subsets of $\reals^d$ such that $\diam(C_j) < \delta$ and $A \subseteq \cup_{j = 1}^{\infty}C_j$. Integration against $\Haus^m$ is defined via Carathéodory's Theorem following the classical measure-theoretic literature. The Hausdorff dimension $\dim_{\Haus}(A)$ of $A$ is defined by $\dim_{\Haus}(A) = \inf\{t \geq 0: \Haus^t(A) = 0\}$. A set $A \subseteq \reals^d$ is said to be $k$-rectifiable if $A$ is of Hausdorff dimension $k$, and there exist a countable collection $\{f_i\}$ of continuously differentiable maps $f_i: \reals^k \to \reals^d$ such that $\Haus^{k}(A \setminus \cup_{i = 0}^{\infty} f_i(\reals^k)) = 0$. $B$ is a \emph{rectifiable curve} if there exists a Lipschitz continuous function $\gamma:[0,1] \to \reals^d$ such that $B=\gamma([0,1])$. We define the curve length function of $B$ to be $\mathfrak{L}({B}) = \sup_{\pi \in \Pi} s(\pi, \gamma)$, where $\Pi = \left\{(t_0, t_1, \ldots, t_N): N \in \mathbb{N}, 0 \leq t_0 < t_1 < \ldots \leq t_N \leq 1\right\}$ and $s(\pi,\gamma) = \sum_{i = 0}^{N-1}\norm{\gamma(t_{i}) - \gamma(t_{i+1})}_2$ for $\pi = (t_0, t_1, \ldots, t_N)$.

    \item \textit{Bounds and Asymptotics}. For real sequences, we write $a_n = o(b_n)$ if $\limsup \frac{|a_n|}{|b_n|} = 0$, and $a_n \lesssim b_n$ if there exist constants $C$ and $N > 0$ such that $n > N$ implies $|a_n| \leq C |b_n|$. For sequences of random variables, we write $a_n = o_{\P}(b_n)$ if $\operatorname{plim}_{n \to \infty}\frac{|a_n|}{|b_n|} = 0$, and $|a_n| \lessp |b_n|$ if $\limsup_{M \to \infty} \limsup_{n \to \infty} \P[|\frac{a_n}{b_n}| \geq M] = 0$.

\end{enumerate}

Let $\Phi(\cdot)$ be the cumulative distribution function of a standard univariate Gaussian random variable.
For a bandwidth vector $\bh=(h_1,\ldots,h_d)^\top\in\mathbb{R}_{++}^d$, let $\bH=\diag(\bh)$, $\underline h=\min_{1\leq l\leq d}h_l$, $\bar{h}=\max_{1\leq l\leq d}h_l$, and $K_{\bH}(\bu)=|\bH|^{-1}K(\bH^{-1}\bu)$. All limits are taken such that $\bar{h}\to0$, $\bar{h}/\underline h\lesssim1$, and the bandwidths vanish as $n\to\infty$. Most of our results hold with fixed but small enough bandwidths. We do not make this distinction explicit to avoid overly-complex statements and technical discussions.


%%%%%%%%%%%%%%%%%%%%%%%%%%%%%%%%%%%%%%%%%%%%%%%%%%%%%%%%%%%%%%%%%%%%%%
\section{Sharp Design Setup}
%%%%%%%%%%%%%%%%%%%%%%%%%%%%%%%%%%%%%%%%%%%%%%%%%%%%%%%%%%%%%%%%%%%%%%

We partition $\mathcal{X}$ into two areas, $\A_t\subset\mathbb{R}^d$ with $t\in\{0,1\}$, which represent the control and treatment regions, respectively. That is, $\mathcal{X} = \A_0 \cup \A_1$, where $\A_0$ and $\A_1$ are two disjoint regions in $\mathbb{R}^d$. The observed outcome is $Y_i = \Indicator(\bX_i \in \A_0) Y_i(0) + \Indicator(\bX_i \in \A_1) Y_i(1)$. The ``boundary'' is $\B = \setbd(\A_0) \cap \setbd(\A_1) \subseteq\mathbb{R}^{d}$, where $\setbd(\cdot)$ denotes the topological boundary of a set, and has effective dimension $d-1$. We assume that $\B$ belongs to $\A_1$, that is, $\B \subseteq \A_1$ and $\B \cap \A_0 = \emptyset$.

%%%%%%%%%%%%%%%%%%%%%%%%%%%%%%%%%%%%%%%%%%%%%%%%%%%%%%%%%%%%%%%%%%%%%%
\subsection{Assumptions}
%%%%%%%%%%%%%%%%%%%%%%%%%%%%%%%%%%%%%%%%%%%%%%%%%%%%%%%%%%%%%%%%%%%%%%

Assumption 1 generalizes as follows.

\begin{assumption}[Data Generating Process]\label{sa-assump: Sharp DGP}
    Let $t\in\{0,1\}$.
    \begin{enumerate}[label=\normalfont(\roman*),noitemsep,leftmargin=*]

    \item $(Y_1(t), \bX_1^\top)^\top,\ldots, (Y_n(t), \bX_n^\top)^\top$ are independent and identically distributed random vectors.

    \item The distribution of $\bX_i$ has compact support $\X\subseteq\mathbb{R}^d$ with nonempty interior and a Lebesgue density $f_X(\bx)$ that is continuous and bounded away from zero on $\X$.

    \item $\mu_t(\bx) = \E[Y_i(t)| \bX_i = \bx]$ has a $(p+1)$-times continuously differentiable extension to an open neighborhood of $\X$.

    \item $\sigma^2_t(\bx) = \Var[Y_i(t)|\bX_i = \bx]$ is bounded away from zero and continuous on $\X$.

    \item $\sup_{\bx \in \X}\E[|Y_i(t)|^{2+v}|\bX_i = \bx] < \infty$ for some $v \geq 2$.
    \end{enumerate}
\end{assumption}

Assumption 2 generalizes as follows.

\begin{assumption}[Boundary and Kernel]\label{sa-assump: Boundary and Kernel} $ $
    \begin{enumerate}[label=\normalfont(\roman*)]
    \item $\B\subset\setint(\X)$, and $\B$ is compact $(d-1)$-rectifiable, with $\Haus^{d-1}(\B)$ positive and finite.
    \item $K: \reals^d \to [0,\infty)$ is compactly supported and Lipschitz continuous, or $K(\bu) = \Indicator(\bu \in [-1,1]^d)$.
    \item For each $t \in \{0,1\}$, there exists a Lebesgue-measurable set $U \subseteq \reals^d$ with $0<\leb(U)<\infty$, such that $K(\bu) \geq \kappa > 0$ for all $\bu \in U$, $\lambda_{\min} (\int_U \br_p(\bz) \br_p(\bz)^{\top} \dif \bz) > 0$, and
    \begin{align*}
        \liminf_{\bar{h} \downarrow 0}\inf_{\bx \in \B} \int_{U} K(\bu) \Indicator(\bx + \bH \bu \in \A_t) \dif \bu \gtrsim 1,
    \end{align*}
    uniformly over diagonal bandwidth matrices satisfying $\bar{h}/\underline h\lesssim1$.
    \end{enumerate}
\end{assumption}

For $d = 2$, Assumption~\ref{sa-assump: Boundary and Kernel}(i) holds whenever $\B$ is a compact rectifiable curve with positive and finite length. For the special case in the main paper, where $\B$ is piecewise linear with finitely many line segments and $K$ is positive in a neighborhood of the origin, Assumption~\ref{sa-assump: Boundary and Kernel}(iii) also holds: locally, each side of the boundary contains a half-plane or polygonal cone of positive aperture, and the finite number of line segments gives a uniform lower bound over $\B$. More generally, Assumption~\ref{sa-assump: Boundary and Kernel}(iii) is a local nondegeneracy condition linking the geometry of $\B$ to the support and shape of the kernel $K$. The density $f_X(\bx)$ does not appear explicitly in this condition because Assumption~\ref{sa-assump: Sharp DGP}(ii) already requires it to be bounded away from zero on $\X$.

More precisely, Assumption~\ref{sa-assump: Boundary and Kernel}(iii) rules out boundary-kernel configurations in which, after local rescaling by $\bH$, one side of the boundary receives too little kernel-weighted mass to guarantee the local polynomial regression is well-defined uniformly over $\B$. Thus, the condition ensures that the local design is sufficiently rich on each side of the boundary. Lemma~\ref{sa-lem: Invertibility} shows that Assumptions~\ref{sa-assump: Sharp DGP}(ii) and \ref{sa-assump: Boundary and Kernel} imply that the population Gram matrix is full rank for $\bar h$ small enough, uniformly in $\bx\in\B$. This guarantees that the treatment effect estimators are well-defined in large samples under the conditions imposed. Assumption~\ref{sa-assump: Boundary and Kernel}(iii) is stated uniformly over $\B$ in anticipation of our uniform estimation and inference results, but it can be relaxed to hold only pointwise in $\bx\in\B$ for pointwise results.

Assumption 3 generalizes as follows.

\begin{assumption}[Weight Function and Boundary]\label{sa-assump: Weight Function and Boundary}
    Let $w:\B \to (0,\infty)$ satisfy $\sup_{\bx\in\B} w(\bx) < \infty$, $\inf_{\bx\in\B} w(\bx) >0$, and $\int_{\B} w(\bx) \dif\Haus^{d-1}(\bx) < \infty$.
\end{assumption}


%%%%%%%%%%%%%%%%%%%%%%%%%%%%%%%%%%%%%%%%%%%%%%%%%%%%%%%%%%%%%%%%%%%%%%
\subsection{Parameters}
%%%%%%%%%%%%%%%%%%%%%%%%%%%%%%%%%%%%%%%%%%%%%%%%%%%%%%%%%%%%%%%%%%%%%%

The multidimensional generalizations of the three causal parameters studied in the paper are:
\begin{enumerate}
    \item \emph{Boundary Average Treatment Effect Curve} (BATEC):
    \begin{align*}
        \tau(\bx) = \E[Y_i(1) - Y_i(0)|\bX_i = \bx], \qquad \bx\in\B.
    \end{align*}

    \item \emph{Weighted Boundary Average Treatment Effect} (WBATE):
    \begin{align*}
        \tau_{\mathtt{WBATE}} = \frac{\int_{\B} \tau(\bx) w(\bx) \dif\Haus^{d-1}(\bx)}
                                     {\int_{\B} w(\bx) \dif\Haus^{d-1}(\bx)}.
    \end{align*}

    \item \emph{Largest Boundary Average Treatment Effect} (LBATE):
    \begin{align*}
        \tau_{\mathtt{LBATE}} = \sup_{\bx\in\B} \tau(\bx).
    \end{align*}

\end{enumerate}
More generally, this supplemental appendix considers the derivatives of the BATEC parameter
\begin{align*}
    \tau^{(\bnu)}(\bx) = \mu^{(\bnu)}_1(\bx) - \mu^{(\bnu)}_0(\bx), \qquad \bx \in \B,
\end{align*}
and the corresponding weighted and largest derivative parameters, respectively given by
\begin{align*}
    \tau_{\mathtt{WBATE}}^{(\bnu)} = \frac{\int_{\B} \tau^{(\bnu)}(\bx) w(\bx) \dif\Haus^{d-1}(\bx)}
                                     {\int_{\B} w(\bx) \dif\Haus^{d-1}(\bx)}
    \qquad\text{and}\qquad
    \tau_{\mathtt{LBATE}}^{(\bnu)} = \sup_{\bx\in\B} \tau^{(\bnu)}(\bx),
\end{align*}
where $|\bnu| \leq p$.

%%%%%%%%%%%%%%%%%%%%%%%%%%%%%%%%%%%%%%%%%%%%%%%%%%%%%%%%%%%%%%%%%%%%%%
\subsection{Estimators and Other Related Quantities}
%%%%%%%%%%%%%%%%%%%%%%%%%%%%%%%%%%%%%%%%%%%%%%%%%%%%%%%%%%%%%%%%%%%%%%

The estimator of the (derivative of the) BATEC parameters
\begin{align*}
    \widehat{\tau}^{(\bnu)}(\bx)
    = \widehat{\mu}^{(\bnu)}_1(\bx) - \widehat{\mu}^{(\bnu)}_0(\bx)
    \qquad \bx\in\B,
\end{align*}
where $\widehat{\mu}_t^{(\bnu)}(\bx) = \bnu! \be_{\bnu}^\top \widehat{\bbeta}_t(\bx)$ for $t\in\{0,1\}$ with
\begin{align}\label{sa-eq: locpoly estimator}
    \widehat{\bbeta}_t(\bx)
    = \argmin_{\bbeta \in \mathbb{R}^{\mathfrak{p}_p}}
      \En \big[ \big(Y_i - \br_p(\bX_i - \bx)^{\top}\bbeta \big)^2 K_{\bH}(\bX_i - \bx)\Indicator(\bX_i \in \A_t) \big], \qquad \bx \in \B,
\end{align}
with $K_{\bH}(\bu) = |\bH|^{-1} K(\bH^{-1}\bu)$ for a $d$-variate kernel function $K(\cdot)$ and a diagonal bandwidth matrix $\bH = \diag(\bh)$ with $\bh = (h_1, \ldots, h_d)^\top\in\mathbb{R}^d_{++}$. Thus,
\begin{align*}
    |\bH| = \prod_{l = 1}^d h_l
    \qquad\text{and}\qquad
    \bH^{-1} = \diag(h_1^{-1}, \ldots, h_d^{-1}).
\end{align*}
Furthermore, if all bandwidths are proportional to each other, that is, $h_l = c_l h$ for fixed constants $c_l > 0$, $l=1,\ldots,d$, and common bandwidth $h > 0$, then $\bH = h \diag(\bc)$ with $\bc=(c_1,\ldots,c_d)^\top$, and hence
\begin{align*}
    |\bH| = h^d \prod_{l = 1}^d c_l,
    \qquad\text{and}\qquad
    \bH^{-1} = h^{-1} \diag(c_1^{-1}, \ldots, c_d^{-1}).
\end{align*}
Therefore, $|\bH|\asymp h^d$ and $\bH^{-1} \asymp h^{-1} \bI_d$.

Under the assumptions imposed,
\begin{align*}
    \widehat{\bbeta}_{t}(\bx)
    = \bS^{-1} \widehat{\bGamma}_{t,\bx}^{-1} \En \big[\br_p(\bH^{-1}(\bX_i-\bx)) K_{\bH}(\bX_i - \bx) Y_i \Indicator(\bX_i \in \A_t) \big],
    \qquad
    \bS = \diag(\br_p(\bh)),
\end{align*}
and
\begin{align*}
    \widehat{\bGamma}_{t,\bx} = \En \big[ \br_p(\bH^{-1}(\bX_i-\bx)) \br_p(\bH^{-1}(\bX_i-\bx))^{\top} K_{\bH}(\bX_i - \bx) \Indicator(\bX_i \in \A_t)\big],
\end{align*}
where its population analog is
\begin{align*}
    \Gammax = \E \big[\br_p(\bH^{-1}(\bX_i-\bx)) \br_p(\bH^{-1}(\bX_i-\bx))^{\top} K_{\bH}(\bX_i - \bx) \Indicator(\bX_i \in \A_t) \big].
\end{align*}
For use in the proofs, let $\bbeta_t(\bx)$ denote the population Taylor coefficient vector satisfying
\begin{align*}
    \bbeta_t(\bx)^\top\br_p(\bu)
    = \sum_{|\bomega|\leq p}\frac{\mu_t^{(\bomega)}(\bx)}{\bomega!}\bu^{\bomega},
    \qquad \bu\in\reals^d.
\end{align*}
Note that $\norm{\be_{\bnu}^\top \bS^{-1}}_2 = \norm{\be_{\bnu}^\top \bS^{-1}}_{\infty} = \bh^{-\bnu}$. In addition, define
\begin{align*}
    \bQ_{t,\bx} = \En \big[\br_p(\bH^{-1}(\bX_i-\bx)) K_{\bH}(\bX_i - \bx) \Indicator(\bX_i \in \A_t) \varepsilon_i\big],
\end{align*}
where $\varepsilon_i = Y_i - [\Indicator(\bX_i \in \A_0) \mu_0(\bX_i) + \Indicator(\bX_i \in \A_1) \mu_1(\bX_i)] = Y_i - \E[Y_i|\bX_i]$.

For $\bx_1, \bx_2 \in \B$ and $t \in \{0,1\}$, we introduce the following population quantities
\begin{align*}
    \Omega_{\bx_1,\bx_2}^{(\bnu)}
    = \Omega_{0,\bx_1,\bx_2}^{(\bnu)} + \Omega_{1,\bx_1,\bx_2}^{(\bnu)},
    \qquad
    \Omega_{t,\bx_1,\bx_2}^{(\bnu)}
    = \frac{(\bnu!)^2}{n |\bH| \bh^{2\bnu}} \be_{\bnu}^{\top}  \bGamma_{t,\bx_1}^{-1} \bSigma_{t,\bx_1,\bx_2} \bGamma_{t,\bx_2}^{-1} \be_{\bnu},
\end{align*}
\begin{align*}
    \bSigma_{t, \bx_1,\bx_2}
    = |\bH| \E \big[\br_p(\bH^{-1}(\bX_i-\bx_1)) \br_p(\bH^{-1}(\bX_i-\bx_2))^{\top} K_{\bH}(\bX_i - \bx_1) K_{\bH}(\bX_i - \bx_2) \sigma_t^2(\bX_i) \Indicator(\bX_i \in \A_t)\big],
\end{align*}
and their sample analogs
\begin{align*}
    \widehat{\Omega}_{\bx_1,\bx_2}^{(\bnu)}
    = \widehat{\Omega}_{0,\bx_1,\bx_2}^{(\bnu)} + \widehat{\Omega}_{1,\bx_1, \bx_2}^{(\bnu)},
    \qquad
    \widehat{\Omega}_{t,\bx_1,\bx_2}^{(\bnu)}
    = \frac{(\bnu!)^2}{n |\bH| \bh^{2\bnu}} \be_{\bnu}^{\top} \widehat{\bGamma}_{t,\bx_1}^{-1} \widehat{\bSigma}_{t,\bx_1,\bx_2} \widehat{\bGamma}_{t,\bx_2}^{-1} \be_{\bnu},
\end{align*}
\begin{align*}
    \widehat{\bSigma}_{t,\bx_1,\bx_2}
    = |\bH| \En \big[\br_p(\bH^{-1}(\bX_i-\bx_1)) \br_p(\bH^{-1}(\bX_i-\bx_2))^{\top} K_{\bH}(\bX_i - \bx_1) K_{\bH}(\bX_i - \bx_2) \widehat{\varepsilon}_i(\bx_1) \widehat{\varepsilon}_i(\bx_2) \Indicator(\bX_i \in \A_t)\big],
\end{align*}
where $\widehat{\varepsilon}_{i}(\bx) = Y_i - \br_p(\bX_i - \bx)^\top[\Indicator(\bX_i \in \A_0) \widehat{\bbeta}_0(\bx) + \Indicator(\bX_i \in \A_1) \widehat{\bbeta}_1(\bx)]$.

Finally, we introduce more notation for the leading bias and variance of the estimator of the (derivative of the) BATEC parameter. For $\bx \in \B$ and $t \in \{0,1\}$, define
\begin{align*}
    B_{t,\bx}^{(\bnu)}(\bh)
    = \bnu!\be_{\bnu}^{\top} \bGamma_{t,\bx}^{-1} \sum_{|\bomega| = p + 1} \frac{\mu_t^{(\bomega)}(\bx)}{\bomega!}\bh^{\bomega-\bnu}
      \E \big[\br_p(\bH^{-1}(\bX_i-\bx)) (\bH^{-1}(\bX_i-\bx))^{\bomega} K_{\bH}(\bX_i - \bx) \Indicator(\bX_i \in \A_t)\big],
\end{align*}
where $B_{\bx}^{(\bnu)}(\bh) = B_{1,\bx}^{(\bnu)}(\bh) - B_{0,\bx}^{(\bnu)}(\bh)$. When $\bh = \bc h$ for a fixed vector $\bc$ and $h$ is a common bandwidth, $B_{\bx}^{(\bnu)}(\bh)=h^{p+1-|\bnu|}\widetilde B_{\bx}^{(\bnu)}(\bc)$ for some $\widetilde B_{\bx}^{(\bnu)}(\bc)$. Also, define
\begin{align*}
    V_{\bx}^{(\bnu)}  = V_{0,\bx}^{(\bnu)} + V_{1,\bx}^{(\bnu)},
\end{align*}
with
\begin{align*}
    V_{t,\bx}^{(\bnu)}
    = (\bnu!)^2\be_{\bnu}^{\top} \bGamma_{t,\bx}^{-1} \bSigma_{t,\bx,\bx} \bGamma_{t,\bx}^{-1}\be_{\bnu}
    = n |\bH| \bh^{2\bnu} \Omega_{t,\bx,\bx}^{(\bnu)}.
\end{align*}
The corresponding empirical analogs of these leading terms, leaving the unknown derivatives explicit, are
\begin{align*}
    \widehat{B}_{t,\bx}^{(\bnu)}(\bh)
    = \bnu!\be_{\bnu}^{\top} \widehat{\bGamma}_{t,\bx}^{-1} \sum_{|\bomega| = p + 1} \frac{\mu_t^{(\bomega)}(\bx)}{\bomega!}\bh^{\bomega-\bnu}
      \En \big[\br_p(\bH^{-1}(\bX_i-\bx)) (\bH^{-1}(\bX_i-\bx))^{\bomega} K_{\bH}(\bX_i - \bx) \Indicator(\bX_i \in \A_t)\big],
\end{align*}
and $\widehat{B}_{\bx}^{(\bnu)}(\bh) = \widehat{B}_{1,\bx}^{(\bnu)}(\bh) - \widehat{B}_{0,\bx}^{(\bnu)}(\bh)$. Also, $\widehat{V}_{\bx}^{(\bnu)} = \widehat{V}_{0,\bx}^{(\bnu)} + \widehat{V}_{1,\bx}^{(\bnu)}$ with $\widehat{V}_{t,\bx}^{(\bnu)} = n |\bH| \bh^{2\bnu} \widehat{\Omega}_{t,\bx,\bx}^{(\bnu)}$.

%%%%%%%%%%%%%%%%%%%%%%%%%%%%%%%%%%%%%%%%%%%%%%%%%%%%%%%%%%%%%%%%%%%%%%
\subsection{Preliminary Lemmas}
%%%%%%%%%%%%%%%%%%%%%%%%%%%%%%%%%%%%%%%%%%%%%%%%%%%%%%%%%%%%%%%%%%%%%%

Let $\bX = (\bX_1^{\top}, \cdots, \bX_n^{\top})$, and recall that $t \in \{0,1\}$.

\begin{lem}[Invertibility]\label{sa-lem: Invertibility}
    Suppose Assumptions \ref{sa-assump: Sharp DGP}(i,ii) and \ref{sa-assump: Boundary and Kernel} hold. Then,
    \begin{align*}
        \liminf_{\bar{h} \to 0} \inf_{\bx \in \B}\lambda_{\min}(\bGamma_{t,\bx}) > 0.
    \end{align*}
\end{lem}

\begin{lem}[Gram]\label{sa-lem: Gram Matrix}
    Suppose Assumptions \ref{sa-assump: Sharp DGP}(i,ii) and \ref{sa-assump: Boundary and Kernel} hold. If $\frac{\log(1/\bar{h})}{n |\bH|} = o(1)$, then
    \begin{align*}
        \sup_{\bx \in \B} \big\|\hGamma - \Gammax\big\|  \lessp \sqrt{\frac{\log(1/\bar{h})}{n |\bH|}},
        \qquad
        \sup_{\bx \in \B} \big\|\hGamma^{-1} - \Gammax^{-1}\big\| \lessp \sqrt{\frac{\log(1/\bar{h})}{n |\bH|}}.
    \end{align*}
    If, in addition, $\bar{h} = o(1)$, then $\inf_{\bx \in \B}\lambda_{\min}(\hGamma) \gtrsim_{\P} 1$, $\sup_{\bx \in \B}\|\hGamma\|\lessp 1$, and $\sup_{\bx \in \B}\|\hGamma^{-1}\|\lessp 1$.
\end{lem}

\begin{lem}[Linear Approximation]\label{sa-lem: Linear Approximation}
    Suppose Assumptions \ref{sa-assump: Sharp DGP}(i,ii,iv,v) and \ref{sa-assump: Boundary and Kernel} hold. If $\frac{\log(1/\bar{h})}{n |\bH|} = o(1)$, then
    \begin{align*}
        \sup_{\bx \in \B} \big\|\bQ_{t,\bx} \big\| \lessp \sqrt{\frac{\log(1/\bar{h})}{n |\bH|}} + \frac{\log(1/\bar{h})}{n^{\frac{1+v}{2+v}}|\bH|}.
    \end{align*}
    If, in addition,  $\bar{h} = o(1)$, then
    \begin{align*}
        \sup_{\bx \in \B} \big|\widehat{\mu}_t^{(\bnu)}(\bx)   - \E\big[\widehat{\mu}_t^{(\bnu)}(\bx)\big|\bX\big] - \bnu!\be_{\bnu}^{\top}\bS^{-1}\bGamma_{t,\bx}^{-1}\bQ_{t,\bx}\big| \lessp \bh^{-\bnu} \sqrt{\frac{\log(1/\bar{h})}{n |\bH|}} \bigg(\sqrt{\frac{\log(1/\bar{h})}{n |\bH|}} + \frac{\log(1/\bar{h})}{n^{\frac{1+v}{2+v}}|\bH|} \bigg).
    \end{align*}
\end{lem}

\begin{lem}[Covariance]\label{sa-lem: covariance}
    Suppose Assumptions \ref{sa-assump: Sharp DGP} and \ref{sa-assump: Boundary and Kernel} hold. If $\frac{\log(1/\bar{h})}{n^{\frac{v}{2+v}} |\bH|} = o(1)$ and $\bar{h}=o(1)$, then
    \begin{align*}
        \sup_{\bx_1,\bx_2 \in \B} \big\|\widehat{\bSigma}_{t, \bx_1,\bx_2} - \bSigma_{t, \bx_1,\bx_2}\big\|
        \lessp \sqrt{\frac{\log(1/\bar{h})}{n |\bH|}} + \frac{\log(1/\bar{h})}{n^{\frac{v}{2+v}}|\bH|} + \bar{h}^{p+1},
    \end{align*}
    \begin{align*}
        \sup_{\bx_1,\bx_2 \in \B} \big|\widehat{\Omega}_{\bx_1,\bx_2}^{(\bnu)} - \Omega_{\bx_1,\bx_2}^{(\bnu)}\big|
        \lessp (n |\bH|\bh^{2\bnu})^{-1}\bigg(\sqrt{\frac{\log(1/\bar{h})}{n |\bH|}} + \frac{\log(1/\bar{h})}{n^{\frac{v}{2+v}}|\bH|}+ \bar{h}^{p+1}\bigg),
    \end{align*}
    and
    \begin{align*}
        \sup_{\bx \in \B} \big|(\widehat{\Omega}_{\bx,\bx}^{(\bnu)})^{-1/2} - (\Omega_{\bx,\bx}^{(\bnu)})^{-1/2} \big|
        \lessp \sqrt{n |\bH|\bh^{2\bnu}}\bigg(\sqrt{\frac{\log(1/\bar{h})}{n |\bH|}} + \frac{\log(1/\bar{h})}{n^{\frac{v}{2+v}}|\bH|} + \bar{h}^{p+1}\bigg).
    \end{align*}
\end{lem}

\begin{lem}[Bias]\label{sa-lem: bias}
    Suppose Assumptions \ref{sa-assump: Sharp DGP}(i,ii,iii) and \ref{sa-assump: Boundary and Kernel} hold. If $\frac{\log(1/\bar{h})}{n |\bH|} = o(1)$ and $\bar{h} = o(1)$, then
    \begin{align*}
        \sup_{\bx \in \B} \big|\E[\widehat{\mu}_t^{(\bnu)}(\bx)|\bX] - \mu_t^{(\bnu)}(\bx)\big| \lessp \bar{h}^{p+1}\bh^{-\bnu},
    \end{align*}
    implying
    \begin{align*}
        \sup_{\bx \in \B} \big|\E[\widehat{\mu}_t^{(\bnu)}(\bx)|\bX] - \mu_t^{(\bnu)}(\bx) - \widehat{B}_{t,\bx}^{(\bnu)}(\bh)\big| = o_{\P}(\bar{h}^{p+1}\bh^{-\bnu}),
    \end{align*}
    with $\sup_{\bx \in \B} \big|\widehat{B}_{t,\bx}^{(\bnu)}(\bh) - B_{t,\bx}^{(\bnu)}(\bh) \big| \lessp \bar{h}^{p+1}\bh^{-\bnu}\sqrt{\frac{\log(1/\bar{h})}{n |\bH|}}$, and hence $\sup_{\bx \in \B}|B_{t,\bx}^{(\bnu)}(\bh)| \lesssim \bar{h}^{p+1}\bh^{-\bnu}$.
\end{lem}


%%%%%%%%%%%%%%%%%%%%%%%%%%%%%%%%%%%%%%%%%%%%%%%%%%%%%%%%%%%%%%%%%%%%%%
\section{Boundary Average Treatment Effect Curve}
%%%%%%%%%%%%%%%%%%%%%%%%%%%%%%%%%%%%%%%%%%%%%%%%%%%%%%%%%%%%%%%%%%%%%%

%%%%%%%%%%%%%%%%%%%%%%%%%%%%%%%%%%%%%%%%%%%%%%%%%%%%%%%%%%%%%%%%%%%%%%
\subsection{Point Estimation and MSE Expansions}
%%%%%%%%%%%%%%%%%%%%%%%%%%%%%%%%%%%%%%%%%%%%%%%%%%%%%%%%%%%%%%%%%%%%%%

\begin{thm}[Convergence Rates: Sharp BATEC]\label{sa-thm: Convergence Rates: Sharp BATEC}
    Suppose Assumptions \ref{sa-assump: Sharp DGP} and \ref{sa-assump: Boundary and Kernel} hold. If $\log(1/\bar{h})/(n |\bH|) = o(1)$ and $\bar{h} = o(1)$, then
    \begin{align*}
        \big|\widehat{\tau}^{(\bnu)}(\bx) - \tau^{(\bnu)}(\bx) \big|
        \lessp \bh^{-\bnu} \Big(\frac{1}{\sqrt{n |\bH|}} + \frac{1}{n^{\frac{1+v}{2+v}}|\bH|} + \bar{h}^{p+1}\Big)
    \end{align*}
    for $\bx \in \B$, and
    \begin{align*}
        \sup_{\bx \in \B} \big|\widehat{\tau}^{(\bnu)}(\bx) - \tau^{(\bnu)}(\bx)\big|
        \lessp \bh^{-\bnu} \Big(\sqrt{\frac{\log(1/\bar{h})}{ n |\bH|}} + \frac{\log(1/\bar{h})}{n^{\frac{1+v}{2+v}}|\bH|} + \bar{h}^{p+1}\Big).
    \end{align*}
\end{thm}

The conditional mean-squared error (MSE) is
\begin{align*}
    & \operatorname{MSE}^{(\bnu)}(\bx) = \E \big[(\widehat{\tau}^{(\bnu)}(\bx) - \tau^{(\bnu)}(\bx))^2 \big| \bX \big]
\end{align*}
for $\bx\in \B$, and the conditional integrated MSE (IMSE) is
\begin{align*}
    & \operatorname{IMSE}^{(\bnu)} = \int_{\B}  \operatorname{MSE}^{(\bnu)}(\bx) w(\bx) \dif\Haus^{d-1}(\bx),
\end{align*}
where $w(\bx)$ satisfies Assumption \ref{sa-assump: Weight Function and Boundary}.

\begin{thm}[MSE Expansions: Sharp BATEC]\label{sa-thm: MSE Expansions: Sharp BATEC}
    Suppose Assumptions \ref{sa-assump: Sharp DGP}, \ref{sa-assump: Boundary and Kernel} and \ref{sa-assump: Weight Function and Boundary} hold. If $\log(1/\bar{h})/(n^{\frac{v}{2+v}} |\bH|) = o(1)$ and $\bar{h} = o(1)$, then
    \begin{align*}
        \operatorname{MSE}^{(\bnu)}(\bx)
        = \big(B_{\bx}^{(\bnu)}(\bh)\big)^2
          + \frac{1}{n|\bH|\bh^{2\bnu}} V_{\bx}^{(\bnu)}
          + o_{\P} \big(\bar{h}^{2p+2}\bh^{-2\bnu} + (n|\bH|\bh^{2\bnu})^{-1}\big)
    \end{align*}
    for $\bx \in \B$, and
    \begin{align*}
        \operatorname{IMSE}^{(\bnu)}
        = \int_{\B} \big[ \big(B_{\bx}^{(\bnu)}(\bh)\big)^2
           + \frac{1}{n|\bH|\bh^{2\bnu}} V_{\bx}^{(\bnu)} \big] w(\bx) \dif\Haus^{d-1}(\bx)
           + o_{\P} \big(\bar{h}^{2p+2}\bh^{-2\bnu} + (n|\bH|\bh^{2\bnu})^{-1}\big).
    \end{align*}
\end{thm}

Theorem \ref{sa-thm: MSE Expansions: Sharp BATEC} can be used to develop feasible bandwidth selectors for diagonal bandwidth matrices by minimizing the displayed criterion over $\bh\in\mathbb{R}_{++}^d$. If $\bh = \bc h$ for a fixed vector $\bc$ and a common bandwidth $h$, the problem reduces to scalar bandwidth selection. If $\widetilde B_{\bx}^{(\bnu)}(\bc)\neq0$, the asymptotic MSE-optimal common scale is
\begin{align*}
    h_{\operatorname{MSE},\bnu}(\bx)
    = \left(\frac{(d + 2|\bnu|) V_{\bx}^{(\bnu)}}
                 {(2p+2 - 2 |\bnu|) (\widetilde B_{\bx}^{(\bnu)}(\bc))^2} \frac{1}{n}\right)^{\frac{1}{2p+d+2}}
\end{align*}
for $\bx \in \B$. Similarly, if $\int_{\B} (\widetilde B_{\bx}^{(\bnu)}(\bc))^2 w(\bx)\dif\Haus^{d-1}(\bx) \neq 0$, the asymptotic IMSE-optimal common scale is
\begin{align*}
    h_{\operatorname{IMSE},\bnu}
    = \left(\frac{(d + 2 |\bnu|) \int_{\B} V_{\bx}^{(\bnu)} w(\bx)\dif\Haus^{d-1}(\bx)}
                 {(2p+2 - 2|\bnu|) \int_{\B} (\widetilde B_{\bx}^{(\bnu)}(\bc))^2 w(\bx) \dif\Haus^{d-1}(\bx)} \frac{1}{n}\right)^{\frac{1}{2p+d+2}}.
\end{align*}

In practice, the unknown bias and variance quantities can be replaced with (consistent) estimators thereof. For example, the unknown functions $\mu_t^{(\bomega)}(\bx)$ can be estimated using higher-order local polynomial estimators in $\widehat{B}_{\bx}^{(\bnu)}(\bh)$, and
$\widehat{V}_{\bx}^{(\bnu)} = \widehat{V}_{0,\bx}^{(\bnu)} + \widehat{V}_{1,\bx}^{(\bnu)}$
with
$\widehat{V}_{t,\bx}^{(\bnu)} = (\bnu!)^2\be_{\bnu}^{\top} \widehat{\bGamma}_{t,\bx}^{-1} \widehat{\bSigma}_{t,\bx,\bx} \widehat{\bGamma}_{t,\bx}^{-1}\be_{\bnu}$, which corresponds to a standard variance estimator (which is also used for asymptotic inference as discussed below).

Finally, some pointwise statements admit weaker side conditions when the empirical-process classes are restricted to a fixed boundary point. We state the stronger uniform side conditions throughout this section to simplify the exposition.

%%%%%%%%%%%%%%%%%%%%%%%%%%%%%%%%%%%%%%%%%%%%%%%%%%%%%%%%%%%%%%%%%%%%%%
\subsection{Distributional Approximation and Inference}
%%%%%%%%%%%%%%%%%%%%%%%%%%%%%%%%%%%%%%%%%%%%%%%%%%%%%%%%%%%%%%%%%%%%%%

Let $\bV = ((Y_1, \bX_1^{\top}), \cdots, (Y_n, \bX_n^{\top}))^\top$, and recall that $t \in \{0,1\}$. For $|\bnu| \leq p$, define the feasible $t$-statistic
\begin{align*}
    \widehat{\Tstat}^{(\bnu)}(\bx)
    = \frac{\widehat{\tau}^{(\bnu)}(\bx) - \tau^{(\bnu)}(\bx)}{\sqrt{\widehat{\Omega}_{\bx,\bx}^{(\bnu)}}},
    \qquad \bx \in \B.
\end{align*}
The associated $100(1-\alpha)\%$ confidence interval estimator is
\begin{align*}
    \widehat{\CI}^{(\bnu)}_{\alpha}(\bx)
    = \bigg[\;\widehat{\tau}^{(\bnu)}(\bx) - \q_{\alpha} \sqrt{\widehat{\Omega}_{\bx,\bx}^{(\bnu)}}
            \; , \;
              \widehat{\tau}^{(\bnu)}(\bx) + \q_{\alpha} \sqrt{\widehat{\Omega}_{\bx,\bx}^{(\bnu)}}\;\bigg],
\end{align*}
where $\q_{\alpha}$ denotes an appropriate quantile depending on the desired confidence level $\alpha\in(0,1)$, and coverage objective (pointwise versus uniform over $\B$).

\begin{thm}[Confidence Intervals: Sharp BATEC]\label{sa-thm: Confidence Intervals: Sharp BATEC}
    Suppose Assumptions \ref{sa-assump: Sharp DGP} and \ref{sa-assump: Boundary and Kernel} hold. If $n |\bH| \to \infty$ and $n |\bH| \bar{h}^{2(p+1)} =o(1)$, then
    \begin{align*}
        \sup_{u \in \reals} \big|\P \big(\widehat{\Tstat}^{(\bnu)}(\bx) \leq u \big) - \Phi(u)\big| = o(1),
        \qquad \bx \in \B,
    \end{align*}
    and
    \begin{align*}
        \P \big(\tau^{(\bnu)}(\bx) \in \widehat{\CI}_{\alpha}^{(\bnu)}(\bx) \big) = 1 - \alpha + o(1),
        \qquad \bx \in \B,
    \end{align*}
    provided that $\q_{\alpha} = \inf \{c > 0: \P( |\widehat{Z}| \geq c | \bV) \leq \alpha \}$ with $\widehat{Z}|\bV \sim \mathsf{Normal}(0, 1)$.
\end{thm}

For uniform inference, we rely on a new strong approximation result established in Section \ref{sa-sec: Gaussian Strong Approximation}. First, we simplify the statistic $\widehat{\Tstat}^{(\bnu)}$, which is not a sum of independent random variables. Let
\begin{align*}
    \overline{\Tstat}^{(\bnu)}(\bx)
    = \E_n \Big[(\Omega_{\bx,\bx}^{(\bnu)})^{-1/2} \bnu!\be_{\bnu}^{\top} \bS^{-1} \big[\Indicator(\bX_i \in \A_1) \Gammap^{-1} - \Indicator(\bX_i \in \A_0) \Gamman^{-1}\big] \br_p(\bH^{-1}(\bX_i - \bx)) K_{\bH}(\bX_i - \bx)  \varepsilon_i \Big],
\end{align*}
where recall that $\varepsilon_i = Y_i - \sum_{t \in \{0,1\}}\Indicator(\bX_i \in \A_t) \mu_t(\bX_i) = Y_i -  \E[Y_i | \bX_i]$.

\begin{thm}[Stochastic Linearization: Sharp BATEC]\label{sa-thm: Stochastic Linearization: Sharp BATEC}
    Suppose Assumptions \ref{sa-assump: Sharp DGP} and \ref{sa-assump: Boundary and Kernel} hold. If $\log(1/\bar{h})/(n^{v/(2+v)}|\bH|) = o(1)$ and $\bar{h} = o(1)$, then
    \begin{align*}
        \sup_{\bx \in \B} \big|\widehat{\Tstat}^{(\bnu)}(\bx) - \overline{\Tstat}^{(\bnu)}(\bx) \big|
        \lessp \bar{h}^{p+1}\sqrt{n |\bH|} + \sqrt{\log(1/\bar{h})} \Big(\sqrt{\frac{\log(1/\bar{h})}{n |\bH|}} + \frac{\log(1/\bar{h})}{n^{\frac{v}{2+v}}|\bH|}\Big).
    \end{align*}
\end{thm}

We can now exploit the linear structure of $(\overline{\Tstat}^{(\bnu)}(\bx): \bx \in \B)$, that is, an average of i.n.i.d.\ random vectors. Define the following functions indexed by $\bx \in \B$:
\begin{align*}
    g_{\bx}(\bu) = \Indicator(\bu \in \A_1)  \K^{(\bnu)}_1(\bu;\bx) - \Indicator(\bu \in \A_0) \K^{(\bnu)}_0(\bu;\bx),
    \qquad \bu \in \X,
\end{align*}
and
\begin{align*}
    \K^{(\bnu)}_t(\bu;\bx) = n^{-1/2} (\Omega_{\bx,\bx}^{(\bnu)})^{-1/2}\bnu!\be_{\bnu}^{\top} \bS^{-1} \bGamma_{t,\bx}^{-1} \br_p(\bH^{-1}(\bu - \bx)) K_{\bH}(\bu - \bx), \qquad \bu \in \X,\quad t \in \{0,1\}.
\end{align*}
Define the associated class of functions $\G = \{g_{\bx}: \bx \in \B\}$ and $\R = \{\operatorname{Id}\}$, where $\operatorname{Id}(x) = x$, for all $x \in \reals$. Then, the \emph{residual-based empirical process} is
\begin{align*}
    R_n(g,r) = n^{-1/2}\sum_{i = 1}^n \Big[g(\bX_i)r(Y_i) - g(\bX_i) \E[r(Y_i)|\bX_i] \Big], \qquad g \in \G, r \in \R,
\end{align*}
and therefore
\begin{align*}
    \overline{\Tstat}^{(\bnu)}(\bx) = R_n(g_{\bx},\operatorname{Id}), \qquad \bx \in \B.
\end{align*}

Theorem~\ref{sa-thm: Residual Empirical Process Strong Approximation} gives a new Gaussian strong approximation that can be applied to our current setup. Specifically, our new theorem allows polynomial moment bounds on the conditional distribution of $Y_i|\bX_i$.

\begin{thm}[Gaussian Strong Approximation: $\overline{\Tstat}^{(\bnu)}$]\label{sa-thm: Gaussian Strong Approximation: Sharp Tstat}
    Suppose Assumptions \ref{sa-assump: Sharp DGP} and \ref{sa-assump: Boundary and Kernel} hold, and $\X=\prod_{l=1}^d[a_l,b_l]$, with $-\infty<a_l<b_l<\infty$ for $l=1,\ldots,d$. Suppose also that there exists a constant $C > 0$ such that for $t \in \{0,1\}$ and for any $\bx \in \B$, the De Giorgi perimeter of the set $E_{t,\bx} = \{\by \in \A_t: \bH^{-1}(\by - \bx) \in \supp(K)\}$ satisfies  $\mathcal{L}(E_{t,\bx}) \leq C \bar{h}^{d-1}$. For the associated classes $\G$ and $\R=\{\operatorname{Id}\}$ defined above, suppose the singleton regularity condition in part (2) of Theorem~\ref{sa-thm: Residual Empirical Process Strong Approximation} holds. If $\liminf_{n \to \infty} \log \bar{h}/\log n > - \infty$ and $n |\bH| \to \infty$, then (on a possibly enlarged probability space) there exists a mean-zero Gaussian process $Z^{(\bnu)}$ indexed by $\B$ with almost surely continuous sample path such that
    \begin{align*}
        \E\Big[\sup_{\bx \in \B} \big|\overline{\Tstat}^{(\bnu)}(\bx) - Z^{(\bnu)}(\bx) \big| \Big]
        \lesssim (\log n)^{\frac{3}{2}} \bigg(\frac{1}{n |\bH|}\bigg)^{\frac{1}{2d+2}\cdot \frac{v}{v + 2}} + \log (n) \left(\frac{1}{n^{\frac{v}{2+v}}|\bH|}\right)^{\frac{1}{2}},
    \end{align*}
    where $\lesssim$ is up to a universal constant, and $Z^{(\bnu)}$ has the same covariance structure as $\overline{\Tstat}^{(\bnu)}$; that is, $\Cov[\overline{\Tstat}^{(\bnu)}(\bx_1), \overline{\Tstat}^{(\bnu)}(\bx_2)] = \Cov[Z^{(\bnu)}(\bx_1), Z^{(\bnu)}(\bx_2)]$ for all $\bx_1, \bx_2 \in \B$.
\end{thm}

The rectangular support condition in Theorem~\ref{sa-thm: Gaussian Strong Approximation: Sharp Tstat} is imposed only for the process-level strong approximation. It permits an affine rescaling of the score support to $[0,1]^d$; combined with the bounded and bounded-away-from-zero density in Assumption~\ref{sa-assump: Sharp DGP}(ii), $\P_X$ can then be used as the surrogate measure and the Rosenblatt transform provides the normalizing transformation with bounded constants. See \citet[Section 3.1]{Cattaneo-Yu_2025_AOS} for details. This normalization simplifies the verification of that stronger coupling and is not used in the pointwise, aggregate, or supremum-based confidence-band results. Because the local-polynomial estimators use only shrinking neighborhoods of $\B\subset\setint(\X)$, the condition is not practically restrictive for the process-level result in settings where the relevant score support can be represented by a compact rectangle after rescaling.

The De Giorgi perimeter condition in Theorem~\ref{sa-thm: Gaussian Strong Approximation: Sharp Tstat} is used only to verify the total variation requirement of the residual empirical process strong approximation theorem (Theorem~\ref{sa-thm: Residual Empirical Process Strong Approximation}) for the class of boundary-indexed local-polynomial weighting functions. Even when the kernel is smooth, the factors $\Indicator(\bu\in\A_t)$ introduce discontinuities along the assignment boundary. The localized perimeter bound controls the total variation of these discontinuities uniformly over $\bx\in\B$, because the relevant jumps occur on the sets $E_{t,\bx}=\{\by \in \A_t: \bH^{-1}(\by-\bx)\in\supp(K)\}$. The separate singleton regularity condition concerns the truncated conditional means introduced by the polynomial-moment truncation step. Informally, these additional technical restrictions avoid an overly ``wiggly'' assignment boundary and irregular truncation-induced conditional means.

The following lemma shows that the simpler boundary assumptions imposed in the paper are sufficient to verify the De Giorgi perimeter condition for standard compactly supported kernels.

\begin{lem}[Localized Perimeter: Piecewise-Linear Boundary]\label{sa-lem: localized perimeter piecewise linear}
    Suppose $d=2$, $\bH=h\mathbf{I}_2$, $\B$ is piecewise linear with finitely many line segments, $\B$ lies in the interior of $\X$, and $\supp(K)$ is bounded. If either $K(\bu)=\Indicator(\bu\in[-1,1]^2)$ or $\supp(K)$ has finite De Giorgi perimeter, then, for $t\in\{0,1\}$,
    \begin{align*}
        \sup_{\bx\in\B}\mathcal{L}\big(\{\by\in\A_t:h^{-1}(\by-\bx)\in \supp(K)\}\big)
        \lesssim h.
    \end{align*}
\end{lem}

Theorem \ref{sa-thm: Gaussian Strong Approximation: Sharp Tstat} can be used to construct confidence bands for $(\tau^{(\bnu)}(\bx):\bx\in\B)$. Let $(\widehat{Z}^{(\bnu)}(\bx):\bx \in \B)$ be a (conditionally on $\bV$) mean-zero Gaussian process with feasible (conditional) covariance function
\begin{align*}
    \Cov \big[\widehat{Z}^{(\bnu)}(\bx_1),\widehat{Z}^{(\bnu)}(\bx_2) \big|\bV \big]
    = \frac{\widehat{\Omega}_{\bx_1,\bx_2}^{(\bnu)}}
           {\sqrt{\widehat{\Omega}_{\bx_1,\bx_1}^{(\bnu)} \widehat{\Omega}_{\bx_2,\bx_2}^{(\bnu)}}},
    \qquad \bx_1, \bx_2 \in \B.
\end{align*}

\begin{thm}[Confidence Bands: Sharp BATEC; process approximation]\label{sa-thm: Confidence Bands: Sharp BATEC -- process}
    Suppose the assumptions and conditions in Theorem \ref{sa-thm: Gaussian Strong Approximation: Sharp Tstat} hold. If $\liminf_{n \to \infty} \log \bar{h}/\log n > - \infty$, $(\log n)^{3}/(n^{v/(2+v)}|\bH|) = o(1)$ and $n |\bH| \bar{h}^{2(p+1)}\log n = o(1)$, then
    \begin{align*}
        \sup_{u \in \reals} \Big|\P \Big(\sup_{\bx \in \B} \big|\widehat{\Tstat}^{(\bnu)}(\bx)\big| \leq u \Big)
                            - \P \Big(\sup_{\bx \in \B} \big|\widehat{Z}^{(\bnu)}(\bx)\big| \leq u \Big| \bV \Big) \Big| = o_{\P}(1)
    \end{align*}
    and
    \begin{align*}
        \P\Big[\tau^{(\bnu)}(\bx) \in \widehat{\CI}_{\alpha}^{(\bnu)}(\bx), \text{ for all } \bx \in \B \Big] = 1 - \alpha + o(1),
    \end{align*}
    provided that $\q_{\alpha} = \inf \big\{c > 0: \P \big(\sup_{\bx \in \B} \big|\widehat{Z}^{(\bnu)}(\bx)\big|\geq c \big| \bV \big) \leq \alpha \big\}$.
\end{thm}

Theorem~\ref{sa-thm: Gaussian Strong Approximation: Sharp Tstat} gives a strong approximation of the entire process $\overline{\Tstat}^{(\bnu)}$ by a Gaussian process $Z^{(\bnu)}$. This process-level coupling is stronger than what is needed for pointwise normal approximation, and also stronger than the supremum approximation needed for confidence bands. Indeed, as demonstrated in the following theorem, approximating only the supremum of the process using the results in \cite{Chernozhukov-Chetverikov-Kato_2014b_AoS} removes the need for the De Giorgi perimeter condition, and allows for a wider class of boundaries $\B$.

\begin{thm}[Gaussian Approximation for $\sup_{\bx \in \B} |\overline{\Tstat}^{(\bnu)}(\bx)|$]\label{sa-thm: Gaussian Approximation Suprema -- Sharp}
Suppose Assumptions \ref{sa-assump: Sharp DGP} and \ref{sa-assump: Boundary and Kernel} hold. Then for $n$ large enough, on a possibly enriched probability space, there exists a Gaussian process $Z^{(\bnu)}(\bx)$, $\bx \in \B$, with almost sure continuous trajectories such that
\begin{align*}
    \Cov[Z^{(\bnu)}(\bx_1), Z^{(\bnu)}(\bx_2)]
    =
    \Cov[\overline{\Tstat}^{(\bnu)}(\bx_1), \overline{\Tstat}^{(\bnu)}(\bx_2)],
    \qquad \bx_1,\bx_2\in\B,
\end{align*}
and
\begin{align*}
    \P\Big[ \big|\sup_{\bx \in \B} |\overline{\Tstat}^{(\bnu)}(\bx)| - \sup_{\bx \in \B} |Z^{(\bnu)}(\bx)|\big| > \frac{\log n}{(n |\bH|)^{1/6}} + \frac{(\log n)^{5/4}}{(n |\bH|)^{1/4}} + \frac{(\log n)^{3/2}}{(n^{\frac{v}{2+v}} |\bH|)^{1/2}}\Big] \lesssim \frac{1}{\log n} + \frac{\log n}{n}.
\end{align*}
\end{thm}

Using this result, we obtain the following theorem for confidence bands that does not require the De Giorgi perimeter condition.

\begin{thm}[Confidence Bands: Sharp BATEC; suprema approximation]\label{sa-thm: Confidence Bands: Sharp BATEC -- Suprema}
    Suppose Assumptions \ref{sa-assump: Sharp DGP} and \ref{sa-assump: Boundary and Kernel} hold. If $\liminf_{n \to \infty} \log \bar{h}/\log n > -\infty$, $(\log n)^9/(n^{v/(2+v)}|\bH|)=o(1)$, and $\bar{h}^{2(p+1)}n|\bH|\log n=o(1)$, then
    \begin{align*}
        \sup_{u \in \reals} \Big|\P \Big(\sup_{\bx \in \B} \big|\widehat{\Tstat}^{(\bnu)}(\bx)\big| \leq u \Big)
                            - \P \Big(\sup_{\bx \in \B} \big|\widehat{Z}^{(\bnu)}(\bx)\big| \leq u \Big| \bV \Big) \Big| = o_{\P}(1)
    \end{align*}
    and
    \begin{align*}
        \P\Big[\tau^{(\bnu)}(\bx) \in \widehat{\CI}_{\alpha}^{(\bnu)}(\bx), \text{ for all } \bx \in \B \Big] = 1 - \alpha + o(1),
    \end{align*}
    provided that $\q_{\alpha} = \inf \big\{c > 0: \P \big(\sup_{\bx \in \B} \big|\widehat{Z}^{(\bnu)}(\bx)\big|\geq c \big| \bV \big) \leq \alpha \big\}$.
\end{thm}

%%%%%%%%%%%%%%%%%%%%%%%%%%%%%%%%%%%%%%%%%%%%%%%%%%%%%%%%%%%%%%%%%%%%%%
\subsection{Covariate Adjustment}
%%%%%%%%%%%%%%%%%%%%%%%%%%%%%%%%%%%%%%%%%%%%%%%%%%%%%%%%%%%%%%%%%%%%%%

Following \cite{Calonico-Cattaneo-Farrell-Titiunik_2019_RESTAT} and \cite{Calonico-Cattaneo-Farrell-Palomba-Titiunik_2026_wp}, we can incorporate covariate adjustment into our framework. To conserve space, we only give a brief description of the basic idea. Suppose that $\bZ_1,\ldots,\bZ_n$ are pre-intervention covariates of dimension $d_Z\geq1$.

For efficiency improvements, the covariate-adjusted BATEC estimator is
\begin{align*}
    \widetilde{\tau}(\bx) = \be_{\mathfrak{p}_p+1}^{\top}\widetilde{\bgamma}(\bx)
\end{align*}
with
\begin{align*}
    \widetilde{\bgamma}(\bx)
    = \argmin_{\bgamma \in \mathbb{R}^{2\mathfrak{p}_p+d_Z}}
      \En \big[ \big(Y_i - \widetilde{\br}_p(\bX_i-\bx,T_i,\bZ_i)^\top \bgamma \big)^2
                K_h(\bX_i - \bx) \big],
\end{align*}
where $T_i=\Indicator(\bX_i \in \A_1)$, and
\begin{align*}
    \widetilde{\br}_p(\bX_i-\bx,T_i,\bZ_i) = [\br_p(\bX_i - \bx)^{\top},T_i\cdot\br_p(\bX_i - \bx)^{\top},\bZ_i^\top]^\top
\end{align*}
contains the full polynomial basis function for each treatment group but the pre-intervention covariates are not interacted with the treatment indicator.

For heterogeneity analysis, the covariate-adjusted BATEC estimator is
\begin{align*}
    \check{\kappa}(\bx,\bz)
    = \be_{\mathfrak{p}_p+1}^{\top}\check{\bgamma}(\bx)
    + \bz^\top
      \begin{bmatrix}
        \mathbf{0}_{d_Z \times (2+d_Z)\mathfrak{p}_p}
        & \mathbf{I}_{d_Z} & \mathbf{0}_{d_Z \times (\mathfrak{p}_p-1)d_Z}
      \end{bmatrix}
      \check{\bgamma}(\bx)
\end{align*}
with
\begin{align*}
    \check{\bgamma}(\bx)
    = \argmin_{\bgamma \in \mathbb{R}^{2\mathfrak{p}_p(d_Z+1)}}
      \En \big[ \big(Y_i - \check{\br}_p(\bX_i-\bx,T_i,\bZ_i)^\top \bgamma \big)^2
                K_h(\bX_i - \bx) \big],
\end{align*}
where $\mathbf{I}_d$ denotes the $(d\times d)$ identity matrix, $\mathbf{0}_{d_1 \times d_2}$ denotes the $(d_1 \times d_2)$ matrix of zeros, $\bz$ takes values on the support of $\bZ_i$, and
\begin{align*}
    \check{\br}_p(\bX_i-\bx,T_i,\bZ_i)
     = [\br_p(\bX_i - \bx)^{\top},T_i\cdot\br_p(\bX_i - \bx)^{\top}, (\br_p(\bX_i - \bx)\otimes \bZ_i)^\top, T_i \cdot (\br_p(\bX_i - \bx)\otimes \bZ_i)^\top]^\top
\end{align*} contains the full interaction between the polynomial basis function, the treatment assignment indicator, and the pre-intervention covariates.

Employing the theoretical results in the supplemental appendix, it is not difficult to develop valid pointwise and uniform estimation and inference results for these covariate-adjusted BATEC estimators, transformations thereof, and their fuzzy counterparts. We plan to formalize these results in future research.


%%%%%%%%%%%%%%%%%%%%%%%%%%%%%%%%%%%%%%%%%%%%%%%%%%%%%%%%%%%%%%%%%%%%%%
\section{Weighted Boundary Average Treatment Effect}
%%%%%%%%%%%%%%%%%%%%%%%%%%%%%%%%%%%%%%%%%%%%%%%%%%%%%%%%%%%%%%%%%%%%%%

Without loss of generality, we set $\int_{\B} w(\bb) \dif\Haus^{d-1} (\bb) = 1$. For $|\bnu|\leq p$, the sharp (derivative) WBATE parameter is
\begin{align*}
    \tau_{\mathtt{WBATE}}^{(\bnu)} = \int_{\B} \tau^{(\bnu)}(\bb) w(\bb) \dif\Haus^{d-1}(\bb),
\end{align*}
where the weight function satisfies Assumption \ref{sa-assump: Weight Function and Boundary}. The original WBATE parameter is the special case $\bnu=\mathbf{0}$.

The sharp (derivative) WBATE estimator is
\begin{align*}
    \widehat{\tau}_{\mathtt{WBATE}}^{(\bnu)} = \int_\B \widehat{\tau}^{(\bnu)}(\bb) w(\bb) \dif\Haus^{d-1}(\bb),
\end{align*}

Our first lemma in this section studies the conditional bias of $\tau_{\mathtt{WBATE}}^{(\bnu)}$. Let
\begin{align*}
    B_{\mathtt{WBATE}}^{(\bnu)}(\bh) = B_{1,\mathtt{WBATE}}^{(\bnu)}(\bh) - B_{0,\mathtt{WBATE}}^{(\bnu)}(\bh),
    \qquad
    B_{t,\mathtt{WBATE}}^{(\bnu)}(\bh) = \int_{\B} B^{(\bnu)}_{t,\bb}(\bh) w(\bb) \dif\Haus^{d-1}(\bb),
\end{align*}
for $t\in\{0,1\}$.

\begin{lem}[Bias: Sharp WBATE]\label{sa-lem: Bias: Sharp WBATE}
    Suppose Assumption~\ref{sa-assump: Sharp DGP}(i)-(iii), \ref{sa-assump: Boundary and Kernel} and \ref{sa-assump: Weight Function and Boundary} hold. If $\frac{\log(1/\bar{h})}{n |\bH|} = o(1)$ and $\bar{h} = o(1)$, then
    \begin{align*}
        \E[\widehat{\tau}_{\mathtt{WBATE}}^{(\bnu)}|\bX] - \tau_{\mathtt{WBATE}}^{(\bnu)} = B_{\mathtt{WBATE}}^{(\bnu)}(\bh) + o_{\P}(\bar{h}^{p+1}\bh^{-\bnu}).
    \end{align*}
    If, in addition, $\bh = \bc h$ for a fixed vector $\bc$ and common bandwidth $h$, then $B_{\mathtt{WBATE}}^{(\bnu)}(\bh) = h^{p+1-|\bnu|}\widetilde{B}_{\mathtt{WBATE}}^{(\bnu)}(\bc) + o(h^{p+1-|\bnu|})$.
\end{lem}

The next lemma studies the conditional variance of $\tau_{\mathtt{WBATE}}^{(\bnu)}$, and a plug-in estimator thereof. Let
\begin{align*}
    \Omega_{\mathtt{WBATE}}^{(\bnu)} = \Omega_{1,\mathtt{WBATE}}^{(\bnu)} + \Omega_{0,\mathtt{WBATE}}^{(\bnu)},
    \qquad
    \Omega_{t,\mathtt{WBATE}}^{(\bnu)} = \int_{\B} \int_{\B} \Omega^{(\bnu)}_{t,\bb_1,\bb_2} w(\bb_1) w(\bb_2) \dif\Haus^{d-1}(\bb_1) \dif\Haus^{d-1}(\bb_2)
\end{align*}
and impose the aggregate variance nondegeneracy condition
\begin{align}
    \Omega_{\mathtt{WBATE}}^{(\bnu)}
    \gtrsim \frac{\bar{h}^{d-1}}{n|\bH|\bh^{2\bnu}}.
    \label{sa-eq: sharp WBATE aggregate nondegeneracy}
\end{align}
This scalar condition rules out cancellation of the local covariance kernel after integrating the weighted boundary process; it is automatically distinct from the pointwise variance lower bound when derivative aggregates are considered.
For the aggregate variance and inference results below, we also use the local boundary-measure regularity condition
\begin{align}
    \sup_{\bx\in\B}\Haus^{d-1}(\B\cap B(\bx,r)) \leq C_{\B}r^{d-1},
    \qquad 0<r\leq r_0,
    \label{sa-eq: local boundary measure regularity}
\end{align}
for some constants $C_{\B}<\infty$ and $r_0>0$. In the proofs below, write $\nu_{\B,k}$ for the $k$-fold product measure $(\Haus^{d-1})^k$ on $\B^k$.
Also let
\begin{align*}
    \widehat{\Omega}_{\mathtt{WBATE}}^{(\bnu)} = \widehat{\Omega}_{1,\mathtt{WBATE}}^{(\bnu)} + \widehat{\Omega}_{0,\mathtt{WBATE}}^{(\bnu)},
    \qquad
    \widehat{\Omega}_{t,\mathtt{WBATE}}^{(\bnu)} = \int_{\B} \int_{\B} \widehat{\Omega}^{(\bnu)}_{t,\bb_1, \bb_2} w(\bb_1) w(\bb_2) \dif\Haus^{d-1}(\bb_1) \dif\Haus^{d-1}(\bb_2),
\end{align*}
for $t\in\{0,1\}$.

\begin{lem}[Variance: Sharp WBATE]\label{sa-lem: Variance: Sharp WBATE}
    Suppose Assumptions \ref{sa-assump: Sharp DGP}, \ref{sa-assump: Boundary and Kernel} and \ref{sa-assump: Weight Function and Boundary} hold, and \eqref{sa-eq: sharp WBATE aggregate nondegeneracy} and \eqref{sa-eq: local boundary measure regularity} hold. If $\frac{\log(1/\bar{h})}{n |\bH|} = o(1)$ and $\bar{h} = o(1)$, then
    \begin{align*}
        \Var[\widehat{\tau}_{\mathtt{WBATE}}^{(\bnu)}|\bX]
        = \Omega_{\mathtt{WBATE}}^{(\bnu)} + o_{\P}\Big(\frac{\bar{h}^{d-1}}{n|\bH|\bh^{2\bnu}}\Big),
    \end{align*}
    where
    \begin{align*}
        \frac{\bar{h}^{d-1}}{n|\bH|\bh^{2\bnu}} \lesssim \Omega_{\mathtt{WBATE}}^{(\bnu)} \lesssim \frac{\bar{h}^{d-1}}{n|\bH|\bh^{2\bnu}}.
    \end{align*}

    Furthermore,
    \begin{align*}
        \Var[\widehat{\tau}_{\mathtt{WBATE}}^{(\bnu)}|\bX] & = \widehat{\Omega}_{\mathtt{WBATE}}^{(\bnu)} + o_{\P}\Big(\frac{\bar{h}^{d-1}}{n|\bH|\bh^{2\bnu}}\Big).
    \end{align*}
\end{lem}

\begin{thm}[MSE Expansion: Sharp WBATE]\label{sa-thm: MSE Expansion: Sharp WBATE}
    Suppose Assumptions \ref{sa-assump: Sharp DGP}, \ref{sa-assump: Boundary and Kernel} and \ref{sa-assump: Weight Function and Boundary} hold, and \eqref{sa-eq: sharp WBATE aggregate nondegeneracy} and \eqref{sa-eq: local boundary measure regularity} hold. If $\log(1/\bar{h})/(n^{v/(2+v)}|\bH|) = o(1)$ and $\bar{h} = o(1)$, then
    \begin{align*}
        \E [(\widehat{\tau}_{\mathtt{WBATE}}^{(\bnu)} - \tau_{\mathtt{WBATE}}^{(\bnu)})^2|\bX]
        = \Omega_{\mathtt{WBATE}}^{(\bnu)} + \big(B_{\mathtt{WBATE}}^{(\bnu)}(\bh)\big)^2
          + o_{\P}\Big(\frac{\bar{h}^{d-1}}{n|\bH|\bh^{2\bnu}}\Big) + o_{\P}(\bar{h}^{2p+2}\bh^{-2\bnu}).
    \end{align*}
\end{thm}

MSE-optimal diagonal bandwidth selection follows directly from Theorem \ref{sa-thm: MSE Expansion: Sharp WBATE}; scalar bandwidth selectors are obtained as the proportional-bandwidth special case $h_j=c_jh$.

For inference, we consider the feasible $t$-statistics
\begin{align*}
    \widehat{\Tstat}_{\mathtt{WBATE}}^{(\bnu)} = \frac{\widehat{\tau}_{\mathtt{WBATE}}^{(\bnu)} - \tau_{\mathtt{WBATE}}^{(\bnu)}}{\sqrt{\widehat{\Omega}_{\mathtt{WBATE}}^{(\bnu)}}}.
\end{align*}

\begin{thm}[Asymptotic Normality: Sharp WBATE]\label{sa-thm: Asymptotic Normality: Sharp WBATE}
    Suppose Assumptions \ref{sa-assump: Sharp DGP}, \ref{sa-assump: Boundary and Kernel} and \ref{sa-assump: Weight Function and Boundary} hold, and \eqref{sa-eq: sharp WBATE aggregate nondegeneracy} and \eqref{sa-eq: local boundary measure regularity} hold. If $\log(1/\bar{h})/(n^{v/(2+v)}|\bH|) = o(1)$ and $n|\bH|\bar{h}^{2p+2}/\bar{h}^{d-1} = o(1)$, then
    \begin{align*}
        \sup_{u \in \reals} \big|\P (\widehat{\Tstat}_{\mathtt{WBATE}}^{(\bnu)} \leq u )  - \Phi(u) \big| = o(1).
    \end{align*}
\end{thm}

Condition~\eqref{sa-eq: local boundary measure regularity} controls the amount of boundary measure in bandwidth-sized neighborhoods, which is what gives the $\bar h^{d-1}$ effective integration factor in aggregate variances. The piecewise-linear boundary and scalar bandwidth assumptions used in the main paper imply this condition.

%%%%%%%%%%%%%%%%%%%%%%%%%%%%%%%%%%%%%%%%%%%%%%%%%%%%%%%%%%%%%%%%%%%%%%
\section{Largest Boundary Average Treatment Effect}
%%%%%%%%%%%%%%%%%%%%%%%%%%%%%%%%%%%%%%%%%%%%%%%%%%%%%%%%%%%%%%%%%%%%%%

For $|\bnu|\leq p$, we consider the sharp (derivative) LBATE estimand and estimator
\begin{align*}
    \tau_{\mathtt{LBATE}}^{(\bnu)} = \sup_{\bb \in \B} \tau^{(\bnu)}(\bb)
    \qquad\text{and}\qquad
    \widehat{\tau}_{\mathtt{LBATE}}^{(\bnu)} = \sup_{\bb \in \B} \widehat{\tau}^{(\bnu)}(\bb).
\end{align*}
The original LBATE parameter is the special case $\bnu=\mathbf{0}$.

\begin{thm}[Convergence Rate: Sharp LBATE]\label{sa-thm: Convergence Rate: Sharp LBATE}
    Suppose Assumptions \ref{sa-assump: Sharp DGP} and \ref{sa-assump: Boundary and Kernel} hold. If $\log(1/\bar{h})/(n |\bH|) = o(1)$ and $\bar{h} = o(1)$, then
    \begin{align*}
        \big|\widehat{\tau}_{\mathtt{LBATE}}^{(\bnu)} - \tau_{\mathtt{LBATE}}^{(\bnu)} \big|
        \lessp \bh^{-\bnu}\Big(\sqrt{\frac{\log(1/\bar{h})}{ n |\bH|}} + \frac{\log(1/\bar{h})}{n^{\frac{1+v}{2+v}}|\bH|} + \bar{h}^{p+1}\Big).
    \end{align*}
\end{thm}

Recall from Theorem \ref{sa-thm: Confidence Bands: Sharp BATEC -- Suprema} that $(\widehat{Z}^{(\bnu)}(\bx):\bx \in \B)$ is a (conditionally on $\bV$) mean-zero Gaussian process with feasible (conditional) covariance function
\begin{align*}
    \Cov \Big[\widehat{Z}^{(\bnu)}(\bx_1),\widehat{Z}^{(\bnu)}(\bx_2) \Big|\bV \Big]
    = \frac{\widehat{\Omega}_{\bx_1,\bx_2}^{(\bnu)}}
           {\sqrt{\widehat{\Omega}_{\bx_1,\bx_1}^{(\bnu)} \widehat{\Omega}_{\bx_2,\bx_2}^{(\bnu)}}},
    \qquad \bx_1, \bx_2 \in \B.
\end{align*}
Consider the confidence interval
\begin{align*}
    \widehat{\CI}_{\alpha,\mathtt{LBATE}}^{(\bnu)} = \bigg[\sup_{\bb \in \B} \Big(\widehat{\tau}^{(\bnu)}(\bb)  - \q_{\alpha} \sqrt{\widehat{\Omega}_{\bb,\bb}^{(\bnu)}}\Big), \; \sup_{\bb \in \B} \Big(\widehat{\tau}^{(\bnu)}(\bb) + \q_{\alpha} \sqrt{\widehat{\Omega}_{\bb,\bb}^{(\bnu)}}\Big)\bigg],
\end{align*}
where $\q_{\alpha} = \inf \big\{c > 0: \P \big(\sup_{\bx \in \B} \big|\widehat{Z}^{(\bnu)}(\bx)\big|\geq c \big| \bV \big) \leq \alpha \big\}$.

\begin{thm}[Confidence Interval: Sharp LBATE]\label{sa-thm: Confidence Intervals: Sharp LBATE}
    Suppose the assumptions and conditions in Theorem \ref{sa-thm: Confidence Bands: Sharp BATEC -- Suprema} hold. Then
    \begin{align*}
            \P\Big[\tau_{\mathtt{LBATE}}^{(\bnu)} \in \widehat{\CI}_{\alpha,\mathtt{LBATE}}^{(\bnu)} \Big] \geq 1 - \alpha + o(1).
    \end{align*}
\end{thm}


%%%%%%%%%%%%%%%%%%%%%%%%%%%%%%%%%%%%%%%%%%%%%%%%%%%%%%%%%%%%%%%%%%%%%%
\section{Imperfect Compliance}\label{sa-sec: Imperfect Compliance}
%%%%%%%%%%%%%%%%%%%%%%%%%%%%%%%%%%%%%%%%%%%%%%%%%%%%%%%%%%%%%%%%%%%%%%

We extend our previous results to settings with imperfect treatment compliance, where treatment assignment and treatment status may not be equal for some units. See, e.g., \cite{Hernan-Robins_2020_Book} for a textbook introduction.

We generalize our notation and assumptions to explicitly account for imperfect treatment compliance. Let
\begin{align*}
    W_i = \Indicator(\bX_i \in \A_0) \cdot W_i(0) + \Indicator(\bX_i \in \A_1) \cdot W_i(1)
\end{align*}
be the observed treatment status, where $W_i(t)$ denotes the potential treatment status under treatment assignment $t\in\{0,1\}$ for each unit. Accordingly, the observed outcome now is
\begin{align*}
    Y_i = \Indicator(\bX_i \in \A_0) \cdot Y_i(0,W_i(0)) + \Indicator(\bX_i \in \A_1) \cdot Y_i(1,W_i(1)),
\end{align*}
where the potential outcomes are now functions of two arguments: $Y_i(t,r)$ denotes the potential outcome for unit $i$ when this unit is assigned to treatment $t\in\{0,1\}$ and takes treatment status $r\in\{0,1\}$.

The usual (fuzzy) estimands and estimators in settings with imperfect treatment compliance are, respectively,
\begin{align*}
    \zeta(\bx) = \frac{\tau_{Y}(\bx)}{\tau_{W}(\bx)}
    \qquad\text{and}\qquad
    \widehat{\zeta}(\bx) = \frac{\widehat{\tau}_{Y}(\bx)}{\widehat{\tau}_{W}(\bx)},
\end{align*}
where, for each $\bx \in \B$,
\begin{align*}
    \tau_{Y}(\bx) = \E[Y_i(1,W_i(1)) - Y_i(0,W_i(0)) | \bX_i = \bx]
    \qquad\text{and}\qquad
    \tau_{W}(\bx) = \E[W_i(1) - W_i(0) | \bX_i = \bx],
\end{align*}
and
\begin{align*}
    \widehat{\tau}_{Y}(\bx) = \be_{\mathbf{0}}^{\top}\widehat{\bbeta}_{Y,1}(\bx) - \be_{\mathbf{0}}^{\top}\widehat{\bbeta}_{Y,0}(\bx)
    \qquad\text{and}\qquad
    \widehat{\tau}_{W}(\bx) = \be_{\mathbf{0}}^{\top}\widehat{\bbeta}_{W,1}(\bx) - \be_{\mathbf{0}}^{\top}\widehat{\bbeta}_{W,0}(\bx)
\end{align*}
with $\widehat{\bbeta}_{A,t}(\bx)$ denoting the local polynomial fit \eqref{sa-eq: locpoly estimator} when using outcome variable $A \in \{Y,W\}$ on side $t\in\{0,1\}$.

To recycle the notation and assumptions from the sharp BD design setup, we redefine $Y_i(t) = Y_i(t,W_i(t))$ for each $t\in\{0,1\}$, which enables us to reuse Assumption \ref{sa-assump: Sharp DGP} in the fuzzy BD design setting. We also impose the following analogous assumption for the reduced-form potential outcome/status pairs.

\begin{assumption}[Fuzzy DGP]\label{sa-assump: Fuzzy DGP}
    Let $t\in\{0,1\}$.
    \begin{enumerate}[label=\normalfont(\roman*),noitemsep,leftmargin=*]

    \item $(Y_1(t), W_1(t), \bX_1^\top)^\top,\ldots, (Y_n(t), W_n(t), \bX_n^\top)^\top$ are independent and identically distributed random vectors.

    \item $\mu_{W,t}(\bx) = \E[W_i(t)| \bX_i = \bx]$ has a $(p+1)$-times continuously differentiable extension to an open neighborhood of $\X$.

    \item $\bUpsilon_{t}(\bx) = \Cov[(Y_i(t), W_i(t))^{\top}|\bX_i = \bx]$ has entries that are continuous on $\X$.

    %\item $\sup_{\bx \in \X}\E[|W_i(t)|^{2+v}|\bX_i = \bx] < \infty$ for some $v \geq 2$.
    \end{enumerate}
\end{assumption}

More generally, for $|\bnu|\leq p$, we consider Wald-type ratios of derivative discontinuities of the form
\begin{align*}
    \zeta^{(\bnu)}(\bx) = \frac{\tau^{(\bnu)}_{Y}(\bx)}{\tau^{(\bnu)}_{W}(\bx)}
    \qquad\text{and}\qquad
    \widehat{\zeta}^{(\bnu)}(\bx) = \frac{\widehat{\tau}^{(\bnu)}_{Y}(\bx)}{\widehat{\tau}^{(\bnu)}_{W}(\bx)}.
\end{align*}
For $\bnu=\mathbf{0}$ this is the usual fuzzy Wald curve. For $|\bnu|>0$, it is a ratio of derivative jumps, not the derivative of the level Wald ratio. Here, for each $\bx \in \B$,
\begin{align*}
    \tau^{(\bnu)}_{Y}(\bx) = \mu^{(\bnu)}_{Y,1}(\bx) - \mu^{(\bnu)}_{Y,0}(\bx)
    \qquad\text{and}\qquad
    \tau^{(\bnu)}_{W}(\bx) = \mu^{(\bnu)}_{W,1}(\bx) - \mu^{(\bnu)}_{W,0}(\bx),
\end{align*}
and
\begin{align*}
    \widehat{\tau}^{(\bnu)}_{Y}(\bx)
    = \bnu!\be_{\bnu}^{\top}\widehat{\bbeta}_{Y,1}(\bx) - \bnu!\be_{\bnu}^{\top}\widehat{\bbeta}_{Y,0}(\bx)
    \qquad\text{and}\qquad
    \widehat{\tau}_{W}^{(\bnu)}(\bx)
    = \bnu!\be_{\bnu}^{\top}\widehat{\bbeta}_{W,1}(\bx) - \bnu!\be_{\bnu}^{\top}\widehat{\bbeta}_{W,0}(\bx)
\end{align*}
with $\widehat{\bbeta}_{A,t}(\bx)$ denoting the local polynomial fit \eqref{sa-eq: locpoly estimator} when using outcome variable $A \in \{Y,W\}$ on side $t\in\{0,1\}$. That is, under the assumptions imposed,
\begin{align*}
    \widehat{\bbeta}_{A,t}(\bx)
     = \bS^{-1} \widehat{\bGamma}_{t,\bx}^{-1} \En \big[\br_p(\bH^{-1}(\bX_i-\bx)) K_{\bH}(\bX_i - \bx) A_i \Indicator(\bX_i \in \A_t) \big],
    \qquad A\in\{Y,W\},
\end{align*}
corresponds to the local polynomial regression fit when using either $Y_i$ or $W_i$ as outcome. We also define
\begin{align*}
    \bQ_{A,t,\bx} = \En \big[\br_p(\bH^{-1}(\bX_i - \bx)) K_{\bH}(\bX_i - \bx) \Indicator(\bX_i \in \A_t) \varepsilon_{A,i}\big],
\end{align*}
where $\varepsilon_{A,i} = A_i - [\Indicator(\bX_i \in \A_0) \mu_{A,0}(\bX_i) + \Indicator(\bX_i \in \A_1) \mu_{A,1}(\bX_i)] = A_i - \E[A_i|\bX_i]$, where $\mu_{A,t}(\bX_i) = \E[A_i(t)|\bX_i]$ and $A\in\{Y,W\}$. Let $\bbeta_{A,t}(\bx)$ denote the corresponding population Taylor coefficient vector satisfying
\begin{align*}
    \bbeta_{A,t}(\bx)^\top\br_p(\bu)
    = \sum_{|\bomega|\leq p}\frac{\mu_{A,t}^{(\bomega)}(\bx)}{\bomega!}\bu^{\bomega},
    \qquad \bu\in\reals^d.
\end{align*}
Finally, for $A_1,A_2\in\{Y,W\}$, $\bx_1, \bx_2 \in \B$ and $t \in \{0,1\}$, we introduce the following population quantities
\begin{align*}
    \Omega_{A_1,A_2,\bx_1,\bx_2}^{(\bnu)}
    = \Omega_{A_1,A_2,0,\bx_1,\bx_2}^{(\bnu)} + \Omega_{A_1,A_2,1,\bx_1,\bx_2}^{(\bnu)},
    \qquad
    \Omega_{A_1,A_2,t,\bx_1,\bx_2}^{(\bnu)}
    = \frac{(\bnu!)^2}{n |\bH|\bh^{2\bnu}} \be_{\bnu}^{\top}  \bGamma_{t,\bx_1}^{-1} \bSigma_{A_1,A_2,t,\bx_1,\bx_2} \bGamma_{t,\bx_2}^{-1} \be_{\bnu},
\end{align*}
where
\begin{align*}
    \bSigma_{A_1,A_2,t, \bx_1,\bx_2}
    &= |\bH| \E \big[\br_p(\bH^{-1}(\bX_i-\bx_1)) \br_p(\bH^{-1}(\bX_i-\bx_2))^{\top}\\
    &\hspace{1in} \cdot K_{\bH}(\bX_i - \bx_1) K_{\bH}(\bX_i - \bx_2) \sigma_{A_1,A_2,t}^2(\bX_i) \Indicator(\bX_i \in \A_t)\big],
\end{align*}
with $\sigma_{A_1,A_2,t}^2(\bX_i) = \Cov(A_1(t),A_2(t)|\bX_i)$, and their sample analogs
\begin{align*}
    \widehat{\Omega}_{A_1,A_2,\bx_1,\bx_2}^{(\bnu)}
    = \widehat{\Omega}_{A_1,A_2,0,\bx_1,\bx_2}^{(\bnu)} + \widehat{\Omega}_{A_1,A_2,1,\bx_1, \bx_2}^{(\bnu)},
    \qquad
    \widehat{\Omega}_{A_1,A_2,t,\bx_1,\bx_2}^{(\bnu)}
    = \frac{(\bnu!)^2}{n |\bH|\bh^{2\bnu}} \be_{\bnu}^{\top} \widehat{\bGamma}_{t,\bx_1}^{-1} \widehat{\bSigma}_{A_1,A_2,t,\bx_1,\bx_2} \widehat{\bGamma}_{t,\bx_2}^{-1} \be_{\bnu},
\end{align*}
where
\begin{align*}
    \widehat{\bSigma}_{A_1,A_2,t,\bx_1,\bx_2}
    &= |\bH| \En \big[\br_p(\bH^{-1}(\bX_i-\bx_1)) \br_p(\bH^{-1}(\bX_i-\bx_2))^{\top}\\
    &\hspace{1in} \cdot K_{\bH}(\bX_i - \bx_1) K_{\bH}(\bX_i - \bx_2) \widehat{\varepsilon}_{A_1,i}(\bx_1) \widehat{\varepsilon}_{A_2,i}(\bx_2) \Indicator(\bX_i \in \A_t)\big],
\end{align*}
where $\widehat{\varepsilon}_{A,i}(\bx) = A_i - \br_p(\bH^{-1}(\bX_i - \bx))^\top\bS[\Indicator(\bX_i \in \A_0) \widehat{\bbeta}_{A,0}(\bx) + \Indicator(\bX_i \in \A_1) \widehat{\bbeta}_{A,1}(\bx)]$, with $A\in\{Y,W\}$.

%%%%%%%%%%%%%%%%%%%%%%%%%%%%%%%%%%%%%%%%%%%%%%%%%%%%%%%%%%%%%%%%%%%%%%
\subsection{Second-order Linearization}\label{sec: Second-Order Linearization: Fuzzy BATEC}
%%%%%%%%%%%%%%%%%%%%%%%%%%%%%%%%%%%%%%%%%%%%%%%%%%%%%%%%%%%%%%%%%%%%%%

The results in this section will follow from the following exact second-order linearization:
\begin{align*}
    \widehat{\zeta}^{(\bnu)}(\bx) - \zeta^{(\bnu)}(\bx)
%    &= \frac{1}{\tau^{(\bnu)}_W(\bx)} \big( \widehat{\tau}^{(\bnu)}_Y(\bx) - \tau^{(\bnu)}_Y(\bx) \big)
%    - \frac{\tau^{(\bnu)}_Y(\bx)}{\tau^{(\bnu)}_W(\bx)^{2}} \big( \widehat{\tau}^{(\bnu)}_W(\bx) - \tau^{(\bnu)}_W(\bx) \big)\\
%    + \mathfrak{R}_{\mathtt{F}}(\bx)
    &= \mathfrak{w}(\bx)^\top \mathfrak{T}(\bx) + \mathfrak{R}_{\mathtt{F}}(\bx)
\end{align*}
with
\begin{align*}
    \mathfrak{w}(\bx)
    &= \Big[ \frac{1}{\tau^{(\bnu)}_W(\bx)}, -\frac{\tau^{(\bnu)}_Y(\bx)}{\tau^{(\bnu)}_W(\bx)^{2}} \Big]^\top
    \qquad\text{and}\qquad
    \mathfrak{T}(\bx)
    = \Big[ \widehat{\tau}^{(\bnu)}_Y(\bx) - \tau^{(\bnu)}_Y(\bx), \widehat{\tau}^{(\bnu)}_W(\bx) - \tau^{(\bnu)}_W(\bx) \Big]^\top,
\end{align*}
and second-order error
\begin{align*}
    \mathfrak{R}_{\mathtt{F}}(\bx)
    &= \frac{\tau^{(\bnu)}_Y(\bx)}{\tau^{(\bnu)}_W(\bx)^{2}\widehat{\tau}^{(\bnu)}_W(\bx)}
       \big(\widehat{\tau}^{(\bnu)}_W(\bx) - \tau^{(\bnu)}_W(\bx)\big)^2
     - \frac{1}{\tau^{(\bnu)}_W(\bx)\widehat{\tau}^{(\bnu)}_W(\bx)}
       \big(\widehat{\tau}^{(\bnu)}_Y(\bx) - \tau^{(\bnu)}_Y(\bx) \big)
       \big(\widehat{\tau}^{(\bnu)}_W(\bx) - \tau^{(\bnu)}_W(\bx) \big).
\end{align*}

For example, under the assumptions imposed, if $\inf_{\bx\in\B} |\tau^{(\bnu)}_W(\bx)| > 0$, $\log(1/\bar{h})/(n^{\frac{1+v}{2+v}}|\bH|)=o(1)$, $\bar{h}=o(1)$, and
\begin{align*}
    \bh^{-\bnu}\bigg(\sqrt{\frac{\log(1/\bar{h})}{n|\bH|}}+\frac{\log(1/\bar{h})}{n^{\frac{1+v}{2+v}}|\bH|}+\bar{h}^{p+1}\bigg)=o(1),
\end{align*}
then two applications of Theorem \ref{sa-thm: Convergence Rates: Sharp BATEC} (taking as outcomes $Y$ and $W$) imply
\begin{align}\label{sa-eq: second order term -- fuzzy}
    \nonumber \sup_{\bx\in\B} |\mathfrak{R}_{\mathtt{F}}(\bx)|
    &\lesssim_\P \sup_{\bx\in\B} |\widehat{\tau}^{(\bnu)}_W(\bx) - \tau^{(\bnu)}_W(\bx)|^2
                + \sup_{\bx\in\B} |\widehat{\tau}^{(\bnu)}_Y(\bx) - \tau^{(\bnu)}_Y(\bx)|
                  \sup_{\bx\in\B} |\widehat{\tau}^{(\bnu)}_W(\bx) - \tau^{(\bnu)}_W(\bx)|\\
    &\lesssim_\P \bh^{-2\bnu} \Big(\sqrt{\frac{\log(1/\bar{h})}{ n |\bH|}} + \frac{\log(1/\bar{h})}{n^{\frac{1+v}{2+v}}|\bH|} + \bar{h}^{p+1}\Big)^2,
\end{align}
thereby demonstrating that the second-order term $\mathfrak{R}_{\mathtt{F}}(\bx)$ is (uniformly) higher-order relative to the leading term $\mathfrak{w}(\bx)^\top \mathfrak{T}(\bx)$.


%%%%%%%%%%%%%%%%%%%%%%%%%%%%%%%%%%%%%%%%%%%%%%%%%%%%%%%%%%%%%%%%%%%%%%
\subsection{Fuzzy BATEC}
%%%%%%%%%%%%%%%%%%%%%%%%%%%%%%%%%%%%%%%%%%%%%%%%%%%%%%%%%%%%%%%%%%%%%%

\subsubsection{Point Estimation and AMSE Expansions}

\begin{thm}[Convergence Rates: Fuzzy BATEC]\label{sa-thm: Convergence Rates: Fuzzy BATEC}
    Suppose Assumptions \ref{sa-assump: Sharp DGP}, \ref{sa-assump: Boundary and Kernel}, and \ref{sa-assump: Fuzzy DGP} hold, and $\inf_{\bx\in\B} |\tau^{(\bnu)}_W(\bx)| > 0$. If $\log(1/\bar{h})/(n^{\frac{1+v}{2+v}} |\bH|) = o(1)$, $\bar{h} = o(1)$, and
    \begin{align*}
        \bh^{-\bnu}\bigg(\sqrt{\frac{\log(1/\bar{h})}{n|\bH|}}+\frac{\log(1/\bar{h})}{n^{\frac{1+v}{2+v}}|\bH|}+\bar{h}^{p+1}\bigg)=o(1),
    \end{align*}
    then
    \begin{align*}
        \big| \widehat{\zeta}^{(\bnu)}(\bx) - \zeta^{(\bnu)}(\bx) \big|
        \lessp \bh^{-\bnu} \Big(\frac{1}{\sqrt{n |\bH|}} + \frac{1}{n^{\frac{1+v}{2+v}}|\bH|} + \bar{h}^{p+1}\Big)
    \end{align*}
    for $\bx \in \B$, and
    \begin{align*}
        \sup_{\bx \in \B} \big| \widehat{\zeta}^{(\bnu)}(\bx) - \zeta^{(\bnu)}(\bx) \big|
        \lessp \bh^{-\bnu} \Big(\sqrt{\frac{\log(1/\bar{h})}{ n |\bH|}} + \frac{\log(1/\bar{h})}{n^{\frac{1+v}{2+v}}|\bH|} + \bar{h}^{p+1}\Big).
    \end{align*}
\end{thm}

The approximate (conditional) mean-squared error (AMSE) is
\begin{align*}
    \operatorname{AMSE}_{\mathtt{F}}^{(\bnu)}(\bx)
    = \E \big[(\mathfrak{w}(\bx)^\top \mathfrak{T}(\bx))^2 \big| \bX \big]
    = \mathfrak{w}(\bx)^\top \E \big[\mathfrak{T}(\bx) \mathfrak{T}(\bx)^\top \big| \bX \big] \mathfrak{w}(\bx)
\end{align*}
for $\bx\in \B$, and the approximate (conditional) integrated MSE (AIMSE) is
\begin{align*}
    & \operatorname{AIMSE}_{\mathtt{F}}^{(\bnu)} = \int_{\B}  \operatorname{AMSE}_{\mathtt{F}}^{(\bnu)}(\bx) w(\bx) \dif\Haus^{d-1}(\bx),
\end{align*}
where $w(\bx)$ satisfies Assumption \ref{sa-assump: Weight Function and Boundary}. These constructions disregard the higher-order term $\mathfrak{R}_{\mathtt{F}}(\bx)$.

To state the AMSE expansions with imperfect compliance, we generalize the notation introduced previously for the leading bias and variance. For the bias, letting $A\in\{Y,W\}$, we define
\begin{align*}
    B_{A,\bx}^{(\bnu)}(\bh) = B_{A,1,\bx}^{(\bnu)}(\bh) - B_{A,0,\bx}^{(\bnu)}(\bh),
\end{align*}
\begin{align*}
    B_{A,t,\bx}^{(\bnu)}(\bh)
    = \bnu!\be_{\bnu}^{\top} \bGamma_{t,\bx}^{-1} \sum_{|\bomega| = p + 1} \frac{\mu_{A,t}^{(\bomega)}(\bx)}{\bomega!}\bh^{\bomega-\bnu}
      \E \big[\br_p(\bH^{-1}(\bX_i-\bx)) (\bH^{-1}(\bX_i - \bx))^{\bomega} K_{\bH}(\bX_i - \bx) \Indicator(\bX_i\in\A_t) \big],
\end{align*}
where recall that $\mu_{A,t}(\bx) = \E[A_i(t)|\bX_i=\bx]$. Therefore, the leading bias takes the form:
\begin{align}\label{sa-eq: leading bias fuzzy}
    B_{\mathtt{F},\bx}^{(\bnu)}(\bh)
    = \mathfrak{w}(\bx)^\top \bB_{\bx}^{(\bnu)}(\bh),
    \qquad
    \bB_{\bx}^{(\bnu)}(\bh) = \big( B_{Y,\bx}^{(\bnu)}(\bh), B_{W,\bx}^{(\bnu)}(\bh) \big)^\top.
\end{align}
When $\bh = \bc h$ for a fixed vector $\bc$ and $h$ is a common bandwidth, $B_{\mathtt{F},\bx}^{(\bnu)}(\bh)=h^{p+1-|\bnu|}\widetilde B_{\mathtt{F},\bx}^{(\bnu)}(\bc)$ for some $\widetilde B_{\mathtt{F},\bx}^{(\bnu)}(\bc)$.

For the variance, letting $A_1,A_2\in\{Y,W\}$, we define
\begin{align*}
    V_{A_1,A_2,\bx}^{(\bnu)}  = V_{A_1,A_2,0,\bx}^{(\bnu)} + V_{A_1,A_2,1,\bx}^{(\bnu)},
\end{align*}
where
\begin{align*}
    V_{A_1,A_2,t,\bx}^{(\bnu)}
    = (\bnu!)^2\be_{\bnu}^{\top} \bGamma_{t,\bx}^{-1} \bSigma_{A_1,A_2,t,\bx,\bx} \bGamma_{t,\bx}^{-1}\be_{\bnu}
    = n|\bH|\bh^{2\bnu} \Omega_{A_1,A_2,t,\bx,\bx}^{(\bnu)}.
\end{align*}
Therefore, the leading variance takes the form:
\begin{align*}
    V_{\mathtt{F},\bx}^{(\bnu)}
    = \mathfrak{w}(\bx)^\top \mathbf{C}_{\bx}^{(\bnu)} \mathfrak{w}(\bx),
    \qquad
    \mathbf{C}_{\bx}^{(\bnu)}
    = \begin{pmatrix}
        V_{Y,Y,\bx}^{(\bnu)} & V_{Y,W,\bx}^{(\bnu)} \\
        V_{W,Y,\bx}^{(\bnu)} & V_{W,W,\bx}^{(\bnu)}
      \end{pmatrix}.
\end{align*}
For inference, we impose $\inf_{\bx\in\B}V_{\mathtt{F},\bx}^{(\bnu)}>0$ directly. This scalar nondegeneracy condition applies to the linearized fuzzy Wald score and does not require the reduced-form conditional covariance matrix of $(Y_i(t),W_i(t))$ to be uniformly positive definite.

\begin{thm}[AMSE Expansions: Fuzzy BATEC]\label{sa-thm: MSE Expansions: Fuzzy BATEC}
    Suppose Assumptions \ref{sa-assump: Sharp DGP}, \ref{sa-assump: Boundary and Kernel}, \ref{sa-assump: Weight Function and Boundary} and \ref{sa-assump: Fuzzy DGP} hold, and $\inf_{\bx\in\B}|\tau_W^{(\bnu)}(\bx)|>0$. If $\log(1/\bar{h})/(n^{\frac{v}{2+v}} |\bH|) = o(1)$ and $\bar{h} = o(1)$, then
    \begin{align*}
        \operatorname{AMSE}_{\mathtt{F}}^{(\bnu)}(\bx)
        = \big(B_{\mathtt{F},\bx}^{(\bnu)}(\bh)\big)^2
          + \frac{1}{n|\bH|\bh^{2\bnu}} V_{\mathtt{F},\bx}^{(\bnu)}
          + o_{\P} \big(\bar{h}^{2p+2}\bh^{-2\bnu} + (n|\bH|\bh^{2\bnu})^{-1}\big)
    \end{align*}
    for $\bx \in \B$, and
    \begin{align*}
        \operatorname{AIMSE}_{\mathtt{F}}^{(\bnu)}
        = \int_{\B} \Big[ \big(B_{\mathtt{F},\bx}^{(\bnu)}(\bh)\big)^2
           + \frac{1}{n|\bH|\bh^{2\bnu}} V_{\mathtt{F},\bx}^{(\bnu)} \Big] w(\bx) \dif\Haus^{d-1}(\bx)
           + o_{\P} \big(\bar{h}^{2p+2}\bh^{-2\bnu} + (n|\bH|\bh^{2\bnu})^{-1}\big).
    \end{align*}
\end{thm}

Theorem \ref{sa-thm: MSE Expansions: Fuzzy BATEC} can be used to develop feasible bandwidth selectors for diagonal bandwidth matrices by minimizing the displayed criterion over $\bh\in\mathbb{R}_{++}^d$. If $\bh=\bc h$ for a fixed vector $\bc$ and a common bandwidth $h$, the problem reduces to scalar bandwidth selection. If $\widetilde B_{\mathtt{F},\bx}^{(\bnu)}(\bc) \neq 0$, the approximate MSE-optimal common scale is
\begin{align*}
    h_{\mathtt{F},\operatorname{AMSE},\bnu}(\bx)
    = \left(\frac{(d + 2|\bnu|) V_{\mathtt{F},\bx}^{(\bnu)}}
                 {(2p+2 - 2 |\bnu|) (\widetilde B_{\mathtt{F},\bx}^{(\bnu)}(\bc))^2} \frac{1}{n}\right)^{\frac{1}{2p+d+2}}
\end{align*}
for $\bx \in \B$. Similarly, if $\int_{\B} (\widetilde B_{\mathtt{F},\bx}^{(\bnu)}(\bc))^2 w(\bx)\dif\Haus^{d-1}(\bx) \neq 0$, the approximate IMSE-optimal common scale is
\begin{align*}
    h_{\mathtt{F},\operatorname{AIMSE},\bnu}
    = \left(\frac{(d + 2 |\bnu|) \int_{\B} V_{\mathtt{F},\bx}^{(\bnu)} w(\bx)\dif\Haus^{d-1}(\bx)}
                 {(2p+2 - 2|\bnu|) \int_{\B} (\widetilde B_{\mathtt{F},\bx}^{(\bnu)}(\bc))^2 w(\bx) \dif\Haus^{d-1}(\bx)} \frac{1}{n}\right)^{\frac{1}{2p+d+2}}.
\end{align*}

In practice, the unknown bias and variance quantities can be replaced with (consistent) estimators thereof, as discussed before.


%%%%%%%%%%%%%%%%%%%%%%%%%%%%%%%%%%%%%%%%%%%%%%%%%%%%%%%%%%%%%%%%%%%%%%
\subsubsection{Distributional Approximation and Inference}
%%%%%%%%%%%%%%%%%%%%%%%%%%%%%%%%%%%%%%%%%%%%%%%%%%%%%%%%%%%%%%%%%%%%%%

Let $\bV = ((Y_1,W_1,\bX_1^{\top}), \cdots, (Y_n,W_n,\bX_n^{\top}))^{\top}$, and recall that $t \in \{0,1\}$. For $|\bnu| \leq p$, define the feasible $t$-statistic
\begin{align*}
    \widehat{\Tstat}_\mathtt{F}^{(\bnu)}(\bx)
    = \frac{\widehat{\zeta}^{(\bnu)}(\bx) - \zeta^{(\bnu)}(\bx)}{\sqrt{\widehat{\Omega}_{\mathtt{F},\bx,\bx}^{(\bnu)}}},
    \qquad \bx \in \B,
\end{align*}
where we define
\begin{align*}
    \widehat{\mathfrak{w}}(\bx)
    &= \Big[ \frac{1}{\widehat{\tau}^{(\bnu)}_W(\bx)}, -\frac{\widehat{\tau}^{(\bnu)}_Y(\bx)}{\big(\widehat{\tau}^{(\bnu)}_W(\bx)\big)^2} \Big]^\top,
\end{align*}
and
\begin{align*}
    \widehat{\Omega}_{\mathtt{F},\bx_1,\bx_2}^{(\bnu)}
    = \widehat{\mathfrak{w}}(\bx_1)^\top \widehat{\bTheta}_{\ttF, \bx_1, \bx_2}^{(\bnu)} \widehat{\mathfrak{w}}(\bx_2),
    \qquad
    \widehat{\bTheta}_{\ttF, \bx_1, \bx_2}^{(\bnu)}
    = \begin{pmatrix}
        \widehat{\Omega}_{Y,Y,\bx_1, \bx_2}^{(\bnu)} & \widehat{\Omega}_{Y,W,\bx_1, \bx_2}^{(\bnu)} \\
        \widehat{\Omega}_{W,Y,\bx_1, \bx_2}^{(\bnu)} & \widehat{\Omega}_{W,W,\bx_1, \bx_2}^{(\bnu)}
      \end{pmatrix}, \qquad \bx_1, \bx_2 \in \B.
\end{align*}
We also define the population version of the denominator by
\begin{align*}
    \Omega_{\mathtt{F},\bx_1, \bx_2}^{(\bnu)}
    = \mathfrak{w}(\bx_1)^\top \bTheta_{\ttF, \bx_1, \bx_2}^{(\bnu)} \mathfrak{w}(\bx_2),
    \qquad
    \bTheta_{\ttF, \bx_1, \bx_2}^{(\bnu)}
    = \begin{pmatrix}
        \Omega_{Y,Y,\bx_1, \bx_2}^{(\bnu)} & \Omega_{Y,W,\bx_1, \bx_2}^{(\bnu)} \\
        \Omega_{W,Y,\bx_1, \bx_2}^{(\bnu)} & \Omega_{W,W,\bx_1, \bx_2}^{(\bnu)}
      \end{pmatrix}, \qquad \bx_1, \bx_2 \in \B.
\end{align*}

\begin{lem}[Covariance: Fuzzy BATEC]\label{sa-lem: Covariance: Fuzzy BATEC}
    Suppose Assumptions \ref{sa-assump: Sharp DGP}, \ref{sa-assump: Boundary and Kernel} and \ref{sa-assump: Fuzzy DGP} hold, $\inf_{\bx\in\B}|\tau_W^{(\bnu)}(\bx)|>0$, and $\inf_{\bx\in\B}V_{\mathtt{F},\bx}^{(\bnu)}>0$. If $\frac{\log(1/\bar{h})}{n^{\frac{v}{2+v}} |\bH|} = o(1)$, $\bar{h}=o(1)$, and
    \begin{align*}
        \bh^{-\bnu}\bigg(\sqrt{\frac{\log(1/\bar{h})}{n |\bH|}}+\frac{\log(1/\bar{h})}{n^{\frac{v}{2+v}}|\bH|}+\bar{h}^{p+1}\bigg)=o(1),
    \end{align*}
    then
    \begin{gather*}
        \sup_{\bx_1,\bx_2 \in \B} \big|\widehat{\Omega}_{\ttF, \bx_1,\bx_2}^{(\bnu)} - \Omega_{\ttF, \bx_1,\bx_2}^{(\bnu)}\big|
        \lessp \frac{\bh^{-\bnu}}{n |\bH|\bh^{2\bnu}}\bigg(\sqrt{\frac{\log(1/\bar{h})}{n |\bH|}} + \frac{\log(1/\bar{h})}{n^{\frac{v}{2+v}}|\bH|}+ \bar{h}^{p+1}\bigg), \\
        \sup_{\bx \in \B} \big|(\widehat{\Omega}_{\ttF, \bx,\bx}^{(\bnu)})^{-1/2} - (\Omega_{\ttF, \bx,\bx}^{(\bnu)})^{-1/2} \big|
        \lessp \bh^{-\bnu}\sqrt{n |\bH|\bh^{2\bnu}}\bigg(\sqrt{\frac{\log(1/\bar{h})}{n |\bH|}} + \frac{\log(1/\bar{h})}{n^{\frac{v}{2+v}}|\bH|} + \bar{h}^{p+1}\bigg).
    \end{gather*}
\end{lem}

The associated $100(1-\alpha)\%$ confidence interval estimator is
\begin{align*}
    \widehat{\CI}^{(\bnu)}_{\mathtt{F},\alpha}(\bx)
    = \bigg[\;\widehat{\zeta}^{(\bnu)}(\bx) - \q_{\alpha} \sqrt{\widehat{\Omega}_{\mathtt{F},\bx,\bx}^{(\bnu)}}
            \; , \;
              \widehat{\zeta}^{(\bnu)}(\bx) + \q_{\alpha} \sqrt{\widehat{\Omega}_{\mathtt{F},\bx,\bx}^{(\bnu)}}\;\bigg],
\end{align*}
where $\q_{\alpha}$ denotes an appropriate quantile depending on the desired confidence level $\alpha\in(0,1)$, and coverage objective (pointwise versus uniform over $\B$).

The following theorem establishes pointwise asymptotic normality and validity of confidence intervals.

\begin{thm}[Confidence Intervals: Fuzzy BATEC]\label{sa-thm: Confidence Intervals: Fuzzy BATEC}
    For a fixed boundary point $\bx\in\B$, suppose Assumptions \ref{sa-assump: Sharp DGP}, \ref{sa-assump: Boundary and Kernel} and \ref{sa-assump: Fuzzy DGP} hold, $|\tau^{(\bnu)}_W(\bx)| > 0$, and $V_{\mathtt{F},\bx}^{(\bnu)}>0$. If $n |\bH| \to \infty$, $n|\bH|\bh^{2\bnu}\to\infty$, and $n |\bH| \bar{h}^{2(p+1)} \to 0$, then
    \begin{align*}
        \sup_{u \in \reals} \Big|\P \big(\widehat{\Tstat}_\mathtt{F}^{(\bnu)}(\bx) \leq u \big) - \Phi(u)\Big| = o(1),
    \end{align*}
    and
    \begin{align*}
        \P \big(\zeta^{(\bnu)}(\bx) \in \widehat{\CI}_{\mathtt{F},\alpha}^{(\bnu)}(\bx) \big) = 1 - \alpha + o(1),
    \end{align*}
    provided that $\q_{\alpha} = \inf \{c > 0: \P( |\widehat{Z}| \geq c | \bV) \leq \alpha \}$ with $\widehat{Z}|\bV \sim \mathsf{Normal}(0, 1)$.
\end{thm}

For uniform inference, we rely on the strong approximation result established in Section~\ref{sa-sec: Gaussian Strong Approximation}. Since $\widehat{\Tstat}^{(\bnu)}_{\mathtt{F}}$ is not directly a sum of independent terms, we introduce the linearized statistic
\begin{align*}
    \overline{\Tstat}_{\mathtt{F}}^{(\bnu)}(\bx)
    &=
    \E_n\Big[
    (\Omega_{\mathtt{F},\bx,\bx}^{(\bnu)})^{-1/2}
    \bnu!\be_{\bnu}^{\top}\bS^{-1}
    \big[\Indicator(\bX_i\in\A_1)\bGamma_{1,\bx}^{-1}-\Indicator(\bX_i\in\A_0)\bGamma_{0,\bx}^{-1}\big]
    \br_p(\bH^{-1}(\bX_i-\bx))\\
    &\hspace{1in} \cdot
    K_{\bH}(\bX_i-\bx) \mathfrak{w}(\bx)^\top \boldsymbol{\varepsilon}_{i}
    \Big],
\end{align*}
where $\boldsymbol{\varepsilon}_i = (\varepsilon_{Y,i}, \varepsilon_{W,i})^\top = (Y_i - \E[Y_i|\bX_i], W_i - \E[W_i|\bX_i])^\top$ is the vector of regression errors.


The following theorem quantifies the approximation error of the stochastic linearization:

\begin{thm}[Stochastic Linearization: Fuzzy BATEC]\label{sa-thm: Stochastic Linearization: Fuzzy BATEC}
    Suppose Assumptions \ref{sa-assump: Sharp DGP}, \ref{sa-assump: Boundary and Kernel} and \ref{sa-assump: Fuzzy DGP} hold, $\inf_{\bx \in \B}|\tau^{(\bnu)}_W(\bx)| > 0$, and $\inf_{\bx\in\B}V_{\mathtt{F},\bx}^{(\bnu)}>0$. If $\log(1/\bar{h})/(n^{v/(2+v)}|\bH|) = o(1)$, $\bar{h} = o(1)$, and
    \begin{align*}
        \bh^{-\bnu}\bigg(\sqrt{\frac{\log(1/\bar{h})}{n |\bH|}}+\frac{\log(1/\bar{h})}{n^{\frac{v}{2+v}}|\bH|}+\bar{h}^{p+1}\bigg)=o(1),
    \end{align*}
    then
    \begin{align*}
        \sup_{\bx \in \B} \Big|\widehat{\Tstat}^{(\bnu)}_{\mathtt{F}}(\bx) - \overline{\Tstat}^{(\bnu)}_{\mathtt{F}}(\bx) \Big|
        \lessp \bar{h}^{p+1}\sqrt{n |\bH|} + \bh^{-\bnu}\; \sqrt{\log(1/\bar{h})} \Big(\sqrt{\frac{\log(1/\bar{h})}{n |\bH|}} +  \frac{\log(1/\bar{h})}{n^{\frac{v}{2+v}}|\bH|} + \bar{h}^{p+1}\Big).
    \end{align*}
\end{thm}

Next, we present a coupling result for the supremum of the fuzzy linearized T-statistics. Unlike Theorem~\ref{sa-thm: Gaussian Strong Approximation: Sharp Tstat}, this result uses a Gaussian approximation for suprema and therefore requires only VC-type entropy and moment bounds, not the total-variation/perimeter condition used for the process-level coupling.

\begin{thm}[Gaussian Approximation for $\sup_{\bx \in \B} |\overline{\Tstat}^{(\bnu)}_{\mathtt{F}}(\bx)|$]\label{sa-thm: Gaussian Approximation Suprema -- Fuzzy}
Suppose Assumptions \ref{sa-assump: Sharp DGP}, \ref{sa-assump: Boundary and Kernel} and \ref{sa-assump: Fuzzy DGP} hold, $\inf_{\bx \in \B}|\tau^{(\bnu)}_W(\bx)| > 0$, and $\inf_{\bx\in\B}V_{\mathtt{F},\bx}^{(\bnu)}>0$. Then for $n$ large enough, on a possibly enriched probability space, there exists a Gaussian process $Z^{(\bnu)}_{\mathtt{F}}(\bx), \bx \in \B$ with almost sure continuous trajectories such that
\begin{align*}
    \Cov[Z^{(\bnu)}_{\mathtt{F}}(\bx_1), Z^{(\bnu)}_{\mathtt{F}}(\bx_2)] =
    \Cov[\overline{\Tstat}^{(\bnu)}_{\mathtt{F}}(\bx_1), \overline{\Tstat}^{(\bnu)}_{\mathtt{F}}(\bx_2)], \qquad \bx_1, \bx_2 \in \B,
\end{align*}
and
\begin{align*}
    \P\Big[ \big|\sup_{\bx \in \B} |\overline{\Tstat}^{(\bnu)}_{\mathtt{F}}(\bx)| - \sup_{\bx \in \B} |Z^{(\bnu)}_{\mathtt{F}}(\bx)|\big| > \frac{\log n}{(n |\bH|)^{1/6}} + \frac{(\log n)^{5/4}}{(n |\bH|)^{1/4}} + \frac{(\log n)^{3/2}}{(n^{\frac{v}{2+v}} |\bH|)^{1/2}}\Big] \lesssim \frac{1}{\log n} + \frac{\log n}{n}.
\end{align*}
\end{thm}

To construct a feasible confidence band for $(\zeta^{(\bnu)}(\bx): \bx \in \B)$, let $(\widehat{Z}^{(\bnu)}_{\mathtt{F}}(\bx): \bx \in \B)$ be a mean-zero Gaussian process conditional on $\bV$. Its feasible conditional covariance function is defined as
\begin{equation*}
    \Cov \left[ \widehat{Z}^{(\bnu)}_{\mathtt{F}}(\bx_1), \widehat{Z}^{(\bnu)}_{\mathtt{F}}(\bx_2) \mid \bV \right]
    = \frac{\widehat{\Omega}_{\mathtt{F}, \bx_1, \bx_2}^{(\bnu)}}{\sqrt{\widehat{\Omega}_{\mathtt{F}, \bx_1, \bx_1}^{(\bnu)} \widehat{\Omega}_{\mathtt{F}, \bx_2, \bx_2}^{(\bnu)}}},
    \quad \bx_1, \bx_2 \in \B.
\end{equation*}

\begin{thm}[Confidence Bands: Fuzzy BATEC; suprema approximation]\label{sa-thm: Confidence Bands: Fuzzy BATEC -- suprema}
    Suppose Assumptions \ref{sa-assump: Sharp DGP}, \ref{sa-assump: Boundary and Kernel}, and \ref{sa-assump: Fuzzy DGP} hold, $\inf_{\bx \in \B} |\tau^{(\bnu)}_W(\bx)| > 0$, and $\inf_{\bx\in\B}V_{\mathtt{F},\bx}^{(\bnu)}>0$. If $\liminf_{n \to \infty} \log \bar{h}/\log n > -\infty$, $(\log n)^9/(n^{v/(2+v)} |\bH|) = o(1)$, $\bar{h}^{p+1} \sqrt{n |\bH|\log n} = o(1)$, and
    \begin{align*}
        &\bh^{-\bnu}\sqrt{\log(1/\bar{h})}
        \left(\sqrt{\log(1/\bar{h})/(n|\bH|)}+
        \log(1/\bar{h})/(n^{v/(2+v)}|\bH|)+\bar{h}^{p+1}\right)=o(1),\\
        &\log^2(n)\bh^{-\bnu}
        \left(
        \sqrt{\frac{\log n}{n|\bH|}}
        +\frac{\log n}{n^{v/(2+v)}|\bH|}
        +\bar h^{p+1}
        \right)=o(1),
    \end{align*}
    then
    \begin{equation*}
        \sup_{u \in \reals} \left| \P \left( \sup_{\bx \in \B} \big| \widehat{\Tstat}_{\mathtt{F}}^{(\bnu)}(\bx) \big| \leq u \right) - \P \left( \sup_{\bx \in \B} \big| \widehat{Z}_{\mathtt{F}}^{(\bnu)}(\bx) \big| \leq u \;\middle|\; \bV \right) \right| = o_{\P}(1)
    \end{equation*}
    and
    \begin{equation*}
        \P \left( \zeta^{(\bnu)}(\bx) \in \widehat{\CI}_{\mathtt{F},\alpha}^{(\bnu)}(\bx) \text{ for all } \bx \in \B \right) = 1 - \alpha + o(1),
    \end{equation*}
    provided that $\q_{\alpha} = \inf \left\{ c > 0 : \P \left( \sup_{\bx \in \B} \big| \widehat{Z}_{\mathtt{F}}^{(\bnu)}(\bx) \big| \geq c \;\middle|\; \bV \right) \leq \alpha \right\}$.
\end{thm}

%%%%%%%%%%%%%%%%%%%%%%%%%%%%%%%%%%%%%%%%%%%%%%%%%%%%%%%%%%%%%%%%%%%%%%
\subsection{Fuzzy WBATE}
%%%%%%%%%%%%%%%%%%%%%%%%%%%%%%%%%%%%%%%%%%%%%%%%%%%%%%%%%%%%%%%%%%%%%%

We now consider WBATE summaries of these fuzzy ratios of derivative discontinuities. For $|\bnu|\le p$, the parameter is
\begin{align*}
    \zeta_{\mathtt{WBATE}}^{(\bnu)} = \int_{\B} \zeta^{(\bnu)}(\bb) w(\bb) \dif\Haus^{d-1} (\bb),
\end{align*}
where the weight function $w(\cdot)$ satisfies Assumption \ref{sa-assump: Weight Function and Boundary}. Without loss of generality, we assume that $\int_{\B} w(\bb) \dif\Haus^{d-1} (\bb) = 1$. The corresponding plug-in estimator is
\begin{align*}
    \widehat{\zeta}_{\mathtt{WBATE}}^{(\bnu)} = \int_\B \widehat{\zeta}^{(\bnu)}(\bb) w(\bb) \dif\Haus^{d-1}(\bb).
\end{align*}
The original fuzzy WBATE parameter is the special case $\bnu=\mathbf{0}$.

Our AMSE results for the fuzzy WBATE estimator leverage the leading linear component from the second-order linearization established in Section \ref{sec: Second-Order Linearization: Fuzzy BATEC},
\begin{align*}
    \check{\zeta}_{\mathtt{WBATE}}^{(\bnu)}
    = \int_{\B}\mathfrak{w}(\bb)^\top\mathfrak{T}(\bb) w(\bb) \dif\Haus^{d-1}(\bb).
\end{align*}
The corresponding approximate MSE is
\begin{align*}
    \operatorname{AMSE}_{\mathtt{F,WBATE}}^{(\bnu)}
    = \E\big[\big(\check{\zeta}_{\mathtt{WBATE}}^{(\bnu)}\big)^2\big|\bX\big].
\end{align*}
The leading bias is defined as the integral of the pointwise leading bias of the fuzzy BATEC estimator:
\begin{align*}
    B_{\mathtt{F,WBATE}}^{(\bnu)}(\bh) = \int_{\B} B_{\mathtt{F},\bb}^{(\bnu)}(\bh) \; w(\bb) \dif\Haus^{d-1}(\bb),
\end{align*}
where $B_{\mathtt{F},\bb}^{(\bnu)}(\bh)$ is given in Equation~\eqref{sa-eq: leading bias fuzzy}.

\begin{lem}[Bias: Fuzzy WBATE Linearization]\label{sa-lem: Fuzzy Integral: Bias}
    Suppose Assumptions \ref{sa-assump: Sharp DGP}, \ref{sa-assump: Boundary and Kernel}, \ref{sa-assump: Weight Function and Boundary} and \ref{sa-assump: Fuzzy DGP} hold, and $\inf_{\bx\in\B}|\tau_W^{(\bnu)}(\bx)|>0$. If $\frac{\log(1/\bar{h})}{n^{\frac{1+v}{2+v}}|\bH|} = o(1)$ and $\bar{h} = o(1)$, then
    \begin{align*}
        \E[\check{\zeta}_{\mathtt{WBATE}}^{(\bnu)}|\bX] = B_{\mathtt{F,WBATE}}^{(\bnu)}(\bh) + o_{\P}(\bar{h}^{p+1}\bh^{-\bnu}).
    \end{align*}
    If, in addition, $\bh=\bc h$ for a fixed vector $\bc$ and common bandwidth $h$, then $B_{\mathtt{F,WBATE}}^{(\bnu)}(\bh)=h^{p+1-|\bnu|}\widetilde B_{\mathtt{F,WBATE}}^{(\bnu)}(\bc)$.
\end{lem}

The conditional variance of this linearized fuzzy WBATE component involves the integration of the cross-covariance structure of the reduced-form processes. Let
\begin{align*}
    \Omega_{\mathtt{F,WBATE}}^{(\bnu)} = \int_{\B} \int_{\B} \Omega_{\mathtt{F},\bb_1, \bb_2}^{(\bnu)} w(\bb_1) w(\bb_2) \dif\Haus^{d-1}(\bb_1) \dif\Haus^{d-1}(\bb_2),
\end{align*}
and impose the aggregate variance nondegeneracy condition
\begin{align}
    \Omega_{\mathtt{F,WBATE}}^{(\bnu)}
    \gtrsim \frac{\bar{h}^{d-1}}{n|\bH|\bh^{2\bnu}}.
    \label{sa-eq: fuzzy WBATE aggregate nondegeneracy}
\end{align}
This condition rules out cancellation after weighting and integrating the linearized fuzzy Wald process.
Let its plug-in estimator be
\begin{align*}
    \widehat{\Omega}_{\mathtt{F,WBATE}}^{(\bnu)} = \int_{\B} \int_{\B} \widehat{\Omega}_{\mathtt{F},\bb_1, \bb_2}^{(\bnu)} w(\bb_1) w(\bb_2) \dif\Haus^{d-1}(\bb_1) \dif\Haus^{d-1}(\bb_2).
\end{align*}

\begin{lem}[Variance: Fuzzy WBATE Linearization]\label{sa-lem: Fuzzy Integral: Variance}
    Suppose the assumptions and conditions of Lemma \ref{sa-lem: Covariance: Fuzzy BATEC} and Assumption \ref{sa-assump: Weight Function and Boundary} hold, and \eqref{sa-eq: fuzzy WBATE aggregate nondegeneracy} and \eqref{sa-eq: local boundary measure regularity} hold. Then
    \begin{align*}
        \Var[\check{\zeta}_{\mathtt{WBATE}}^{(\bnu)}|\bX]
        = \Omega_{\mathtt{F,WBATE}}^{(\bnu)} + o_{\P}\Big(\frac{\bar{h}^{d-1}}{n|\bH|\bh^{2\bnu}}\Big),
    \end{align*}
    where
    \begin{align*}
        \frac{\bar{h}^{d-1}}{n|\bH|\bh^{2\bnu}}
         \lesssim \Omega_{\mathtt{F,WBATE}}^{(\bnu)}
         \lesssim \frac{\bar{h}^{d-1}}{n|\bH|\bh^{2\bnu}}.
    \end{align*}

    In addition,
    \begin{align*}
        \Var[\check{\zeta}_{\mathtt{WBATE}}^{(\bnu)}|\bX]
        & = \widehat{\Omega}_{\mathtt{F,WBATE}}^{(\bnu)} + o_{\P}\Big(\frac{\bar{h}^{d-1}}{n|\bH|\bh^{2\bnu}}\Big).
    \end{align*}
\end{lem}

\begin{thm}[AMSE Expansion: Fuzzy WBATE]\label{sa-thm: MSE Expansion: Fuzzy WBATE}
    Suppose the assumptions of Lemma \ref{sa-lem: Fuzzy Integral: Variance} hold, including \eqref{sa-eq: local boundary measure regularity}. If $\log(1/\bar{h})/(n^{v/(2+v)}|\bH|) = o(1)$ and $\bar{h} = o(1)$, then
    \begin{align*}
        \operatorname{AMSE}_{\mathtt{F,WBATE}}^{(\bnu)}
        = \Omega_{\mathtt{F,WBATE}}^{(\bnu)} + \big(B_{\mathtt{F,WBATE}}^{(\bnu)}(\bh)\big)^2
        + o_{\P}\Big(\frac{\bar{h}^{d-1}}{n|\bH|\bh^{2\bnu}} + \bar{h}^{2p+2}\bh^{-2\bnu}\Big).
    \end{align*}
\end{thm}

For inference, we define the feasible $t$-statistic for the aggregate effect:
\begin{align*}
    \widehat{\Tstat}_{\mathtt{F,WBATE}}^{(\bnu)} = \frac{\widehat{\zeta}_{\mathtt{WBATE}}^{(\bnu)} - \zeta_{\mathtt{WBATE}}^{(\bnu)}}{\sqrt{\widehat{\Omega}_{\mathtt{F,WBATE}}^{(\bnu)}}}.
\end{align*}

\begin{thm}[Asymptotic Normality: Fuzzy WBATE]\label{sa-thm: Asymptotic Normality: Fuzzy WBATE}
    Suppose Assumptions \ref{sa-assump: Sharp DGP}, \ref{sa-assump: Boundary and Kernel}, \ref{sa-assump: Weight Function and Boundary} and \ref{sa-assump: Fuzzy DGP} hold, $\inf_{\bx\in\B}|\tau_W^{(\bnu)}(\bx)|>0$, $\inf_{\bx\in\B}V_{\mathtt{F},\bx}^{(\bnu)}>0$, and \eqref{sa-eq: fuzzy WBATE aggregate nondegeneracy} and \eqref{sa-eq: local boundary measure regularity} hold. If $\log(1/\bar{h})/(n^{v/(2+v)}|\bH|) = o(1)$, $n|\bH|\bar{h}^{2p+2}/\bar{h}^{d-1} = o(1)$, and
    \begin{align*}
        \frac{\bh^{-2\bnu}\log^2(1/\bar{h})}{n|\bH|\bar{h}^{d-1}}=o(1),
        \qquad
        \frac{\bh^{-2\bnu}\log^4(1/\bar{h})}{n^{\frac{4v}{2+v}-1}|\bH|^3\bar{h}^{d-1}}=o(1),
    \end{align*}
    then
    \begin{align*}
        \sup_{u \in \reals} \big|\P (\widehat{\Tstat}_{\mathtt{F,WBATE}}^{(\bnu)} \leq u )  - \Phi(u) \big| = o(1).
    \end{align*}
\end{thm}

The local boundary-measure condition in \eqref{sa-eq: local boundary measure regularity} is used here for the same reason as in the sharp WBATE results: it controls the effective size of the boundary pairs that can have nonzero local-kernel covariance. The more specialized piecewise-linear boundary conditions in the main paper imply this condition.


\subsection{Fuzzy LBATE}

For $|\bnu|\leq p$, we consider the LBATE summary of the fuzzy ratio of derivative discontinuities, with estimand and estimator
\begin{align*}
    \zeta_{\mathtt{LBATE}}^{(\bnu)} = \sup_{\bb \in \B} \zeta^{(\bnu)}(\bb)
    \qquad\text{and}\qquad
    \widehat{\zeta}_{\mathtt{LBATE}}^{(\bnu)} = \sup_{\bb \in \B} \widehat{\zeta}^{(\bnu)}(\bb).
\end{align*}
The original fuzzy LBATE parameter is the special case $\bnu=\mathbf{0}$.

\begin{thm}[Convergence Rate: Fuzzy LBATE]\label{sa-thm: Convergence Rate: Fuzzy LBATE}
    Suppose the assumptions and conditions of Theorem~\ref{sa-thm: Convergence Rates: Fuzzy BATEC} hold. Then
    \begin{align*}
        \big|\widehat{\zeta}_{\mathtt{LBATE}}^{(\bnu)} - \zeta_{\mathtt{LBATE}}^{(\bnu)} \big|
        \lessp \bh^{-\bnu}\Big(\sqrt{\frac{\log(1/\bar{h})}{ n |\bH|}} + \frac{\log(1/\bar{h})}{n^{\frac{1+v}{2+v}}|\bH|} + \bar{h}^{p+1}\Big).
    \end{align*}
\end{thm}

Recall from Theorem~\ref{sa-thm: Confidence Bands: Fuzzy BATEC -- suprema} that $(\widehat{Z}_{\mathtt{F}}^{(\bnu)}(\bx):\bx \in \B)$ is a (conditionally on $\bV$) mean-zero Gaussian process with feasible (conditional) covariance function
\begin{align*}
    \Cov \Big[\widehat{Z}_{\mathtt{F}}^{(\bnu)}(\bx_1),\widehat{Z}_{\mathtt{F}}^{(\bnu)}(\bx_2) \Big|\bV \Big]
    = \frac{\widehat{\Omega}_{\mathtt{F},\bx_1,\bx_2}^{(\bnu)}}
           {\sqrt{\widehat{\Omega}_{\mathtt{F},\bx_1,\bx_1}^{(\bnu)} \widehat{\Omega}_{\mathtt{F},\bx_2,\bx_2}^{(\bnu)}}},
    \qquad \bx_1, \bx_2 \in \B.
\end{align*}
Consider the confidence interval
\begin{align*}
    \widehat{\CI}_{\alpha,\mathtt{F,LBATE}}^{(\bnu)}
    =
    \bigg[\sup_{\bb \in \B} \Big(\widehat{\zeta}^{(\bnu)}(\bb)  - \q_{\alpha} \sqrt{\widehat{\Omega}_{\mathtt{F},\bb,\bb}^{(\bnu)}}\Big), \;
    \sup_{\bb \in \B} \Big(\widehat{\zeta}^{(\bnu)}(\bb) + \q_{\alpha} \sqrt{\widehat{\Omega}_{\mathtt{F},\bb,\bb}^{(\bnu)}}\Big)\bigg],
\end{align*}
where $\q_{\alpha} = \inf \big\{c > 0: \P \big(\sup_{\bx \in \B} \big|\widehat{Z}_{\mathtt{F}}^{(\bnu)}(\bx)\big|\geq c \big| \bV \big) \leq \alpha \big\}$.

\begin{thm}[Confidence Interval: Fuzzy LBATE]\label{sa-thm: Confidence Intervals: Fuzzy LBATE}
    Suppose the assumptions and conditions in Theorem~\ref{sa-thm: Confidence Bands: Fuzzy BATEC -- suprema} hold. Then
    \begin{align*}
            \P\Big[\zeta_{\mathtt{LBATE}}^{(\bnu)} \in \widehat{\CI}_{\alpha,\mathtt{F,LBATE}}^{(\bnu)} \Big] \geq 1 - \alpha + o(1).
    \end{align*}
\end{thm}



%%%%%%%%%%%%%%%%%%%%%%%%%%%%%%%%%%%%%%%%%%%%%%%%%%%%%%%%%%%%%%%%%%%%%%
\section{Gaussian Strong Approximation}\label{sa-sec: Gaussian Strong Approximation}
%%%%%%%%%%%%%%%%%%%%%%%%%%%%%%%%%%%%%%%%%%%%%%%%%%%%%%%%%%%%%%%%%%%%%%

We present a Gaussian strong approximation theorem, which is the key technical tool behind Theorem~\ref{sa-thm: Gaussian Strong Approximation: Sharp Tstat}. The theorem builds on and generalizes the results in \cite{Cattaneo-Yu_2025_AOS}. Consider the \emph{residual-based empirical process} given by
\begin{align*}
    R_n[g,r] = \frac{1}{\sqrt{n}} \sum_{i = 1}^n \Big[g(\bx_i) r(y_i) - \E[g(\bx_i) r(y_i)|\bx_i] \Big], \qquad g \in \G, r \in \R,
\end{align*}
where $\G$ and $\R$ are classes of functions satisfying certain regularity conditions.

\subsection{Definitions for Function Spaces}

Let $\F$ be a class of measurable functions from a probability space $(\reals^q, \B(\reals^q), \P)$ to $\reals$. We introduce several definitions that capture properties of $\F$.

\begin{enumerate}[label=(\roman*)]
    \item  $\F$ is pointwise measurable if it contains a countable subset $\G$ such that for any $f \in \F$, there exists a sequence $(g_m:m\geq1) \subseteq \G$ such that $\lim_{m \to \infty} g_m(\bu) = f(\bu)$ for all $\bu \in \reals^q$.
    \item  Let $\supp(\F) = \cup_{f \in \F}\supp(f)$. A probability measure $\bbQ_\F$ on $(\reals^q,\B(\reals^q))$ is a surrogate measure for $\P$ with respect to $\F$ if
    \begin{enumerate}[label=(\roman*)]
        \item $\bbQ_\F$ agrees with $\P$ on $\supp(\P) \cap \supp(\F)$.
        \item $\bbQ_\F(\supp(\F) \setminus \supp(\P)) = 0$.
    \end{enumerate}
    Let $\mathcal{Q}_\F=\supp(\bbQ_\F)$.
    \item For $q=1$ and an interval $\mathcal{I}\subseteq\reals$, the pointwise total variation of $\F$ over $\mathcal{I}$ is
    \begin{align*}
        \mathtt{pTV}_{\F,\mathcal{I}} = \sup_{f \in \F} \sup_{P\geq1}\sup_{\calP_P \in \mathcal{I}} \sum_{i = 1}^{P-1}|f(a_{i+1}) - f(a_i)|,
    \end{align*}
    where $\calP_P=\{(a_1,\dots,a_P):a_1 \leq \cdots \leq a_P\}$ denotes the collection of all partitions of $\mathcal{I}$.
    \item For a non-empty $\mathcal{C} \subseteq \reals^q$, the total variation of $\F$ over $\mathcal{C}$ is
    \begin{align*}
        \ttTV_{\F, \mathcal{C}} = \inf_{\mathcal{U} \in \mathcal{O}(\mathcal{C})}\sup_{f \in \F} \sup_{\phi \in \mathscr{D}_{q}(\mathcal{U})} \int_{\reals^q} f(\bu)\operatorname{div}(\phi)(\bu) \dif \bu / \infnorm{\norm{\phi}_2},
    \end{align*}
    where $\mathcal{O}(\mathcal{C})$ denotes the collection of all open sets that contains $\mathcal{C}$, and $\mathscr{D}_{q}(\mathcal{U})$ denotes the space of infinitely differentiable functions from $\reals^q$ to $\reals^q$ with compact support contained in $\mathcal{U}$.
    \item For a non-empty $\mathcal{C} \subseteq \reals^q$, the local total variation constant of $\F$ over $\mathcal{C}$, is a positive number $\ttK_{\F,\mathcal{C}}$ such that for any cube $\mathcal{D} \subseteq \reals^q$ with edges of length $\ell$ parallel to the coordinate axes,
    \begin{align*}
        \ttTV_{\F, \mathcal{D} \cap \mathcal{C}}  \leq \ttK_{\F, \mathcal{C}} \ell^{q-1}.
    \end{align*}
    \item For a non-empty $\mathcal{C} \subseteq \reals^q$, the envelopes of $\F$ over $\mathcal{C}$ are
    \begin{align*}
        \ttM_{\F,\C} = \sup_{\bu \in \C }M_{\F,\C}(\bu),
        \qquad M_{\F,\C}(\bu) = \sup_{f \in \F}|f(\bu)|,
        \qquad \bu \in \C.
    \end{align*}
    \item  For a non-empty $\mathcal{C} \subseteq \reals^q$, the Lipschitz constant of $\F$ over $\C$ is
    \begin{align*}
        \ttL_{\F,\C} = \sup_{f \in \F}\sup_{\bu_1, \bu_2 \in \C} \frac{|f(\bu_1) - f(\bu_2)|}{\|\bu_1 - \bu_2\|_\infty}.
    \end{align*}
    \item For a non-empty $\mathcal{C} \subseteq \reals^q$, the $L_1$ bound of $\F$ over $\C$ is
    \begin{align*}
        \mathtt{E}_{\F,\C} = \sup_{f \in \F} \int_{\C} |f| \dif\P.
    \end{align*}
    \item   For a non-empty $\mathcal{C} \subseteq \reals^q$, the uniform covering number of $\F$ with envelope $M_{\F,\C}$ over $\C$ is
    \begin{align*}
        \ttN_{\F,\C}(\delta,M_{\F,\C}) = \sup_{\mu} N(\F,\norm{\cdot}_{\mu,2},\delta \norm{M_{\F,\C}}_{\mu,2}),
        \qquad \delta \in (0, \infty),
    \end{align*}
    where the supremum is taken over all finite discrete measures on $(\C, \B(\C))$. We assume that $M_{\F,\C}(\bu)$ is finite for every $\bu \in \C$.
    \item  For a non-empty $\mathcal{C} \subseteq \reals^q$, the uniform entropy integral of $\F$ with envelope $M_{\F,\C}$ over $\C$ is
    \begin{align*}
        J_\C(\delta, \F, M_{\F,\C}) = \int_0^{\delta} \sqrt{1 + \log \ttN_{\F,\C}(\varepsilon,M_{\F,\C})} \dif \varepsilon,
    \end{align*}
    where it is assumed that $M_{\F,\C}(\bu)$ is finite for every $\bu \in \C$.
    \item For a non-empty $\mathcal{C} \subseteq \reals^q$, $\F$ is a VC-type class with envelope $M_{\F,\C}$ over $\C$ if (i) $M_{\F,\C}$ is measurable and $M_{\F,\C}(\bu)$ is finite for every $\bu \in \C$, and (ii) there exist $\ttc_{\F,\C}>0$ and $\mathtt{d}_{\F,\C}>0$ such that
    \begin{align*}
        \ttN_{\F,\C}(\varepsilon,M_{\F,\C}) \leq \ttc_{\F,\C} \varepsilon^{-\mathtt{d}_{\F,\C}},
        \qquad \varepsilon\in(0,1).
    \end{align*}

\end{enumerate}

If a surrogate measure $\bbQ_\F$ for $\P$ with respect to $\F$ has been assumed, and it is clear from the context, we drop the dependence on $\C = \mathcal{Q}_{\F}$ for all quantities in the previous definitions. That is, to save notation, we set $\ttTV_{\F}=\ttTV_{\F,\mathcal{Q}_{\F}}$, $\ttK_{\F}=\ttK_{\F,\mathcal{Q}_{\F}}$, $\ttM_{\F}=\ttM_{\F,\mathcal{Q}_{\F}}$, $M_{\F}(\bu)=M_{\F,\mathcal{Q}_{\F}}(\bu)$, $\ttL_{\F}=\ttL_{\F,\mathcal{Q}_{\F}}$, and so on, whenever there is no confusion.

\subsection{Residual-based Empirical Process}

The following theorem generalizes \citet[Theorem 2]{Cattaneo-Yu_2025_AOS} by requiring only bounded polynomial moments for $y_i$ conditional on $\bx_i$.

\begin{thm}[Strong Approximation for Residual-based Empirical Processes]\label{sa-thm: Residual Empirical Process Strong Approximation}
    Suppose $(\bz_i=(\bx_i, y_i): 1 \leq i \leq n)$ are i.i.d.\ random vectors taking values in $(\reals^{d+1}, \B(\reals^{d+1}))$ with common law $\P_Z$, where $\bx_i$ has distribution $\P_X$ supported on $\X\subseteq\reals^d$, $y_i$ has distribution $\P_Y$ supported on $\Y\subseteq\reals$, $\sup_{\bx \in \X}\E[|y_i|^{2 + v}|\bx_i = \bx] \leq C_y<\infty$ for some $v > 0$, and the following conditions hold:

    \begin{enumerate}[label=\emph{(\roman*)}]
        \item $\G$ is a real-valued pointwise measurable class of functions on $(\reals^d, \B(\reals^d), \P_X)$.
        \item There exists a surrogate measure $\bbQ_\G$ for $\P_X$ with respect to $\G$ such that $\bbQ_\G = \mathfrak{m} \circ \phi_\G$, where the \textit{normalizing transformation} $\phi_{\G}: \Q_\G \to [0,1]^d$ is a diffeomorphism.
        \item $\G$ is a VC-type class with envelope $\ttM_{\G}$ over $\Q_\G$ with $\ttc_{\G} \geq e$ and $\mathtt{d}_{\G} \geq 1$.
        \item $\R$ is a real-valued pointwise measurable class of functions on $(\reals, \mathsf{Borel}(\reals),\P_Y)$.
        \item $\R$ is a VC-type class with envelope $M_{\R,\Y}$ over $\Y$ with $\ttc_{\R,\Y}\geq e$ and $\mathtt{d}_{\R,\Y}\geq 1$, where $M_{\R,\Y}(y) + \mathtt{pTV}_{\R,(-|y|,|y|)} \leq \mathtt{v} (1 + |y|)$ for all $y \in \Y$, for some $\mathtt{v}>0$.
        \item There exists a constant $\mathtt{k}$ such that $|\log_2 \mathtt{E}_{\G}| + |\log_2 \ttTV| + |\log_2 \ttM_{\G}| \leq \mathtt{k} \log_2 n$, where the constant $\ttTV = \max \{\ttTV_{\G}, \ttTV_{\G \times \mathscr{U}_{\R},\Q_\G}\}$ with $\mathscr{U}_{\R} = \{\theta(\cdot,r, \tau) : r \in \R, \tau \in (0,\infty]\}$, and $\theta(\bx,r,\tau) = \E[r(y_i) \Indicator(|y_i| \leq \tau)|\bx_i = \bx]$ for $\bx \in \X$.
    \end{enumerate}
    Define the residual-based empirical process
    \begin{align*}
        R_n(g,r) = \frac{1}{\sqrt{n}} \sum_{i = 1}^n g(\bx_i)(r(y_i) - \E[r(y_i)|\bx_i]), \qquad g \in \G, r \in \R.
    \end{align*}
    Then (1) on a possibly enlarged probability space, there exists a sequence of mean-zero Gaussian processes $(Z_n^R(g,r): g\in\G, r \in \R)$ with almost sure continuous trajectories such that:
    \begin{itemize}
        \item $\E[R_n(g_1, r_1) R_n(g_2, r_2)] = \E[Z^R_n(g_1, r_1) Z^R_n(g_2, r_2)]$ for all $(g_1, r_1), (g_2, r_2) \in \G \times \R$, and
        \item $\E\big[\norm{R_n - Z_n^R}_{\G \times \R}\big] \leq C \mathtt{v} \ttd \log(\ttc n) \rho_n$,
    \end{itemize}
    with
    \begin{align*}
        \rho_n = \sqrt{\ttd \log (\ttc n)} \; \mathtt{r}_n^{\frac{v}{v +2}}(\sqrt{\ttM_{\G}\ttE_{\G}})^{\frac{2}{v+2}} + \ttM_{\G} n^{-\frac{v/2}{2+v}} +  \frac{\ttM_{\G}}{\sqrt{n}}  \Big(\frac{\sqrt{\ttM_{\G} \ttE_{\G}}}{\mathtt{r}_n}\Big)^{\frac{2}{v+2}},
    \end{align*}
    where $C$ is a positive constant that depends only on $v$ and $C_y$, $\ttc = \ttc_{\G} \ttc_{\R,\Y} \ttk$, $\ttd = \ttd_{\G} + \ttd_{\R,\Y} + \ttk$, and
    \begin{gather*}
        \mathtt{r}_n = \min\Big\{\frac{(\mathtt{c}_1^d \ttM_{\G}^{d+1} \ttTV^d \mathtt{E}_{\G})^{1/(2d+2)} }{n^{1/(2d+2)}}, \frac{(\mathtt{c}_1^{d/2} \mathtt{c}_2^{d/2} \ttM_{\G} \ttTV^{d/2} \mathtt{E}_{\G} \ttL^{d/2})^{1/(d+2)}}{n^{1/(d+2)}} \Big\}, \\
        \ttc_1 = d \sup_{\bx \in \mathcal{Q}_{\G}} \prod_{j = 1}^{d-1} \sigma_j(\nabla \phi_{\G}(\bx)), \qquad
        \ttc_2 = \sup_{\bx \in \mathcal{Q}_{\G}} \frac{1}{\sigma_{d}(\nabla \phi_{\G}(\bx))}, \qquad
        \ttL = \max\{\ttL_{\G}, \ttL_{\G \times \mathscr{U}_{\R},\Q_{\G}}\};
    \end{gather*}
    and (2) if $\R$ is a singleton and the truncated conditional-mean class $\mathscr{U}_{\R}$ satisfies
    \[
        \ttTV_{\G\times\mathscr{U}_{\R},\Q_{\G}}\lesssim \ttTV_{\G\times\V_{\R},\Q_{\G}},
        \qquad
        \ttL_{\G\times\mathscr{U}_{\R},\Q_{\G}}\lesssim \ttL_{\G\times\V_{\R},\Q_{\G}},
    \]
    then we can replace $\ttTV$ and $\ttL$ in the previous conditions and statements by $\ttTV_{\text{sing}} = \max\{\ttTV_{\G}, \ttTV_{\G \times \V_{\R}, \Q_{\G}}\}$, and $\ttL_{\text{sing}} = \max\{\ttL_{\G}, \ttL_{\G \times \V_{\R},\Q_{\G}}\}$, respectively, with $\V_{\R} = \{\theta(\cdot,r) : r \in \R\}$, and $\theta(\bx,r) = \E[r(y_i) |\bx_i = \bx]$ for $\bx \in \X$.
\end{thm}

\begin{remark}
The class $\mathscr{U}_{\mathcal R}$ comprises truncated conditional means at all truncation levels. Its Lipschitz and total-variation constants can be bounded, for example, if $f(y \mid \bx)$ is Lipschitz in $\bx$ uniformly over $(\bx,y)$ in the support of $(\bx_i,y_i)$. When $\mathcal R$ is a singleton, regularity of the untruncated conditional mean class $\V_{\mathcal R}$ is enough whenever the truncated conditional means inherit the same bounds, as stated in part (2) of the theorem; this is automatic under bounded support for the residual component and holds under standard smooth conditional-density conditions.
\end{remark}


%%%%%%%%%%%%%%%%%%%%%%%%%%%%%%%%%%%%
\subsection{Proof of Theorem~\ref{sa-thm: Residual Empirical Process Strong Approximation}}\label{sa-sec: proof of residual empirical process strong approximation}
%%%%%%%%%%%%%%%%%%%%%%%%%%%%%%%%%%%%

We use a truncation argument. Let $\kappa_n \geq 1$ be the level of truncation. For each $r \in \R$, define
\begin{align*}
    \tilde{r}(y) = r(y) \Indicator(|y| \leq \kappa_n), \qquad y \in \reals,
\end{align*}
and define the class $\tilde{\R} = \{\tilde{r}: r \in \R\}$. Suppose $Z_n^R$ is a mean-zero Gaussian process indexed by $\G \times \R \cup \G \times \tilde{\R}$, whose construction is given below. Then
\begin{align*}
    R_n(g,r) - Z_n^R(g,r)
    & = \big[R_n(g,\tilde{r}) - Z_n^R(g,\tilde{r})\big] + \big[R_n(g,r) - R_n(g,\tilde{r})\big] + \big[Z_n^R(g,\tilde{r}) - Z_n^R(g,r)\big].
\end{align*}

\subsubsection*{Part 1: Strong approximation for truncated residual empirical process.}

The truncation class $\tilde{\R}$ is pointwise measurable whenever $\R$ is, and it is VC-type with envelope $M_{\tilde{\R},\Y} = M_{\R,\Y} \Indicator(|\cdot| \leq \kappa_n)$ over $\Y$ and constants $\ttc_{\R,\Y}$ and $\ttd_{\R,\Y}$ up to universal factors. Moreover, the envelope and pointwise total variation satisfy
\begin{align*}
    M_{\tilde{\R},\Y}(y) + \mathtt{pTV}_{\tilde{\R},(-|y|,|y|)} \lesssim \ttv \kappa_n, \qquad y \in \Y,
\end{align*}
because truncating $r$ only adds jumps at $\pm\kappa_n$. Therefore \citet[Theorem 2]{Cattaneo-Yu_2025_AOS}, applied with $\alpha=0$ and with the bounded-outcome constant proportional to $\ttv\kappa_n$, implies on a possibly enlarged probability space that there exists a sequence of mean-zero Gaussian processes $(Z_n^R(g,r): (g,r)\in \G\times \tilde{\R})$ with almost sure continuous trajectories on $(\G \times \tilde{\R}, \rho_{\P})$ such that $\E[R_n(g_1, r_1) R_n(g_2, r_2)] = \E[Z^R_n(g_1, r_1) Z^R_n(g_2, r_2)]$ for all $(g_1, r_1), (g_2, r_2) \in \G \times \tilde{\R}$, and, after integrating the tail bound in that theorem,
\begin{align*}
    & \E[\lVert R_n(g,\tilde{r}) - Z_n^R(g,\tilde{r}) \rVert_{\G \times \R}] \\
    & \leq C_1 \ttv \kappa_n \bigg(\sqrt{d} \min\Big\{\frac{(\ttc_1^d \ttM_{\G}^{d+1} \ttTV^d \mathtt{E}_{\G})^{\frac{1}{2d+2}} }{n^{1/(2d+2)}},
                     \frac{(\ttc_1^{\frac{d}{2}} \ttc_2^{\frac{d}{2}}\ttM_{\G} \ttTV^{\frac{d}{2}} \mathtt{E}_{\G} \ttL^{\frac{d}{2}})^{\frac{1}{d+2}}}{n^{1/(d+2)}} \Big\} (\ttd\log (\ttc n))^{3/2}
        + \frac{\ttd\log (\ttc n)}{\sqrt{n}} \ttM_{\G} \bigg) \\
    & = C_1 \ttv \kappa_n \bigg(\sqrt{d} \mathtt{r}_n (\ttd\log (\ttc n))^{\frac{3}{2}}
        + \frac{\ttd\log (\ttc n)}{\sqrt{n}} \ttM_{\G}\bigg),
\end{align*}
where $C_1$ is a positive universal constant. The total-variation constant $\ttTV$ used in the theorem controls the truncated class because
\begin{align*}
    \V_{\tilde{\R}}
    = \{\theta(\cdot,r,\kappa_n): r\in\R\}
    \subseteq \mathscr{U}_{\R}.
\end{align*}
Thus $\ttTV = \max \{\ttTV_{\G}, \ttTV_{\G \times \mathscr{U}_{\R},\Q_\G}\}$ is an upper bound for $\max \{\ttTV_{\G}, \ttTV_{\G \times \V_{\tilde{\R}},\Q_\G}\}$. The same argument applies to $\ttL$.

In the special case that $\R = \{r_{\ast}\}$ is a singleton, take $\tilde{y}_i = r_{\ast}(y_i) \Indicator(|y_i| \leq \kappa_n)/(\ttv \kappa_n)$. Then $\E[\exp(|\tilde{y}_i|)] \leq 2$, and $\tilde{y}_i$ is supported on $\tilde{\Y} = [-1,1]$. Moreover,
\begin{align*}
    \frac{1}{\ttv \kappa_n}R_n(g, \tilde{r}_{\ast}) = \frac{1}{\sqrt{n}} \sum_{i = 1}^n g(\bx_i) (\tilde{y}_i - \E[\tilde{y}_i|\bx_i]), \qquad g \in \G.
\end{align*}
The right-hand side is a residual empirical process based on the sample $(\bx_i, \tilde{y}_i)$, $1 \leq i \leq n$, indexed by $\G \times \{\operatorname{Id}\}$, where $\operatorname{Id}: \reals \to \reals$ is the identity function. Under the singleton domination condition in part (2) of the theorem, applying \citet[Theorem 2]{Cattaneo-Yu_2025_AOS} with $\alpha = 0$ to this empirical process gives the same upper bound with $\ttTV$ and $\ttL$ replaced by $\ttTV_{\text{sing}}$ and $\ttL_{\text{sing}}$.

\subsubsection*{Part 2: Truncation error for the empirical process --- $\lVert R_n(g,r) - R_n(g, \widetilde{r})\rVert_{\G \times \R}$}

Consider the class of differences due to truncation, $\Delta \R = \{r - \tilde{r}: r \in \R\}$. The assumptions imply $\G \times \Delta \R$ is VC-type in the sense that for all $0 <\varepsilon < 1$,
\begin{align*}
    \sup_{\Q} N(\G \times \Delta \R, \norm{\cdot}_{\Q,2}, \varepsilon \lVert \ttM_{\G} (M_{\R,\Y} - M_{\tilde{\R},\Y}) \rVert_{\Q,2}) \leq \ttc_{\G} \ttc_{\R,\Y} (\varepsilon^2/4)^{-\ttd_{\G} - \ttd_{\R,\Y}},
\end{align*}
where the supremum is over all finite discrete measures on $\reals^{d+1}$, and $M_{\tilde{\R},\Y}(y) = M_{\R,\Y}(y) \Indicator(|y| \leq \kappa_n)$. The function $\ttM_{\G} (M_{\R,\Y} - M_{\tilde{\R},\Y})$ is an envelope for $\G \times \Delta \R$, since all functions in $\Delta\R$ vanish on $[-\kappa_n,\kappa_n]$. Denote $\bX = (\bx_i)_{1 \leq i \leq n}$. The linear envelope condition on $\R$ and the conditional $(2+v)$ moment bound on $y_i$ imply
\begin{align*}
    \E \Big[\max_{1 \leq i \leq n}\ttM_{\G}^2 (M_{\R,\Y}(y_i) - M_{\tilde{\R},\Y}(y_i))^2\Big|\bX\Big]^{\frac{1}{2}}
    & \lesssim \ttM_{\G} \E \Big[ \big( \max_{1 \leq i \leq n}M_{\R,\Y}(y_i)\big)^2\Big| \bX \Big]^{\frac{1}{2}} \lesssim  \ttM_{\G} n^{\frac{1}{2+v}}, \\
    \sup_{(g,r)  \in \G \times \R} \E[g(\bx_i)^2 r(y_i)^2 \Indicator(|y_i| \geq \kappa_n)]^{\frac{1}{2}}
    & \lesssim \sup_{(g,r)  \in \G \times \R} \E \left[g(\bx_i)^2 \E[|r(y_i)|^{2+v}|\bx_i]^{\frac{2}{2+v}} \P(|y_i| \geq \kappa_n| \bx_i)^{\frac{v}{2+v}}\right]^{\frac{1}{2}} \\
    & \lesssim \sqrt{\ttM_{\G} \ttE_{\G} \kappa_n^{-v}}.
\end{align*}
By Jensen's inequality, we also have
\begin{align*}
    & \E \Big[\max_{1 \leq i \leq n}\ttM_{\G}^2 (\E[M_{\R,\Y}(y_i) - M_{\tilde{\R},\Y}(y_i)|\bx_i])^2\Big|\bX\Big]^{\frac{1}{2}} \lesssim \ttM_{\G} n^{\frac{1}{2+v}}, \\
    & \sup_{(g,r)  \in \G \times \R} \E[g(\bx_i)^2 \E[r(y_i) - \tilde{r}(y_i)|\bx_i]^2]^{\frac{1}{2}} \lesssim \sqrt{\ttM_{\G} \ttE_{\G} \kappa_n^{-v}}, \\
    & \E[\ttM_{\G}^2(M_{\R,\Y}(y_i) - M_{\tilde{\R},\Y}(y_i))^2]^{1/2} \lesssim \ttM_{\G} \kappa_n^{-v/2}.
\end{align*}
Denote $A =  (\ttc_{\G} \ttc_{\R,\Y})^{\frac{1}{2 \ttd_{\G} + 2 \ttd_{\R,\Y}}}/4$ and $D = 2 \ttd_{\G} + 2 \ttd_{\R,\Y}$. \citet[Corollary 5.1]{Chernozhukov-Chetverikov-Kato_2014b_AoS} gives
\begin{align*}
    \E \left[\lVert R_n(g, r) - R_n(g,\widetilde{r})\rVert_{\G \times \R}\right] &
    \lesssim \E \left[\sup_{g \in \G} \sup_{h \in \Delta \R} \frac{1}{\sqrt{n}} \sum_{i = 1}^n g(\bx_i)(h(y_i) - \E[h(y_i)|\bx_i])\right]\\
    & \lesssim \sqrt{D \ttM_{\G} \ttE_{\G}\kappa_n^{-v} \log (A \sqrt{\ttM_{\G}/\ttE_{\G}})} + \frac{D \ttM_{\G} n^{\frac{1}{2+v}}}{\sqrt{n}} \log (A \sqrt{\ttM_{\G}/\ttE_{\G}})\\
    & \lesssim \sqrt{D \log (A \sqrt{\ttM_{\G}/\ttE_{\G}})} \sqrt{\ttM_{\G} \ttE_{\G}} \kappa_n^{-v/2} + \frac{D \log (A \sqrt{\ttM_{\G}/\ttE_{\G}}) \ttM_{\G}}{\sqrt{n^{\frac{v}{2+v}}}}.
\end{align*}

\subsubsection*{Part 3: Truncation error for the Gaussian process --- $\lVert Z_n^R(g,r) - Z_n^R(g,\tilde{r}) \rVert_{\G \times \R}$}

The class $\G \times \tilde{\R} \cup \G \times \R$ is VC-type with respect to envelope function $2 \ttM_{\G} \ttM_{\R,\Y}$ in the sense that for all $0 <\varepsilon < 1$,
\begin{align*}
    \sup_{Q} N(\G \times \R \cup \G \times \tilde{\R}, \norm{\cdot}_{Q,2}, 2 \varepsilon \lVert \ttM_{\G} \ttM_{\R,\Y} \rVert_{Q,2}) \leq \ttc_{\G} \ttc_{\R} (\varepsilon^2/4)^{-\ttd_{\G}-\ttd_{\R}},
\end{align*}
where the supremum is over all finite discrete measures on $\reals^{d+1}$. Here $\ttc_{\R}$ and $\ttd_{\R}$ are VC-type constants for $\R\cup\tilde{\R}$; truncation by the fixed indicator $\Indicator(|y|\leq\kappa_n)$ preserves the entropy exponent, so these constants can be chosen with $\ttc_{\R}\geq \ttc_{\R,\Y}$ and $\ttd_{\R}=\ttd_{\R,\Y}$ up to universal constants. Hence $\G \times \tilde{\R} \cup \G \times \R$ is pre-Gaussian, and on some probability space there exists a mean-zero Gaussian process $\bar{Z}_n^R$ indexed by $\F = \G \times \tilde{\R} \cup \G \times \R$ with the same covariance structure as $R_n$ and almost surely continuous paths with respect to the metric $\rho$, given by
\begin{align*}
    \rho((g_1,r_1),(g_2,r_2)) =  \E[(\bar{Z}_n^R(g_1, r_1) - \bar{Z}_n^R(g_2, r_2))^2]^{\frac{1}{2}} = \E[(R_n(g_1, r_1) - R_n(g_2, r_2))^2]^{\frac{1}{2}},
\end{align*}
for $(g_1, r_1), (g_2, r_2) \in \F$. Recall the definition of $\Delta\R$ in Part 2. The increment diameter between each original function and its truncated counterpart satisfies
\begin{align*}
    \sigma \equiv \sup_{g \in \G, r \in \R}\rho((g,r),(g,\tilde r))
    \leq \sup_{g \in \G, h \in \Delta\R} \E[g(\bx_i)^2 h(y_i)^2]^{1/2}
    \lesssim \sqrt{\ttM_{\G} \ttE_{\G} \kappa_n^{-v}}.
\end{align*}
The entropy bound above also implies, for all $0 < \varepsilon < 1$,
\begin{align*}
    N(\G \times \R \cup \G \times \tilde{\R}, \rho, 2\varepsilon\lVert \ttM_{\G} M_{\R,\Y}\rVert_{\P,2})
    \leq \ttc_{\G} \ttc_{\R,\Y}(\varepsilon^2/4)^{-\ttd_{\G}-\ttd_{\R,\Y}}.
\end{align*}
Denote $A =  (\ttc_{\G} \ttc_{\R,\Y})^{\frac{1}{2 \ttd_{\G} + 2 \ttd_{\R,\Y}}}/4$ and $D = 2 \ttd_{\G} + 2 \ttd_{\R,\Y}$. By \citet[Corollary 2.2.8]{van-der-Vaart-Wellner_1996_Book}, choosing any $(g_0, r_0) \in \G \times \R$,
\begin{align*}
    \E\bigg[\lVert \bar{Z}_n^R(g,r) - \bar{Z}_n^R(g,\tilde{r}) \rVert_{\G \times \R}\bigg] & \lesssim \E \left[\left|\bar{Z}_n^R(g_0, r_0) - \bar{Z}_n^R(g_0, \tilde{r}_0)\right|\right] + \int_{0}^{\sigma} \sqrt{\log\Big( \ttc_{\G} \ttc_{\R}\big(\frac{\ttM_{\G}}{\varepsilon}\big)^{\ttd_{\G} + \ttd_{\R}}\Big)} \dif \varepsilon \\
    & \leq \sqrt{D \log (A \sqrt{\ttM_{\G}/\ttE_{\G}})} \sqrt{\ttM_{\G} \ttE_{\G}} \kappa_n^{-v/2} \\
    & \lesssim \sqrt{(\ttd_{\G} + \ttd_{\R,\Y}) \log(\ttc_{\G} \ttc_{\R,\Y} \ttk n)} \sqrt{\ttM_{\G} \ttE_{\G}} \kappa_n^{-v/2}.
\end{align*}
Since the restriction of $(\bar{Z}_n^R(g,r): (g,r)\in\F)$ to $\G\times\tilde{\R}$ has the same distribution as the Gaussian process constructed in Part 1, the Vorob'ev--Berkes--Philipp theorem \citep[Theorem 1.31]{dudley2014uniform} allows us to construct $\bar{Z}_n^R$ on the same probability space as $(\bx_i,y_i)_{1 \leq i \leq n}$ and $Z_n^R$, with $\bar{Z}_n^R$ and $Z_n^R$ coinciding on $\G \times \tilde{\R}$. By an abuse of notation, call this extended Gaussian process $Z_n^R$.

\subsubsection*{Part 4: Putting Together}

Combining the previous three parts, choose $\kappa_n$ so that
\begin{align*}
    \mathtt{r}_n \kappa_n \asymp \sqrt{\ttM_{\G} \ttE_{\G}} \kappa_n^{-v/2},
    \qquad\text{that is}\qquad
    \kappa_n \asymp \left(\frac{\sqrt{\ttM_{\G}\ttE_{\G}}}{\mathtt{r}_n}\right)^{\frac{2}{v+2}}.
\end{align*}
Then
\begin{align*}
    \E \big[\lVert R_n - Z_n^R \rVert_{\G \times \R} \big] & \lesssim \ttv(\ttd\log (\ttc n))^{3/2} \mathtt{r}_n^{\frac{v}{v +2}}(\sqrt{\ttM_{\G}\ttE_{\G}})^{\frac{2}{v+2}} + \ttv\ttd\log(\ttc n) \ttM_{\G} n^{-\frac{v/2}{2+v}} \\
    & \qquad +  \ttv\ttd\log(\ttc n) \ttM_{\G} n^{-1/2} \Big(\frac{\sqrt{\ttM_{\G} \ttE_{\G}}}{\mathtt{r}_n}\Big)^{\frac{2}{v+2}} \\
    & \lesssim \ttv \ttd\log(\ttc n)\rho_n,
\end{align*}
where $\ttd = \ttd_{\G} + \ttd_{\R,\Y} + \ttk$, and $\ttc = \ttc_{\G} \ttc_{\R,\Y} \ttk$. This proves the stated bound.
\qed


%%%%%%%%%%%%%%%%%%%%%%%%%%%%%%%%%%%%%%%%%%%%%%%%%%%%%%%%%%%%%%%%%%%%%%
\section{Proofs}\label{sa-sec: Proofs}
%%%%%%%%%%%%%%%%%%%%%%%%%%%%%%%%%%%%%%%%%%%%%%%%%%%%%%%%%%%%%%%%%%%%%%

%%%%%%%%%%%%%%%%%%%%%%%%%%%%%%%%%%%%
\subsection{Proof of Lemma~\ref{sa-lem: Invertibility}}
%%%%%%%%%%%%%%%%%%%%%%%%%%%%%%%%%%%%

Because $\B$ is compact and $\B\subset\setint(\X)$, compact support of $K$ implies that $\bx+\bH\supp(K)\subseteq\X$ uniformly over $\bx\in\B$ for all sufficiently small $\bar h$. Assumption~\ref{sa-assump: Sharp DGP}(ii) and uniform continuity of $f_X$ on the compact set $\X$ imply
\begin{align*}
    \bGamma_{t,\bx} & = \E \Big[\br_p(\bH^{-1}(\bX_i-\bx)) \br_p(\bH^{-1}(\bX_i-\bx))^{\top} K_{\bH}(\bX_i - \bx) \Indicator(\bX_i \in \A_t) \Big] \\
    & = \int_{\A_t} \br_p(\bH^{-1}(\bu - \bx)) \br_p(\bH^{-1}(\bu - \bx))^{\top} K_{\bH}(\bu - \bx) f_X(\bu) \dif \bu \\
    & = f_X(\bx)\int_{\A_t} \br_p(\bH^{-1}(\bu - \bx)) \br_p(\bH^{-1}(\bu - \bx))^{\top} K_{\bH}(\bu - \bx) \dif \bu + o(1),
\end{align*}
where in the last line we have used $\sup_{\bu\in\bx+\bH\supp(K)}|f_X(\bu)-f_X(\bx)|=o(1)$ uniformly over $\bx\in\B$ and $\int_{\A_t} (\bH^{-1}(\bu - \bx))^{\bv} K_{\bH}(\bu - \bx) \dif \bu = O(1)$ for any multi-index $\bv$ from a standard change of variable argument.

\begin{center}
\textbf{I. Polynomial Representation of Minimum Eigenvalue}
\end{center}
For each bandwidth matrix $\bH$, define
\begin{align*}
    \bA_{t,\bx}(\bH) = \int_{\A_t} \br_p(\bH^{-1}(\bu - \bx)) \br_p(\bH^{-1}(\bu - \bx))^{\top} K_{\bH}(\bu - \bx) \dif \bu.
\end{align*}
The argument below establishes a uniform positive lower bound for $\lambda_{\min}(\bA_{t,\bx}(\bH))$ for all sufficiently small $\bar{h}$; no pointwise limit of $\bA_{t,\bx}(\bH)$ is needed.
A change of variable gives
\begin{align*}
    \bA_{t,\bx}(\bH) = \int \br_p(\bz) \br_p(\bz)^{\top} K(\bz) \Indicator(\bx + \bH \bz \in \A_t) \dif \bz.
\end{align*}
Let $\ba \in \reals^{\mathfrak{p}_p}$, where $\mathfrak{p}_p = \frac{(d + p)!}{d! p !}$.
Then the equivalent representation of minimum eigenvalue gives
\begin{align}\label{sa-eq: min eigenvalue}
    \nonumber \lambda_{\min}(\bA_{t,\bx}(\bH))
    & = \min_{\norm{\ba} = 1} \int (\ba^{\top} \br_p(\bz))^2 K(\bz) \Indicator(\bx + \bH \bz \in \A_t) \dif \bz \\
    & \geq \kappa \min_{\norm{\ba} = 1} \int_{U} (\ba^{\top} \br_p(\bz))^2  \Indicator(\bx + \bH \bz \in \A_t) \dif \bz,
\end{align}
where in the last line we have used $K(\bu) \geq \kappa$ for all $\bu \in U$.

\begin{center}
\textbf{II. Mass Retaining Ratio in Treatment/Control Region}
\end{center}

Denote $E_{\bH}(\bx,t) = \{\bz \in U: \bx + \bH \bz \in \A_t \}$. Assumption~\ref{sa-assump: Boundary and Kernel} (ii) implies there is some upper bound $\Lambda > 0$ of $K(\cdot)$. Hence for $c_0 = 1/2 \; \liminf_{\bar{h} \downarrow 0}\inf_{\bx \in \B} \int_{U} K(\bu) \Indicator(\bx + \bH \bu \in \A_t) \dif \bu$, we have
\begin{align*}
    \Lambda \leb(E_{\bH}(\bx,t))
    & \geq \int_{U} K(\bu) \Indicator(\bx + \bH \bu \in \A_t) \dif \bu
    \geq c_0
\end{align*}
for small enough $\bar{h}$, which implies
\begin{align}\label{sa-eq: mass relation}
      \leb(E_{\bH}(\bx,t))  \geq \alpha \leb(U), \qquad \alpha = \frac{c_0}{\Lambda \leb(U)}.
\end{align}

\begin{center}
\textbf{III. $L_2$ Integral of Polynomials in Full Versus Treatment/Control Regions}
\end{center}

Consider $S = \{f \in \mathcal{P}_{p}: \int_U f(\bu)^2 \dif \bu = 1\}$, where $\mathcal{P}_{p}$ is the collection of all polynomials in $\bu$ with total degree at most $p$. Let $(\phi_j, 1 \leq j \leq \mathfrak{p}_p)$ be an orthonormal basis of $(\mathcal{P}_{p}, \lVert \cdot \rVert_{L_2})$. Then $T(\mathbf{a}) = \sum_{j = 1}^{\mathfrak{p}_p} a_j \phi_j$ is an isometry. Since $T(S) = \{\mathbf{a} \in \reals^{\mathfrak{p}_p}: \norm{\ba} = 1\}$ is compact, $S$ is also compact in $(\mathcal{P}_{p}, \lVert \cdot \rVert_{L_2})$. Since $\mathcal{P}_{p}$ is $\mathfrak{p}_p$-dimensional, equivalence of norms implies that $S$ is also compact in $(\mathcal{P}_{p}, \lVert \cdot \rVert_{L_{\infty}})$. Now consider
\begin{align*}
    \Phi_q(\varepsilon) = \leb(\{\bu \in U: |q(\bu)| < \varepsilon \}), \qquad q \in S, \varepsilon > 0,
\end{align*}
and
\begin{align*}
    \psi(q) = \sup \Big\{\varepsilon > 0: \Phi_q(\varepsilon) \leq \frac{\alpha}{2} \leb(U) \Big\}.
\end{align*}
Since $\int_U q^2 = 1$ and $q$ is polynomial, $\lim_{\varepsilon \downarrow 0} \Phi_q(\varepsilon) = 0$ and $\Phi_q(\lVert q\rVert_{\infty}) = \leb(U)$. Continuity and Lipschitzness of $q \in S$ imply $\psi(q) > 0$ for all $q \in S$.

Next, we want to show $\psi$ is lower-semicontinuous on $(\mathcal{P}_{p}, \lVert \cdot \rVert_{L_{\infty}})$. Suppose $q_n \to q$ uniformly on $U$. For every $\varepsilon_0 \in (0, \psi(q))$, there exists $\eta > 0$ such that $\Phi_q(\varepsilon_0) \leq \frac{\alpha}{2} \leb(U) - \eta$. Continuity of polynomials and the fact that level sets of nonzero polynomials have zero Lebesgue measure imply $\Indicator_{\{|q_n| < \varepsilon_0\}}(\cdot) \to \Indicator_{\{|q| < \varepsilon_0\}}(\cdot)$ almost surely. By Dominated Convergence Theorem, $\Phi_{q_n}(\varepsilon_0) \to \Phi_q(\varepsilon_0)$. Hence for large enough $n$, $\Phi_{q_n}(\varepsilon_0) \leq \frac{\alpha}{2}\leb(U)$, which implies $\varepsilon_0 \leq \psi(q_n)$. This implies $\liminf_{n \to \infty} \psi(q_n) \geq \varepsilon_0$. Since $\varepsilon_0$ is arbitrary in $(0, \psi(q))$, we have $\liminf_{n \to \infty} \psi(q_n) \geq \psi(q)$.

Compactness of $S$ and lower-semicontinuity of $\psi$ implies $\psi$ attains its minimum on $S$. Since $\psi(q) > 0$ for all $q \in S$, we know $\varepsilon_* = \inf_{q \in S} \psi(q) > 0$. Then for every $q \in S$,
\begin{align*}
    \int_{E_{\bH}(\bx,t)} q^2
    & \geq \varepsilon_*^2 \; \leb \Big(E_{\bH}(\bx,t) \setminus \{|q| \leq \varepsilon_*\} \Big) \\
    & \geq \varepsilon_*^2 \; \Big(\leb(E_{\bH}(\bx,t)) - \leb(\{|q| \leq \varepsilon_*\}) \Big) \\
    & \geq \varepsilon_*^2 \; \frac{\alpha}{2} \leb(U).
\end{align*}
Scaling $q$ from $S$ gives
\begin{align}\label{sa-eq: integral relation}
    \int_{E_{\bH}(\bx,t)} q^2 \geq \varepsilon_*^2 \; \frac{\alpha}{2} \int_{U} q^2, \qquad q \in \mathcal{P}_{p}.
\end{align}

\begin{center}
\textbf{IV. Lower Bound of Minimum Eigenvalue}
\end{center}

Equations~\eqref{sa-eq: min eigenvalue}, \eqref{sa-eq: mass relation} and \eqref{sa-eq: integral relation} together give for small enough $\bar{h}$,
\begin{align*}
    \inf_{\bx \in \B}\lambda_{\min}(\bA_{t,\bx}(\bH))
    & \geq \kappa \inf_{\bx \in \B}\min_{\norm{\ba} = 1} \int_{E_{\bH}(\bx,t)} (\ba^{\top} \br_p(\bz))^2   \dif \bz, \\
    & \geq \kappa \varepsilon_*^2 \; \frac{\alpha}{2}  \min_{\norm{\ba} = 1} \int_{U} (\ba^{\top} \br_p(\bz))^2 \dif \bz \\
    & \geq \kappa \varepsilon_*^2 \; \frac{\alpha}{2}  \lambda_{\min} \Big(\int_U \br_p(\bz) \br_p(\bz)^{\top} \dif \bz \Big),
\end{align*}
which implies $\liminf_{\bar{h} \to 0} \inf_{\bx \in \B}\lambda_{\min}(\bA_{t,\bx}(\bH)) > 0$. Since Assumption~\ref{sa-assump: Sharp DGP}(ii) gives $\inf_{\bx\in\X}f_X(\bx)>0$ and the preceding expansion of $\bGamma_{t,\bx}$ is uniform over $\bx\in\B$, it follows that
\begin{align*}
    \liminf_{\bar h\to0}\inf_{\bx\in\B}\lambda_{\min}(\bGamma_{t,\bx})
    \geq
    \inf_{\bx\in\X}f_X(\bx)\,
    \liminf_{\bar h\to0}\inf_{\bx\in\B}\lambda_{\min}(\bA_{t,\bx}(\bH))
    >0.
\end{align*}


%%%%%%%%%%%%%%%%%%%%%%%%%%%%%%%%%%%%
\subsection{Proof of Lemma~\ref{sa-lem: Gram Matrix}}
%%%%%%%%%%%%%%%%%%%%%%%%%%%%%%%%%%%%

Since $\widehat{\bGamma}_{t,\bx}$ is a finite dimensional matrix, it suffices to show the stated rate of convergence for each entry. Let $\bv$ be a multi-index such that $|\bv| \leq 2 p$. Define
\begin{align*}
    g_n(\xi,\bx) = (\bH^{-1}(\xi - \bx))^{\bv} K_{\bH}(\xi - \bx)\Indicator(\xi \in \A_t), \qquad \xi \in \mathcal{X}, \bx \in \B.
\end{align*}
and $\F = \{g_n(\cdot, \bx): \bx \in \B\}$. We will show $\F$ is a VC-type class. In order to do this, we study the following quantities.

\medskip\textit{Constant Envelope Function}. We assume $K$ is continuous and has compact support, or $K = \Indicator(\cdot \in [-1,1]^d)$. Hence there exists a constant $C_1$ such that for all $l \in \mathcal{F}$ and all $\xi \in \X$, $|l(\xi)| \leq C_1 |\bH|^{-1} = F$.

\medskip\textit{Diameter of $\F$ in $L_2$}. $ \sup_{l \in \mathcal{F}} \norm{l}_{\P,2}
= \sup_{\bx \in \B}(\int_{\bH^{-1}(\A_t - \bx)} |\bH|^{-1}\by^{2\bv}K(\by)^2 f_X(\bx + \bH \by) \dif \by)^{1/2} \leq C_2 |\bH|^{-1/2}$ for some constant $C_2$. We can take $C_1$ large enough so that $\sigma = C_2 |\bH|^{-1/2} \leq F = C_1 |\bH|^{-1}$.

\medskip\textit{Ratio}. For some constant $C_3$, $\delta = \frac{\sigma}{F}  = C_3 \sqrt{|\bH|}$.

\medskip\textit{Covering Numbers}. Case 1: $K$ is Lipschitz. Let $\bx,\bx' \in \B$. Then, for generic evaluation points $\bx = (x_1,\ldots,x_d)^\top$ and $\bx' = (x_1',\ldots,x_d')^\top$,
\begin{align*}
    \sup_{\xi \in \X} \left|g_n(\xi,\bx) - g_n(\xi, \px) \right|
    & \leq  \left|(\bH^{-1}(\xi-\bx))^{\bv}  - (\bH^{-1}(\xi-\px))^{\bv} \right| K_{\bH}(\xi - \bx) \\
    & \qquad + |(\bH^{-1}(\xi-\px))^{\bv}| \Big|K_{\bH}(\xi - \bx) - K_{\bH}(\xi - \px) \Big| \\
    & \lesssim |\bH|^{-1}\underline h^{-1} \infnorm{\bx - \px},
\end{align*}
since we have assumed that $K$ has compact support and is Lipschitz continuous. Hence, for any $\varepsilon \in (0,1]$ and for any finitely supported measure $Q$ and metric $\norm{\cdot}_{Q,2}$ based on $L_2(Q)$,
\begin{align*}
    N(\F, \norm{\cdot}_{Q,2}, \varepsilon \norm{F}_{Q,2})
    \leq N(\X, \infnorm{\cdot}, \varepsilon \norm{F}_{Q,2} |\bH|\underline h)
    \stackrel{(i)}{\lesssim} \left(\frac{\diam(\X)}{ \varepsilon \norm{F}_{Q,2} |\bH|\underline h}\right)^d
    \lesssim \left(\frac{\diam(\X)}{\varepsilon \underline h}\right)^d,
\end{align*}
where in ($i$) we used compactness of $\X$ and the fact that $\varepsilon \norm{F}_{Q,2}|\bH|\underline h \lesssim \varepsilon \underline h \lesssim 1$. Hence, $\F$ forms a VC-type class, and taking $A_1 = \diam(\X)/\underline h$ and $A_2 = d$, $\sup_Q N(\F, \norm{\cdot}_{Q,2}, \varepsilon \norm{F}_{Q,2}) \lesssim (A_1/ \varepsilon)^{A_2}$, $\varepsilon \in (0,1]$, and where the supremum is over all finite discrete measures.

Case 2: $K = \Indicator(\cdot \in [-1,1]^d)$. Consider
\begin{align*}
    m_n(\xi, \bx)
    = (\bH^{-1}(\xi - \bx))^{\bv} |\bH|^{-1} \Indicator(\xi \in \A_t),
    \qquad \xi, \bx \in \mathcal{X},
\end{align*}
$\M = \{m_n(\cdot, \bx): \bx \in \B\}$ and the constant envelope function $M = C_4 |\bH|^{-1}\underline h^{-|\bv|}$, for some constant $C_4$ only depending on diameter of $\X$. The same argument as before shows that for any discrete measure $Q$, we have
\begin{align*}
    N(\M, \norm{\cdot}_{Q,2}, \varepsilon \norm{M}_{Q,2})
    & \leq N(\X, \infnorm{\cdot}, \varepsilon \norm{M}_{Q,2} |\bH|\underline h^{|\bv|+1})
    \lesssim \Big(\frac{\diam(\X)}{\varepsilon \norm{M}_{Q,2} |\bH|\underline h^{|\bv|+1}}\Big)^d
    \lesssim \Big(\frac{\diam(\X)}{\varepsilon \underline h}\Big)^d.
\end{align*}
The class $\G = \{\Indicator(\bH^{-1}(\cdot - \bx) \in [-1,1]^d): \bx \in \B\}$ has VC dimension no greater than $2d$ \cite[Example 2.6.1]{van-der-Vaart-Wellner_1996_Book}, and by \citet[Theorem 2.6.4]{van-der-Vaart-Wellner_1996_Book}, for any discrete measure $Q$, $N(\G, \norm{\cdot}_{Q,2}, \varepsilon) \leq 2d (4 e)^{2d} \varepsilon^{-4d}$, $0 < \varepsilon \leq 1$. The class $\F$ is contained in the product class $\M\cdot\G$. Since the envelope ratio satisfies $\norm{F}_{Q,2}/\norm{M}_{Q,2}\gtrsim\underline h^{|\bv|}$ uniformly in $Q$, the standard product-covering inequality gives, for a constant $c>0$,
\begin{align*}
    N(\F, \norm{\cdot}_{Q,2}, \varepsilon \norm{F}_{Q,2})
    &\leq
    N(\M, \norm{\cdot}_{Q,2}, c\varepsilon \underline h^{|\bv|} \norm{M}_{Q,2})
    N(\G, \norm{\cdot}_{Q,2}, c\varepsilon \underline h^{|\bv|})\\
    &\lesssim
    \left(\frac{\diam(\X)}{\varepsilon \underline h^{|\bv|+1}}\right)^d
    (\varepsilon\underline h^{|\bv|})^{-4d}
    \lesssim
    \varepsilon^{-5d}\underline h^{-d(5|\bv|+1)}.
\end{align*}
Hence, taking $A_1 = C\underline h^{-(|\bv|+1)}$ and $A_2 = 5d$ for a large enough constant $C$, $\sup_Q N(\F, \norm{\cdot}_{Q,2}, \varepsilon \norm{F}_{Q,2}) \lesssim (A_1/ \varepsilon)^{A_2}$, $\varepsilon \in (0,1]$, where the supremum is over all finite discrete measures.

\medskip\textit{Maximal Inequality}. By Corollary 5.1 in \cite{Chernozhukov-Chetverikov-Kato_2014b_AoS} for the empirical process on class $\F$,
\begin{align*}
     \E \Big[\sup_{\bx \in \B} \big|\En [g_n(\bX_i, \bx)]  - \E[g_n(\bX_i, \bx)]\big| \Big]
    & \lesssim \frac{\sigma}{\sqrt{n}}\sqrt{A_2\log(A_1/\delta)} + \frac{\pnorm{F} A_2\log(A_1/\delta)}{n} \\
    & \lesssim \sqrt{\frac{\log(1/\bar{h})}{n |\bH|}} + \frac{\log(1/\bar{h})}{n |\bH|},
\end{align*}
where $A_1, A_2, \sigma, F, \delta$ are all given previously. Assuming $\frac{\log(1/\bar{h})}{n |\bH|} \to 0$, we conclude that $\sup_{\bx \in \B} \big\|\widehat{\bGamma}_{t,\bx} - \bGamma_{t,\bx}\big\| \lessp \sqrt{\frac{\log(1/\bar{h})}{n |\bH|}}$. By Weyl's Theorem, $\sup_{\bx \in \B}|\lambda_{\min}(\widehat{\bGamma}_{t,\bx}) - \lambda_{\min}(\bGamma_{t,\bx})| \leq \sup_{\bx \in \B} \big\|\widehat{\bGamma}_{t,\bx} - \bGamma_{t,\bx}\big\| \lessp \sqrt{\frac{\log(1/\bar{h})}{n |\bH| }}$. Lemma~\ref{sa-lem: Invertibility} gives $\lambda_{\min}(\bGamma_{t,\bx}) \gtrsim 1$ uniformly over $\bx\in\B$. Therefore, $\inf_{\bx \in \B}\lambda_{\min}(\widehat{\bGamma}_{t,\bx}) \gtrsim_{\P} 1$, $\sup_{\bx \in \B} \big\|\widehat{\bGamma}_{t,\bx}\big\| \lessp 1$, and $\sup_{\bx \in \B} \big\|\widehat{\bGamma}_{t,\bx}^{-1}\big\| \lessp 1$. Hence $\sup_{\bx \in \B} \big\|\widehat{\bGamma}_{t,\bx}^{-1} - \bGamma_{t,\bx}^{-1}\big\| \leq \sup_{\bx \in \B} \big\|\bGamma_{t,\bx}^{-1}\big\| \big\|\bGamma_{t,\bx} -\widehat{\bGamma}_{t,\bx}\big\| \big\|\widehat{\bGamma}_{t,\bx}^{-1}\big\| \lessp \sqrt{\frac{\log(1/\bar{h})}{n |\bH|}}$.


%%%%%%%%%%%%%%%%%%%%%%%%%%%%%%%%%%%%
\subsection{Proof of Lemma~\ref{sa-lem: Linear Approximation}}
%%%%%%%%%%%%%%%%%%%%%%%%%%%%%%%%%%%%

The proof is similar to the proof of Lemma~\ref{sa-lem: Gram Matrix}. Let $\bv$ be a multi-index such that $ 0 \leq |\bv| \leq p$. Let
\begin{align*}
    g_n(\xi,\bx) = (\bH^{-1}(\xi - \bx))^{\bv} K_{\bH}(\xi - \bx)\Indicator(\xi\in \A_t), \qquad \xi , \bx \in \X.
\end{align*}
Define the class of functions $\F = \{(\xi, u) \in \X \times \reals \mapsto g_n(\xi,\bx)u: \bx \in \B\}$.

\medskip\textit{Envelope Function}. Since $K$ is continuous on its compact support, there exists a constant $C_1 > 0$ such that $|g_n(\xi,\bx)u| \leq C_1 |\bH|^{-1} |u|$, for $\xi, \bx \in \X$ and $ u \in \reals$. We define the envelope function $F(\xi, u) = C_1 |\bH|^{-1}|u|$, for $\xi \in \X$ and $ u \in \reals$. Moreover, by Assumption~\ref{sa-assump: Sharp DGP}(v), let $M = \max_{1 \leq i \leq n} F(\bX_i, \varepsilon_i)$, then
\begin{align*}
    \E[M^2]^{1/2}
    \lesssim |\bH|^{-1} \E \big[\max_{1 \leq i \leq n} |\varepsilon_i|^2 \big]^{1/2}
    \lesssim |\bH|^{-1} \E \big[\max_{1 \leq i \leq n} |\varepsilon_i|^{2+v}\big]^{1/(2+v)}
    \lesssim n^{1/(2+v)} |\bH|^{-1}.
\end{align*}

\medskip\textit{Diameter of $\F$ in $L_2$}. Recall we denote $\varepsilon_i = Y_i - \E[Y_i|\bX_i]$, then
\begin{align*}
    \sup_{l \in \mathcal{F}} \E[l(\bX_i, \varepsilon_i)^2]^{1/2}
    \leq \sup_{\xi \in \X} \E[\varepsilon_i^2 | \bX_i = \xi]^{1/2} \sup_{\xi \in \X} \E[g_n(\bX_i, \xi)^2]^{1/2} \leq C_3 |\bH|^{-1/2}
    = \sigma.
\end{align*}

\medskip\textit{Ratio}. We set $\delta = \frac{\sigma}{\norm{F}_{\P,2}} \lesssim |\bH|^{1/2}$.

\medskip\textit{Covering Numbers}. Case 1: $K$ is Lipschitz. Let $\mathbb{Q}$ be a finite distribution on $(\X \times \reals, \B(\X) \otimes \mathsf{Borel}(\reals))$. Let $\bx, \bx' \in \X$. In the proof of Lemma~\ref{sa-lem: Gram Matrix}, we showed that $\sup_{\xi \in \mathcal{X}} \sup_{\bx, \bx' \in \X } \frac{|g_n(\xi,\bx) - g_n(\xi, \bx')|}{\infnorm{\bx - \bx'}} \lesssim |\bH|^{-1}\underline h^{-1}$. Hence,
\begin{align*}
    \| g_n(\bX_i, \bx) \varepsilon_i - g_n(\bX_i, \bx')\varepsilon_i \|_{\mathbb{Q},2}
    \leq \infnorm{g_n(\cdot, \bx) - g_n(\cdot,\bx')} \norm{\varepsilon_i}_{\mathbb{Q},2}
    \lesssim \underline h^{-1} \norm{F}_{\mathbb{Q},2} \infnorm{\bx- \px}.
\end{align*}
It follows that $\sup_{Q} N(\F, \lVert \cdot \rVert_{Q,2}, \epsilon \| F \|_{Q,2}) \lesssim (\frac{\diam(\X)}{\epsilon \underline h})^d$, where $\sup$ is over all finite probability distributions on $(\X \times \reals, \B(\X) \otimes \mathsf{Borel}(\reals))$. Letting $A_1 = \frac{\diam(\X)}{\underline h}$ and $ A_2 = d$, we conclude that
\begin{align*}
    \sup_Q N(\F, \norm{\cdot}_{Q,2}, \epsilon \norm{F}_{Q,2}) \lesssim (A_1/\epsilon)^{A_2}, \qquad \epsilon \in (0,1].
\end{align*}

Case 2: $K$ is the uniform kernel. Let
\begin{align*}
    m_n(\xi, \bx) =  (\bH^{-1}(\xi - \bx))^{\bv} |\bH|^{-1} \Indicator(\xi \in \A_t), \qquad \xi, \bx \in \mathcal{X},
\end{align*}
with $\M = \{(\xi,u) \in \X \times \reals \to m_n(\xi,\bx)u: \bx \in \B\}$ and envelope function $M(\xi,u) = C_1 |\bH|^{-1}\underline h^{-|\bv|} |u|$, for a positive constant $C_1$ depending only on $K$. By similar arguments as Case 1 and the proof of Lemma~\ref{sa-lem: Gram Matrix}, it follows that $\sup_Q N(\M, \norm{\cdot}_{Q,2}, \varepsilon \|M\|_{Q,2}) \lesssim \big(\frac{\diam(\X)}{\varepsilon \underline h}\big)^d$, where the supremum is taken over all finite discrete measures. Taking $\G = \{\Indicator(\bH^{-1}(\cdot - \bx) \in [-1,1]^d): \bx \in \B\}$, the proof of Lemma~\ref{sa-lem: Gram Matrix} shows that
\begin{align*}
    \sup_Q N(\G, \norm{\cdot}_{Q,2}, \varepsilon) \leq 2d (4 e)^{2d} \varepsilon^{-4d}, \qquad \varepsilon \in (0,1],
\end{align*}
where the supremum is taken over all finite discrete measures. The class $\F$ is contained in $\M\cdot\G$. Since $\norm{F}_{Q,2}/\norm{M}_{Q,2}\gtrsim\underline h^{|\bv|}$, the product-covering inequality gives, for a constant $c>0$,
\begin{align*}
    N(\F, \norm{\cdot}_{Q,2}, \varepsilon \norm{F}_{Q,2})
    &\leq
    N(\M, \norm{\cdot}_{Q,2}, c\varepsilon\underline h^{|\bv|}\norm{M}_{Q,2})
    N(\G, \norm{\cdot}_{Q,2}, c\varepsilon\underline h^{|\bv|})\\
    &\lesssim
    \varepsilon^{-5d}\underline h^{-d(5|\bv|+1)}.
\end{align*}
Taking $A_1 = C\underline h^{-(|\bv|+1)}$ and $A_2 = 5d$ for a large enough constant $C$, we have
\begin{align*}
    \sup_Q N(\F, \norm{\cdot}_{Q,2}, \varepsilon \norm{F}_{Q,2}) \lesssim (A_1/ \varepsilon)^{A_2}, \qquad \varepsilon \in (0,1],
\end{align*}
the supremum is over all finite discrete measures.


\medskip\textit{Maximal Inequality}. By Corollary 5.1 in \cite{Chernozhukov-Chetverikov-Kato_2014b_AoS},
\begin{align*}
    \E \Big[\sup_{\bx \in \B} \big|\En [g_n(\bX_i,\bx)\varepsilon_i] \big| \Big]
    & \lesssim \frac{\sigma}{\sqrt{n}}\sqrt{A_2\log(A_1/\delta)} + \frac{\pnorm{M} A_2\log(A_1/\delta)}{n}\\
    & \lesssim \sqrt{\frac{\log(1/\bar{h})}{n |\bH|}} +\frac{\log(1/\bar{h})}{n^{\frac{1+v}{2+v}} |\bH|}.
\end{align*}
Since $\bQ_{t,\bx}$ is finite-dimensional, entry-wise convergence implies convergence in norm with the same rate. Hence, $\sup_{\bx \in \B} \big\|\bQ_{t,\bx}\big\| \lessp \sqrt{\frac{\log(1/\bar{h})}{n |\bH|}} + \frac{\log(1/\bar{h})}{n^{\frac{1+v}{2+v}}|\bH|}$. By Lemma~\ref{sa-lem: Gram Matrix},
\begin{align*}
    \sup_{\bx \in \B} \big|\widehat{\mu}_t^{(\bnu)}(\bx) - \E [\widehat{\mu}_t^{(\bnu)}(\bx) | \bX ]
                           - \bnu!\be_{\bnu}^{\top} \bS^{-1} \bGamma_{t,\bx}^{-1} \bQ_{t,\bx} \big|
    &= \sup_{\bx \in \B} \big| \bnu!\be_{\bnu}^{\top} \bS^{-1} \big(\widehat{\bGamma}_{t,\bx}^{-1} - \bGamma_{t,\bx}^{-1} \big) \bQ_{t,\bx}\big| \\
    & \lessp \bh^{-\bnu} \sqrt{\frac{\log(1/\bar{h})}{n |\bH|}} \Big(\sqrt{\frac{\log(1/\bar{h})}{n |\bH|}} + \frac{\log(1/\bar{h})}{n^{\frac{1+v}{2+v}}|\bH|} \Big),
\end{align*}
and
\begin{align*}
    \sup_{\bx \in \B} \big|\widehat{\mu}_t^{(\bnu)}(\bx) - \E[\widehat{\mu}_t^{(\bnu)}(\bx) | \bX ] \big|
    \lessp \bh^{-\bnu} \Big(\sqrt{\frac{\log(1/\bar{h})}{n |\bH|}} + \frac{\log(1/\bar{h})}{n^{\frac{1+v}{2+v}}|\bH|}\Big),
\end{align*}
which completes the proof.
\qed


%%%%%%%%%%%%%%%%%%%%%%%%%%%%%%%%%%%%
\subsection{Proof of Lemma~\ref{sa-lem: covariance}}
%%%%%%%%%%%%%%%%%%%%%%%%%%%%%%%%%%%%

Let $\eta_i(\bx) = \sum_{t \in \{0,1\}} \Indicator(\bX_i \in \A_t) (\mu_t(\bX_i) - \widehat{\bbeta}_t(\bx)^\top \br_p(\bX_i - \bx))$. Then, for all $\bx, \by \in \B$, the difference between the estimated and true variance matrices is
\begin{align*}
    \widehat{\bSigma}_{t,\bx,\by} - \bSigma_{t,\bx,\by} = \bM_{1,\bx,\by} + \bM_{2,\bx,\by} + \bM_{3,\bx,\by} + \bM_{4,\bx,\by}
\end{align*}
where
\begin{align*}
    \bM_{1,\bx,\by}
    &= \En \Big[\br_p(\bH^{-1}(\bX_i - \bx)) \br_p(\bH^{-1}(\bX_i - \by))^{\top} |\bH| K_{\bH}(\bX_i - \bx)K_{\bH}(\bX_i - \by)\eta_i(\bx) \eta_i(\by) \Indicator(\bX_i \in  \A_t)\Big],\\
    \bM_{2,\bx,\by}
    &= \En \Big[\br_p(\bH^{-1}(\bX_i - \bx)) \br_p(\bH^{-1}(\bX_i - \by))^{\top} |\bH| K_{\bH}(\bX_i - \bx)K_{\bH}(\bX_i - \by) (\eta_i(\bx) + \eta_i(\by)) \varepsilon_i \Indicator(\bX_i \in \A_t)\Big],\\
    \bM_{3,\bx,\by}
    &= \En \Big[\br_p(\bH^{-1}(\bX_i - \bx))\br_p(\bH^{-1}(\bX_i - \by))^{\top}|\bH|K_{\bH}(\bX_i - \bx)K_{\bH}(\bX_i - \by)(\varepsilon_i^2 - \sigma_t(\bX_i)^2)\Indicator(\bX_i \in \A_t)\Big], \\
    \bM_{4,\bx,\by}
    &= \En \Big[\br_p(\bH^{-1}(\bX_i - \bx))\br_p(\bH^{-1}(\bX_i - \by))^{\top}|\bH|K_{\bH}(\bX_i - \bx)K_{\bH}(\bX_i - \by)\sigma_t(\bX_i)^2\Indicator(\bX_i \in  \A_t)\Big]\\
    &\qquad - \E \Big[\br_p(\bH^{-1}(\bX_i - \bx))\br_p(\bH^{-1}(\bX_i - \by))^{\top}|\bH|K_{\bH}(\bX_i - \bx)K_{\bH}(\bX_i - \by)\sigma_t(\bX_i)^2 \Indicator(\bX_i \in  \A_t)\Big].
\end{align*}

For $\bu$ and $\bv$ multi-indices, let $g_n(\bX_i; \bx, \by) = |\bH|K_{\bH}(\bX_i - \bx)K_{\bH}(\bX_i - \by)(\bH^{-1}(\bX_i - \bx))^{\bu}(\bH^{-1}(\bX_i - \by))^{\bv} \Indicator(\bX_i \in  \A_t)$. Set
\begin{align*}
    \mathfrak{R}_{\mathtt{EP}} = \sqrt{\frac{\log (1/\bar{h})}{n |\bH|}} + \frac{\log (1/\bar{h})}{n^{\frac{1+v}{2+v}}|\bH|}.
\end{align*}

First, we present a bound on $\max_{1 \leq i \leq n} |\eta_i(\bx)| \Indicator(\bH^{-1}(\bX_i - \bx) \in \supp(K))$. By Lemma~\ref{sa-lem: bias} and Lemma~\ref{sa-lem: Linear Approximation}, for any multi-index $\boldsymbol{\alpha}$ such that $|\boldsymbol{\alpha}| \leq p$,
\begin{align*}
    \sup_{\bx \in \B}|\be_{\boldsymbol{\alpha}}^\top (\widehat{\bbeta}_t(\bx) - \bbeta_t(\bx))| \lessp \bh^{-\boldsymbol{\alpha}} (\bar{h}^{p+1} + \mathfrak{R}_{\mathtt{EP}}).
\end{align*}
Since $K$ is compactly supported, we have
\begin{align*}
     \max_{1 \leq i \leq n} \Big|\sum_{t \in \{0,1\}}\Indicator(\bX_i \in \A_t) (\widehat{\bbeta}_t(\bx) - \bbeta_t(\bx))^\top \br_p(\bX_i - \bx) \Indicator(\bH^{-1}(\bX_i - \bx) \in \supp(K))\Big|
     \lessp \bar{h}^{p+1} + \mathfrak{R}_{\mathtt{EP}}.
\end{align*}
Since $\mu_t$ is $p+1$ times continuously differentiable,
\begin{align*}
    \max_{1 \leq i \leq n} \Big|\sum_{t \in \{0,1\}}\Indicator(\bX_i \in \A_t) (\mu_t(\bX_i)
    - \bbeta_t(\bx)^\top \br_p(\bX_i - \bx)) \Indicator(\bH^{-1}(\bX_i - \bx) \in \supp(K))\Big|
    \lesssim \bar{h}^{p+1}.
\end{align*}
It follows that
\begin{align*}
    \sup_{\bx \in \B}\max_{1 \leq i \leq n} |\eta_i(\bx)| \Indicator(\bH^{-1}(\bX_i - \bx) \in \supp(K))
    \lessp \bar{h}^{p+1} + \mathfrak{R}_{\mathtt{EP}}.
\end{align*}

\medskip\textit{Term $\bM_{1,\bx,\by}$}. From the proof for Lemma~\ref{sa-lem: Gram Matrix}, applied to the absolute-value class generated by $g_n$, $\sup_{\bx,\by \in \B}\En[ |g_n(\bX_i; \bx,\by)|]\lessp 1$. Thus,
\begin{align*}
    &\sup_{\bx, \by \in \B} \big|\En[g_n(\bX_i; \bx, \by) \eta_i(\bx) \eta_i(\by)] \big|\\
    & \leq
      \Big(\sup_{\bx \in \B}\max_{1 \leq i \leq n} |\eta_i(\bx)| \Indicator(\bH^{-1}(\bX_i - \bx) \in \supp(K))\Big)^2
      \sup_{\bx, \by \in \B}\En[|g_n(\bX_i; \bx, \by)|] \\
    & \lessp (\bar{h}^{p+1} + \mathfrak{R}_{\mathtt{EP}})^2,
\end{align*}
where we have used the preceding local-support bound on $\eta_i(\bx)$. Finite dimensionality of $\bM_{1,\bx,\by}$ then implies
\begin{align*}
    \sup_{\bx, \by \in \B}\norm{\bM_{1,\bx,\by}} \lessp (\bar{h}^{p+1} + \mathfrak{R}_{\mathtt{EP}})^2.
\end{align*}

\medskip\textit{Term $\bM_{2,\bx,\by}$}. The same maximal-inequality argument used for Lemma~\ref{sa-lem: Linear Approximation}, applied to the absolute-value class $|g_n(\cdot;\bx,\by)u|$, gives $\sup_{\bx,\by \in \B}\En[|g_n(\bX_i; \bx, \by)\varepsilon_i|]\lessp 1$. Thus,
\begin{align*}
    &\sup_{\bx, \by \in \B} \big|\En[g_n(\bX_i; \bx, \by)(\eta_i(\bx) + \eta_i(\by)) \varepsilon_i] \big|\\
    & \leq
      2\sup_{\bx \in \B}\max_{1 \leq i \leq n} |\eta_i(\bx)|\Indicator(\bH^{-1}(\bX_i-\bx)\in\supp(K))
      \sup_{\bx, \by \in \B} \En[|g_n(\bX_i; \bx, \by) \varepsilon_i|]\\
    &\lessp \bar{h}^{p+1} + \mathfrak{R}_{\mathtt{EP}},
\end{align*}
which implies that
\begin{align*}
    \sup_{\bx, \by \in \B}\norm{\bM_{2,\bx,\by}} \lessp \bar{h}^{p+1} + \mathfrak{R}_{\mathtt{EP}}.
\end{align*}

\medskip\textit{Term $\bM_{3,\bx,\by}$}. Define $l_n(\cdot, \cdot;\bx,\by): \X \times \reals \to \reals$ as
\begin{align*}
    l_n(\xi, \varepsilon;\bx, \by)
    = |\bH|K_{\bH}(\xi - \bx)K_{\bH}(\xi - \by)(\bH^{-1}(\xi - \bx))^{\bu}(\bH^{-1}(\xi - \by))^{\bv} \Indicator(\xi \in \A_t) (\varepsilon^2 - \sigma_t^2(\xi)),
\end{align*}
and consider the function class $\mathcal{L} = \{l_n(\cdot,\cdot;\bx,\by): \bx,\by \in \B\}$. Let $L: \X \times \reals \to \reals$ be $L(\xi, \varepsilon) = c|\bH|^{-1}|\varepsilon^2 - \sigma_t^2(\xi)|$ with $c = \sup_{\bx, \by \in \B} \big|(\bH^{-1}(\xi - \bx))^{\bu}(\bH^{-1}(\xi - \by))^{\bv}K(\bH^{-1}(\xi - \bx)) K(\bH^{-1}(\xi - \by))\big|$. By similar argument as in the proof for Lemma~\ref{sa-lem: Linear Approximation}, we can show $\mathcal{L}$ is a VC-type class such that $\E [l_n(\bX_i,\varepsilon_i; \bx, \by)] = 0$, for all $\bx,\by \in \B$,
\begin{align*}
    \sup_{\bx,\by \in \B}\E[l_n(\bX_i,\varepsilon_i;\bx,\by)^2]^{\frac{1}{2}}
    \lesssim \sup_{\bx, \by \in \B} \E[g_n(\bX_i; \bx, \by)^2]^{\frac{1}{2}}
    \sup_{\xi \in \X}\E[(\varepsilon_i^2-\sigma_t^2(\xi))^2|\bX_i = \xi]^{1/2}
    \lesssim |\bH|^{-1/2}
\end{align*}
because Assumption~\ref{sa-assump: Sharp DGP}(v), with $v\ge2$, gives a uniformly bounded conditional fourth moment for $\varepsilon_i$.
and
\begin{align*}
    \E\big[\max_{1 \leq i \leq n} L(\bX_i,\varepsilon_i)^2\big]^{\frac{1}{2}}
    \lesssim |\bH|^{-1} \E\big[\max_{1 \leq i \leq n}\varepsilon_i^4\big]^{1/2}
    \lesssim |\bH|^{-1} \E\big[\max_{1 \leq i \leq n} |\varepsilon_i|^{2+v}\big]^{\frac{2}{2+v}}
    \lesssim |\bH|^{-1}n^{\frac{2}{2+v}}.
\end{align*}
Applying Corollary 5.1 in \cite{Chernozhukov-Chetverikov-Kato_2014b_AoS}, we obtain
\begin{align*}
    \sup_{\bx, \by \in \B} \big|\En[l_n(\bX_i,\varepsilon_i;\bx,\by)] \big|
    \lessp \sqrt{\frac{\log (1/\bar{h})}{n |\bH|}} + \frac{\log (1/\bar{h})}{n^{\frac{v}{2+v}}|\bH|}
\end{align*}
and
\begin{align*}
    \sup_{\bx, \by \in \B} \norm{\bM_{3,\bx,\by}}
    \lessp \sqrt{\frac{\log (1/\bar{h})}{n |\bH|}} + \frac{\log (1/\bar{h})}{n^{\frac{v}{2+v}}|\bH|}.
\end{align*}

\medskip\textit{Term $\bM_{4,\bx,\by}$}. Notice that $\{g_n(\cdot; \bx, \by)\sigma_t^2(\cdot): \bx, \by \in \B\}$ is a VC-type class with constant envelope function $C |\bH|^{-1}$ for some positive constant $C$, where
\begin{align*}
    \sup_{\bx, \by \in \B} \sup_{\xi \in \X}|g_n(\xi; \bx, \by)\sigma_t^2(\xi)| \lesssim |\bH|^{-1}
\end{align*}
and
\begin{align*}
    \sup_{\bx, \by \in \B} \E[g_n(\bX_i; \bx, \by)^2 \sigma_t(\bX_i)^2]^{\frac{1}{2}} \lesssim |\bH|^{-1/2}.
\end{align*}
Then, similar to the proof of $\bM_{1,\bx,\by}$, we conclude that
\begin{align*}
     \sup_{\bx, \by \in \B} \big|\En[g_n(\bX_i; \bx, \by)\sigma_t^2(\bX_i)] -\E[g_n(\bX_i; \bx, \by)\sigma_t^2(\bX_i)] \big|
     \lessp \sqrt{\frac{\log (1/\bar{h})}{n |\bH|}}
\end{align*}
and
\begin{align*}
     \sup_{\bx, \by \in \B} \norm{\bM_{4,\bx,\by}}
     \lessp \sqrt{\frac{\log (1/\bar{h})}{n |\bH|}}.
\end{align*}

\medskip\textit{Final result}. Combining the upper bounds of the four terms,
\begin{align*}
    \sup_{\bx,\by \in \B} \big\|\widehat{\bSigma}_{t, \bx,\by} - \bSigma_{t, \bx,\by} \big\|
    \lessp \bar{h}^{p+1} + \sqrt{\frac{\log (1/\bar{h})}{n |\bH|}} + \frac{\log (1/\bar{h})}{n^{\frac{v}{2+v}}|\bH|},
\end{align*}
which implies $\sup_{\bx, \by \in \B} \|\widehat{\bSigma}_{t,\bx,\by}\| \lessp 1$. It follows that, for each $t\in\{0,1\}$,
\begin{align*}
    \sup_{\bx, \by \in \B}  |\widehat{\Omega}_{t,\bx,\by}^{(\bnu)} - \Omega_{t,\bx,\by}^{(\bnu)}|
    & \leq  \frac{1}{n |\bH|\bh^{2\bnu}}  \Big( \sup_{\bx, \by \in \B} \big\|\widehat{\bGamma}_{t,\bx}^{-1} - \bGamma_{t,\bx}^{-1} \big\| \big\|\widehat{\bSigma}_{t,\bx,\by}\big\| \big\|\widehat{\bGamma}_{t,\by}^{-1}\big\|\\
    &\qquad +  \sup_{\bx, \by \in \B} \big\|\bGamma_{t,\bx}^{-1}\big\|  \big\|\widehat{\bSigma}_{t,\bx,\by} - \bSigma_{t,\bx,\by}\big\| \big\|\widehat{\bGamma}_{t,\by}^{-1}\big\|\\
    &\qquad +  \sup_{\bx, \by \in \B} \big\|\bGamma_{t,\bx}^{-1}\big\| \big\|\bSigma_{t,\bx,\by}\big\| \big\|\widehat{\bGamma}_{t,\by}^{-1} - \bGamma_{t,\by}^{-1}\big\| \Big) \\
    & \lesssim_\P \frac{1}{n |\bH|\bh^{2\bnu}} \Big(\bar{h}^{p+1} + \sqrt{\frac{\log (1/\bar{h})}{n |\bH|}} + \frac{\log (1/\bar{h})}{n^{\frac{v}{2+v}}|\bH|}\Big).
\end{align*}
Summing the last display over $t\in\{0,1\}$ gives the asserted bound for $\widehat{\Omega}_{\bx,\by}^{(\bnu)} - \Omega_{\bx,\by}^{(\bnu)}$. By Assumption~\ref{sa-assump: Sharp DGP}(iv) and Assumption~\ref{sa-assump: Boundary and Kernel}(ii), $\inf_{\bx \in \B}\Omega_{\bx,\bx}^{(\bnu)} \gtrsim (n |\bH|\bh^{2\bnu})^{-1}$. Therefore, $\inf_{\bx \in \B}\widehat{\Omega}_{\bx,\bx}^{(\bnu)} \gtrsim_{\P} (n |\bH|\bh^{2\bnu})^{-1}$. Furthermore,
\begin{align*}
    & \sup_{\bx \in \B} \Big|\sqrt{\widehat{\Omega}_{\bx,\bx}^{(\bnu)}} - \sqrt{\Omega_{\bx,\bx}^{(\bnu)}}\Big|
    \lessp \sup_{\bx \in \B}\sqrt{n |\bH|\bh^{2\bnu}} \Big|\widehat{\Omega}_{\bx,\bx}^{(\bnu)} - \Omega_{\bx,\bx}^{(\bnu)}\Big|
    \lessp \frac{1}{\sqrt{n |\bH|\bh^{2\bnu}}} \Big(\bar{h}^{p+1} + \sqrt{\frac{\log (1/\bar{h})}{n |\bH|}} + \frac{\log (1/\bar{h})}{n^{\frac{v}{2+v}}|\bH|}\Big)
\end{align*}
and
\begin{align*}
    \sup_{\bx \in \B} \left|\frac{\bh^{-\bnu}}{\sqrt{\widehat{\Omega}_{\bx,\bx}^{(\bnu)}}} - \frac{\bh^{-\bnu}}{\sqrt{\Omega_{\bx,\bx}^{(\bnu)}}}\right|
    = \bh^{-\bnu}\sup_{\bx \in \B} \left|\frac{\sqrt{\widehat{\Omega}_{\bx,\bx}^{(\bnu)}} - \sqrt{\Omega_{\bx,\bx}^{(\bnu)}}}{\sqrt{\widehat{\Omega}_{\bx,\bx}^{(\bnu)} \Omega_{\bx,\bx}^{(\bnu)}}} \right|
    \lessp \sqrt{n |\bH|} \Big(\bar{h}^{p+1} + \sqrt{\frac{\log (1/\bar{h})}{n |\bH|}} + \frac{\log (1/\bar{h})}{n^{\frac{v}{2+v}}|\bH|}\Big),
\end{align*}
which completes the proof.
\qed

%%%%%%%%%%%%%%%%%%%%%%%%%%%%%%%%%%%%
\subsection{Proof of Lemma~\ref{sa-lem: bias}}
%%%%%%%%%%%%%%%%%%%%%%%%%%%%%%%%%%%%

Define
\begin{align*}
    \boldsymbol{\chi}_{t, \bx}
    = \En \Big[\br_p(\bH^{-1}(\bX_i-\bx)) K_{\bH}(\bX_i - \bx)\Indicator(\bX_i \in \A_t)\fr_t(\bX_i;\bx) \Big],
    \quad
    \fr_t(\xi;\bx) = \mu_t(\xi) - \sum_{0 \leq |\bomega| \leq p}\frac{\mu^{(\bomega)}_t(\bx)}{\bomega!}(\xi - \bx)^{\bomega}.
\end{align*}
Since $\mu_t$ is $(p+1)$-times continuously differentiable, Taylor's theorem and compact support of $K$ imply
\begin{align*}
    \sup_{\bx\in\B}\max_{1\leq i\leq n}
    \big|\fr_t(\bX_i;\bx)\big|\Indicator\{K_{\bH}(\bX_i-\bx)\neq0\}
    \lesssim \bar{h}^{p+1}.
\end{align*}
Since $\frac{\log(1/\bar{h})}{n|\bH|} = o(1)$, the same argument as the proof of Lemma~\ref{sa-lem: Gram Matrix} shows
\begin{align*}
    \sup_{\bx\in\B}
    \En \Big[\big\lVert \br_p(\bH^{-1}(\bX_i-\bx)) \big\rVert K_{\bH}(\bX_i - \bx)\Indicator(\bX_i \in \A_t)\Big] \lessp 1.
\end{align*}
Therefore, $\sup_{\bx\in\B}\norm{\boldsymbol{\chi}_{t,\bx}}\lessp \bar{h}^{p+1}$.
It then follows from Lemma~\ref{sa-lem: Gram Matrix} that
\begin{align*}
    \sup_{\bx \in \B} \big|\E[\widehat{\mu}_t^{(\bnu)}(\bx)|\bX] - \mu_t^{(\bnu)}(\bx)\big|
    = \sup_{\bx \in \B} \big|\bnu!\be_{\bnu}^{\top} \bS^{-1} \widehat{\bGamma}_{t,\bx}^{-1} \boldsymbol{\chi}_{t,\bx}\big|
    \lessp \bar{h}^{p + 1}\bh^{-\bnu}.
\end{align*}

Now also assume that $\bar{h} = o(1)$. By Assumption~\ref{sa-assump: Sharp DGP}(iii), for all $\bx \in \B$ and $\xi \in \X$,
\begin{align*}
    \fr_t(\xi;\bx)
    = \sum_{|\bv|=p+1}\frac{\mu_t^{(\bv)}(\bx)}{\bv!}(\xi-\bx)^{\bv}
      + \sum_{|\bv|=p+1} r_{t,\bv}(\xi;\bx)(\xi-\bx)^{\bv},
\end{align*}
where
\begin{align*}
    \max_{|\bv|=p+1}\sup_{\bx\in\B,\xi\in\X}
    |r_{t,\bv}(\xi;\bx)|\Indicator\{K_{\bH}(\xi-\bx)\neq0\}=o(1).
\end{align*}
Letting
\begin{align*}
    \widetilde{\boldsymbol{\chi}}_{t,\bx}
    = \En \Big[\br_p(\bH^{-1}(\bX_i-\bx)) K_{\bH}(\bX_i - \bx)\Indicator(\bX_i \in \A_t) ( \sum_{|\bv| = p+1} \frac{\mu_t^{(\bv)}(\bx)}{\bv!}(\bX_i - \bx)^{\bv})\Big],
\end{align*}
we conclude that
\begin{align*}
   \sup_{\bx \in \B} \big\| \boldsymbol{\chi}_{t,\bx} - \widetilde{\boldsymbol{\chi}}_{t,\bx} \big\|
   &\lesssim M_n \sup_{\bx \in \B} \Big\| \En \Big[\br_p(\bH^{-1}(\bX_i-\bx)) K_{\bH}(\bX_i - \bx) \Indicator(\bX_i \in \A_t) \Big(\sum_{|\bv| = p+1} \frac{\left|\bX_i - \bx\right|^{\bv}}{\bv!} \Big) \Big] \Big\|\\
   &= o_{\P}(\bar{h}^{p+1}),
\end{align*}
where the last equality employs the same arguments as in the proof of Lemma~\ref{sa-lem: Gram Matrix}. Hence,
\begin{align*}
    \sup_{\bx \in \B} \Big|\E[\widehat{\mu}_t^{(\bnu)}(\bx)|\bX] - \mu_t^{(\bnu)}(\bx) - \widehat{B}_{t,\bx}^{(\bnu)}(\bh)\Big|
    &= \sup_{\bx \in \B} \Big|\bnu!\be_{\bnu}^{\top} \bS^{-1} \widehat{\bGamma}_{t,\bx}^{-1} \boldsymbol{\chi}_{t,\bx} - \bnu!\be_{\bnu}^{\top} \bS^{-1} \widehat{\bGamma}_{t,\bx}^{-1} \widetilde{\boldsymbol{\chi}}_{t,\bx}\Big|\\
    &= o_{\P}(\bar{h}^{p+1}\bh^{-\bnu}).
\end{align*}
Using Lemma~\ref{sa-lem: Gram Matrix} and the maximal inequality as in the proof of Lemma~\ref{sa-lem: Gram Matrix}, we conclude that
\begin{align*}
    \max_{t \in \{0,1\}} \sup_{\bx \in \B} \big|\widehat{B}_{t,\bx}^{(\bnu)}(\bh) - B_{t,\bx}^{(\bnu)}(\bh)\big|
    \lesssim_{\P} \bar{h}^{p+1}\bh^{-\bnu}\sqrt{\frac{\log(1/\bar{h})}{n |\bH|}}.
\end{align*}
Since $\max_{t \in \{0,1\}}\sup_{\bx \in \B}|B_{t,\bx}^{(\bnu)}(\bh)| \lesssim \bar{h}^{p+1}\bh^{-\bnu}$, it follows that $\max_{t \in \{0,1\}} \sup_{\bx \in \B}|\widehat{B}_{t,\bx}^{(\bnu)}(\bh)| \lessp \bar{h}^{p+1}\bh^{-\bnu}$.
\qed


%%%%%%%%%%%%%%%%%%%%%%%%%%%%%%%%%%%%
\subsection{Proof of Theorem~\ref{sa-thm: Convergence Rates: Sharp BATEC}}
%%%%%%%%%%%%%%%%%%%%%%%%%%%%%%%%%%%%

For a fixed $\bx\in\B$, the pointwise versions of Lemma~\ref{sa-lem: Linear Approximation} and Lemma~\ref{sa-lem: bias} decompose
$\widehat{\mu}^{(\bnu)}_t(\bx)-\mu_t^{(\bnu)}(\bx)$ into a stochastic component and a conditional bias component of the stated pointwise orders. Applying this decomposition to both $t=0$ and $t=1$ and using the triangle inequality gives the pointwise rate for $\widehat{\tau}^{(\bnu)}(\bx)-\tau^{(\bnu)}(\bx)$. Taking suprema over $\bx\in\B$ and using the displayed uniform bounds in the two lemmas gives the uniform rate.
\qed

%%%%%%%%%%%%%%%%%%%%%%%%%%%%%%%%%%%%
\subsection{Proof of Theorem~\ref{sa-thm: MSE Expansions: Sharp BATEC}}
%%%%%%%%%%%%%%%%%%%%%%%%%%%%%%%%%%%%

For the conditional bias, by Lemma~\ref{sa-lem: bias},
\begin{align*}
    & \sup_{\bx \in \B} \big|\E \big[\widehat{\tau}^{(\bnu)}(\bx) - \tau^{(\bnu)}(\bx) \big| \bX \big]^2 - (B_{\bx}^{(\bnu)}(\bh))^2 \big| \\
    & \leq \sup_{\bx \in \B} \big|\E \big[\widehat{\tau}^{(\bnu)}(\bx) - \tau^{(\bnu)}(\bx) \big| \bX \big] - B_{\bx}^{(\bnu)}(\bh) \big| \cdot \sup_{\bx \in \B} \big|\E \big[\widehat{\tau}^{(\bnu)}(\bx) - \tau^{(\bnu)}(\bx) \big| \bX \big] + B_{\bx}^{(\bnu)}(\bh) \big| \\
    & = o_{\P}(\bar{h}^{2p+2}\bh^{-2\bnu}).
\end{align*}

For the conditional variance, by Lemma~\ref{sa-lem: covariance},
\begin{align*}
    \sup_{\bx \in \B} \big|\Var\big[\widehat{\tau}^{(\bnu)}(\bx) \big|\bX \big] -  (n|\bH|\bh^{2\bnu})^{-1} V_{\bx}^{(\bnu)} \big|
    = o_{\P}((n|\bH|\bh^{2\bnu})^{-1}).
\end{align*}

The pointwise MSE expansion follows directly. For the IMSE expansion, notice that
\begin{align*}
    & \Big|\operatorname{IMSE}^{(\bnu)} - \int_{\B} \big[(B_{\bx}^{(\bnu)}(\bh))^2 + (n|\bH|\bh^{2\bnu})^{-1} V_{\bx}^{(\bnu)}\big] w(\bx) \dif\Haus^{d-1}(\bx)\Big|\\
    & \leq \int_{\B} |w(\bx)| \dif\Haus^{d-1}(\bx) \cdot \sup_{\bx \in \B} \big|\operatorname{MSE}^{(\bnu)}(\bx) -  (B_{\bx}^{(\bnu)}(\bh))^2 - (n|\bH|\bh^{2\bnu})^{-1} V_{\bx}^{(\bnu)}\big|\\
    & = o_{\P} \big(\bar{h}^{2p+2}\bh^{-2\bnu} + (n|\bH|\bh^{2\bnu})^{-1} \big),
\end{align*}
which completes the proof.
\qed

%%%%%%%%%%%%%%%%%%%%%%%%%%%%%%%%%%%%
\subsection{Proof of Theorem~\ref{sa-thm: Confidence Intervals: Sharp BATEC}}
%%%%%%%%%%%%%%%%%%%%%%%%%%%%%%%%%%%%

We have $\overline{\Tstat}^{(\bnu)}(\bx) = \sum_{i = 1}^n \chi_i$ with
\begin{align*}
    \chi_i = \sum_{t \in \{0,1\}} s_t n^{-1} (\Omega_{\bx,\bx}^{(\bnu)})^{-1/2} \bnu!\be_{\bnu}^{\top} \bS^{-1} \bGamma_{t,\bx}^{-1} \br_p(\bH^{-1}(\bX_i - \bx)) K_{\bH}(\bX_i - \bx) \Indicator(\bX_i \in \A_t) \varepsilon_i,
\end{align*}
where $s_1=1$, $s_0=-1$, $\E[\chi_i] = 0$, and $\Var[\chi_i] = n^{-1}$. By the Berry-Esseen Theorem,
\begin{align*}
    \sup_{u \in \reals} \Big|\P \big(\overline{\Tstat}^{(\bnu)}(\bx) \leq u\big) - \Phi(u) \Big|
    \lesssim B_n^{-1} \sum_{i = 1}^n \E[|\chi_i|^3],
\end{align*}
where $B_n  = \sum_{i = 1}^n \Var[\chi_i] = 1$. Moreover,
\begin{align*}
    \sum_{i = 1}^n \E[|\chi_i|^3]
    & =  n^{-3} (\Omega_{\bx,\bx}^{(\bnu)})^{-3/2} \sum_{i = 1}^n \E \Big[\Big|\sum_{t \in \{0,1\}} s_t\bnu!\be_{\bnu}^{\top} \bS^{-1} \bGamma_{t,\bx}^{-1} \br_p(\bH^{-1}(\bX_i - \bx)) K_{\bH}(\bX_i - \bx) \Indicator(\bX_i \in \A_t) \varepsilon_i \Big|^3 \Big]\\
    & \lesssim n^{-3} (\Omega_{\bx,\bx}^{(\bnu)})^{-3/2} \sum_{i = 1}^n \E \Big[\Big|\sum_{t \in \{0,1\}}s_t\bnu!\be_{\bnu}^{\top} \bS^{-1} \bGamma_{t,\bx}^{-1} \br_p(\bH^{-1}(\bX_i - \bx)) K_{\bH}(\bX_i - \bx) \Indicator(\bX_i \in \A_t)\Big|^3 \Big]\\
    & \lesssim n^{-2} \bh^{-\bnu}|\bH|^{-1} (\Omega_{\bx,\bx}^{(\bnu)})^{-3/2} \E \Big[\Big|\sum_{t \in \{0,1\}}s_t\bnu!\be_{\bnu}^{\top} \bS^{-1} \bGamma_{t,\bx}^{-1} \br_p(\bH^{-1}(\bX_i - \bx)) K_{\bH}(\bX_i - \bx) \Indicator(\bX_i \in \A_t)\Big|^2 \Big]\\
    & \lesssim n^{-1} \bh^{-\bnu}|\bH|^{-1}  (\Omega_{\bx,\bx}^{(\bnu)})^{-1/2}\\
    & \lesssim (n |\bH|)^{-1/2},
\end{align*}
where the second line uses Assumption~\ref{sa-assump: Sharp DGP}(v): since $v\ge2$, the conditional third moment of $\varepsilon_i$ is uniformly bounded, so conditioning on $\bX_i$ absorbs the factor $|\varepsilon_i|^3$ into a constant. The third line uses
\begin{align*}
    \bigg|\sum_{t \in \{0,1\}}s_t\bnu!\be_{\bnu}^{\top} \bS^{-1} \bGamma_{t,\bx}^{-1} \br_p(\bH^{-1}(\bX_i - \bx)) K_{\bH}(\bX_i - \bx) \Indicator(\bX_i \in \A_t)\bigg| \lesssim \bh^{-\bnu}|\bH|^{-1}
\end{align*}
and the fourth line uses the definition of $\Omega_{\bx,\bx}^{(\bnu)}$.

Finally, although Lemma~\ref{sa-lem: Gram Matrix} through Lemma~\ref{sa-lem: covariance} establish uniform convergence over $\bx \in \B$, the corresponding pointwise results for a fixed $\bx$ follow by restricting the function classes in those proofs to singletons indexed by $\bx$. In particular, when handling the variance terms in Lemma~\ref{sa-lem: covariance} (specifically $\bM_{3,\bx,\by}$), we bound terms involving $\varepsilon_i^2$ using either the Fuk--Nagaev inequality (under $2+v$ moments) or Markov's inequality. Thus, for each fixed $\bx\in\B$,
\begin{align*}
    \|\widehat{\bGamma}_{t,\bx}^{-1}-\bGamma_{t,\bx}^{-1}\|
    &=O_{\P}\big((n|\bH|)^{-1/2}\big),\\
    \|\widehat{\bSigma}_{t,\bx,\bx}-\bSigma_{t,\bx,\bx}\|
    &=O_{\P}\big((n|\bH|)^{-1/2}+\bar h^{p+1}\big),
\end{align*}
and therefore
\begin{align*}
    \left|\frac{\widehat{\Omega}_{\bx,\bx}^{(\bnu)}}{\Omega_{\bx,\bx}^{(\bnu)}}-1\right|
    &=O_{\P}\big((n|\bH|)^{-1/2}+\bar h^{p+1}\big),\\
    \left|(\widehat{\Omega}_{\bx,\bx}^{(\bnu)})^{-1/2}
          -(\Omega_{\bx,\bx}^{(\bnu)})^{-1/2}\right|
    &\lessp \sqrt{n|\bH|\bh^{2\bnu}}
      \big((n|\bH|)^{-1/2}+\bar h^{p+1}\big).
\end{align*}
Together with
\begin{align*}
    \big|\bnu!\be_{\bnu}^{\top}\bS^{-1}\bGamma_{t,\bx}^{-1}\bQ_{t,\bx}\big|
    =O_{\P}\big(\bh^{-\bnu}(n|\bH|)^{-1/2}\big)
\end{align*}
and the bias bound
\begin{align*}
    \big|\E[\widehat{\tau}^{(\bnu)}(\bx)|\bX]-\tau^{(\bnu)}(\bx)\big|
    =O_{\P}(\bar h^{p+1}\bh^{-\bnu}),
\end{align*}
this yields
\begin{align}\label{sa-eq: lin error 2d}
\Big|\widehat{\Tstat}^{(\bnu)}(\bx) - \overline{\Tstat}^{(\bnu)}(\bx)\Big|
\lessp
\bar{h}^{p+1}\sqrt{n |\bH|} + \frac{1}{\sqrt{n |\bH|}},
\end{align}
provided that $\bar{h} \to 0$ and $n |\bH| \to \infty$.

The final results follow by weak convergence to a Gaussian distribution, and properties of the distribution function.
\qed

%%%%%%%%%%%%%%%%%%%%%%%%%%%%%%%%%%%%
\subsection{Proof of Theorem~\ref{sa-thm: Stochastic Linearization: Sharp BATEC}}
%%%%%%%%%%%%%%%%%%%%%%%%%%%%%%%%%%%%

For all $\bx \in \B$, we have $\widehat{\Tstat}^{(\bnu)}(\bx) = \overline{\Tstat}^{(\bnu)}(\bx) + G_1^{(\bnu)}(\bx) +  G_2^{(\bnu)}(\bx)$, where
\begin{align*}
  G_1^{(\bnu)}(\bx)
  = \Big(\E\big[\widehat{\tau}^{(\bnu)}(\bx) \big| \bX \big] - \tau^{(\bnu)}(\bx)\Big) (\widehat{\Omega}_{\bx,\bx}^{(\bnu)})^{-1/2},
\end{align*}
and
\begin{align*}
  G_2^{(\bnu)}(\bx)
  = \bnu!\be_{\bnu}^{\top}\bS^{-1} \Big[\big(\hGammap^{-1} \bQ_{1, \bx}- \hGamman^{-1} \bQ_{0, \bx}\big) (\widehat{\Omega}_{\bx,\bx}^{(\bnu)})^{-\frac{1}{2}} - \big(\Gammap^{-1} \bQ_{1, \bx}- \Gamman^{-1} \bQ_{0, \bx}\big)(\Omega_{\bx,\bx}^{(\bnu)})^{-\frac{1}{2}}\Big].
\end{align*}
By Lemma~\ref{sa-lem: bias} and Lemma~\ref{sa-lem: covariance},
\begin{align*}
    \sup_{\bx \in \B} \big|G_1^{(\bnu)}(\bx)\big| \lessp \bar{h}^{p+1}\bh^{-\bnu} (n |\bH|\bh^{2\bnu})^{1/2} \lesssim \bar{h}^{p+1}\sqrt{n |\bH|}.
\end{align*}
By Lemma~\ref{sa-lem: Gram Matrix}, Lemma~\ref{sa-lem: Linear Approximation} and Lemma~\ref{sa-lem: covariance},
\begin{align*}
    \sup_{\bx \in \B} \Big|\bnu!\be_{\bnu}^{\top} \bS^{-1} \big[\widehat{\bGamma}_{t,\bx}^{-1} - \bGamma_{t,\bx}^{-1}\big] \bQ_{t,\bx} (\widehat{\Omega}_{\bx,\bx}^{(\bnu)})^{-1/2}\Big|
    \lessp \sqrt{\log(1/\bar{h})} \bigg(\sqrt{\frac{\log(1/\bar{h})}{n |\bH|}} + \frac{\log(1/\bar{h})}{n^{\frac{1 + v}{2+v}}|\bH|}\bigg)
\end{align*}
and
\begin{align*}
    & \sup_{\bx \in \B} \Big|\bnu!\be_{\bnu}^{\top} \bS^{-1} \bGamma_{t,\bx}^{-1} \bQ_{t,\bx} \Big[(\widehat{\Omega}_{\bx,\bx}^{(\bnu)})^{-1/2} - (\Omega_{\bx,\bx}^{(\bnu)})^{-1/2} \Big] \Big| \\
    & \lessp \bh^{-\bnu} \cdot \bigg(\sqrt{\frac{\log(1/\bar{h})}{n |\bH|}}  + \frac{\log(1/\bar{h})}{n^{\frac{1+v}{2+v}}|\bH|} \bigg) \cdot \sqrt{n|\bH|\bh^{2\bnu}} \bigg(\sqrt{\frac{\log(1/\bar{h})}{n |\bH|}}  + \frac{\log(1/\bar{h})}{n^{\frac{v}{2+v}}|\bH|} + \bar{h}^{p+1}\bigg) \\
    & \lesssim \frac{\log(1/\bar{h})}{\sqrt{n |\bH|}} + \frac{(\log(1/\bar{h}))^{3/2}}{n^{\frac{v}{2+v}}|\bH|}
      + \bar{h}^{p+1}\sqrt{n|\bH|}\bigg(\sqrt{\frac{\log(1/\bar{h})}{n |\bH|}}  + \frac{\log(1/\bar{h})}{n^{\frac{1+v}{2+v}}|\bH|} \bigg) \\
    & \lesssim \frac{\log(1/\bar{h})}{\sqrt{n |\bH|}} + \frac{(\log(1/\bar{h}))^{3/2}}{n^{\frac{v}{2+v}}|\bH|} + o(\bar{h}^{p+1}\sqrt{n |\bH|}).
\end{align*}
The last term is absorbed by the bound for $G_1^{(\bnu)}$, and the two displayed logarithmic terms give the second component of the rate in the theorem.
The result now follows from combining the bounds above.
\qed

%%%%%%%%%%%%%%%%%%%%%%%%%%%%%%%%%%%%
\subsection{Proof of Theorem~\ref{sa-thm: Gaussian Strong Approximation: Sharp Tstat}}
%%%%%%%%%%%%%%%%%%%%%%%%%%%%%%%%%%%%

We verify the high-level conditions of Theorem~\ref{sa-thm: Residual Empirical Process Strong Approximation}. We will employ the following technical lemma.

\begin{lem}[VC Class to VC2 Class]\label{sa-lem: VC Class to VC2 Class}
    Assume $\F$ is a VC class on a measure space $(\X, \mathcal{B})$: there exists an envelope function $F$ and positive constants $c(\F), d(\F)$ such that for all $\varepsilon\in(0,1)$,
    \begin{align*}
        \sup_{Q}N(\mathcal{F}, \norm{\cdot}_{Q,1}, \varepsilon \norm{F}_{Q,1}) \leq c(\F)\varepsilon^{-d(\F)},
    \end{align*}
    where the supremum is taken over all finite discrete measures. Then, $\F$ is also VC2 class: for all $\varepsilon\in(0,1)$,
    \begin{align*}
        \sup_{Q} N(\mathcal{F}, \norm{\cdot}_{Q,2}, \varepsilon \norm{F}_{Q,2}) \leq c(\F) (\varepsilon^2/2)^{-d(\mathcal{F})},
    \end{align*}
    where the supremum is taken over all finite discrete measures.
\end{lem}

\noindent\textbf{Proof of Lemma~\ref{sa-lem: VC Class to VC2 Class}}. Let $Q$ be a finite discrete probability measure. Let $f,g \in \mathcal{F}$. Then, $\int |f - g|^2 \dif Q \leq 2 \int |f - g| |F| \dif Q$. Define another probability measure $\tilde{Q}(c_k) = F(c_k) Q(c_k) / \norm{F}_{Q,1}$ on the support of $Q$, denoted by $\{c_1, \ldots, c_k, \ldots\}$. Then,
\begin{align*}
    \int |f - g|^2 \dif Q
    \leq 2 \norm{F}_{Q,1} \int |f - g| \dif \tilde{Q}
    \leq 2 \norm{F}_{Q,1} \norm{f - g}_{\tilde{Q},1}.
\end{align*}
Hence, if we take an $\varepsilon^2/2$-net in $(\mathcal{F}, \norm{\cdot}_{\tilde{Q},1})$ with cardinality no greater than $c(\mathcal{F})(\varepsilon^2/2)^{- d (\mathcal{F})}$, then for any $f \in \mathcal{F}$, there exists a $g \in \mathcal{F}$ such that $\norm{f - g}_{\tilde{Q},1} \leq \varepsilon^2/2 \norm{F}_{\tilde{Q},1}$, and hence
\begin{align*}
    \norm{f - g}_{Q,2}^2 \leq 2 \varepsilon^2/2 \norm{F}_{Q,1} \norm{F}_{\tilde{Q},1} \leq \varepsilon^2 \norm{F}_{Q,2}^2,
\end{align*}
which gives the result.
\qed

By the rectangular-support condition imposed in Theorem~\ref{sa-thm: Gaussian Strong Approximation: Sharp Tstat}, after a componentwise affine rescaling we can take $\X=[0,1]^d$. Since $f_X$ is continuous and bounded away from zero on this compact rectangle, \citet[Section 3.1]{Cattaneo-Yu_2025_AOS} implies that $\Q_{\F_t}=\P_X$ is a valid surrogate measure for $\P_X$ with respect to $\F_t$, and that the Rosenblatt transform $\phi_{\F_t}=T_{\P_X}$ is a valid normalizing transformation. The constants $\ttc_1$ and $\ttc_2$ from Theorem~\ref{sa-thm: Residual Empirical Process Strong Approximation} are therefore finite and uniformly bounded by constants depending only on the density bounds and the support geometry.

Consider first the class of functions $\F_t = \{\K_t^{(\bnu)}(\cdot; \bx): \bx \in \B\}$, for $t \in \{0,1\}$.

\medskip\textit{Envelope Function}. By Lemma~\ref{sa-lem: Invertibility}, Assumption~\ref{sa-assump: Sharp DGP}(iv), and Assumption~\ref{sa-assump: Boundary and Kernel}, uniformly in $\bx \in \B$,
\begin{align*}
    \Omega_{\bx,\bx}^{(\bnu)} \asymp (n|\bH|\bh^{2\bnu})^{-1}.
\end{align*}
Together with the compact support of $K$ and $\|\bS^{-1}\be_{\bnu}\|\asymp \bh^{-\bnu}$, this gives
\begin{align*}
    \sup_{\bx \in \B} \sup_{\xi \in \X} \big|\K_t^{(\bnu)}(\xi; \bx)\big|
    \lesssim n^{-1/2} (n|\bH|\bh^{2\bnu})^{1/2} \bh^{-\bnu} |\bH|^{-1}
    \lesssim |\bH|^{-1/2}.
\end{align*}
Hence, there exists a constant $C_1 > 0$ such that $\ttM_{\F_t} = C_1|\bH|^{-1/2}$ is a constant envelope function.

\medskip\textit{$L_1$ Bound}. Since $\K_t^{(\bnu)}(\cdot;\bx)$ is supported on $\bx+\bH\supp(K)$, whose Lebesgue measure is of order $|\bH|$,
\begin{align*}
    \ttE_{\F_t} = \sup_{\bx \in \B} \E[|\K_t^{(\bnu)}(\bX_i; \bx)|] \lesssim |\bH|^{1/2}.
\end{align*}

\medskip\textit{Uniform Variation}. Case 1: $K$ is Lipschitz. By Assumption \ref{sa-assump: Sharp DGP}(iv) and Assumption \ref{sa-assump: Boundary and Kernel},
\begin{align*}
    \ttL_{\F_t}
    = \sup_{\bx \in \B}\sup_{\xi, \xi' \in \X} \frac{|\K_t^{(\bnu)}(\xi; \bx) - \K_t^{(\bnu)}(\xi'; \bx)|}{\infnorm{\xi - \xi'}}
    \lesssim |\bH|^{-1/2}\underline h^{-1}.
\end{align*}
Each entry of $\bGamma_{t,\bx}$ and $\bSigma_{t,\bx,\bx}$ is a smooth local integral with argument $\bH^{-1}(\xi-\bx)$ and is therefore $\underline h^{-1}$-Lipschitz in $\bx$. Similarly,
\begin{align*}
    \big|\Omega_{t,\bx,\bx}^{(\bnu)} - \Omega_{t,\px,\px}^{(\bnu)}\big|
    \lesssim (n|\bH|\bh^{2\bnu}\underline h)^{-1}\infnorm{\bx-\px},
\end{align*}
which, using $\Omega_{t,\bx,\bx}^{(\bnu)}\asymp(n|\bH|\bh^{2\bnu})^{-1}$, implies
\begin{align*}
    \big|(\Omega_{t,\bx,\bx}^{(\bnu)})^{-1/2} - (\Omega_{t,\px,\px}^{(\bnu)})^{-1/2}\big|
    \lesssim \underline h^{-1}(n|\bH|\bh^{2\bnu})^{1/2}\infnorm{\bx-\px}.
\end{align*}
It follows that the class is Lipschitz with respect to the point of evaluation:
\begin{align*}
    \mathtt{l}_{\F_t}
    = \sup_{\xi \in \X}\sup_{\bx, \px \in \B} \frac{\left|\K_t^{(\bnu)}(\xi; \bx) - \K_t^{(\bnu)}(\xi; \px) \right|}{\infnorm{\bx - \px}}
    \lesssim |\bH|^{-1/2}\underline h^{-1}.
\end{align*}
Since $\K_t^{(\bnu)}(\cdot;\bx)$ is supported on $\bx+\bH\supp(K)$, we have
\begin{align*}
    \ttTV_{\F_t}
    \lesssim |\bH| \, |\bH|^{-1/2}\underline h^{-1}
    \lesssim |\bH|^{1/2}\underline h^{-1}.
\end{align*}

Case 2: $K = \Indicator(\cdot \in [-1,1]^d)$. Define the smooth part
\begin{align*}
    \tilde{\K}^{(\bnu)}_t(\bu;\bx)
    = n^{-1/2} (\Omega_{\bx,\bx}^{(\bnu)})^{-1/2}\bnu!\be_{\bnu}^{\top} \bS^{-1} \bGamma_{t,\bx}^{-1} \br_p(\bH^{-1}(\bu - \bx)) |\bH|^{-1},
    \qquad \bu \in \X,\quad t \in \{0,1\}.
\end{align*}
Then, $\K_t^{(\bnu)}(\bu;\bx) = \tilde{\K}_t^{(\bnu)}(\bu;\bx)\Indicator(\bH^{-1}(\bu-\bx)\in[-1,1]^d)$, and the preceding argument gives $\ttTV_{\tilde{\F}_t}\lesssim |\bH|^{1/2}\underline h^{-1}$. For $\mathscr{L}=\{\Indicator(\bH^{-1}(\cdot-\bx)\in[-1,1]^d):\bx\in\B\}$, the boundary of each rectangle has perimeter $O(|\bH|\underline h^{-1})$, so the product rule yields
\begin{align*}
    \ttTV_{\F_t}
    \leq \ttTV_{\tilde{\F}_t}\ttM_{\mathscr{L}} + \ttM_{\tilde{\F}_t}\ttTV_{\mathscr{L}}
    \lesssim |\bH|^{1/2}\underline h^{-1}.
\end{align*}

\medskip\textit{VC-type Class}. Case 1: $K$ is Lipschitz. We apply \citet[Lemma 7]{Cattaneo-Chandak-Jansson-Ma_2024_Bernoulli} to the rescaled maps
\begin{align*}
    f_{\bx}(\cdot)
    = \frac{\bnu!}{\sqrt{n \Omega_{\bx,\bx}^{(\bnu)}}}\be_{\bnu}^{\top} \bS^{-1} \bGamma_{t,\bx}^{-1} \br_p(\cdot)K(\cdot),
    \qquad \bx \in \B,
\end{align*}
indexed by $\mathcal{H}=\{f_{\bx}(\bH^{-1}(\cdot-\bx)):\bx\in\B\}$. The boundedness and compact support conditions are unchanged; the Lipschitz constant in the index is now of order $\underline h^{-1}$. Therefore, for $0<\varepsilon\leq1$,
\begin{align*}
    \sup_{Q} N\big(\F_t,\norm{\cdot}_{Q,2},(2c+1)^{d+1}\varepsilon\ttM_{\F_t}\big)
    \leq \bc'2^{2d+4}\varepsilon^{-2d-4}+1,
\end{align*}
where the supremum is taken over all finite discrete measures.

\emph{Case 2: Suppose $K=\Indicator(\cdot\in[-1,1]^d)$.} The smooth class $\tilde{\F}_t$ satisfies the same entropy bound. The indicator class $\mathscr{L}$ is a VC class of axis-aligned rectangles with side lengths $2h_j$, so
\begin{align*}
    \sup_Q N(\mathscr{L},\norm{\cdot}_{Q,2},\varepsilon)\leq C\varepsilon^{-4d},\qquad 0<\varepsilon\leq1.
\end{align*}
The product entropy bound then gives, for a constant $c>0$,
\begin{align*}
    \sup_Q N(\F_t,\norm{\cdot}_{Q,2},\varepsilon C_1\ttM_{\tilde{\F}_t})
    &\leq
    \sup_Q N(\tilde{\F}_t,\norm{\cdot}_{Q,2},c\varepsilon\ttM_{\tilde{\F}_t})
    \sup_Q N(\mathscr{L},\norm{\cdot}_{Q,2},c\varepsilon)\\
    &\leq C_2\varepsilon^{-(6d+4)},
\end{align*}
for constants depending only on $d$.

\bigskip
Consider next the class of functions $\G = \{g_{\bx}: \bx \in \B\}$, where $g_{\bx}(\bu)=\Indicator(\bu\in\A_1)\K_1^{(\bnu)}(\bu;\bx)-\Indicator(\bu\in\A_0)\K_0^{(\bnu)}(\bu;\bx)$. We have immediately that $\ttM_{\G}\lesssim|\bH|^{-1/2}$, $\ttE_{\G}\lesssim|\bH|^{1/2}$, and
\begin{align*}
    \sup_Q N(\G,\norm{\cdot}_{Q,2},\varepsilon(2c+1)^{d+1}\ttM_{\G})\leq 2\bc'\varepsilon^{-4d-4}+2.
\end{align*}

\medskip\textit{Total Variation}. Observe that $\Indicator(\bu\in\A_t)\K_t^{(\bnu)}(\bu;\bx)\neq0$ implies $\bu\in E_{t,\bx}=\{\by\in\A_t:\bH^{-1}(\by-\bx)\in\supp(K)\}$. By assumption, $\mathcal{L}(E_{t,\bx})\leq C\bar{h}^{d-1}$. Using the product rule and the bandwidth comparability condition,
\begin{align*}
    \ttTV_{\G}
    &\leq \sup_{\bx\in\B}\sum_{t\in\{0,1\}}\ttTV_{\{\K_t^{(\bnu)}(\cdot;\bx)\}}+\ttM_{\F_t}\ttTV_{\{\Indicator_{E_{t,\bx}}\}}\\
    &\lesssim |\bH|^{1/2}\underline h^{-1}+|\bH|^{-1/2}\bar{h}^{d-1}
    \lesssim |\bH|^{1/2}\underline h^{-1}.
\end{align*}

The imposed singleton regularity condition verifies the remaining truncated conditional-mean requirement in Theorem~\ref{sa-thm: Residual Empirical Process Strong Approximation}. We have therefore verified all high-level sufficient conditions, which immediately give the result.
\qed

%%%%%%%%%%%%%%%%%%%%%%%%%%%%%%%%%%%%
\subsection{Proof of Lemma~\ref{sa-lem: localized perimeter piecewise linear}}
%%%%%%%%%%%%%%%%%%%%%%%%%%%%%%%%%%%%

Let $\mathcal{N}_{\bx,h}=\bx+h\,\supp(K)$. Since $\supp(K)$ is bounded, there is a constant $R<\infty$ such that $\mathcal{N}_{\bx,h}\subseteq B(\bx,Rh)$ for all $\bx\in\B$ and all $h>0$. Since $\B$ is a finite union of line segments, say $\B=\cup_{j=1}^{J}L_j$, uniformly over $\bx\in\B$,
\begin{align*}
    \Haus^1(\B\cap \mathcal{N}_{\bx,h})
    \leq \Haus^1(\B\cap B(\bx,Rh))
    \leq \sum_{j=1}^{J}\Haus^1(L_j\cap B(\bx,Rh))
    \leq 2JRh
    \lesssim h.
\end{align*}
Because $\B$ lies in the interior of $\X$, for all small enough $h$ the window $\mathcal{N}_{\bx,h}$ does not intersect $\setbd(\X)$ uniformly over $\bx\in\B$. Write $\mathcal{L}(\A_t;\mathcal{N}_{\bx,h})$ for the perimeter of $\A_t$ relative to $\mathcal{N}_{\bx,h}$. The relative boundary of $\A_t$ inside $\mathcal{N}_{\bx,h}$ is contained in $\B\cap \mathcal{N}_{\bx,h}$, up to finitely many endpoints, and therefore
\begin{align*}
    \mathcal{L}(\A_t;\mathcal{N}_{\bx,h})\lesssim \Haus^1(\B\cap \mathcal{N}_{\bx,h})\lesssim h.
\end{align*}

If $K(\bu)=\Indicator(\bu\in[-1,1]^2)$, then $\mathcal{N}_{\bx,h}$ is a square with perimeter $8h$. More generally, if $\supp(K)$ has finite De Giorgi perimeter, the scaling property of perimeter gives $\mathcal{L}(\mathcal{N}_{\bx,h})=h\mathcal{L}(\supp(K))\lesssim h$. By the standard subadditivity/product rule for finite-perimeter sets,
\begin{align*}
    \mathcal{L}(\A_t\cap \mathcal{N}_{\bx,h})
    \leq \mathcal{L}(\A_t;\mathcal{N}_{\bx,h})+\mathcal{L}(\mathcal{N}_{\bx,h})
    \lesssim h,
\end{align*}
uniformly over $\bx\in\B$ and $t\in\{0,1\}$. Since $d=2$ and $\bar h=h$, this is the claimed localized perimeter bound. \qed

%%%%%%%%%%%%%%%%%%%%%%%%%%%%%%%%%%%%
\subsection{Proof of Theorem~\ref{sa-thm: Confidence Bands: Sharp BATEC -- process}}
%%%%%%%%%%%%%%%%%%%%%%%%%%%%%%%%%%%%

The proof is divided into three theorem-specific technical lemmas. We first record a generic comparison lemma that will be used to translate a uniform process approximation into a Kolmogorov distance bound for suprema.

\begin{lem}[Generic Suprema Comparison]\label{sa-lem: generic suprema comparison}
Let $\mathbb{T}_j=(\mathbb{T}_j(\bx):\bx\in\B)$, $j=1,2$, be real-valued processes defined on a common probability space. Write
\[
    M_j=\sup_{\bx\in\B}|\mathbb{T}_j(\bx)|,\quad j=1,2,
    \qquad
    \Delta_{12}=\sup_{\bx\in\B}|\mathbb{T}_1(\bx)-\mathbb{T}_2(\bx)|.
\]
For any $\eta>0$,
\begin{align*}
    \sup_{u\in\reals}\left|\P(M_1\le u)-\P(M_2\le u)\right|
    \le
    \mathsf{AC}_2(\eta)
    +\P(\Delta_{12}>\eta),
\end{align*}
where
\[
    \mathsf{AC}_2(\eta)
    =
    \sup_{u\in\reals}\P(|M_2-u|\le \eta).
\]
If, in addition, $\mathbb{T}_2$ is a separable mean-zero Gaussian process and the Gaussian anti-concentration bound applies to $M_2$, then
\begin{align*}
    \sup_{u\in\reals}\left|\P(M_1\le u)-\P(M_2\le u)\right|
    \le
    4\eta\{\E[M_2]+1\}
    +\P(\Delta_{12}>\eta).
\end{align*}
Consequently, if $\E[\Delta_{12}]\le r_n$, then the last probability can be further bounded by $r_n/\eta$.
\end{lem}

\noindent\textbf{Proof of Lemma~\ref{sa-lem: generic suprema comparison}.}
On the event $\{\Delta_{12}\le\eta\}$, the inequalities
\[
    M_1\le u \implies M_2\le u+\eta,
    \qquad
    M_2\le u-\eta \implies M_1\le u
\]
hold for every $u\in\reals$. Therefore,
\begin{align*}
    \P(M_1\le u)-\P(M_2\le u)
    &\le \P(u<M_2\le u+\eta)
       +\P(\Delta_{12}>\eta),\\
    \P(M_2\le u)-\P(M_1\le u)
    &\le \P(u-\eta<M_2\le u)
       +\P(\Delta_{12}>\eta).
\end{align*}
Taking absolute values and then the supremum over $u$ gives the first result. The second display follows from the Gaussian anti-concentration inequality of \citet[Theorem 2.1]{Chernozhukov-Chetverikov-Kato_2014a_AoS}, applied to the symmetrized Gaussian process whose supremum is $M_2$. The final claim follows from Markov's inequality.
\qed


For the theorem-specific lemmas below, define the strong-approximation rate
\[
    \mathfrak{R}_{\mathtt{sa}}
    =
    (\log n)^{\frac{3}{2}}\Big(\frac{1}{n |\bH|}\Big)^{\frac{1}{2d+2}\cdot \frac{v}{v+2}}
    + \log(n)\sqrt{\frac{1}{n^{\frac{v}{2+v}}|\bH|}},
\]
the stochastic-linearization remainder
\[
    \mathfrak{R}_{\mathtt{lin}}
    =
    \sqrt{\log(1/\bar{h})}
    \Big(\sqrt{\frac{\log(1/\bar{h})}{n|\bH|}}
    + \frac{\log(1/\bar{h})}{n^{\frac{v}{2+v}}|\bH|} \Big)
    + \bar{h}^{p+1} \sqrt{n|\bH|},
\]
and the covariance-estimation rate
\[
    \mathfrak{R}_{\mathtt{cov}}
    =
    \sqrt{\frac{\log n}{n|\bH|}}
    + \frac{\log n}{n^{\frac{v}{2+v}}|\bH|}.
\]
Also set the covariance-comparison rate $\mathfrak{R}_{\mathtt{cc}}=\mathfrak{R}_{\mathtt{cov}}+\bar h^{p+1}$.

\begin{lem}[KS Distance Between $\overline{\Tstat}^{(\bnu)}$ and $Z^{(\bnu)}$]\label{sa-lem: infeasible gaussian to bahadur representation}
Suppose the conditions of Theorem~\ref{sa-thm: Confidence Bands: Sharp BATEC -- process} hold. Then,
\begin{align*}
    \sup_{u \in \reals} \Big|\P \Big(\sup_{\bx \in \B} \big|\overline{\Tstat}^{(\bnu)}(\bx)\big| \leq u\Big)
     - \P \Big(\sup_{\bx \in \B} \big|Z^{(\bnu)}(\bx)\big| \leq u\Big)\Big|
    \lesssim (\log n)^{1/4}\mathfrak{R}_{\mathtt{sa}}^{1/2}.
\end{align*}
\end{lem}


\noindent\textbf{Proof of Lemma~\ref{sa-lem: infeasible gaussian to bahadur representation}.} Let $a_n$ be a positive sequence to be determined below. Applying Lemma~\ref{sa-lem: generic suprema comparison} with $\mathbb{T}_1=\overline{\Tstat}^{(\bnu)}$, $\mathbb{T}_2=Z^{(\bnu)}$, and $\eta=a_n$, and then using the coupling tail bound in Theorem~\ref{sa-thm: Gaussian Strong Approximation: Sharp Tstat}, gives
\begin{align*}
    &\sup_{u\in\reals}\Big|
    \P \Big(\sup_{\bx \in \B} \big|\overline{\Tstat}^{(\bnu)}(\bx)\big| \leq u\Big)
    -\P \Big(\sup_{\bx \in \B} \big|Z^{(\bnu)}(\bx)\big| \leq u\Big)
    \Big|\\
    &\qquad\le
    4a_n\left(\E\left[\sup_{\bx\in\B}|Z^{(\bnu)}(\bx)|\right]+1\right)
    + \frac{C\mathfrak{R}_{\mathtt{sa}}}{a_n}.
\end{align*}

For the next displays, write $\K(\bu,\bx)=g_{\bx}(\bu)$ and $\sigma^2(\bu)=\Var(Y_i|\bX_i=\bu)$. Notice that $Z^{(\bnu)}(\bx), \bx \in \B$ is a mean-zero Gaussian process satisfying
\begin{align*}
    \E \big[\big(Z^{(\bnu)}(\bx) - Z^{(\bnu)}(\by)\big)^2\big]^{\frac{1}{2}}
    &= \E \big[(\K(\bX_i, \bx) - \K(\bX_i, \by))^2 \sigma(\bX_i)^2\big]^{\frac{1}{2}}  \\
    &\leq C' l_{n,2} \|\bx - \by\|_{\infty},
\end{align*}
where $C'$ is a constant and $l_{n,2} \asymp \underline h^{-1}$. Hence the canonical metric
$$
\rho(\bx,\by)
:= \E \big[\big(Z^{(\bnu)}(\bx) - Z^{(\bnu)}(\by)\big)^2\big]^{1/2}
\lesssim \underline h^{-1}\|\bx-\by\|_{\infty}.
$$
Moreover,
$$
\sup_{\bx\in\B} \E \big[ Z^{(\bnu)}(\bx)^2 \big]
= \sup_{\bx\in\B} \E \big[\K(\bX_i,\bx)^2 \sigma^2(\bX_i)\big]
\lesssim 1.
$$

Let $N(\varepsilon,\B,\rho)$ denote the covering number of $\B$ under the metric $\rho$.
Since $\rho(\bx,\by)\lesssim \underline h^{-1}\|\bx-\by\|_{\infty}$ and $\B$ is compact in $\R^d$,
$$
N(\varepsilon,\B,\rho)
\lesssim
N(\underline h\varepsilon,\B,\|\cdot\|_{\infty})
\lesssim
\Big(\frac{1}{\underline h\varepsilon}\Big)^d,
\qquad 0<\varepsilon\le 1,
$$
and therefore
$$
\log N(\varepsilon,\B,\rho)
\lesssim
d\log\!\Big(\frac{1}{\underline h\varepsilon}\Big).
$$

Applying Corollary 2.2.8 in \cite{van-der-Vaart-Wellner_1996_Book} yields
\begin{align*}
\E\!\left[\sup_{\bx\in\B}\big|Z^{(\bnu)}(\bx)\big|\right]
&\lesssim
\int_0^{1}
\sqrt{\log N(\varepsilon,\B,\rho)}\dif \varepsilon  \\
&\lesssim
\int_0^{1}
\sqrt{\log\!\Big(\frac{1}{\underline h\varepsilon}\Big)}\dif \varepsilon
\lesssim
\sqrt{\log(1/\underline h)} \lesssim \sqrt{\log n}.
\end{align*}

Choosing $a_n \asymp (\mathfrak{R}_{\mathtt{sa}}/\sqrt{\log n})^{1/2}$, the result now follows.
\qed

\begin{lem}[KS Distance Between $\widehat{\Tstat}^{(\bnu)}$ and $\overline{\Tstat}^{(\bnu)}$]\label{sa-lem: t stats to bahadur}
Suppose the conditions in Theorem~\ref{sa-thm: Confidence Bands: Sharp BATEC -- process} hold. Then, for any multi-index $|\bnu| \leq p$,
\begin{align*}
    \sup_{u \in \reals} \Big| \P \Big(\sup_{\bx \in \B} \big|\widehat{\Tstat}^{(\bnu)}(\bx)\big| \leq u\Big)
    - \P \Big(\sup_{\bx \in \B} \big|\overline{\Tstat}^{(\bnu)}(\bx)\big| \leq u \Big)\Big| = o(1).
\end{align*}
\end{lem}

\noindent\textbf{Proof of Lemma~\ref{sa-lem: t stats to bahadur}.} Choose a deterministic sequence $a_n$ such that $\mathfrak{R}_{\mathtt{lin}}=o(a_n)$ and $a_n\sqrt{\log n}=o(1)$; such a sequence exists under the side conditions of Theorem~\ref{sa-thm: Confidence Bands: Sharp BATEC -- process}. Then,
\[
    \Delta_{\widehat T,\overline T}
    :=
    \sup_{\bx \in \B} \big|\overline{\Tstat}^{(\bnu)}(\bx) - \widehat{\Tstat}^{(\bnu)}(\bx)\big|
    =
    o_{\P}(a_n).
\]
Define
\[
    M_{\widehat T}=\sup_{\bx\in\B}|\widehat{\Tstat}^{(\bnu)}(\bx)|,
    \qquad
    M_{\overline T}=\sup_{\bx\in\B}|\overline{\Tstat}^{(\bnu)}(\bx)|,
    \qquad
    M_Z=\sup_{\bx\in\B}|Z^{(\bnu)}(\bx)|.
\]
Applying the first part of Lemma~\ref{sa-lem: generic suprema comparison} with $\mathbb{T}_1=\widehat{\Tstat}^{(\bnu)}$ and $\mathbb{T}_2=\overline{\Tstat}^{(\bnu)}$ gives
\[
    \sup_{u\in\reals}\big|\P(M_{\widehat T}\le u)-\P(M_{\overline T}\le u)\big|
    \le
    \sup_{u\in\reals}\P(|M_{\overline T}-u|\le a_n)
    +\P(\Delta_{\widehat T,\overline T}>a_n).
\]
The anti-concentration term for $M_{\overline T}$ can be bounded through the Gaussian approximation in Lemma~\ref{sa-lem: infeasible gaussian to bahadur representation}: for every $u\in\reals$,
\begin{align*}
    \P(|M_{\overline T}-u|\le a_n)
    &\le
    \P(|M_Z-u|\le a_n)
    +2\sup_{v\in\reals}\big|\P(M_{\overline T}\le v)-\P(M_Z\le v)\big| \\
    &\le
    4a_n\{\E[M_Z]+1\}
    +2\sup_{v\in\reals}\big|\P(M_{\overline T}\le v)-\P(M_Z\le v)\big|.
\end{align*}
By Lemma~\ref{sa-lem: infeasible gaussian to bahadur representation}, the last Kolmogorov distance is $O((\log n)^{1/4}\mathfrak{R}_{\mathtt{sa}}^{1/2})=o(1)$. From the proof of that lemma, $\E[M_Z]\lesssim\sqrt{\log n}$. Since $a_n\sqrt{\log n}=o(1)$ and $\P(\Delta_{\widehat T,\overline T}>a_n)=o(1)$, the result follows.
\qed


\begin{lem}[KS Distance Between $Z^{(\bnu)}$ and $\widehat{Z}^{(\bnu)}$]\label{sa-lem: feasible gaussian to infeasible gaussian}
Suppose the conditions for Theorem~\ref{sa-thm: Confidence Bands: Sharp BATEC -- process} hold. Then, for any multi-index $|\bnu| \leq p$,
\begin{align*}
    \sup_{u \in \reals} \Big|\P \Big(\sup_{\bx \in \B} \big|Z^{(\bnu)}(\bx)\big| \leq u \Big)
                             - \P \Big(\sup_{\bx \in \B} \big|\widehat{Z}^{(\bnu)}(\bx) \big| \leq u \big| \bV \Big) \Big|
    \lessp \log(n)\mathfrak{R}_{\mathtt{cc}}^{1/2}.
\end{align*}
\end{lem}

\noindent \textbf{Proof of Lemma~\ref{sa-lem: feasible gaussian to infeasible gaussian}.} In this proof, we suppress the superscript $(\bnu)$ on $\Omega$ and $\widehat{\Omega}$. First, using Lemma~\ref{sa-lem: covariance}, we provide an upper bound between covariance functions of the feasible Gaussian process and the infeasible Gaussian process. Letting $\bPi_{\bx, \by} = \Omega_{\bx, \by} / \sqrt{\Omega_{\bx,\bx} \Omega_{\by,\by}}$ and $\widehat{\bPi}_{\bx,\by} = \widehat{\Omega}_{\bx,\by} / \sqrt{\widehat{\Omega}_{\bx,\bx} \widehat{\Omega}_{\by,\by}}$,
\begin{align*}
    \sup_{\bx,\by \in \B} \big|\bPi_{\bx, \by} - \widehat{\bPi}_{\bx,\by}\big|
    = \sup_{\bx, \by \in \B}
      \Big|\frac{\Omega_{\bx, \by} - \widehat{\Omega}_{\bx, \by}}{\sqrt{\Omega_{\bx,\bx} \Omega_{\by,\by}}}
         + \frac{\widehat{\Omega}_{\bx, \by}}{\sqrt{\widehat{\Omega}_{\bx,\bx} \widehat{\Omega}_{\by,\by}}} \Big(\sqrt{\frac{\widehat{\Omega}_{\bx,\bx} \widehat{\Omega}_{\by,\by}}{\Omega_{\bx,\bx} \Omega_{\by,\by}}} - 1\Big)\Big|
\end{align*}
From Lemma~\ref{sa-lem: covariance} and the fact that $|\sqrt{x} - \sqrt{y}| \leq (x \wedge y)^{-1/2} |x - y|/2$ for $x, y > 0$,
\begin{align*}
    \sup_{\bx,\by \in \B}
    \frac{\big|\big(\widehat{\Omega}_{\bx,\bx} \widehat{\Omega}_{\by,\by}\big)^{1/2} - \big(\Omega_{\bx,\bx} \Omega_{\by,\by}\big)^{1/2}\big|}
         {\big(\Omega_{\bx,\bx} \Omega_{\by,\by}\big)^{1/2}}
    \lesssim \frac{\sup_{\bx, \by \in \B} \big|\widehat{\Omega}_{\bx,\bx} \widehat{\Omega}_{\by,\by}  - \Omega_{\bx,\bx} \Omega_{\by,\by}\big|}
                  {\inf_{\bx,\by\in\B} \widehat{\Omega}_{\bx,\bx} \widehat{\Omega}_{\by,\by} \wedge \inf_{\bx,\by\in\B} \Omega_{\bx,\bx} \Omega_{\by,\by}}
    \lessp \bar{h}^{p+1} + \mathfrak{R}_{\mathtt{cov}}
\end{align*}
and
\begin{align*}
    \sup_{\bx, \by \in \B} \frac{\big|\Omega_{\bx,\by} - \widehat{\Omega}_{\bx,\by}\big|}{\sqrt{\Omega_{\bx,\bx} \Omega_{\by,\by}}}
    \lessp \bar{h}^{p+1} + \mathfrak{R}_{\mathtt{cov}}.
\end{align*}
Therefore, $\sup_{\bx, \by \in \B} |\bPi_{\bx, \by} - \widehat{\bPi}_{\bx,\by}| \lessp \mathfrak{R}_{\mathtt{cc}}$. Then, we bound the KS distance between the maximum of $Z^{(\bnu)}$ and $\widehat{Z}^{(\bnu)}$ on a $\delta_n$-net of $\B$, denoted by $\B_{\delta_n}$: for all $\bx \in \B$, there exists $\bz \in \B_{\delta_n}$ such that $\infnorm{\bx - \bz} \leq \delta_n$. Since $\B$ is compact, we can assume $M := \operatorname{Card}\left(\B_{\delta_n}\right) \lesssim \delta_n^{-d}$. Denote $\bZ^{\delta_n}$ and $\widehat{\bZ}^{\delta_n}$ as the restrictions of $Z^{(\bnu)}$ and $\widehat{Z}^{(\bnu)}$ to $\B_{\delta_n}$, respectively. Then, by \citet[Theorem 2.1]{Chernozhukov-Chetverikov-Kato-Koike_2022_AoS},
\begin{align*}
    \sup_{\by \in \reals^{M}} \big|\P(\bZ^{\delta_n}\leq \by) - \P(\widehat{\bZ}^{\delta_n} \leq \by | \bV) \big|
    \lesssim \log(M) \sup_{\bx,\by \in \B} \big|\bPi_{\bx, \by} - \widehat{\bPi}_{\bx,\by}\big|^{\frac{1}{2}}
    \lessp \log(M) \mathfrak{R}_{\mathtt{cc}}^{\frac{1}{2}},
\end{align*}
and hence
\begin{align*}
    \sup_{x \in \reals} \big|\P(\infnorm{\bZ^{\delta_n}} \leq x ) - \P(\infnorm{\widehat{\bZ}^{\delta_n}} \leq x | \bV)\big|
    &\leq \sup_{x \in \reals} \big|\P(- x\mathbf{1} \leq \bZ^{\delta_n} \leq x \mathbf{1}) - \P(- x\mathbf{1} \leq \widehat{\bZ}^{\delta_n}  \leq x \mathbf{1} | \bV)\big| \\
    & \lessp  \log(M) \mathfrak{R}_{\mathtt{cc}}^{\frac{1}{2}} =: \mathfrak{R}_M.
\end{align*}

Finally, we extend the comparison from the finite net to the full boundary. Let $\delta_n=n^{-s}$ for a fixed constant $s>0$ chosen large enough that
\begin{align*}
    \sqrt{\log n}\,|\bH|^{-1/2}\underline h^{-1}\delta_n=o(\mathfrak{R}_{\mathtt{cc}}),
    \qquad
    \sqrt{\log n}\,\underline h^{-1}\delta_n=o(\mathfrak{R}_{\mathtt{cc}}).
\end{align*}
Such an $s$ exists under the polynomial lower bound on $\bar h$ and the bounded ratio $\bar h/\underline h$. Put $b_n=\mathfrak{R}_{\mathtt{cc}}^{1/2}$ and define the modulus probabilities
\begin{align*}
    \Psi_{\delta_n}(b_n)
    &= \P \Big(\sup_{\substack{\bx,\by \in \B \\ \infnorm{\bx - \by} \leq \delta_n}} \big|Z^{(\bnu)}(\bx) - Z^{(\bnu)}(\by)\big| \geq b_n \Big),\\
    \widehat{\Psi}_{\delta_n}(b_n)
    &= \P \Big(\sup_{\substack{\bx,\by \in \B \\ \infnorm{\bx - \by} \leq \delta_n}} \big|\widehat{Z}^{(\bnu)}(\bx) - \widehat{Z}^{(\bnu)}(\by)\big| \geq b_n \Big| \bV \Big).
\end{align*}
The usual chaining argument for the local-polynomial Gaussian metrics gives
\begin{align*}
    \E \Big[\sup_{\substack{\bx,\by \in \B \\ \infnorm{\bx - \by}\leq \delta_n}} \big|Z^{(\bnu)}(\bx) - Z^{(\bnu)}(\by)\big|\Big]
    &\lesssim \sqrt{\log n}\,\underline h^{-1}\delta_n,\\
    \E \Big[\sup_{\substack{\bx,\by \in \B \\ \infnorm{\bx - \by}\leq \delta_n}} \big|\widehat{Z}^{(\bnu)}(\bx) - \widehat{Z}^{(\bnu)}(\by)\big| \Big| \bV \Big]
    &\lessp \sqrt{\log n}\,|\bH|^{-1/2}\underline h^{-1}\delta_n.
\end{align*}
The second display follows from Lemma~\ref{sa-lem: covariance}, Lemma~\ref{sa-lem: Gram Matrix}, and the Lipschitz bound for the feasible local-polynomial weights; the first display follows from the corresponding population bound. Therefore Markov's inequality implies
\begin{align*}
    \Psi_{\delta_n}(b_n)+\widehat{\Psi}_{\delta_n}(b_n)=o_{\P}(\mathfrak{R}_{\mathtt{cc}}^{1/2}).
\end{align*}
Moreover, the same entropy calculation gives
\[
    \E\Big[\sup_{\bx\in\B}|Z^{(\bnu)}(\bx)|\Big]\lesssim\sqrt{\log n},
    \qquad
    \E\Big[\sup_{\bx\in\B}|\widehat Z^{(\bnu)}(\bx)|\Big|\bV\Big]\lessp\sqrt{\log n}.
\]
For every $t>0$, the net approximation, the finite-dimensional comparison, and Gaussian anti-concentration give
\begin{align*}
    &\P \Big( \sup_{\bx \in \B} \big|Z^{(\bnu)}(\bx)\big| \leq t \Big)\\
    &\leq \P \Big( \sup_{\bx \in \B} \big|\widehat{Z}^{(\bnu)}(\bx)\big| \leq t \Big| \bV \Big)
       +4b_n\Big(\E\Big[\sup_{\bx\in\B}|\widehat Z^{(\bnu)}(\bx)|\Big|\bV\Big]+1\Big)
       +\Psi_{\delta_n}(b_n)+\widehat\Psi_{\delta_n}(b_n)+\mathfrak R_M,
\end{align*}
and the same argument with the roles of $Z^{(\bnu)}$ and $\widehat Z^{(\bnu)}$ reversed gives the corresponding lower bound. Since $\log M\lesssim \log n$, $\mathfrak R_M\lessp \log(n)\mathfrak{R}_{\mathtt{cc}}^{1/2}$; the anti-concentration term is $O_{\P}(\sqrt{\log n}\mathfrak{R}_{\mathtt{cc}}^{1/2})$; and the two modulus probabilities are $o_{\P}(\mathfrak{R}_{\mathtt{cc}}^{1/2})$. Combining the upper and lower bounds gives the stated result.
\qed


\bigskip
The proof of Theorem \ref{sa-thm: Confidence Bands: Sharp BATEC -- process} now follows directly from Lemma~\ref{sa-lem: infeasible gaussian to bahadur representation}, Lemma~\ref{sa-lem: t stats to bahadur} and Lemma~\ref{sa-lem: feasible gaussian to infeasible gaussian}. Furthermore, by definition of $\widehat{\CI}_{\alpha}^{(\bnu)}(\bx)$,
\begin{align*}
    \P\Big[\tau^{(\bnu)}(\bx) \in \widehat{\CI}_{\alpha}^{(\bnu)}(\bx), \text{ for all } \bx \in \B \Big]
    & = \P \Big[\sup_{\bx \in \B} \big|\widehat{\Tstat}^{(\bnu)}(\bx)\big| \leq \q_{\alpha}\Big] \\
    & = \P \Big[\sup_{\bx \in \B} \big|\widehat{Z}^{(\bnu)}(\bx)\big|\leq \q_{\alpha}\Big] + o(1) \\
    & = \E \Big[\P \Big[\sup_{\bx \in \B} \big|\widehat{Z}^{(\bnu)}(\bx)\big| \leq \q_{\alpha}\Big| \bV \Big]\Big] + o(1) \\
    & = 1 - \alpha + o(1),
\end{align*}
which completes the proof of the theorem.
\qed

%%%%%%%%%%%%%%%%%%%%%%%%%%%%%%%%%%%%
\subsection{Proof of Theorem~\ref{sa-thm: Gaussian Approximation Suprema -- Sharp}}
%%%%%%%%%%%%%%%%%%%%%%%%%%%%%%%%%%%%

We employ the following result established by \citet[Corollary 2.2]{Chernozhukov-Chetverikov-Kato_2014b_AoS}.

\begin{lem}[Gaussian Approximation of Suprema]\label{sa-lem: suprema coupling lemma}
    Let $X_1, \dots, X_n$ be i.i.d.\ random variables in a measurable space $(S, \mathcal{S})$ with law $P$, where $n \ge 3$. Let $\mathcal{F}$ be a class of measurable functions $f: S \to \mathbb{R}$ that is $P$-centered, i.e., $Pf = 0$ for all $f \in \mathcal{F}$. Define the empirical process $\mathbb{G}_n$ indexed by $\mathcal{F}$ as $\mathbb{G}_n f = \frac{1}{\sqrt{n}} \sum_{i=1}^n f(X_i)$ for $f \in \mathcal{F}$. Suppose that:
    \begin{enumerate}
        \item The class $\mathcal{F}$ is pointwise measurable, that is, it contains a countable subset $\mathcal{G}$ such that for every $f \in \mathcal{F}$ there exists a sequence $g_m \in \mathcal{G}$ with $g_m(x) \to f(x)$ for every $x \in S$.
        \item There exists a measurable envelope $F$ such that $|f| \le F$ for all $f \in \mathcal{F}$, and constants $A \ge e$ and $\upsilon \ge 1$ such that $\sup_Q N(\mathcal{F}, e_Q, \epsilon \|F\|_{Q,2}) \le (A/\epsilon)^{\upsilon}$ for all $0 < \epsilon \le 1$, where the supremum is taken over all finite discrete probability measures on $(S, \mathcal{S})$.
        \item For some $b \ge \sigma > 0$ and $q \in [4, \infty]$, we have $\sup_{f \in \mathcal{F}} P|f|^k \le \sigma^2 b^{k-2}$ for $k = 2, 3, 4$ and $\|F\|_{P,q} \le b$.
    \end{enumerate}
    Let $Z = \sup_{f \in \mathcal{F}} \mathbb{G}_n f$. Then for every $\gamma \in (0, 1)$, there exists, on a possibly extended probability space, a tight Gaussian process $(\mathbb{G}_P f)_{f \in \mathcal{F}}$ with mean zero, covariance function $\E[\mathbb{G}_P(f)\mathbb{G}_P(g)] = P(fg)$, and almost surely uniformly $e_P$-continuous trajectories, such that for $\tilde{Z} = \sup_{f \in \mathcal{F}} \mathbb{G}_P f$:
    $$ P\left\{ |Z - \tilde{Z}| > \frac{b K_n}{\gamma^{1/q} n^{1/2-1/q}} + \frac{(b \sigma)^{1/2} K_n^{3/4}}{\gamma^{1/2} n^{1/4}} + \frac{(b \sigma^2 K_n^2)^{1/3}}{\gamma^{1/3} n^{1/6}} \right\} \le C \left( \gamma + \frac{\log n}{n} \right) $$
    where $K_n = \upsilon(\log n \vee \log(A b/\sigma))$, and $C > 0$ is a constant that depends only on $q$.
\end{lem}


Consider $\bz_i=(\bX_i,\varepsilon_i)$, $1\leq i\leq n$, where $\varepsilon_i=Y_i-\E[Y_i|\bX_i]$. For each $\bx\in\B$, let $g_{\bx}$ be the boundary-indexed local-polynomial weight function defined before Theorem~\ref{sa-thm: Gaussian Strong Approximation: Sharp Tstat}, and consider the class
\begin{align*}
    \F=\{(\bu,u)\in\reals^d\times\reals\mapsto g_{\bx}(\bu)u:\bx\in\B\},
    \qquad
    \F_{\pm}=\F\cup(-\F).
\end{align*}
Then
\begin{align*}
    \sup_{\bx\in\B}|\overline{\Tstat}^{(\bnu)}(\bx)|
    =
    \sup_{f\in\F_{\pm}}
    \frac{1}{\sqrt n}\sum_{i=1}^n\{f(\bz_i)-\E[f(\bz_i)]\}.
\end{align*}
We apply Lemma~\ref{sa-lem: suprema coupling lemma} to this empirical process. First, the fixed indicators $\Indicator(\bu\in\A_t)$ are absorbed into $g_{\bx}$ and do not increase $L_2(Q)$ covering numbers beyond a constant factor, because multiplication by a fixed indicator is a contraction in $L_2(Q)$. Therefore, using the same local-polynomial entropy calculation as in the proofs of Lemmas~\ref{sa-lem: Gram Matrix} and \ref{sa-lem: Linear Approximation}, but without the total-variation step, $\{g_{\bx}:\bx\in\B\}$ is pointwise measurable and has an envelope $\ttM_{\G}\lesssim|\bH|^{-1/2}$ such that, for constants $\bc,\bc'>0$ depending only on the compact set $\X$,
\begin{align*}
    \sup_Q N(\G,\norm{\cdot}_{Q,2},\varepsilon(2\bc+1)^{d+1}\ttM_{\G})
    \leq 2\bc'\varepsilon^{-4d-4}+2,
    \qquad 0<\varepsilon<1,
\end{align*}
where the supremum is taken over all finite discrete measures. Since the residual class is the singleton $\{\operatorname{Id}\}$, the class $\F_{\pm}$ satisfies the same VC-type entropy bound up to a universal constant; centering the functions changes only the envelope by a constant factor.

For the moment condition, take $q=2+v\geq4$ in Lemma~\ref{sa-lem: suprema coupling lemma}. Then, Assumption~\ref{sa-assump: Sharp DGP} implies $\sup_{\bx\in\X}\E[|\varepsilon_i|^{2+v}|\bX_i=\bx]<\infty$. The envelope $F(\bu,u)=\ttM_{\G}|u|$ satisfies $\|F\|_{\P,2+v}\lesssim|\bH|^{-1/2}$. The same local-polynomial moment calculation used in the proof of Lemma~\ref{sa-lem: Linear Approximation} gives, for $\sigma=1$ and $b=|\bH|^{-1/2}$, that $\sup_{f\in\F_{\pm}}\P|f|^k\leq\sigma^2b^{k-2}$ for $k=2,3,4$ and $\|F\|_{\P,2+v}\leq b$. Taking $\gamma=(\log n)^{-1}$ in Lemma~\ref{sa-lem: suprema coupling lemma} yields
\begin{align*}
    \P\left\{ \left|\sup_{\bx\in\B}|\overline{\Tstat}^{(\bnu)}(\bx)|-\sup_{\bx\in\B}|Z^{(\bnu)}(\bx)|\right|>
    \frac{\log n}{(n|\bH|)^{1/6}}
    +\frac{(\log n)^{5/4}}{(n|\bH|)^{1/4}}
    +\frac{(\log n)^{3/2}}{(n^{\frac{v}{2+v}}|\bH|)^{1/2}}
    \right\}
    \leq C\left(\frac{1}{\log n}+\frac{\log n}{n}\right),
\end{align*}
for a constant $C>0$. This proves the theorem.
\qed

%%%%%%%%%%%%%%%%%%%%%%%%%%%%%%%%%%%%
\subsection{Proof of Theorem~\ref{sa-thm: Confidence Bands: Sharp BATEC -- Suprema}}
%%%%%%%%%%%%%%%%%%%%%%%%%%%%%%%%%%%%
Define
\begin{equation*}
    \mathfrak{R}_{\mathtt{sup}}
    =
    \frac{\log n}{(n|\bH|)^{1/6}}
    +\frac{(\log n)^{5/4}}{(n|\bH|)^{1/4}}
    +\frac{(\log n)^{3/2}}{(n^{\frac{v}{2+v}}|\bH|)^{1/2}}.
\end{equation*}
In this proof, let $\mathfrak{R}_{\mathtt{lin}}$ denote the stochastic-linearization remainder defined at the beginning of the proof of Theorem~\ref{sa-thm: Confidence Bands: Sharp BATEC -- process}.
The entropy calculation in the proof of Theorem~\ref{sa-thm: Gaussian Approximation Suprema -- Sharp} gives
\[
    \E\left[\sup_{\bx\in\B}|Z^{(\bnu)}(\bx)|\right]\lesssim\sqrt{\log n}.
\]
The side condition $(\log n)^9/(n^{v/(2+v)}|\bH|)=o(1)$ implies $\mathfrak{R}_{\mathtt{sup}}\sqrt{\log n}=o(1)$. Let
\[
    M_{\overline T}=\sup_{\bx\in\B}|\overline{\Tstat}^{(\bnu)}(\bx)|,
    \qquad
    M_Z=\sup_{\bx\in\B}|Z^{(\bnu)}(\bx)|.
\]
Applying Lemma~\ref{sa-lem: generic suprema comparison} with $\mathbb{T}_1=M_{\overline T}$ and $\mathbb{T}_2=M_Z$, viewed as singleton-indexed processes, and $\eta=\mathfrak{R}_{\mathtt{sup}}$, gives
\[
    \sup_{u\in\reals}\left|\P(M_{\overline T}\le u)-\P(M_Z\le u)\right|
    \le
    \sup_{u\in\reals}\P(|M_Z-u|\le \mathfrak{R}_{\mathtt{sup}})
    +\P(|M_{\overline T}-M_Z|>\mathfrak{R}_{\mathtt{sup}}).
\]
The first term is $O(\mathfrak{R}_{\mathtt{sup}}\sqrt{\log n})=o(1)$ by Gaussian anti-concentration, and the second term is $o(1)$ by Theorem~\ref{sa-thm: Gaussian Approximation Suprema -- Sharp}. Therefore,
\begin{equation}\label{sa-eq: t-to-z sharp suprema}
    \sup_{u\in\reals}
    \left|
    \P\left(\sup_{\bx\in\B}|\overline{\Tstat}^{(\bnu)}(\bx)|\leq u\right)
    -
    \P\left(\sup_{\bx\in\B}|Z^{(\bnu)}(\bx)|\leq u\right)
    \right|
    =o(1).
\end{equation}

Next, by Theorem~\ref{sa-thm: Stochastic Linearization: Sharp BATEC}, the stochastic-linearization error is $O_{\P}(\mathfrak{R}_{\mathtt{lin}})$. The polynomial lower bound on $\bar h$, together with $(\log n)^9/(n^{v/(2+v)}|\bH|)=o(1)$ and $\bar h^{p+1}\sqrt{n|\bH|\log n}=o(1)$, implies $\mathfrak{R}_{\mathtt{lin}}\sqrt{\log n}=o(1)$. Choose a deterministic sequence $a_n$ such that $\mathfrak{R}_{\mathtt{lin}}=o(a_n)$ and $a_n\sqrt{\log n}=o(1)$. Applying Lemma~\ref{sa-lem: generic suprema comparison} with $\mathbb{T}_1=\widehat{\Tstat}^{(\bnu)}$, $\mathbb{T}_2=\overline{\Tstat}^{(\bnu)}$, and $\eta=a_n$, and then bounding the anti-concentration of $M_{\overline T}$ through Equation~\eqref{sa-eq: t-to-z sharp suprema}, gives
\begin{equation}\label{sa-eq: t-to-t sharp suprema}
    \sup_{u\in\reals}
    \left|
    \P\left(\sup_{\bx\in\B}|\widehat{\Tstat}^{(\bnu)}(\bx)|\leq u\right)
    -
    \P\left(\sup_{\bx\in\B}|\overline{\Tstat}^{(\bnu)}(\bx)|\leq u\right)
    \right|
    =o(1).
\end{equation}

Finally, set
\[
    \mathfrak{R}_{\mathtt{cc}}
    =
    \sqrt{\frac{\log n}{n|\bH|}}
    +\frac{\log n}{n^{\frac{v}{2+v}}|\bH|}
    +\bar h^{p+1}.
\]
This is the same covariance-comparison rate used in Lemma~\ref{sa-lem: feasible gaussian to infeasible gaussian}. The covariance-comparison and net-modulus arguments in that lemma do not use the De Giorgi perimeter condition. Under the present side conditions, they give
\begin{equation}\label{sa-eq: z-to-z sharp suprema}
    \sup_{u\in\reals}
    \left|
    \P\left(\sup_{\bx\in\B}|Z^{(\bnu)}(\bx)|\leq u\right)
    -
    \P\left(\sup_{\bx\in\B}|\widehat Z^{(\bnu)}(\bx)|\leq u\;\middle|\;\bV\right)
    \right|
    \lessp \log(n)\mathfrak{R}_{\mathtt{cc}}^{1/2}
    =o_{\P}(1).
\end{equation}
The first conclusion follows from the triangle inequality applied to Equations~\eqref{sa-eq: t-to-z sharp suprema}, \eqref{sa-eq: t-to-t sharp suprema}, and \eqref{sa-eq: z-to-z sharp suprema}.

For the coverage statement, by definition of $\widehat{\CI}_{\alpha}^{(\bnu)}(\bx)$ and the first conclusion,
\begin{align*}
    \P\Big[\tau^{(\bnu)}(\bx)\in\widehat{\CI}_{\alpha}^{(\bnu)}(\bx),\text{ for all }\bx\in\B\Big]
    &=
    \P\left(\sup_{\bx\in\B}|\widehat{\Tstat}^{(\bnu)}(\bx)|\leq\q_{\alpha}\right)\\
    &=
    \E\left[
    \P\left(\sup_{\bx\in\B}|\widehat Z^{(\bnu)}(\bx)|\leq\q_{\alpha}\;\middle|\;\bV\right)
    \right]+o(1)\\
    &=1-\alpha+o(1),
\end{align*}
which completes the proof.
\qed

%%%%%%%%%%%%%%%%%%%%%%%%%%%%%%%%%%%%
\subsection{Proof of Lemma~\ref{sa-lem: Bias: Sharp WBATE}}
%%%%%%%%%%%%%%%%%%%%%%%%%%%%%%%%%%%%

By Lemma~\ref{sa-lem: bias},
\begin{align*}
    \sup_{\bx\in\B}
    \big|\E[\widehat{\tau}^{(\bnu)}(\bx)|\bX]-\tau^{(\bnu)}(\bx)-B_{\bx}^{(\bnu)}(\bh)\big|
    = o_{\P}(\bar{h}^{p+1}\bh^{-\bnu}).
\end{align*}
Integrating this uniform expansion against $w(\bx)\dif\Haus^{d-1}(\bx)$ gives the displayed WBATE bias expansion because $\int_{\B}|w(\bx)|\dif\Haus^{d-1}(\bx)<\infty$. When $\bh=\bc h$, the pointwise leading bias is $h^{p+1-|\bnu|}\widetilde{B}_{\bx}^{(\bnu)}(\bc)$; integrating this leading term over $\B$ gives $\widetilde{B}_{\mathtt{WBATE}}^{(\bnu)}(\bc)$, and the integrated remainder remains $o(h^{p+1-|\bnu|})$ by the same integrability argument.
\qed

%%%%%%%%%%%%%%%%%%%%%%%%%%%%%%%%%%%%
\subsection{Proof of Lemma~\ref{sa-lem: Variance: Sharp WBATE}}
%%%%%%%%%%%%%%%%%%%%%%%%%%%%%%%%%%%%

Since $\Var[\widehat{\tau}_{\mathtt{WBATE}}^{(\bnu)}|\bX] = \Var[\int_{\B}\widehat{\mu}^{(\bnu)}_0(\bb) w(\bb) \dif\Haus^{d-1}(\bb)|\bX] + \Var[\int_{\B}\widehat{\mu}^{(\bnu)}_1(\bb) w(\bb) \dif\Haus^{d-1}(\bb)|\bX]$, it is enough to consider only one treatment assignment group $t \in \{0,1\}$. The conditional cross-covariance between the two side-specific integrals is zero because the side indicators are disjoint observation by observation, and observations are conditionally independent given $\bX$. In addition,
\begin{align*}
    \Var \Big[\int_{\B}\widehat{\mu}^{(\bnu)}_t(\bb) w(\bb) \dif\Haus^{d-1}(\bb) \Big| \bX \Big]
    = \int_{\B} \int_{\B} \Cov \big[\widehat{\mu}^{(\bnu)}_t(\bb_1), \widehat{\mu}^{(\bnu)}_t(\bb_2) \big| \bX \big] w(\bb_1) w(\bb_2) \dif\Haus^{d-1}(\bb_1) \dif\Haus^{d-1}(\bb_2)
\end{align*}
and
\begin{align*}
    \Omega_{t,\mathtt{WBATE}}^{(\bnu)}
    = \int_{\B} \int_{\B} \Omega_{t,\bb_1, \bb_2}^{(\bnu)} w(\bb_1) w(\bb_2) \dif\Haus^{d-1}(\bb_1) \dif\Haus^{d-1}(\bb_2).
\end{align*}

Proceeding as in the proof of Lemma \ref{sa-lem: covariance}, we have
\begin{align*}
    \sup_{\bb_1, \bb_2 \in \B} \big| \Cov[\widehat{\mu}^{(\bnu)}_t(\bb_1), \widehat{\mu}^{(\bnu)}_t(\bb_2)| \bX ] - \Omega_{t, \bb_1, \bb_2}^{(\bnu)} \big|
    \lessp \frac{1}{n|\bH|\bh^{2\bnu}}\sqrt{\frac{\log(1/\bar{h})}{n|\bH|}}.
\end{align*}

Since $K$ is supported on a compact set, let $R\in(0,\infty)$ denote the diameter of the support, and define the ``effective domain'' $\mathcal{E}(\bH) = \{(\bx, \by) \in \B \times \B: \lVert \bx - \by \rVert \leq \bar{h} R\}$. Let $\nu_{\B,2}=\Haus^{d-1}\times\Haus^{d-1}$ on $\B\times\B$. By \eqref{sa-eq: local boundary measure regularity}, $\nu_{\B,2}(\mathcal{E}(\bH)) \lesssim \bar{h}^{d-1}$. Therefore,
\begin{align*}
    &\Big| \Var \Big[\int_{\B}\widehat{\mu}^{(\bnu)}_t(\bb) w(\bb) \dif\Haus^{d-1}(\bb) \Big| \bX \Big] - \Omega_{t,\mathtt{WBATE}}^{(\bnu)} \Big|\\
    &= \Big| \int_{\B} \int_{\B} \big( \Cov[\widehat{\mu}^{(\bnu)}_t(\bb_1), \widehat{\mu}^{(\bnu)}_t(\bb_2)| \bX] - \Omega_{t,\bb_1, \bb_2}^{(\bnu)} \big) w(\bb_1) w(\bb_2) \dif\Haus^{d-1}(\bb_1) \dif\Haus^{d-1}(\bb_2)\Big|\\
    & \lesssim \sup_{\bb_1, \bb_2 \in \B} \big| \Cov[\widehat{\mu}^{(\bnu)}_t(\bb_1), \widehat{\mu}^{(\bnu)}_t(\bb_2)|\bX] -
    \Omega_{t, \bb_1, \bb_2}^{(\bnu)}\big| \int_{\B} \int_{\B} \Indicator((\bb_1, \bb_2) \in \mathcal{E}(\bH)) w(\bb_1) w(\bb_2) \dif\Haus^{d-1}(\bb_1) \dif\Haus^{d-1}(\bb_2) \\
    &= o_{\P}\Big(\frac{\bar{h}^{d-1}}{n|\bH|\bh^{2\bnu}}\Big),
\end{align*}
because $\frac{\log(1/\bar{h})}{n|\bH|} = o(1)$ and $\bar{h}=o(1)$. This proves the first claim. Next, for each $t\in\{0,1\}$,
\begin{align*}
    \Omega_{t,\mathtt{WBATE}}^{(\bnu)}
    & = \int_{\B} \int_{\B} \Omega_{t,\bb_1, \bb_2}^{(\bnu)} w(\bb_1) w(\bb_2) \dif\Haus^{d-1}(\bb_1) \dif\Haus^{d-1}(\bb_2)\\
    & \leq \sup_{\bb_1, \bb_2 \in \B} \big|\Omega_{t,\bb_1, \bb_2}^{(\bnu)}\big| \int_{\B} \int_{\B} \Indicator((\bb_1, \bb_2) \in \mathcal{E}(\bH)) w(\bb_1) w(\bb_2) \dif\Haus^{d-1}(\bb_1) \dif\Haus^{d-1}(\bb_2) \\
    & \lesssim (n|\bH|\bh^{2\bnu})^{-1} \nu_{\B,2}(\mathcal{E}(\bH)) \lesssim \frac{\bar{h}^{d-1}}{n|\bH|\bh^{2\bnu}},
\end{align*}
which verifies the upper bound after summing over $t$. The lower bound follows from the aggregate variance nondegeneracy condition \eqref{sa-eq: sharp WBATE aggregate nondegeneracy}. This condition is needed because the covariance kernel can contain negative regions after local-polynomial weighting, so pointwise variance lower bounds alone do not rule out cancellation in the weighted boundary integral.

The third and final claim of the lemma follows from Lemma~\ref{sa-lem: covariance} and the same analysis as above.
\qed

%%%%%%%%%%%%%%%%%%%%%%%%%%%%%%%%%%%%
\subsection{Proof of Theorem~\ref{sa-thm: MSE Expansion: Sharp WBATE}}
%%%%%%%%%%%%%%%%%%%%%%%%%%%%%%%%%%%%

Lemma~\ref{sa-lem: Bias: Sharp WBATE} gives
\begin{align*}
    \E[\widehat{\tau}_{\mathtt{WBATE}}^{(\bnu)}|\bX]-\tau_{\mathtt{WBATE}}^{(\bnu)}
    = B_{\mathtt{WBATE}}^{(\bnu)}(\bh)+o_{\P}(\bar{h}^{p+1}\bh^{-\bnu}),
\end{align*}
so the squared conditional bias equals $(B_{\mathtt{WBATE}}^{(\bnu)}(\bh))^2+o_{\P}(\bar{h}^{2p+2}\bh^{-2\bnu})$. Lemma~\ref{sa-lem: Variance: Sharp WBATE} gives the conditional variance expansion with remainder $o_{\P}(\bar{h}^{d-1}/(n|\bH|\bh^{2\bnu}))$. Adding the squared-bias and variance expansions gives the stated MSE expansion.
\qed

%%%%%%%%%%%%%%%%%%%%%%%%%%%%%%%%%%%%
\subsection{Proof of Theorem~\ref{sa-thm: Asymptotic Normality: Sharp WBATE}}
%%%%%%%%%%%%%%%%%%%%%%%%%%%%%%%%%%%%

Since $\widehat{\tau}_{\mathtt{WBATE}}^{(\bnu)} - \tau_{\mathtt{WBATE}}^{(\bnu)} = (\widehat{\mu}_{1,\mathtt{WBATE}}^{(\bnu)} - \mu_{1,\mathtt{WBATE}}^{(\bnu)}) - (\widehat{\mu}_{0,\mathtt{WBATE}}^{(\bnu)} - \mu_{0,\mathtt{WBATE}}^{(\bnu)})$, it is enough to start with only one treatment assignment group $t \in \{0,1\}$. The same argument applies to both side-specific components, and the two components combine by summing their variances because the treatment- and control-side kernel-weighted summands have disjoint side indicators observation by observation. Furthermore,
\begin{align*}
    \widehat{\mu}_{t,\mathtt{WBATE}}^{(\bnu)} - \mu_{t,\mathtt{WBATE}}^{(\bnu)}
    &= \int_{\B} (\widehat{\mu}^{(\bnu)}_{t}(\bb) - \mu^{(\bnu)}_t(\bb) ) w(\bb) \dif\Haus^{d-1}(\bb)\\
    &= \int_{\B} \bnu!\be_{\bnu}^{\top}\bS^{-1} \bGamma_{t,\bb}^{-1} \bQ_{t,\bb} w(\bb) \dif\Haus^{d-1}(\bb)\\
    &\qquad  + \int_{\B} \bnu!\be_{\bnu}^{\top}\bS^{-1}(\widehat{\bGamma}_{t,\bb}^{-1} - \bGamma_{t,\bb}^{-1}) \bQ_{t,\bb} w(\bb) \dif\Haus^{d-1}(\bb)
      + O_\P(\bar{h}^{p+1}\bh^{-\bnu})
\end{align*}
using Lemma~\ref{sa-lem: bias} to bound the approximation error.
The side condition $n|\bH|\bar{h}^{2p+2}/\bar{h}^{d-1}=o(1)$ implies that the bias term is $o_{\P}((\Omega_{\mathtt{WBATE}}^{(\bnu)})^{1/2})$.

For the second integral, let
\begin{align*}
    \overline{\bSigma}_{t,\bx_1,\bx_2}
    = \En \Big[\br_p(\bH^{-1}(\bX_i - \bx_1)) \br_p(\bH^{-1}(\bX_i - \bx_2))^{\top} |\bH| K_{\bH}(\bX_i - \bx_1) K_{\bH}(\bX_i - \bx_2) \sigma_t^2(\bX_i) \Indicator(\bX_i \in \A_t)\Big],
\end{align*}
and since $\overline{\bSigma}_{t,\bb_1, \bb_2} = 0$ if $\lVert \bb_1-\bb_2\rVert$ is larger than a constant multiple of $\bar{h}$, let
\[
    \mathcal{E}(\bH)=\{(\bb_1,\bb_2)\in\B\times\B:\lVert\bb_1-\bb_2\rVert\lesssim\bar h\}.
\]
By \eqref{sa-eq: local boundary measure regularity}, $\nu_{\B,2}(\mathcal{E}(\bH))\lesssim \bar h^{d-1}$. Then
\begin{align*}
    &\E \bigg[ \bigg(\int_{\B} \bnu!\be_{\bnu}^{\top}\bS^{-1}(\widehat{\bGamma}_{t,\bb}^{-1} - \bGamma_{t,\bb}^{-1}) \bQ_{t,\bb} w(\bb) \dif\Haus^{d-1}(\bb) \bigg)^2\bigg|\bX\bigg]\\
    &= \int_{\B} \int_{\B} (\bnu!)^2\be_{\bnu}^{\top}\bS^{-1}(\widehat{\bGamma}_{t,\bb_1}^{-1} - \bGamma_{t,\bb_1}^{-1}) (n|\bH|)^{-1} \overline{\bSigma}_{t,\bb_1, \bb_2} (\widehat{\bGamma}_{t,\bb_2}^{-1} - \bGamma_{t,\bb_2}^{-1}) \bS^{-1}\be_{\bnu} w(\bb_1) w(\bb_2) \dif\Haus^{d-1}(\bb_1) \dif\Haus^{d-1}(\bb_2),\\
    &\leq \bh^{-2\bnu}\sup_{\bb \in \B} \|\widehat{\bGamma}_{t,\bb}^{-1} - \bGamma_{t,\bb}^{-1}\|^2 \sup_{\bb_1, \bb_2 \in \B} \lVert \overline{\bSigma}_{t,\bb_1, \bb_2}  \rVert \sup_{\bb \in \B} |w(\bb)| (n|\bH|)^{-1} \nu_{\B,2}(\mathcal{E}(\bH)) \\
    &= o_\P\Big(\frac{\bar{h}^{d-1}}{n|\bH|\bh^{2\bnu}}\Big),
\end{align*}
and hence $\int_{\B} \bnu!\be_{\bnu}^{\top}\bS^{-1}(\widehat{\bGamma}_{t,\bb}^{-1} - \bGamma_{t,\bb}^{-1}) \bQ_{t,\bb} w(\bb) \dif\Haus^{d-1}(\bb) = o_\P\big((\bar{h}^{d-1}/(n|\bH|\bh^{2\bnu}))^{1/2}\big)$.

Next, combine the two side-specific expansions with signs $s_1=1$ and $s_0=-1$. Define the signed linearized statistic
\begin{align*}
    \overline{\Tstat}_{\mathtt{WBATE}}^{(\bnu)}
    =
    (\Omega_{\mathtt{WBATE}}^{(\bnu)})^{-1/2}
    \sum_{t\in\{0,1\}}s_t
    \int_{\B} \bnu!\be_{\bnu}^{\top}\bS^{-1}\bGamma_{t,\bb}^{-1} \bQ_{t,\bb} w(\bb) \dif\Haus^{d-1}(\bb).
\end{align*}
Using Lemma~\ref{sa-lem: Variance: Sharp WBATE} and the previous side-specific bounds,
\begin{align*}
    \widehat{\Tstat}_{\mathtt{WBATE}}^{(\bnu)} - \overline{\Tstat}_{\mathtt{WBATE}}^{(\bnu)}
    = \big((\widehat{\Omega}_{\mathtt{WBATE}}^{(\bnu)})^{-1/2} - (\Omega_{\mathtt{WBATE}}^{(\bnu)})^{-1/2}\big)
      \sum_{t\in\{0,1\}}s_t
      \int_{\B} \bnu!\be_{\bnu}^{\top}\bS^{-1}\bGamma_{t,\bb}^{-1} \bQ_{t,\bb} w(\bb) \dif\Haus^{d-1}(\bb) + o_{\P}(1)
    = o_{\P}(1),
\end{align*}

Finally, apply the Berry--Esseen lemma to this signed treatment-control statistic. Write $\overline{\Tstat}_{\mathtt{WBATE}}^{(\bnu)} = \sum_{i = 1}^n \chi_i$, where
\begin{align*}
    \chi_i
    &= n^{-1}(\Omega_{\mathtt{WBATE}}^{(\bnu)})^{-1/2}
       \sum_{t\in\{0,1\}}s_t
       \int_{\B} \bnu!\be_{\bnu}^{\top}\bS^{-1}\bGamma_{t,\bb}^{-1}
       \br_p(\bH^{-1}(\bX_i - \bb)) K_{\bH}(\bX_i - \bb) \\
    &\qquad\qquad\qquad\qquad\qquad\qquad\qquad
       \cdot \Indicator(\bX_i \in \A_t) \varepsilon_i w(\bb) \dif\Haus^{d-1}(\bb).
\end{align*}
This satisfies $\E[\chi_i] = 0$. The definition of $\Omega_{\mathtt{WBATE}}^{(\bnu)}$ implies that $\sum_{i = 1}^n\Var[\chi_i] = 1$. Since the indicators $\Indicator(\bX_i\in\A_t)$ are disjoint observation by observation, the Lyapunov bound is the same as the side-specific bound up to a fixed constant. Hence, it remains to bound
\begin{align*}
    &\sum_{i = 1}^n \E[|\chi_i|^3]\\
    &\lesssim n^{-3} (\Omega_{\mathtt{WBATE}}^{(\bnu)})^{-3/2} \sum_{i = 1}^n\sum_{t\in\{0,1\}} \E \bigg[\bigg| \int_{\B} \bnu!\be_{\bnu}^{\top}\bS^{-1}\bGamma_{t,\bb}^{-1} \br_p(\bH^{-1}(\bX_i - \bb)) K_{\bH}(\bX_i - \bb) \Indicator(\bX_i \in \A_t) \varepsilon_i w(\bb) \dif\Haus^{d-1}(\bb) \bigg|^3\bigg].
\end{align*}

Let $R$ denote the diameter of the (compact) support of $K$, and define $\mathcal{E}_3(\bH) = \{(\bb_1, \bb_2, \bb_3) \in \B^3: \lVert \bb_i - \bb_j\rVert \leq R\bar{h}, i,j = 1,2,3\}$. Let $\nu_{\B,3}$ denote the product measure $\Haus^{d-1}\times\Haus^{d-1}\times\Haus^{d-1}$ on $\B^3$. Iterating \eqref{sa-eq: local boundary measure regularity} gives $\nu_{\B,3}(\mathcal{E}_3(\bH)) \lesssim \bar{h}^{2d-2}$. Then,
\begin{align*}
    & \E \bigg[\bigg| \int_{\B} \bnu!\be_{\bnu}^{\top}\bS^{-1}\bGamma_{t,\bb}^{-1} \br_p(\bH^{-1}(\bX_i - \bb)) K_{\bH}(\bX_i - \bb) \Indicator(\bX_i \in \A_t) \varepsilon_i w(\bb) \dif\Haus^{d-1}(\bb) \bigg|^3\bigg] \\
    & \leq \E \bigg[\int_{\bb_1 \in \B} \int_{\bb_2 \in \B} \int_{\bb_3 \in \B} |G(\bb_1, \bb_2, \bb_3)| w(\bb_1) w(\bb_2) w(\bb_3) \dif\Haus^{d-1}(\bb_1) \dif\Haus^{d-1}(\bb_2) \dif\Haus^{d-1}(\bb_3)\bigg] \\
    & \lesssim \nu_{\B,3}(\mathcal{E}_3(\bH))  \sup_{\bb_1, \bb_2, \bb_3 \in \B} \E[|G(\bb_1, \bb_2, \bb_3)|],
\end{align*}
where $G(\bb_1, \bb_2, \bb_3) = g(\bX_i, \varepsilon_i, \bb_1) g(\bX_i, \varepsilon_i, \bb_2) g(\bX_i, \varepsilon_i, \bb_3)$ with
\begin{align*}
    g(\bX_i, \varepsilon_i, \bb)
    = \bnu!\be_{\bnu}^{\top}\bS^{-1}\bGamma_{t,\bb}^{-1} \br_p(\bH^{-1}(\bX_i - \bb)) K_{\bH}(\bX_i - \bb) \Indicator(\bX_i \in \A_t) \varepsilon_i.
\end{align*}

Proceeding as in the proof of Lemma~\ref{sa-lem: Gram Matrix} and Lemma~\ref{sa-lem: Variance: Sharp WBATE}, it can be shown that
\begin{align*}
    \sup_{\bb_1, \bb_2, \bb_3 \in \B} \E[|G(\bb_1, \bb_2, \bb_3)|] \lesssim \bh^{-3\bnu}|\bH|^{-2}
\end{align*}
provided that $\frac{\log(1/\bar{h})}{n|\bH|} = o(1)$. Therefore, together with the rate of $\Omega_{\mathtt{WBATE}}^{(\bnu)}$ from Lemma~\ref{sa-lem: Variance: Sharp WBATE}, we have $\sum_{i = 1}^n \E[|\chi_i|^3] \lesssim (\bar{h}^{d-1}/(n|\bH|))^{1/2} \asymp (n\bar{h})^{-1/2}$, and the result follows.
\qed

\subsection{Proof of Theorem~\ref{sa-thm: Convergence Rate: Sharp LBATE}}

Follows by Theorem \ref{sa-thm: Convergence Rates: Sharp BATEC} after noting that
\begin{align*}
    |\widehat{\tau}_{\mathtt{LBATE}}^{(\bnu)} - \tau_{\mathtt{LBATE}}^{(\bnu)}|
    \leq \sup_{\bx\in\B} |\widehat{\tau}^{(\bnu)}(\bx) - \tau^{(\bnu)}(\bx)|.
\end{align*}
\qed

%%%%%%%%%%%%%%%%%%%%%%%%%%%%%%%%%%%%
\subsection{Proof of Theorem~\ref{sa-thm: Confidence Intervals: Sharp LBATE}}
%%%%%%%%%%%%%%%%%%%%%%%%%%%%%%%%%%%%

Consider the event $E = \Big\{\sup_{\bb \in \B} \frac{|\widehat{\tau}^{(\bnu)}(\bb) - \tau^{(\bnu)}(\bb)|}{(\widehat{\Omega}_{\bb, \bb}^{(\bnu)})^{1/2}} \leq \q_{\alpha}\Big\}$. Theorem~\ref{sa-thm: Confidence Bands: Sharp BATEC -- Suprema} implies that $\P(E) = 1 - \alpha + o(1)$. On the event $E$, we also have
\begin{align*}
    \widehat{\tau}^{(\bnu)}(\bb) - \q_{\alpha}(\widehat{\Omega}_{\bb,\bb}^{(\bnu)})^{1/2} \leq \tau^{(\bnu)}(\bb)
    \leq \widehat{\tau}^{(\bnu)}(\bb) + \q_{\alpha}(\widehat{\Omega}_{\bb,\bb}^{(\bnu)})^{1/2}, \qquad \forall \; \bb \in \B,
\end{align*}
which implies
\begin{align*}
    \sup_{\bb \in \B}  \widehat{\tau}^{(\bnu)}(\bb) - \q_{\alpha}(\widehat{\Omega}_{\bb,\bb}^{(\bnu)})^{1/2}
    \leq \sup_{\bb \in \B} \tau^{(\bnu)}(\bb)
    \leq \sup_{\bb \in \B} \widehat{\tau}^{(\bnu)}(\bb) + \q_{\alpha}(\widehat{\Omega}_{\bb,\bb}^{(\bnu)})^{1/2}.
\end{align*}
The stated result then follows.
\qed

%%%%%%%%%%%%%%%%%%%%%%%%%%%%%%%%%%%%
\subsection{Proof of Theorem~\ref{sa-thm: Convergence Rates: Fuzzy BATEC}}
%%%%%%%%%%%%%%%%%%%%%%%%%%%%%%%%%%%%

Theorem \ref{sa-thm: Convergence Rates: Sharp BATEC} applies directly to the reduced-form estimator $\widehat{\tau}^{(\bnu)}_Y$. For the treatment-status estimator $\widehat{\tau}^{(\bnu)}_W$, the same convergence-rate proof applies with $W$ in place of $Y$: Assumption~\ref{sa-assump: Fuzzy DGP}(ii) gives the required smoothness of $\mu_{W,t}$, and $W_i(t)\in\{0,1\}$ gives the required moment bounds. The variance lower bound in Assumption~\ref{sa-assump: Sharp DGP}(iv) is not used for these consistency rates. The strengthened bandwidth condition makes the uniform reduced-form convergence rate vanish, so $\inf_{\bx\in\B}|\widehat{\tau}_W^{(\bnu)}(\bx)|$ is bounded away from zero with probability approaching one. Since $\inf_{\bx\in\B}|\tau_W^{(\bnu)}(\bx)|>0$, the displayed rates follow from the exact expansion
\begin{align*}
    \widehat{\zeta}^{(\bnu)}(\bx) - \zeta^{(\bnu)}(\bx)
    = \frac{\widehat{\tau}^{(\bnu)}_Y(\bx)-\tau^{(\bnu)}_Y(\bx)}{\tau_W^{(\bnu)}(\bx)}
      - \frac{\tau_Y^{(\bnu)}(\bx)}{\tau_W^{(\bnu)}(\bx)\widehat{\tau}_W^{(\bnu)}(\bx)}
        \big(\widehat{\tau}^{(\bnu)}_W(\bx)-\tau^{(\bnu)}_W(\bx)\big),
\end{align*}
and the same argument uniformly over $\bx\in\B$.
\qed

%%%%%%%%%%%%%%%%%%%%%%%%%%%%%%%%%%%%


\subsection{Proof of Theorem~\ref{sa-thm: MSE Expansions: Fuzzy BATEC}}

\begin{lem}[Covariance for Fuzzy Designs]\label{sa-lem: covariance for fuzzy designs}
    Suppose Assumptions \ref{sa-assump: Sharp DGP}, \ref{sa-assump: Boundary and Kernel} and \ref{sa-assump: Fuzzy DGP} hold. If $\frac{\log(1/\bar{h})}{n^{\frac{v}{2+v}} |\bH|} = o(1)$ and $\bar{h}=o(1)$, then for $t = 0,1$, and $A_1, A_2 \in \{Y, W\}$, we have
    \begin{align*}
        \sup_{\bx_1,\bx_2 \in \B} \big\|\widehat{\bSigma}_{A_1,A_2,t,\bx_1,\bx_2} - \bSigma_{A_1,A_2,t,\bx_1,\bx_2}\big\|
        \lessp \sqrt{\frac{\log(1/\bar{h})}{n |\bH|}} + \frac{\log(1/\bar{h})}{n^{\frac{v}{2+v}}|\bH|} + \bar{h}^{p+1},
    \end{align*}
    \begin{align*}
        \sup_{\bx_1,\bx_2 \in \B} \big|\widehat{\Omega}_{A_1, A_2, \bx_1,\bx_2}^{(\bnu)} - \Omega_{A_1, A_2, \bx_1,\bx_2}^{(\bnu)}\big|
        \lessp (n |\bH|\bh^{2\bnu})^{-1}\bigg(\sqrt{\frac{\log(1/\bar{h})}{n |\bH|}} + \frac{\log(1/\bar{h})}{n^{\frac{v}{2+v}}|\bH|}+ \bar{h}^{p+1}\bigg),
    \end{align*}
\end{lem}

\begin{proof}
The argument parallels the proof of Lemma~\ref{sa-lem: covariance}.
The only additional moment input is for the treatment-status residuals and cross-products: $W_i$ is binary, hence $\varepsilon_{W,i}$ is bounded, while Assumption~\ref{sa-assump: Sharp DGP}(v) gives the required conditional fourth-moment control for $\varepsilon_{Y,i}$; mixed terms with $A_1,A_2\in\{Y,W\}$ then follow by Cauchy--Schwarz.
We begin with the decomposition
\begin{align*}
    \widehat{\varepsilon}_{A_1,i}\widehat{\varepsilon}_{A_2,i}
    = (\varepsilon_{A_1,i} + \widehat{\varepsilon}_{A_1,i}-\varepsilon_{A_1,i})
      (\varepsilon_{A_2,i} + \widehat{\varepsilon}_{A_2,i}-\varepsilon_{A_2,i}),
\end{align*}
which allows us to control the difference between estimated and population covariances through uniform bounds on
$\widehat{\varepsilon}_{A,i}-\varepsilon_{A,i}$ for $A\in\{Y,W\}$. Fix $A\in\{Y,W\}$. For each $i$,
\begin{align*}
    &|\widehat{\varepsilon}_{A,i}-\varepsilon_{A,i}|\\
    &\le
    \sup_{\bx\in\B}
    \Big|
        \sum_{s\in\{0,1\}}\Indicator(\bX_i\in\A_s)\mu_{A,s}(\bX_i)\\
    &\hspace{1in}
        - \br_p(\bH^{-1}(\bX_i-\bx))^\top\bS
        \big[
            \Indicator(\bX_i\in\A_0)\widehat{\bbeta}_{A,0}(\bx)
            + \Indicator(\bX_i\in\A_1)\widehat{\bbeta}_{A,1}(\bx)
        \big]
    \Big|
    \Indicator(K_{\bH}(\bX_i-\bx)\neq0).
\end{align*}
Taking the maximum over $i$ and decomposing the difference into approximation and estimation components yields
\begin{align*}
    & \max_{1\le i\le n}|\widehat{\varepsilon}_{A,i}-\varepsilon_{A,i}| \\
    &\le
    \max_{1\le i\le n}\sup_{\bx\in\B}
    \Big|
        \sum_{s\in\{0,1\}}\Indicator(\bX_i\in\A_s)\mu_{A,s}(\bX_i)\\
    &\hspace{1in}
        - \br_p(\bH^{-1}(\bX_i-\bx))^\top\bS
        \big[
            \Indicator(\bX_i\in\A_0)\widehat{\bbeta}_{A,0}(\bx)
            + \Indicator(\bX_i\in\A_1)\widehat{\bbeta}_{A,1}(\bx)
        \big]
    \Big|
    \Indicator(K_{\bH}(\bX_i-\bx)\neq 0) \\
    &\le
    \sup_{\bx\in\B}\sup_{\by\in\X}
    \Big|
        \sum_{s\in\{0,1\}}\Indicator(\by\in\A_s)\mu_{A,s}(\by)\\
    &\hspace{1in}
        - \br_p(\bH^{-1}(\by-\bx))^\top\bS
        \big[
            \Indicator(\by\in\A_0)\bbeta_{A,0}(\bx)
            + \Indicator(\by\in\A_1)\bbeta_{A,1}(\bx)
        \big]
    \Big|
    \Indicator(K_{\bH}(\by-\bx)\neq 0) \\
    &\quad
    + \max_{t=0,1}\sup_{\bx\in\B}\sup_{\by\in\X}
    \Big|
        \br_p(\bH^{-1}(\by-\bx))^\top\bS
        \big(\widehat{\bbeta}_{A,t}(\bx)-\bbeta_{A,t}(\bx)\big)
    \Big|
    \Indicator(K_{\bH}(\by-\bx)\neq 0).
\end{align*}

The first term is the usual local polynomial approximation error and is of order $\bar{h}^{p+1}$ uniformly over $\bx$.
The second term is controlled by the uniform convergence rate of
$\widehat{\bbeta}_{A,t}(\bx)$, giving
\begin{align*}
    \max_{1\le i\le n}
    |\widehat{\varepsilon}_{A,i}-\varepsilon_{A,i}|
    \lessp
    \bar{h}^{p+1}
    + \sqrt{\frac{\log(1/\bar{h})}{n |\bH|}}
    + \frac{\log(1/\bar{h})}{n^{\frac{1+v}{2+v}}|\bH|}.
\end{align*}

Substituting this bound into the covariance decomposition and proceeding exactly as in Lemma~\ref{sa-lem: covariance} yields the stated uniform bounds for
$\widehat{\bSigma}_{A_1,A_2,t,\bx_1,\bx_2}$,
$\widehat{\Omega}_{A_1,A_2,\bx_1,\bx_2}^{(\bnu)}$.
\end{proof}

\noindent\textbf{Completion of the AMSE proof.}
Applying Lemma~\ref{sa-lem: bias} with $A \in \{Y,W\}$ yields
\begin{align*}
    \sup_{\bx \in \B}
    \Big|
        \E[\widehat{\tau}_{A}^{(\bnu)}(\bx)\mid\bX]
        - \tau_{A}^{(\bnu)}(\bx)
        - B_{A,\bx}^{(\bnu)}(\bh)
    \Big|
    = o_{\P}\!\left(\bar{h}^{p+1}\bh^{-\bnu}\right).
\end{align*}
Premultiplying the reduced-form bias vector by $\mathfrak{w}(\bx)$ gives $B_{\mathtt{F},\bx}^{(\bnu)}(\bh)$ in \eqref{sa-eq: leading bias fuzzy}. Combining this bias expansion with Lemma~\ref{sa-lem: covariance for fuzzy designs} gives the AMSE expansion for the leading linear component, and integrating the same pointwise expansion over $\B$ establishes the AIMSE result. \qed



%%%%%%%%%%%%%%%%%%%%%%%%%%%%%%%%%%%
\subsection{Proof of Lemma~\ref{sa-lem: Covariance: Fuzzy BATEC}}
%%%%%%%%%%%%%%%%%%%%%%%%%%%%%%%%%%%
Let
\begin{align*}
    a_n = \sqrt{\frac{\log(1/\bar{h})}{n |\bH|}}
          + \frac{\log(1/\bar{h})}{n^{\frac{v}{2+v}}|\bH|}
          + \bar{h}^{p+1}.
\end{align*}
Lemma~\ref{sa-lem: covariance for fuzzy designs} gives
\begin{align*}
    \sup_{\bx_1,\bx_2\in\B}
    \big\|\widehat{\bTheta}_{\ttF,\bx_1,\bx_2}^{(\bnu)}
          -\bTheta_{\ttF,\bx_1,\bx_2}^{(\bnu)}\big\|
    \lessp (n|\bH|\bh^{2\bnu})^{-1} a_n .
\end{align*}
The reduced-form convergence rates, together with $\inf_{\bx\in\B}|\tau_W^{(\bnu)}(\bx)|>0$, imply
\begin{align*}
    \sup_{\bx\in\B}\|\widehat{\mathfrak{w}}(\bx)-\mathfrak{w}(\bx)\|
    \lessp \bh^{-\bnu}a_n.
\end{align*}
The displayed rate condition in Lemma~\ref{sa-lem: Covariance: Fuzzy BATEC} makes the right-hand side $o_{\P}(1)$, so the feasible weights are uniformly bounded with probability approaching one.
Because $\mathfrak{w}(\cdot)$ is bounded and
$\sup_{\bx_1,\bx_2\in\B}\|\bTheta_{\ttF,\bx_1,\bx_2}^{(\bnu)}\|
\lesssim (n|\bH|\bh^{2\bnu})^{-1}$, the feasible plug-in covariance satisfies
\begin{align*}
    \sup_{\bx_1,\bx_2 \in \B} \big|\widehat{\Omega}_{\ttF, \bx_1,\bx_2}^{(\bnu)} - \Omega_{\ttF, \bx_1,\bx_2}^{(\bnu)}\big|
    \lessp \frac{\bh^{-\bnu}}{n|\bH|\bh^{2\bnu}}a_n.
\end{align*}
Finally, $\Omega_{\ttF,\bx,\bx}^{(\bnu)}=V_{\mathtt{F},\bx}^{(\bnu)}/(n|\bH|\bh^{2\bnu})$, so $\inf_{\bx\in\B}V_{\mathtt{F},\bx}^{(\bnu)}>0$ gives the required lower bound on the population denominator. The inverse-square-root bound follows from the displayed covariance bound and a mean-value expansion of $x\mapsto x^{-1/2}$ on this bounded-away-from-zero neighborhood. \qed

%%%%%%%%%%%%%%%%%%%%%%%%%%%%%%%%%%%%
\subsection{Proof of Theorem~\ref{sa-thm: Confidence Intervals: Fuzzy BATEC}}
%%%%%%%%%%%%%%%%%%%%%%%%%%%%%%%%%%%%
For the fixed boundary point $\bx$ in the theorem, a similar argument as the proof for Theorem~\ref{sa-thm: Confidence Intervals: Sharp BATEC} gives the following stochastic linearization,
\begin{align*}
    \Big|\widehat{\Tstat}_{\mathtt{F}}^{(\bnu)}(\bx) - \overline{\Tstat}_{\mathtt{F}}^{(\bnu)}(\bx)\Big|
    \lessp
    \bar{h}^{p+1}\sqrt{n |\bH|} + \frac{1}{\sqrt{n |\bH|}},
\end{align*}
provided that $\bar{h} \to 0$ and $n |\bH| \to \infty$. To see this, the pointwise reduced-form rates give
\begin{align*}
    |\widehat{\tau}_{A}^{(\bnu)}(\bx)-\tau_A^{(\bnu)}(\bx)|
    =O_{\P}\big(\bh^{-\bnu}((n|\bH|)^{-1/2}+\bar h^{p+1})\big),
    \qquad A\in\{Y,W\}.
\end{align*}
Since $|\tau_W^{(\bnu)}(\bx)|>0$, this implies $|\widehat{\tau}_W^{(\bnu)}(\bx)|^{-1}=O_{\P}(1)$ and
\begin{align*}
    \|\widehat{\mathfrak w}(\bx)-\mathfrak w(\bx)\|
    =O_{\P}\big(\bh^{-\bnu}((n|\bH|)^{-1/2}+\bar h^{p+1})\big).
\end{align*}
The exact second-order expansion in Section~\ref{sec: Second-Order Linearization: Fuzzy BATEC} gives
\begin{align*}
    |\mathfrak R_{\mathtt F}(\bx)|
    =O_{\P}\big(\bh^{-2\bnu}((n|\bH|)^{-1/2}+\bar h^{p+1})^2\big),
\end{align*}
and hence
\begin{align*}
    \frac{|\mathfrak R_{\mathtt F}(\bx)|}{\sqrt{\Omega_{\mathtt F,\bx,\bx}^{(\bnu)}}}
    =O_{\P}\!\left(
        \bh^{-\bnu}\sqrt{n|\bH|}
        ((n|\bH|)^{-1/2}+\bar h^{p+1})^2
      \right)=o_{\P}(1)
\end{align*}
under the theorem's bandwidth conditions. Combining these bounds with the pointwise Gram and covariance bounds displayed in the proof of Theorem~\ref{sa-thm: Confidence Intervals: Sharp BATEC} yields the displayed linearization. Here
\begin{align*}
    \overline{\Tstat}_{\mathtt{F}}^{(\bnu)}(\bx) = \sum_{i=1}^n \chi_i,
\end{align*}
with $\mathfrak{w}(\bx) = (\mathfrak{w}_1(\bx), \mathfrak{w}_2(\bx))$ and
\begin{align*}
    \chi_i
    &= n^{-1} (\Omega_{\mathtt{F}, \bx,\bx}^{(\bnu)})^{-1/2}
    \bnu!\be_{\bnu}^{\top} \bS^{-1}
    \Big[\Indicator(\bX_i\in\A_1)\bGamma_{1,\bx}^{-1}-\Indicator(\bX_i\in\A_0)\bGamma_{0,\bx}^{-1}\Big]
    \br_p(\bH^{-1}(\bX_i - \bx))\\
    &\hspace{1in} \cdot K_{\bH}(\bX_i - \bx)
    \big(\mathfrak{w}_1(\bx) \varepsilon_{Y,i} + \mathfrak{w}_2(\bx) \varepsilon_{W,i}\big).
\end{align*}
In particular, we have $\E[\chi_i] = 0$ and $\Var\big[\sum_{i=1}^n \chi_i\big] = 1$. Since $\mathfrak w(\bx)$ is fixed at the point $\bx$, $\varepsilon_{W,i}$ is bounded, and Assumption~\ref{sa-assump: Sharp DGP}(v) gives a uniformly bounded conditional third moment for $\varepsilon_{Y,i}$,
\begin{align*}
    \sum_{i=1}^n\E[|\chi_i|^3]\lesssim (n|\bH|\bh^{2\bnu})^{-1/2}=o(1).
\end{align*}
Berry--Esseen gives the normal approximation for $\overline{\Tstat}_{\mathtt F}^{(\bnu)}(\bx)$, and the stochastic linearization transfers it to $\widehat{\Tstat}_{\mathtt F}^{(\bnu)}(\bx)$. \qed

%%%%%%%%%%%%%%%%%%%%%%%%%%%%%%%%%%%%
\subsection{Proof of Theorem~\ref{sa-thm: Stochastic Linearization: Fuzzy BATEC}}
%%%%%%%%%%%%%%%%%%%%%%%%%%%%%%%%%%%%
Use the exact second-order expansion in Section~\ref{sec: Second-Order Linearization: Fuzzy BATEC}. The leading linear component is a fixed linear combination, with weights $\mathfrak{w}(\bx)$, of the reduced-form stochastic expansions for $Y$ and $W$. Applying Theorem~\ref{sa-thm: Stochastic Linearization: Sharp BATEC} to these two reduced-form estimators and then using Lemma~\ref{sa-lem: Covariance: Fuzzy BATEC} to replace $\Omega_{\mathtt{F},\bx,\bx}^{(\bnu)}$ by $\widehat{\Omega}_{\mathtt{F},\bx,\bx}^{(\bnu)}$ gives
\begin{align*}
    &\sup_{\bx \in \B}
    \left|\frac{\mathfrak{w}(\bx)^\top\mathfrak{T}(\bx)}
    {\sqrt{\widehat{\Omega}_{\mathtt{F},\bx,\bx}^{(\bnu)}}}
    - \overline{\Tstat}^{(\bnu)}_{\mathtt{F}}(\bx)\right| \\
    &\qquad \lessp
    \bar{h}^{p+1}\sqrt{n |\bH|}
    + \bh^{-\bnu}\sqrt{\log(1/\bar{h})}
      \left(\sqrt{\frac{\log(1/\bar{h})}{n|\bH|}}
      +\frac{\log(1/\bar{h})}{n^{\frac{v}{2+v}}|\bH|}
      +\bar{h}^{p+1}\right).
\end{align*}
Equation~\eqref{sa-eq: second order term -- fuzzy} and the lower bound
$\inf_{\bx\in\B}\Omega_{\mathtt{F},\bx,\bx}^{(\bnu)}\gtrsim(n|\bH|\bh^{2\bnu})^{-1}$ from $\inf_{\bx\in\B}V_{\mathtt{F},\bx}^{(\bnu)}>0$ imply that the normalized second-order remainder is bounded by the same rate. Combining the two displays proves the theorem.
\qed

%%%%%%%%%%%%%%%%%%%%%%%%%%%%%%%%%%%%
\subsection{Proof of Theorem~\ref{sa-thm: Gaussian Approximation Suprema -- Fuzzy}}
%%%%%%%%%%%%%%%%%%%%%%%%%%%%%%%%%%%%
Consider $\bz_i = (\bx_i, \varepsilon_{Y,i}, \varepsilon_{W,i})$, $1 \leq i \leq n$. For each $\bx \in \B$, define
\begin{align*}
    g_{\ttF, \bx}(\bu) = \Indicator(\bu \in \A_1)  \K^{(\bnu)}_{\ttF,1}(\bu;\bx) - \Indicator(\bu \in \A_0) \K^{(\bnu)}_{\ttF,0}(\bu;\bx),
    \qquad \bu \in \X,
\end{align*}
where
\begin{align*}
    \K^{(\bnu)}_{\ttF, t}(\bu;\bx) = n^{-1/2} (\Omega_{\ttF, \bx,\bx}^{(\bnu)})^{-1/2}\bnu!\be_{\bnu}^{\top} \bS^{-1} \bGamma_{t,\bx}^{-1} \br_p(\bH^{-1}(\bu - \bx)) K_{\bH}(\bu - \bx), \qquad \bu \in \X,\quad t \in \{0,1\}.
\end{align*}
Consider the class of functions
\begin{align*}
    \F = \{(\bu, \varepsilon_Y, \varepsilon_W) \in \reals^{d} \times \reals \times \reals \mapsto g_{\ttF, \bx}(\bu) \; (\mathfrak{w}_1(\bx)\varepsilon_{Y}+\mathfrak{w}_2(\bx)\varepsilon_{W}): \bx \in \B\},
\end{align*}
and let $\F_{\pm}=\F\cup(-\F)$. Also define
\begin{gather*}
    \G_{\ttF} = \{g_{\ttF, \bx}: \bx \in \B \}, \qquad
    \R_{\ttF} = \{(\varepsilon_Y, \varepsilon_W) \mapsto \mathfrak{w}_1(\bx)\varepsilon_{Y}+\mathfrak{w}_2(\bx)\varepsilon_{W}: \bx \in \B\}.
\end{gather*}
Then we have
\begin{align*}
    \sup_{\bx \in \B} |\overline{\Tstat}^{(\bnu)}_{\mathtt{F}}(\bx)| = \sup_{f \in \F_{\pm}} \frac{1}{\sqrt{n}} \sum_{i = 1}^n \Big\{f(\bz_i) - \E[f(\bz_i)]\Big\}.
\end{align*}
We use Lemma~\ref{sa-lem: suprema coupling lemma} to approximate the distribution of the right hand side of the above equation. First, we verify the VC-type entropy and moment conditions. The fixed indicators $\Indicator(\bu\in\A_t)$ are absorbed into $\G_{\ttF}$ and do not increase $L_2(Q)$ covering numbers beyond a constant factor, because multiplication by a fixed indicator is a contraction in $L_2(Q)$. Therefore, using the same local-polynomial entropy calculation as in the proofs of Lemmas~\ref{sa-lem: Gram Matrix} and \ref{sa-lem: Linear Approximation}, but without the total-variation step, we obtain that: (i) $\G_{\ttF}$ is pointwise measurable; (ii) there exists a constant envelope function $\ttM_{\G_{\ttF}} \lesssim |\bH|^{-1/2}$ such that for all $0 < \varepsilon < 1$,
\begin{align*}
    \sup_{Q} N(\G_{\ttF}, \norm{\cdot}_{Q,2}, \varepsilon (2 \bc +1)^{d+1} \ttM_{\G_{\ttF}}) \leq 2 \bc' \varepsilon^{-4d-4} + 2,
\end{align*}
where the supremum is taken over all finite discrete measures and some constants $\bc$ and $\bc^{\prime}$ only depending on the compact set $\X$. Moreover, Assumption~\ref{sa-assump: Sharp DGP} and \ref{sa-assump: Fuzzy DGP} imply that $\R_{\ttF}$ is also pointwise measurable and for an envelope function $M_{\R_{\ttF}}(e_1, e_2) = |e_1| + |e_2|$, $(e_1, e_2) \in \reals^2$, we have for all $0 < \varepsilon < 1$,
\begin{align*}
    \sup_{Q} N(\R_{\ttF}, \norm{\cdot}_{Q,2}, \varepsilon \bc M_{\R_{\ttF}}) \leq \bc' \varepsilon^{-2d},
\end{align*}
where the supremum is taken over all finite discrete measures. The last entropy bound follows because Assumption~\ref{sa-assump: Sharp DGP}(iii), Assumption~\ref{sa-assump: Fuzzy DGP}(ii), and $|\bnu|\leq p$ imply that $\tau_Y^{(\bnu)}(\bx)$ and $\tau_W^{(\bnu)}(\bx)$ are continuously differentiable on a neighborhood of $\B$; combined with $\inf_{\bx\in\B}|\tau_W^{(\bnu)}(\bx)|>0$, the map $\bx\mapsto\mathfrak w(\bx)$ is bounded and Lipschitz on the compact set $\B$. Thus the coefficient set $\{\mathfrak w(\bx):\bx\in\B\}\subset\reals^2$ has polynomial covering numbers inherited from $\B$. Since $\F\subseteq \G_{\ttF}\cdot\R_{\ttF}$, the product-covering inequality gives, for a constant $c>0$,
\begin{align*}
    N(\F, \norm{\cdot}_{Q,2}, \varepsilon c \ttM_{\G_{\ttF}}M_{\R_{\ttF}})
    &\leq
    N(\G_{\ttF}, \norm{\cdot}_{Q,2}, c\varepsilon\ttM_{\G_{\ttF}})
    N(\R_{\ttF}, \norm{\cdot}_{Q,2}, c\varepsilon M_{\R_{\ttF}}).
\end{align*}
Hence $\F_{\pm}$ is also pointwise measurable and satisfies the same VC-type bound as $\F$ up to a universal constant: for all $0 < \varepsilon < 1$,
\begin{align*}
    \sup_{Q} N(\F_{\pm}, \norm{\cdot}_{Q,2}, \varepsilon \bc \; \ttM_{\G_{\ttF}} M_{\R_{\ttF}}) \leq 64 \bc'^2 \varepsilon^{-6d-4} + 4 \bc' \varepsilon^{-2d},
\end{align*}
where the supremum is taken over all finite discrete measures. The denominator condition $\inf_{\bx \in \B}|\tau_W^{(\bnu)}(\bx)|>0$, continuity, and compactness of $\B$ imply that the population Wald weights $\mathfrak{w}(\bx)$ are uniformly bounded. For the moment condition, take $q=2+v\geq4$ in Lemma~\ref{sa-lem: suprema coupling lemma}. Since $\varepsilon_{W,i}=W_i-\E[W_i|\bX_i]$ is bounded and Assumption~\ref{sa-assump: Sharp DGP} gives $\sup_{\bx \in \X}\E[|Y_i(t)|^{2+v}|\bX_i = \bx] < \infty$, the envelope
\[
    F(\bx, \varepsilon_Y, \varepsilon_W) = \ttM_{\G_{\ttF}} M_{\R_{\ttF}}(\varepsilon_Y, \varepsilon_W)
\]
satisfies $\|F\|_{\P,2+v}\lesssim |\bH|^{-1/2}$. Together with the local-polynomial moment calculation used in the proof of Lemma~\ref{sa-lem: Linear Approximation}, for
\begin{align*}
    \sigma = 1, \qquad b = |\bH|^{-1/2},
\end{align*}
we have $\sup_{f \in \F_{\pm}} \P|f|^k \le \sigma^2 b^{k-2}$ for $k = 2, 3, 4$ and $\|F\|_{\P,2+v} \le b$.
Taking $\gamma = (\log n)^{-1}$ in Lemma~\ref{sa-lem: suprema coupling lemma},
\begin{align*}
    \P\left\{ |\sup_{\bx \in \B} |\overline{\Tstat}^{(\bnu)}_{\mathtt{F}}(\bx)| - \sup_{\bx \in \B} |Z^{(\bnu)}_{\mathtt{F}}(\bx)|| > \frac{\log n}{(n |\bH|)^{1/6}} + \frac{(\log n)^{5/4}}{(n |\bH|)^{1/4}} + \frac{(\log n)^{3/2}}{(n^{\frac{v}{2+v}} |\bH|)^{1/2}}\right\} \le C \left( \frac{1}{\log n} + \frac{\log n}{n} \right).
\end{align*}
\qed

\begin{lem}[KS Distance Between $Z_{\mathtt{F}}^{(\bnu)}$ and $\widehat{Z}_{\mathtt{F}}^{(\bnu)}$]\label{sa-lem: feasible fuzzy gaussian to infeasible fuzzy gaussian}
Suppose the conditions for Theorem~\ref{sa-thm: Confidence Bands: Fuzzy BATEC -- suprema} hold. Then, for any multi-index $|\bnu|\leq p$,
\begin{align*}
    \sup_{u \in \reals}
    \left|
    \P \left( \sup_{\bx \in \B} \big| \widehat{Z}_{\mathtt{F}}^{(\bnu)}(\bx) \big| \leq u \;\middle|\; \bV \right)
    - \P \left( \sup_{\bx \in \B} \big| Z_{\mathtt{F}}^{(\bnu)}(\bx) \big| \leq u \right)
    \right|
    = o_{\P}(1).
\end{align*}
\end{lem}

\begin{proof}
Let
\begin{align*}
    \mathfrak{R}_{\mathtt{F,cc}}
    =
    \bh^{-\bnu}
    \left(
    \sqrt{\frac{\log n}{n|\bH|}}
    +\frac{\log n}{n^{v/(2+v)}|\bH|}
    +\bar h^{p+1}
    \right).
\end{align*}
For $\bx,\by\in\B$, define
\[
    \bPi_{\mathtt{F},\bx,\by}
    =\frac{\Omega_{\mathtt{F},\bx,\by}^{(\bnu)}}
    {\sqrt{\Omega_{\mathtt{F},\bx,\bx}^{(\bnu)}\Omega_{\mathtt{F},\by,\by}^{(\bnu)}}},
    \qquad
    \widehat{\bPi}_{\mathtt{F},\bx,\by}
    =\frac{\widehat{\Omega}_{\mathtt{F},\bx,\by}^{(\bnu)}}
    {\sqrt{\widehat{\Omega}_{\mathtt{F},\bx,\bx}^{(\bnu)}\widehat{\Omega}_{\mathtt{F},\by,\by}^{(\bnu)}}}.
\]
Lemma~\ref{sa-lem: Covariance: Fuzzy BATEC}, the lower bound
$\inf_{\bx\in\B}V_{\mathtt{F},\bx}^{(\bnu)}>0$, and the mean-value expansion of $x\mapsto x^{-1/2}$ give
\begin{align*}
    \sup_{\bx,\by\in\B}
    \big|\widehat{\bPi}_{\mathtt{F},\bx,\by}
    -\bPi_{\mathtt{F},\bx,\by}\big|
    \lessp \mathfrak{R}_{\mathtt{F,cc}}.
\end{align*}
The Gaussian processes $Z_{\mathtt{F}}^{(\bnu)}$ and $\widehat Z_{\mathtt{F}}^{(\bnu)}$ have the same boundary-indexed local-polynomial structure as the sharp Gaussian processes, with the scalar residual replaced by the linear combination $\mathfrak w(\bx)^\top\boldsymbol{\varepsilon}_i$, whose coefficients are uniformly bounded by the denominator nondegeneracy condition. The entropy and modulus calculations in Lemma~\ref{sa-lem: feasible gaussian to infeasible gaussian} therefore apply with the sharp covariance-comparison rate $\mathfrak{R}_{\mathtt{cc}}$ used there replaced by the present fuzzy rate $\mathfrak{R}_{\mathtt{F,cc}}$. In particular, on a $\delta_n=n^{-s}$ net with $s$ large enough,
\begin{align*}
    \sup_{u \in \reals}
    \left|
    \P \left( \sup_{\bx \in \B} \big| \widehat{Z}_{\mathtt{F}}^{(\bnu)}(\bx) \big| \leq u \;\middle|\; \bV \right)
    - \P \left( \sup_{\bx \in \B} \big| Z_{\mathtt{F}}^{(\bnu)}(\bx) \big| \leq u \right)
    \right|
    \lessp \log(n)\mathfrak{R}_{\mathtt{F,cc}}^{1/2}.
\end{align*}
The second displayed rate condition in Theorem~\ref{sa-thm: Confidence Bands: Fuzzy BATEC -- suprema} gives $\log^2(n)\mathfrak{R}_{\mathtt{F,cc}}=o(1)$, and hence the right-hand side is $o_{\P}(1)$.
\end{proof}


\subsection{Proof of Theorem~\ref{sa-thm: Confidence Bands: Fuzzy BATEC -- suprema}}

Define the convergence rates
\begin{equation*}
    \mathfrak{R}_{\mathtt{sup}}
    =
    \frac{\log n}{(n |\bH|)^{1/6}}
    + \frac{(\log n)^{5/4}}{(n |\bH|)^{1/4}}
    + \frac{(\log n)^{3/2}}{(n^{v/(2+v)} |\bH|)^{1/2}},
\end{equation*}
and
\[
    \mathfrak{R}_{\mathtt{F,lin}}
    =
    \bar{h}^{p+1}\sqrt{n |\bH|}
    + \bh^{-\bnu}\sqrt{\log(1/\bar{h})}
      \left(\sqrt{\frac{\log(1/\bar{h})}{n|\bH|}}
      +\frac{\log(1/\bar{h})}{n^{\frac{v}{2+v}}|\bH|}
      +\bar{h}^{p+1}\right).
\]
First, we establish the coupling between the intermediate processes. Let
\[
    M_{\overline T,\mathtt{F}}=\sup_{\bx\in\B}|\overline{\Tstat}_{\mathtt{F}}^{(\bnu)}(\bx)|,
    \qquad
    M_{Z,\mathtt{F}}=\sup_{\bx\in\B}|Z_{\mathtt{F}}^{(\bnu)}(\bx)|.
\]
Applying Lemma~\ref{sa-lem: generic suprema comparison} with $\mathbb{T}_1=M_{\overline T,\mathtt{F}}$ and $\mathbb{T}_2=M_{Z,\mathtt{F}}$, viewed as singleton-indexed processes, and $\eta=\mathfrak{R}_{\mathtt{sup}}$, gives
\[
    \sup_{u\in\reals}\left|\P(M_{\overline T,\mathtt{F}}\le u)-\P(M_{Z,\mathtt{F}}\le u)\right|
    \le
    \sup_{u\in\reals}\P(|M_{Z,\mathtt{F}}-u|\le \mathfrak{R}_{\mathtt{sup}})
    +\P(|M_{\overline T,\mathtt{F}}-M_{Z,\mathtt{F}}|>\mathfrak{R}_{\mathtt{sup}}).
\]
The first term is $O(\mathfrak{R}_{\mathtt{sup}}\sqrt{\log n})=o(1)$ by Gaussian anti-concentration and the entropy calculation in the proof of Theorem~\ref{sa-thm: Gaussian Approximation Suprema -- Fuzzy}. The second term is $o(1)$ by Theorem~\ref{sa-thm: Gaussian Approximation Suprema -- Fuzzy}. Therefore,
\begin{equation}\label{sa-eq: t-to-z fuzzy}
    \sup_{u \in \reals} \left| \P \left( \sup_{\bx \in \B} \big| \overline{\Tstat}_{\mathtt{F}}^{(\bnu)}(\bx) \big| \leq u \right) - \P \left( \sup_{\bx \in \B} \big| Z_{\mathtt{F}}^{(\bnu)}(\bx) \big| \leq u \right) \right| = o(1).
\end{equation}

Furthermore, Theorem~\ref{sa-thm: Stochastic Linearization: Fuzzy BATEC} gives
\[
    \sup_{\bx\in\B}
    |\widehat{\Tstat}_{\mathtt{F}}^{(\bnu)}(\bx)
    -\overline{\Tstat}_{\mathtt{F}}^{(\bnu)}(\bx)|
    =
    O_{\P}(\mathfrak{R}_{\mathtt{F,lin}}).
\]
The side conditions of the theorem imply $\mathfrak{R}_{\mathtt{F,lin}}\sqrt{\log n}=o(1)$. Choose a deterministic sequence $a_n$ such that $\mathfrak{R}_{\mathtt{F,lin}}=o(a_n)$ and $a_n\sqrt{\log n}=o(1)$. Applying Lemma~\ref{sa-lem: generic suprema comparison} with $\mathbb{T}_1=\widehat{\Tstat}_{\mathtt{F}}^{(\bnu)}$, $\mathbb{T}_2=\overline{\Tstat}_{\mathtt{F}}^{(\bnu)}$, and $\eta=a_n$, and then bounding the anti-concentration of $M_{\overline T,\mathtt{F}}$ through Equation~\eqref{sa-eq: t-to-z fuzzy}, gives
\begin{equation}\label{sa-eq: t-to-t fuzzy}
    \sup_{u \in \reals} \left| \P \left( \sup_{\bx \in \B} \big| \widehat{\Tstat}_{\mathtt{F}}^{(\bnu)}(\bx) \big| \leq u \right) - \P \left( \sup_{\bx \in \B} \big| \overline{\Tstat}_{\mathtt{F}}^{(\bnu)}(\bx) \big| \leq u \right) \right| = o(1),
\end{equation}
and Lemma~\ref{sa-lem: feasible fuzzy gaussian to infeasible fuzzy gaussian} gives
\begin{equation}\label{sa-eq: z-to-z fuzzy}
    \sup_{u \in \reals} \left| \P \left( \sup_{\bx \in \B} \big| \widehat{Z}_{\mathtt{F}}^{(\bnu)}(\bx) \big| \leq u \;\middle|\; \bV \right) - \P \left( \sup_{\bx \in \B} \big| Z_{\mathtt{F}}^{(\bnu)}(\bx) \big| \leq u \right) \right| = o_{\P}(1).
\end{equation}
The distributional approximation follows from the triangle inequality applied to Equations~\eqref{sa-eq: t-to-z fuzzy}, \eqref{sa-eq: t-to-t fuzzy}, and \eqref{sa-eq: z-to-z fuzzy}. For coverage of the feasible confidence band, let
\[
    E_{\mathtt{F},\alpha}
    =
    \left\{
    \zeta^{(\bnu)}(\bx)\in\widehat{\CI}_{\mathtt{F},\alpha}^{(\bnu)}(\bx)
    \text{ for all }\bx\in\B
    \right\}.
\]
By definition of the band, $E_{\mathtt{F},\alpha}=\{\sup_{\bx\in\B}|\widehat{\Tstat}_{\mathtt{F}}^{(\bnu)}(\bx)|\leq \q_\alpha\}$. Hence the preceding approximation and the conditional definition of $\q_\alpha$ give
\begin{align*}
    \P(E_{\mathtt{F},\alpha})
    &= \P\left(\sup_{\bx\in\B}|\widehat{\Tstat}_{\mathtt{F}}^{(\bnu)}(\bx)|\leq \q_\alpha\right)\\
    &= \E\left[
        \P\left(\sup_{\bx\in\B}|\widehat Z_{\mathtt{F}}^{(\bnu)}(\bx)|\leq \q_\alpha\;\middle|\;\bV\right)
       \right]+o(1)\\
    &= 1-\alpha+o(1),
\end{align*}
using continuity of the supremum distribution of the feasible Gaussian process.
\qed


%%%%%%%%%%%%%%%%%%%%%%%%%%%%%%%%%%%%
\subsection{Proof of Lemma~\ref{sa-lem: Fuzzy Integral: Bias}}
%%%%%%%%%%%%%%%%%%%%%%%%%%%%%%%%%%%%

A similar argument to the proof of Lemma~\ref{sa-lem: bias} shows that under the assumptions imposed,
\begin{align*}
    \sup_{\bx \in \B} \big|\E[\mathfrak{w}(\bx)^\top\mathfrak{T}(\bx)|\bX] - B_{\mathtt{F},\bx}^{(\bnu)}(\bh) \big| = o_{\P}(\bar{h}^{p+1}\bh^{-\bnu}).
\end{align*}
Because $\mathfrak{w}(\cdot)$ is bounded under the denominator nondegeneracy condition and the reduced-form bias expansion is uniform in $\bx$, the displayed remainder is uniform over $\B$. Integrating against $w(\bx)\dif\Haus^{d-1}(\bx)$ therefore preserves the order by $\int_{\B} |w(\bx)| \dif\Haus^{d-1}(\bx)<\infty$, and the integral of $B_{\mathtt{F},\bx}^{(\bnu)}(\bh)$ is $B_{\mathtt{F,WBATE}}^{(\bnu)}(\bh)$ by definition. Under the common-bandwidth restriction $\bh=\bc h$, the same integration of the pointwise leading term gives the stated $h^{p+1-|\bnu|}\widetilde B_{\mathtt{F,WBATE}}^{(\bnu)}(\bc)$ form.
\qed

%%%%%%%%%%%%%%%%%%%%%%%%%%%%%%%%%%%%
\subsection{Proof of Lemma~\ref{sa-lem: Fuzzy Integral: Variance}}
%%%%%%%%%%%%%%%%%%%%%%%%%%%%%%%%%%%%

The proof is the fuzzy analog of Lemma~\ref{sa-lem: Variance: Sharp WBATE}. By Fubini's theorem,
\begin{align*}
    \Var[\check{\zeta}_{\mathtt{WBATE}}^{(\bnu)}|\bX]
    = \int_{\B}\int_{\B}
      \Cov[\mathfrak{w}(\bb_1)^\top\mathfrak{T}(\bb_1),\mathfrak{w}(\bb_2)^\top\mathfrak{T}(\bb_2)|\bX]
      w(\bb_1)w(\bb_2)\dif\Haus^{d-1}(\bb_1)\dif\Haus^{d-1}(\bb_2).
\end{align*}
Lemma~\ref{sa-lem: Covariance: Fuzzy BATEC} replaces this covariance uniformly by $\Omega_{\mathtt{F},\bb_1,\bb_2}^{(\bnu)}$ with an integrated error of order $o_{\P}(\bar{h}^{d-1}/(n|\bH|\bh^{2\bnu}))$. Compact support of $K$ implies that $\Omega_{\mathtt{F},\bb_1,\bb_2}^{(\bnu)}=0$ unless $\|\bb_1-\bb_2\|\lesssim \bar{h}$; by \eqref{sa-eq: local boundary measure regularity}, the effective integration set has product Hausdorff measure $O(\bar{h}^{d-1})$. The pointwise covariance is of order $(n|\bH|\bh^{2\bnu})^{-1}$, which gives the stated upper bound. The lower bound follows from the aggregate fuzzy variance nondegeneracy condition \eqref{sa-eq: fuzzy WBATE aggregate nondegeneracy}, which rules out cancellation after applying the fuzzy Wald weights and integrating over $\B$. Replacing $\Omega_{\mathtt{F},\bb_1,\bb_2}^{(\bnu)}$ by $\widehat{\Omega}_{\mathtt{F},\bb_1,\bb_2}^{(\bnu)}$ uses Lemma~\ref{sa-lem: Covariance: Fuzzy BATEC} once more.
\qed

%%%%%%%%%%%%%%%%%%%%%%%%%%%%%%%%%%%%
\subsection{Proof of Theorem~\ref{sa-thm: MSE Expansion: Fuzzy WBATE}}
%%%%%%%%%%%%%%%%%%%%%%%%%%%%%%%%%%%%

By Lemma~\ref{sa-lem: Fuzzy Integral: Bias},
\begin{align*}
    \E[\check{\zeta}_{\mathtt{WBATE}}^{(\bnu)}|\bX]
    = B_{\mathtt{F,WBATE}}^{(\bnu)}(\bh)+o_{\P}(\bar{h}^{p+1}\bh^{-\bnu}).
\end{align*}
Squaring the bias expansion gives $\big(B_{\mathtt{F,WBATE}}^{(\bnu)}(\bh)\big)^2+o_{\P}(\bar{h}^{2p+2}\bh^{-2\bnu})$. Lemma~\ref{sa-lem: Fuzzy Integral: Variance} gives the variance term $\Omega_{\mathtt{F,WBATE}}^{(\bnu)}$ with remainder $o_{\P}(\bar{h}^{d-1}/(n|\bH|\bh^{2\bnu}))$. Adding these two components gives the displayed AMSE expansion.
\qed

%%%%%%%%%%%%%%%%%%%%%%%%%%%%%%%%%%%%
\subsection{Proof of Theorem~\ref{sa-thm: Asymptotic Normality: Fuzzy WBATE}}
%%%%%%%%%%%%%%%%%%%%%%%%%%%%%%%%%%%%

The proof follows the same Lyapunov argument as Theorem~\ref{sa-thm: Asymptotic Normality: Sharp WBATE}, applied to the linearized fuzzy statistic obtained by integrating $\overline{\Tstat}^{(\bnu)}_{\mathtt{F}}(\bb)$ against $w(\bb)$. Lemma~\ref{sa-lem: Fuzzy Integral: Variance} gives the variance rate $\bar{h}^{d-1}/(n|\bH|\bh^{2\bnu})$ and consistency of its plug-in estimator. Lemma~\ref{sa-lem: Fuzzy Integral: Bias} and the side condition $n|\bH|\bar{h}^{2p+2}/\bar{h}^{d-1}=o(1)$ make the linearized bias negligible relative to the standard deviation. The final two displayed rate conditions imply the two stochastic pieces of the uniform reduced-form consistency condition, while $\bh^{-\bnu}\bar h^{p+1}=o(1)$ follows from $|\bnu|\le p$ and $\bar h/\underline h\lesssim1$. Moreover, \eqref{sa-eq: second order term -- fuzzy}, integrability of $w$, and the final displayed rate conditions in the theorem imply
\begin{align*}
    \frac{\left|\int_{\B} \mathfrak{R}_{\mathtt{F}}(\bb) w(\bb)\dif\Haus^{d-1}(\bb)\right|}
    {\left(\bar{h}^{d-1}/(n|\bH|\bh^{2\bnu})\right)^{1/2}}
    =o_{\P}(1).
\end{align*}
The Lyapunov condition follows from the same calculation as in the sharp WBATE proof, with the scalar residual replaced by the linearized fuzzy residual $\mathfrak{w}_1(\bb)\varepsilon_{Y,i}+\mathfrak{w}_2(\bb)\varepsilon_{W,i}$. In particular, if $\chi_i$ denotes the normalized integrated linearized fuzzy summand, then the boundedness of $\mathfrak{w}(\bb)$, the boundedness of $\varepsilon_{W,i}$, and Assumption~\ref{sa-assump: Sharp DGP}(v) imply
\begin{align*}
    \sum_{i=1}^n \E[|\chi_i|^3]
    \lesssim
    \left(\frac{\bar{h}^{d-1}}{n|\bH|}\right)^{1/2}
    = o(1).
\end{align*}
Therefore the feasible statistic $\widehat{\Tstat}_{\mathtt{F,WBATE}}^{(\bnu)}$ is asymptotically standard normal.
\qed

%%%%%%%%%%%%%%%%%%%%%%%%%%%%%%%%%%%%
\subsection{Proof of Theorem~\ref{sa-thm: Convergence Rate: Fuzzy LBATE}}
%%%%%%%%%%%%%%%%%%%%%%%%%%%%%%%%%%%%

Follows by Theorem~\ref{sa-thm: Convergence Rates: Fuzzy BATEC} after noting that
\begin{align*}
    |\widehat{\zeta}_{\mathtt{LBATE}}^{(\bnu)} - \zeta_{\mathtt{LBATE}}^{(\bnu)}|
    \leq \sup_{\bx\in\B} |\widehat{\zeta}^{(\bnu)}(\bx) - \zeta^{(\bnu)}(\bx)|.
\end{align*}
\qed

%%%%%%%%%%%%%%%%%%%%%%%%%%%%%%%%%%%%
\subsection{Proof of Theorem~\ref{sa-thm: Confidence Intervals: Fuzzy LBATE}}
%%%%%%%%%%%%%%%%%%%%%%%%%%%%%%%%%%%%

Consider the event $E_{\mathtt{F}} = \Big\{\sup_{\bb \in \B} \frac{|\widehat{\zeta}^{(\bnu)}(\bb) - \zeta^{(\bnu)}(\bb)|}{(\widehat{\Omega}_{\mathtt{F},\bb,\bb}^{(\bnu)})^{1/2}} \leq \q_{\alpha}\Big\}$. Theorem~\ref{sa-thm: Confidence Bands: Fuzzy BATEC -- suprema} implies that $\P(E_{\mathtt{F}}) = 1 - \alpha + o(1)$. On the event $E_{\mathtt{F}}$, we have
\begin{align*}
    \widehat{\zeta}^{(\bnu)}(\bb) - \q_{\alpha}(\widehat{\Omega}_{\mathtt{F},\bb,\bb}^{(\bnu)})^{1/2}
    \leq \zeta^{(\bnu)}(\bb)
    \leq \widehat{\zeta}^{(\bnu)}(\bb) + \q_{\alpha}(\widehat{\Omega}_{\mathtt{F},\bb,\bb}^{(\bnu)})^{1/2}, \qquad \forall \; \bb \in \B,
\end{align*}
which implies
\begin{align*}
    \sup_{\bb \in \B}  \widehat{\zeta}^{(\bnu)}(\bb) - \q_{\alpha}(\widehat{\Omega}_{\mathtt{F},\bb,\bb}^{(\bnu)})^{1/2}
    \leq \sup_{\bb \in \B} \zeta^{(\bnu)}(\bb)
    \leq \sup_{\bb \in \B} \widehat{\zeta}^{(\bnu)}(\bb) + \q_{\alpha}(\widehat{\Omega}_{\mathtt{F},\bb,\bb}^{(\bnu)})^{1/2}.
\end{align*}
The stated result then follows.
\qed


%%%%%%%%%%%%%%%%%%%%%%%%%%%%%%%%%%%%%%%%%%%%%%%%%%%%%%%%%%%%%%%%%%%%%%
\section{Implementation Details}\label{sa-sec: Implementation Details}
%%%%%%%%%%%%%%%%%%%%%%%%%%%%%%%%%%%%%%%%%%%%%%%%%%%%%%%%%%%%%%%%%%%%%%

This section provides additional details on the practical implementation of the pointwise and uniform inference methods discussed in Section 3.3 of the paper. The main issue for uniform inference is that the target process is indexed by the continuum $\bx\in\B$, whereas computation is necessarily carried out on a finite discretization of the boundary. The following steps describe the default implementation used by the companion software and recommended in applications.

\begin{enumerate}[leftmargin=*]

    \item \textit{Represent the boundary and choose the target set.}
    When $\B$ is piecewise linear, as in the empirical application, it is convenient to parameterize the boundary by arc length. Let $\gamma:[0,\mathfrak{L}(\B)]\to\B$ denote such a parameterization, where $\mathfrak{L}(\B)$ is the total boundary length. If the researcher wants inference only on a subset of the assignment boundary, for example after trimming endpoints, corners, or regions with weak local support, the same construction can be applied to the trimmed target set $\B_0\subseteq\B$. This distinction is important: uniform confidence bands cover the chosen target set, so any trimming changes the estimand from the full boundary to the trimmed boundary.

    \item \textit{Construct a quasi-uniform grid on the boundary.}
    Choose grid points $\bb_1,\ldots,\bb_J\in\B_0$ with approximately equal spacing in arc length. For a single connected component this can be done by setting $\bb_j=\gamma(s_j)$ with $s_j=(j-1)\mathfrak{L}(\B_0)/(J-1)$, $j=1,\ldots,J$. For a piecewise linear boundary with multiple segments, grid points can be allocated across segments in proportion to segment length, while forcing vertices or other substantively important boundary points to be included when desired. In higher-dimensional settings, the analogous requirement is a quasi-uniform mesh on the $(d-1)$-dimensional boundary with respect to Hausdorff measure.

    \item \textit{Choose the mesh size.}
    The grid should be fine enough that the discretized supremum is a good approximation to the continuous supremum. A useful practical rule is to make the maximum spacing between adjacent boundary grid points small relative to the bandwidth scale used for inference. Equivalently, the mesh should be refined until the simulated critical value and the visual shape of the confidence band are stable. In simple two-dimensional applications, values such as $J=40$ or $J=80$ often give similar results, but the appropriate choice depends on the length and geometry of $\B$, the bandwidth, and the degree of treatment effect heterogeneity. We recommend reporting a grid-sensitivity check, for example comparing $J$, $2J$, and, when computationally feasible, $4J$.

    \item \textit{Select bandwidths and polynomial orders.}
    The plug-in selectors described in Section 3.3 of the paper can be implemented pointwise, producing $\widehat{h}_{\mathtt{MSE},\bx}$, or globally, producing $\widehat{h}_{\mathtt{IMSE}}$. For uniform inference and graphical display along the boundary, a global bandwidth is often easier to interpret and tends to produce smoother covariance estimates. Pointwise bandwidths can also be used, provided the covariance matrix of the resulting process is computed using the same bandwidth choices. Robust bias-corrected inference is implemented by using the bandwidth selected for the $p$th-order point estimator and constructing the associated $t$-statistic with a $(p+1)$th-order local polynomial, including the corresponding variance adjustment. When the coordinates of the score have different scales, one may either standardize the score before using a scalar bandwidth or use a diagonal bandwidth matrix $\bH=\diag(h_1,\ldots,h_d)$ as allowed in this supplemental appendix.

    \item \textit{Compute point estimates and covariance estimates on the grid.}
    For each grid point $\bb_j$, compute the local-polynomial estimates, the robust bias-corrected $t$-statistics, and the variance estimates. For uniform inference, also compute the estimated cross-covariances
    \[
        \widehat{\Omega}_{\bb_j,\bb_k}^{(\bnu)},
        \qquad 1\leq j,k\leq J,
    \]
    using the formulas in Sections SA-3 and SA-6. The simulated Gaussian process is then represented by a mean-zero Gaussian vector with correlation matrix
    \[
        \widehat{\bPi}_{j,k}
        =
        \frac{\widehat{\Omega}_{\bb_j,\bb_k}^{(\bnu)}}
        {\sqrt{\widehat{\Omega}_{\bb_j,\bb_j}^{(\bnu)}
               \widehat{\Omega}_{\bb_k,\bb_k}^{(\bnu)}}},
        \qquad 1\leq j,k\leq J.
    \]
    In finite samples, $\widehat{\bPi}$ should be symmetrized before simulation.

    \item \textit{Regularize the covariance matrix if needed.}
    The estimated matrix $\widehat{\bPi}$ may fail to be positive semidefinite because of simulation error, numerical precision, or unstable local covariance estimates. This is a finite-sample numerical issue rather than a change in the target parameter. A simple regularization is to replace $\widehat{\bPi}$ by its nearest positive semidefinite approximation, or equivalently to truncate negative eigenvalues at a small nonnegative tolerance and rescale the diagonal entries to one. Another common approach is to use a small ridge adjustment, $\widehat{\bPi}_{\lambda}=(1-\lambda)\widehat{\bPi}+\lambda \mathbf{I}_J$, with the smallest $\lambda\geq0$ that makes the matrix positive semidefinite. The chosen adjustment should be reported when it is not negligible. See \citet{Cattaneo-Feng-Underwood_2024_JASA} for related discussion and theoretical results on covariance regularization.

    \item \textit{Simulate the critical value.}
    Given the regularized correlation matrix, draw $M$ independent vectors from $N(0,\widehat{\bPi})$ and compute the empirical $(1-\alpha)$ quantile of
    \[
        \max_{1\leq j\leq J}|Z_j|.
    \]
    This quantile is the simulated critical value for the uniform confidence band over the grid. In practice, $M$ should be large enough that simulation error is small relative to the desired reporting precision; increasing $M$ and checking that the critical value is stable is a simple diagnostic.

    \item \textit{Check local support and numerical stability.}
    Uniform inference can be sensitive to boundary regions where one side has little local support or the local design matrix is poorly conditioned. We recommend inspecting, for each grid point, effective local sample sizes or kernel masses on the two sides of the boundary, together with condition numbers of the estimated Gram matrices. Grid points with very small one-sided kernel mass, extremely unbalanced mass across sides, or unstable Gram matrices indicate that the local design is weak at those locations. In such cases, researchers should report the issue, consider sensitivity to bandwidth choices, and, when substantively justified, trim the problematic portion of the boundary and redefine the target set accordingly.

    \item \textit{Report sensitivity analyses.}
    A transparent empirical implementation should report the main estimates together with basic sensitivity checks: alternative grid sizes, alternative bandwidth choices, and whether covariance regularization was needed. For applications with corners or visibly irregular boundaries, it is also useful to report results with and without small neighborhoods of those regions. Researchers should also examine whether the density of the score changes abruptly near the boundary or whether observations appear to bunch on one side of $\B$. In BD designs this type of manipulation or sorting diagnostic is necessarily multivariate and boundary-local: useful checks include scatterplots near $\B$, side-specific kernel masses along the boundary grid, and sensitivity to excluding small neighborhoods where the empirical distribution of $\bX_i$ appears discontinuous or unusually concentrated. These checks help distinguish genuine treatment effect heterogeneity from artifacts of discretization, weak local support, numerical instability, or nonrandom sorting around the assignment boundary.

\end{enumerate}

The companion software package \texttt{rd2d} and article \citep{Cattaneo-Titiunik-Yu_2025_rd2d} discuss additional implementation details.

%%%%%%%%%%%%%%%%%%%%%%%%%%%%%%%%%%%%%%%%%%%%%%%%%%%%%%%%%%%%%%%%%%%%%%
\section{Empirical Application}
%%%%%%%%%%%%%%%%%%%%%%%%%%%%%%%%%%%%%%%%%%%%%%%%%%%%%%%%%%%%%%%%%%%%%%

This section reports numerical results for the SPP empirical application that are omitted from the main paper. Replication codes are available at \url{https://rdpackages.github.io/replication/}. The main paper uses the ITT bandwidth selector throughout; this section reports the corresponding full boundary-grid results and then gives the analogous results using the fuzzy bandwidth selector.


\subsection{ITT Bandwidth Selector}

Due to space limitations, the paper reports selected boundary points and aggregate summaries in Table 1. Here we provide the corresponding full boundary-grid results for all $40$ evaluation points.
\begin{itemize}
    \item Table~\ref{tab:sa-empapp_main_itt} reports the reduced-form, or intention-to-treat (ITT), estimates for college enrollment.
    \item Table~\ref{tab:sa-empapp_main_fs} reports the first-stage estimates for treatment take-up.
    \item Table~\ref{tab:sa-empapp_main_fuzzy} reports the fuzzy estimates for college enrollment.
    \item Table~\ref{tab:sa-empapp_covbal_itt} reports the covariate-balance reduced-form estimates using mother's education as the outcome.
\end{itemize}

The figures in this subsection complement the confidence intervals and bands plots presented in the paper. For each quantity, we report heat maps along the boundary for the point estimates and associated pointwise robust bias-corrected $p$-values. Specifically:
\begin{itemize}
    \item Figure~\ref{fig:sa-heatmaps-itt-main} reports heatmaps for the reduced-form analysis of college enrollment.
    \item Figure~\ref{fig:sa-heatmaps-fs-main} reports heatmaps for the first-stage estimates for treatment take-up.
    \item Figure~\ref{fig:sa-heatmaps-fuzzy-main} reports heatmaps for the fuzzy analysis of college enrollment.
    \item Figure~\ref{fig:sa-heatmaps-covbal} reports heatmaps for the covariate-balance analysis using mother's education as the outcome.
\end{itemize}
These plots make it easier to see where along the boundary the estimated effects are largest and where the evidence against the zero-effect null is strongest.

In addition, we report baseline estimates and treatment effects relative to control-side baselines.
\begin{itemize}
    \item Table~\ref{tab:sa-empapp_main_itt0} reports the control-side baseline estimates $\widehat{\mu}_{Y,0}(\bx)$ for college enrollment.
    \item Table~\ref{tab:sa-empapp_main_percent} reports ITT and Fuzzy treatment effects relative to control-side baseline estimates $\widehat{\mu}_{Y,0}(\bx)$ for college enrollment.
    \item Figure~\ref{fig:sa-empapp_main_itt0} reports baseline estimates and uncertainty quantification.
    \item Figure~\ref{fig:sa-empapp_main_percent} reports treatment effects relative to the estimated control-side baseline.
\end{itemize}

The baseline estimates in Table~\ref{tab:sa-empapp_main_itt0} help interpret the magnitude of the treatment effects.

\clearpage

\begin{table}[p]
    \centering
\begin{tabular}{lccccccc}
  \toprule\toprule
   & $h_1$ & $h_2$ & $N_{\mathrm{Co}}$ & $N_{\mathrm{Tr}}$ & Estimate & $p$-value & 95\% CI \\
  \midrule
   $\bb_{1}$ & 27.6 & 12.2 & 8921 & 3532 & 0.311 & 0.000 & (0.204, 0.423)\\
   $\bb_{2}$ & 26.0 & 11.5 & 8998 & 3677 & 0.313 & 0.000 & (0.217, 0.430)\\
   $\bb_{3}$ & 24.6 & 10.9 & 8510 & 3678 & 0.316 & 0.000 & (0.218, 0.428)\\
   $\bb_{4}$ & 24.5 & 10.9 & 9236 & 3976 & 0.309 & 0.000 & (0.197, 0.399)\\
   $\bb_{5}$ & 23.1 & 10.2 & 9044 & 3997 & 0.295 & 0.000 & (0.164, 0.370)\\
   $\bb_{6}$ & 20.7 & 9.2 & 7176 & 3525 & 0.277 & 0.000 & (0.128, 0.350)\\
   $\bb_{7}$ & 19.0 & 8.4 & 6089 & 3209 & 0.267 & 0.000 & (0.118, 0.356)\\
   $\bb_{8}$ & 19.2 & 8.5 & 6725 & 3422 & 0.283 & 0.000 & (0.137, 0.370)\\
   $\bb_{9}$ & 19.2 & 8.5 & 6803 & 3516 & 0.305 & 0.000 & (0.159, 0.391)\\
   $\bb_{10}$ & 19.3 & 8.6 & 6925 & 3563 & 0.324 & 0.000 & (0.189, 0.417)\\
   $\bb_{11}$ & 21.1 & 9.3 & 8665 & 4206 & 0.332 & 0.000 & (0.219, 0.424)\\
   $\bb_{12}$ & 21.8 & 9.7 & 9017 & 4456 & 0.327 & 0.000 & (0.223, 0.420)\\
   $\bb_{13}$ & 21.1 & 9.3 & 8683 & 4417 & 0.322 & 0.000 & (0.209, 0.412)\\
   $\bb_{14}$ & 21.8 & 9.6 & 9069 & 4646 & 0.318 & 0.000 & (0.194, 0.388)\\
   $\bb_{15}$ & 22.2 & 9.8 & 10063 & 4977 & 0.314 & 0.000 & (0.190, 0.378)\\
   $\bb_{16}$ & 22.5 & 10.0 & 10210 & 5162 & 0.317 & 0.000 & (0.208, 0.390)\\
   $\bb_{17}$ & 21.3 & 9.4 & 9409 & 4444 & 0.323 & 0.000 & (0.251, 0.416)\\
   $\bb_{18}$ & 20.7 & 9.2 & 9182 & 3850 & 0.328 & 0.000 & (0.264, 0.427)\\
   $\bb_{19}$ & 19.3 & 8.5 & 8468 & 3157 & 0.334 & 0.000 & (0.263, 0.439)\\
   $\bb_{20}$ & 19.4 & 8.6 & 8945 & 2680 & 0.327 & 0.000 & (0.231, 0.433)\\
   $\bb_{21}$ & 20.9 & 9.3 & 10677 & 2442 & 0.314 & 0.000 & (0.157, 0.470)\\
   $\bb_{22}$ & 21.5 & 9.5 & 10221 & 2786 & 0.306 & 0.000 & (0.183, 0.418)\\
   $\bb_{23}$ & 20.3 & 9.0 & 7743 & 2679 & 0.298 & 0.000 & (0.186, 0.408)\\
   $\bb_{24}$ & 20.5 & 9.1 & 7097 & 2856 & 0.295 & 0.000 & (0.198, 0.413)\\
   $\bb_{25}$ & 21.0 & 9.3 & 6188 & 3114 & 0.291 & 0.000 & (0.211, 0.423)\\
   $\bb_{26}$ & 23.2 & 10.3 & 6948 & 3635 & 0.284 & 0.000 & (0.211, 0.414)\\
   $\bb_{27}$ & 25.1 & 11.1 & 7232 & 4098 & 0.273 & 0.000 & (0.195, 0.393)\\
   $\bb_{28}$ & 25.2 & 11.2 & 5984 & 4204 & 0.269 & 0.000 & (0.182, 0.391)\\
   $\bb_{29}$ & 25.6 & 11.3 & 5319 & 4343 & 0.266 & 0.000 & (0.171, 0.393)\\
   $\bb_{30}$ & 26.3 & 11.7 & 4698 & 4560 & 0.257 & 0.000 & (0.157, 0.398)\\
   $\bb_{31}$ & 28.5 & 12.6 & 4759 & 4995 & 0.247 & 0.000 & (0.139, 0.377)\\
   $\bb_{32}$ & 28.9 & 12.8 & 4482 & 4655 & 0.240 & 0.000 & (0.123, 0.370)\\
   $\bb_{33}$ & 30.8 & 13.7 & 4559 & 4766 & 0.234 & 0.000 & (0.115, 0.360)\\
   $\bb_{34}$ & 31.2 & 13.8 & 4048 & 4222 & 0.228 & 0.000 & (0.103, 0.357)\\
   $\bb_{35}$ & 28.1 & 12.5 & 2875 & 2950 & 0.219 & 0.000 & (0.077, 0.373)\\
   $\bb_{36}$ & 26.0 & 11.5 & 2214 & 2243 & 0.215 & 0.000 & (0.064, 0.394)\\
   $\bb_{37}$ & 26.8 & 11.9 & 1987 & 2002 & 0.209 & 0.000 & (0.059, 0.391)\\
   $\bb_{38}$ & 28.5 & 12.6 & 2049 & 2068 & 0.200 & 0.000 & (0.052, 0.379)\\
   $\bb_{39}$ & 33.6 & 14.9 & 2758 & 2787 & 0.191 & 0.000 & (0.053, 0.349)\\
   $\bb_{40}$ & 36.3 & 16.1 & 2916 & 2993 & 0.185 & 0.000 & (0.044, 0.335)\\
  \midrule
   $\mathtt{WBATE}$ &  &  &  &  & 0.282 & 0.000 & (0.248, 0.316)\\
  \midrule
   $\mathtt{LBATE}$ &  &  &  &  & 0.334 &  & (0.263, 0.472)\\
  \bottomrule\bottomrule
\end{tabular}

    \caption{Intention-to-Treat Effects (SPP Application, ITT Bandwidth Selector).}
    \label{tab:sa-empapp_main_itt}
\end{table}

\begin{table}[p]
    \centering
\begin{tabular}{lccccccc}
  \toprule\toprule
   & $h_1$ & $h_2$ & $N_{\mathrm{Co}}$ & $N_{\mathrm{Tr}}$ & Estimate & $p$-value & 95\% CI \\
  \midrule
   $\bb_{1}$ & 27.6 & 12.2 & 8921 & 3532 & 0.551 & 0.000 & (0.495, 0.655)\\
   $\bb_{2}$ & 26.0 & 11.5 & 8998 & 3677 & 0.546 & 0.000 & (0.496, 0.651)\\
   $\bb_{3}$ & 24.6 & 10.9 & 8510 & 3678 & 0.544 & 0.000 & (0.494, 0.646)\\
   $\bb_{4}$ & 24.5 & 10.9 & 9236 & 3976 & 0.536 & 0.000 & (0.483, 0.629)\\
   $\bb_{5}$ & 23.1 & 10.2 & 9044 & 3997 & 0.528 & 0.000 & (0.466, 0.612)\\
   $\bb_{6}$ & 20.7 & 9.2 & 7176 & 3525 & 0.517 & 0.000 & (0.441, 0.595)\\
   $\bb_{7}$ & 19.0 & 8.4 & 6089 & 3209 & 0.507 & 0.000 & (0.419, 0.582)\\
   $\bb_{8}$ & 19.2 & 8.5 & 6725 & 3422 & 0.507 & 0.000 & (0.407, 0.570)\\
   $\bb_{9}$ & 19.2 & 8.5 & 6803 & 3516 & 0.512 & 0.000 & (0.406, 0.571)\\
   $\bb_{10}$ & 19.3 & 8.6 & 6925 & 3563 & 0.524 & 0.000 & (0.417, 0.582)\\
   $\bb_{11}$ & 21.1 & 9.3 & 8665 & 4206 & 0.538 & 0.000 & (0.448, 0.599)\\
   $\bb_{12}$ & 21.8 & 9.7 & 9017 & 4456 & 0.542 & 0.000 & (0.463, 0.607)\\
   $\bb_{13}$ & 21.1 & 9.3 & 8683 & 4417 & 0.552 & 0.000 & (0.472, 0.616)\\
   $\bb_{14}$ & 21.8 & 9.6 & 9069 & 4646 & 0.562 & 0.000 & (0.480, 0.617)\\
   $\bb_{15}$ & 22.2 & 9.8 & 10063 & 4977 & 0.564 & 0.000 & (0.481, 0.614)\\
   $\bb_{16}$ & 22.5 & 10.0 & 10210 & 5162 & 0.568 & 0.000 & (0.493, 0.622)\\
   $\bb_{17}$ & 21.3 & 9.4 & 9409 & 4444 & 0.576 & 0.000 & (0.501, 0.636)\\
   $\bb_{18}$ & 20.7 & 9.2 & 9182 & 3850 & 0.577 & 0.000 & (0.498, 0.642)\\
   $\bb_{19}$ & 19.3 & 8.5 & 8468 & 3157 & 0.575 & 0.000 & (0.488, 0.654)\\
   $\bb_{20}$ & 19.4 & 8.6 & 8945 & 2680 & 0.574 & 0.000 & (0.467, 0.674)\\
   $\bb_{21}$ & 20.9 & 9.3 & 10677 & 2442 & 0.573 & 0.000 & (0.424, 0.762)\\
   $\bb_{22}$ & 21.5 & 9.5 & 10221 & 2786 & 0.576 & 0.000 & (0.469, 0.715)\\
   $\bb_{23}$ & 20.3 & 9.0 & 7743 & 2679 & 0.574 & 0.000 & (0.473, 0.696)\\
   $\bb_{24}$ & 20.5 & 9.1 & 7097 & 2856 & 0.574 & 0.000 & (0.472, 0.683)\\
   $\bb_{25}$ & 21.0 & 9.3 & 6188 & 3114 & 0.579 & 0.000 & (0.481, 0.684)\\
   $\bb_{26}$ & 23.2 & 10.3 & 6948 & 3635 & 0.587 & 0.000 & (0.490, 0.679)\\
   $\bb_{27}$ & 25.1 & 11.1 & 7232 & 4098 & 0.593 & 0.000 & (0.494, 0.674)\\
   $\bb_{28}$ & 25.2 & 11.2 & 5984 & 4204 & 0.600 & 0.000 & (0.499, 0.682)\\
   $\bb_{29}$ & 25.6 & 11.3 & 5319 & 4343 & 0.611 & 0.000 & (0.511, 0.697)\\
   $\bb_{30}$ & 26.3 & 11.7 & 4698 & 4560 & 0.623 & 0.000 & (0.522, 0.711)\\
   $\bb_{31}$ & 28.5 & 12.6 & 4759 & 4995 & 0.636 & 0.000 & (0.535, 0.720)\\
   $\bb_{32}$ & 28.9 & 12.8 & 4482 & 4655 & 0.649 & 0.000 & (0.547, 0.738)\\
   $\bb_{33}$ & 30.8 & 13.7 & 4559 & 4766 & 0.661 & 0.000 & (0.564, 0.751)\\
   $\bb_{34}$ & 31.2 & 13.8 & 4048 & 4222 & 0.674 & 0.000 & (0.578, 0.772)\\
   $\bb_{35}$ & 28.1 & 12.5 & 2875 & 2950 & 0.691 & 0.000 & (0.587, 0.811)\\
   $\bb_{36}$ & 26.0 & 11.5 & 2214 & 2243 & 0.710 & 0.000 & (0.601, 0.849)\\
   $\bb_{37}$ & 26.8 & 11.9 & 1987 & 2002 & 0.719 & 0.000 & (0.610, 0.861)\\
   $\bb_{38}$ & 28.5 & 12.6 & 2049 & 2068 & 0.722 & 0.000 & (0.611, 0.863)\\
   $\bb_{39}$ & 33.6 & 14.9 & 2758 & 2787 & 0.721 & 0.000 & (0.618, 0.852)\\
   $\bb_{40}$ & 36.3 & 16.1 & 2916 & 2993 & 0.723 & 0.000 & (0.621, 0.855)\\
  \midrule
   $\mathtt{WBATE}$ &  &  &  &  & 0.592 & 0.000 & (0.567, 0.620)\\
  \midrule
   $\mathtt{LBATE}$ &  &  &  &  & 0.723 &  & (0.624, 0.860)\\
  \bottomrule\bottomrule
\end{tabular}

    \caption{First-Stage Effects (SPP Application, ITT Bandwidth Selector).}
    \label{tab:sa-empapp_main_fs}
\end{table}

\begin{table}[p]
    \centering
\begin{tabular}{lccccccc}
  \toprule\toprule
   & $h_1$ & $h_2$ & $N_{\mathrm{Co}}$ & $N_{\mathrm{Tr}}$ & Estimate & $p$-value & 95\% CI \\
  \midrule
   $\bb_{1}$ & 27.6 & 12.2 & 8921 & 3532 & 0.564 & 0.000 & (0.374, 0.715)\\
   $\bb_{2}$ & 26.0 & 11.5 & 8998 & 3677 & 0.572 & 0.000 & (0.398, 0.730)\\
   $\bb_{3}$ & 24.6 & 10.9 & 8510 & 3678 & 0.581 & 0.000 & (0.403, 0.732)\\
   $\bb_{4}$ & 24.5 & 10.9 & 9236 & 3976 & 0.576 & 0.000 & (0.374, 0.697)\\
   $\bb_{5}$ & 23.1 & 10.2 & 9044 & 3997 & 0.559 & 0.000 & (0.323, 0.666)\\
   $\bb_{6}$ & 20.7 & 9.2 & 7176 & 3525 & 0.535 & 0.000 & (0.268, 0.654)\\
   $\bb_{7}$ & 19.0 & 8.4 & 6089 & 3209 & 0.527 & 0.000 & (0.259, 0.689)\\
   $\bb_{8}$ & 19.2 & 8.5 & 6725 & 3422 & 0.559 & 0.000 & (0.304, 0.734)\\
   $\bb_{9}$ & 19.2 & 8.5 & 6803 & 3516 & 0.595 & 0.000 & (0.350, 0.776)\\
   $\bb_{10}$ & 19.3 & 8.6 & 6925 & 3563 & 0.619 & 0.000 & (0.401, 0.812)\\
   $\bb_{11}$ & 21.1 & 9.3 & 8665 & 4206 & 0.618 & 0.000 & (0.437, 0.791)\\
   $\bb_{12}$ & 21.8 & 9.7 & 9017 & 4456 & 0.602 & 0.000 & (0.433, 0.767)\\
   $\bb_{13}$ & 21.1 & 9.3 & 8683 & 4417 & 0.584 & 0.000 & (0.401, 0.740)\\
   $\bb_{14}$ & 21.8 & 9.6 & 9069 & 4646 & 0.566 & 0.000 & (0.371, 0.692)\\
   $\bb_{15}$ & 22.2 & 9.8 & 10063 & 4977 & 0.556 & 0.000 & (0.362, 0.675)\\
   $\bb_{16}$ & 22.5 & 10.0 & 10210 & 5162 & 0.557 & 0.000 & (0.388, 0.686)\\
   $\bb_{17}$ & 21.3 & 9.4 & 9409 & 4444 & 0.560 & 0.000 & (0.458, 0.715)\\
   $\bb_{18}$ & 20.7 & 9.2 & 9182 & 3850 & 0.570 & 0.000 & (0.484, 0.727)\\
   $\bb_{19}$ & 19.3 & 8.5 & 8468 & 3157 & 0.580 & 0.000 & (0.484, 0.744)\\
   $\bb_{20}$ & 19.4 & 8.6 & 8945 & 2680 & 0.569 & 0.000 & (0.440, 0.723)\\
   $\bb_{21}$ & 20.9 & 9.3 & 10677 & 2442 & 0.548 & 0.000 & (0.338, 0.720)\\
   $\bb_{22}$ & 21.5 & 9.5 & 10221 & 2786 & 0.532 & 0.000 & (0.359, 0.655)\\
   $\bb_{23}$ & 20.3 & 9.0 & 7743 & 2679 & 0.519 & 0.000 & (0.357, 0.659)\\
   $\bb_{24}$ & 20.5 & 9.1 & 7097 & 2856 & 0.513 & 0.000 & (0.373, 0.686)\\
   $\bb_{25}$ & 21.0 & 9.3 & 6188 & 3114 & 0.503 & 0.000 & (0.385, 0.703)\\
   $\bb_{26}$ & 23.2 & 10.3 & 6948 & 3635 & 0.483 & 0.000 & (0.382, 0.689)\\
   $\bb_{27}$ & 25.1 & 11.1 & 7232 & 4098 & 0.461 & 0.000 & (0.352, 0.655)\\
   $\bb_{28}$ & 25.2 & 11.2 & 5984 & 4204 & 0.449 & 0.000 & (0.326, 0.643)\\
   $\bb_{29}$ & 25.6 & 11.3 & 5319 & 4343 & 0.435 & 0.000 & (0.302, 0.632)\\
   $\bb_{30}$ & 26.3 & 11.7 & 4698 & 4560 & 0.412 & 0.000 & (0.272, 0.627)\\
   $\bb_{31}$ & 28.5 & 12.6 & 4759 & 4995 & 0.388 & 0.000 & (0.238, 0.584)\\
   $\bb_{32}$ & 28.9 & 12.8 & 4482 & 4655 & 0.370 & 0.000 & (0.206, 0.561)\\
   $\bb_{33}$ & 30.8 & 13.7 & 4559 & 4766 & 0.354 & 0.000 & (0.188, 0.534)\\
   $\bb_{34}$ & 31.2 & 13.8 & 4048 & 4222 & 0.338 & 0.000 & (0.164, 0.517)\\
   $\bb_{35}$ & 28.1 & 12.5 & 2875 & 2950 & 0.317 & 0.000 & (0.123, 0.521)\\
   $\bb_{36}$ & 26.0 & 11.5 & 2214 & 2243 & 0.303 & 0.000 & (0.100, 0.532)\\
   $\bb_{37}$ & 26.8 & 11.9 & 1987 & 2002 & 0.291 & 0.000 & (0.090, 0.522)\\
   $\bb_{38}$ & 28.5 & 12.6 & 2049 & 2068 & 0.277 & 0.000 & (0.079, 0.506)\\
   $\bb_{39}$ & 33.6 & 14.9 & 2758 & 2787 & 0.265 & 0.000 & (0.079, 0.468)\\
   $\bb_{40}$ & 36.3 & 16.1 & 2916 & 2993 & 0.255 & 0.000 & (0.067, 0.447)\\
  \midrule
   $\mathtt{WBATE}$ &  &  &  &  & 0.487 & 0.000 & (0.435, 0.535)\\
  \midrule
   $\mathtt{LBATE}$ &  &  &  &  & 0.619 &  & (0.487, 0.808)\\
  \bottomrule\bottomrule
\end{tabular}

    \caption{Fuzzy Effects (SPP Application, ITT Bandwidth Selector).}
    \label{tab:sa-empapp_main_fuzzy}
\end{table}

\begin{table}[p]
    \centering
\begin{tabular}{lccccccc}
  \toprule\toprule
   & $h_1$ & $h_2$ & $N_{\mathrm{Co}}$ & $N_{\mathrm{Tr}}$ & Estimate & $p$-value & 95\% CI \\
  \midrule
   $\bb_{1}$ & 27.4 & 12.2 & 8750 & 3469 & -0.027 & 0.170 & (-0.153, 0.058)\\
   $\bb_{2}$ & 25.9 & 11.5 & 8454 & 3560 & -0.032 & 0.115 & (-0.158, 0.050)\\
   $\bb_{3}$ & 24.8 & 11.0 & 8447 & 3647 & -0.031 & 0.115 & (-0.157, 0.050)\\
   $\bb_{4}$ & 24.7 & 11.0 & 9184 & 3947 & -0.022 & 0.213 & (-0.139, 0.058)\\
   $\bb_{5}$ & 24.2 & 10.7 & 9767 & 4156 & -0.015 & 0.398 & (-0.123, 0.069)\\
   $\bb_{6}$ & 22.7 & 10.1 & 8801 & 3951 & -0.008 & 0.739 & (-0.108, 0.086)\\
   $\bb_{7}$ & 21.9 & 9.7 & 8282 & 3901 & -0.001 & 0.984 & (-0.099, 0.098)\\
   $\bb_{8}$ & 21.4 & 9.5 & 8372 & 3993 & -0.002 & 0.678 & (-0.115, 0.087)\\
   $\bb_{9}$ & 21.7 & 9.6 & 8680 & 4130 & -0.010 & 0.367 & (-0.130, 0.071)\\
   $\bb_{10}$ & 21.0 & 9.3 & 8085 & 3981 & -0.011 & 0.272 & (-0.139, 0.065)\\
   $\bb_{11}$ & 20.6 & 9.1 & 7894 & 3964 & -0.011 & 0.444 & (-0.129, 0.077)\\
   $\bb_{12}$ & 21.6 & 9.6 & 8832 & 4358 & -0.009 & 0.931 & (-0.100, 0.095)\\
   $\bb_{13}$ & 21.0 & 9.3 & 8098 & 4270 & -0.001 & 0.851 & (-0.094, 0.106)\\
   $\bb_{14}$ & 20.9 & 9.3 & 8189 & 4348 & 0.002 & 0.849 & (-0.091, 0.104)\\
   $\bb_{15}$ & 20.3 & 9.0 & 8051 & 4329 & 0.000 & 0.909 & (-0.093, 0.100)\\
   $\bb_{16}$ & 21.8 & 9.7 & 9040 & 4809 & 0.001 & 0.560 & (-0.071, 0.105)\\
   $\bb_{17}$ & 22.7 & 10.0 & 10699 & 4678 & 0.002 & 0.516 & (-0.057, 0.088)\\
   $\bb_{18}$ & 20.1 & 8.9 & 8788 & 3746 & -0.002 & 0.472 & (-0.058, 0.095)\\
   $\bb_{19}$ & 18.7 & 8.3 & 7730 & 3014 & -0.005 & 0.815 & (-0.085, 0.073)\\
   $\bb_{20}$ & 20.2 & 8.9 & 9828 & 2825 & -0.009 & 0.613 & (-0.103, 0.073)\\
   $\bb_{21}$ & 19.8 & 8.8 & 9415 & 2203 & -0.004 & 0.503 & (-0.118, 0.185)\\
   $\bb_{22}$ & 21.0 & 9.3 & 9846 & 2645 & 0.001 & 0.350 & (-0.076, 0.144)\\
   $\bb_{23}$ & 17.9 & 7.9 & 5971 & 2207 & 0.012 & 0.065 & (-0.043, 0.176)\\
   $\bb_{24}$ & 17.9 & 7.9 & 4992 & 2339 & 0.018 & 0.037 & (-0.034, 0.180)\\
   $\bb_{25}$ & 18.7 & 8.3 & 4697 & 2598 & 0.020 & 0.058 & (-0.039, 0.168)\\
   $\bb_{26}$ & 20.6 & 9.1 & 4929 & 3087 & 0.019 & 0.112 & (-0.046, 0.148)\\
   $\bb_{27}$ & 22.3 & 9.9 & 5182 & 3532 & 0.015 & 0.233 & (-0.057, 0.130)\\
   $\bb_{28}$ & 22.5 & 10.0 & 4222 & 3677 & 0.013 & 0.312 & (-0.065, 0.130)\\
   $\bb_{29}$ & 23.0 & 10.2 & 3755 & 3831 & 0.009 & 0.412 & (-0.077, 0.134)\\
   $\bb_{30}$ & 24.5 & 10.9 & 3842 & 3945 & 0.005 & 0.556 & (-0.085, 0.126)\\
   $\bb_{31}$ & 29.4 & 13.0 & 5150 & 5113 & 0.004 & 0.678 & (-0.079, 0.105)\\
   $\bb_{32}$ & 30.0 & 13.3 & 4746 & 4961 & 0.000 & 0.961 & (-0.094, 0.097)\\
   $\bb_{33}$ & 27.7 & 12.3 & 3589 & 3723 & -0.004 & 0.798 & (-0.117, 0.099)\\
   $\bb_{34}$ & 26.1 & 11.6 & 2715 & 2794 & -0.004 & 0.665 & (-0.135, 0.101)\\
   $\bb_{35}$ & 23.1 & 10.3 & 1825 & 1842 & -0.006 & 0.639 & (-0.156, 0.114)\\
   $\bb_{36}$ & 21.2 & 9.4 & 1396 & 1393 & -0.009 & 0.528 & (-0.172, 0.113)\\
   $\bb_{37}$ & 21.7 & 9.6 & 1252 & 1223 & -0.012 & 0.423 & (-0.177, 0.103)\\
   $\bb_{38}$ & 22.7 & 10.1 & 1304 & 1274 & -0.014 & 0.372 & (-0.173, 0.094)\\
   $\bb_{39}$ & 29.3 & 13.0 & 1944 & 1969 & -0.010 & 0.434 & (-0.134, 0.079)\\
   $\bb_{40}$ & 27.8 & 12.3 & 1560 & 1511 & -0.011 & 0.516 & (-0.134, 0.087)\\
  \midrule
   $\mathtt{WBATE}$ &  &  &  &  & -0.004 & 0.938 & (-0.030, 0.028)\\
  \midrule
   $\mathtt{LBATE}$ &  &  &  &  & 0.020 &  & (-0.031, 0.182)\\
  \bottomrule\bottomrule
\end{tabular}

    \caption{Covariate Balance for Mother's Education (SPP Application, ITT Bandwidth Selector).}
    \label{tab:sa-empapp_covbal_itt}
\end{table}

\clearpage

\begin{figure}
    \centering
    \begin{subfigure}[b]{0.70\textwidth}
        \centering
        \includegraphics[width=\linewidth]{figures/empapp_main_itt_heatmap_bwitt.png}
        \caption{Point estimates.}
    \end{subfigure}

    \begin{subfigure}[b]{0.70\textwidth}
        \centering
        \includegraphics[width=\linewidth]{figures/empapp_main_itt_heatmap_pval_bwitt.png}
        \caption{Pointwise $p$-values.}
    \end{subfigure}
    \caption{Intention-to-Treat Effects (SPP Application, ITT Bandwidth Selector).}
    \label{fig:sa-heatmaps-itt-main}
\end{figure}

\begin{figure}
    \centering
    \begin{subfigure}[b]{0.70\textwidth}
        \centering
        \includegraphics[width=\linewidth]{figures/empapp_main_fs_heatmap_bwitt.png}
        \caption{Point estimates.}
    \end{subfigure}

    \begin{subfigure}[b]{0.70\textwidth}
        \centering
        \includegraphics[width=\linewidth]{figures/empapp_main_fs_heatmap_pval_bwitt.png}
        \caption{Pointwise $p$-values.}
    \end{subfigure}
    \caption{First-Stage Effects (SPP Application, ITT Bandwidth Selector).}
    \label{fig:sa-heatmaps-fs-main}
\end{figure}

\begin{figure}
    \centering
    \begin{subfigure}[b]{0.70\textwidth}
        \centering
        \includegraphics[width=\linewidth]{figures/empapp_main_fuzzy_heatmap_bwitt.png}
        \caption{Point estimates.}
    \end{subfigure}

    \begin{subfigure}[b]{0.70\textwidth}
        \centering
        \includegraphics[width=\linewidth]{figures/empapp_main_fuzzy_heatmap_pval_bwitt.png}
        \caption{Pointwise $p$-values.}
    \end{subfigure}
    \caption{Fuzzy Effects (SPP Application, ITT Bandwidth Selector).}
    \label{fig:sa-heatmaps-fuzzy-main}
\end{figure}

\begin{figure}
    \centering
    \begin{subfigure}[b]{0.70\textwidth}
        \centering
        \includegraphics[width=\linewidth]{figures/empapp_covbal_itt_heatmap_bwitt.png}
        \caption{Point estimates.}
    \end{subfigure}

    \begin{subfigure}[b]{0.70\textwidth}
        \centering
        \includegraphics[width=\linewidth]{figures/empapp_covbal_itt_heatmap_pval_bwitt.png}
        \caption{Pointwise $p$-values.}
    \end{subfigure}
    \caption{Covariate Balance for Mother's Education (SPP Application, ITT Bandwidth Selector).}
    \label{fig:sa-heatmaps-covbal}
\end{figure}

\clearpage

\begin{table}[p]
    \centering
\begin{tabular}{lccccccc}
  \toprule\toprule
   & $h_1$ & $h_2$ & $N_{\mathrm{Co}}$ & $N_{\mathrm{Tr}}$ & Estimate & $p$-value & 95\% CI \\
  \midrule
   $\bb_{1}$ & 27.6 & 12.2 & 8921 &  & 0.371 & 0.000 & (0.295, 0.459)\\
   $\bb_{2}$ & 26.0 & 11.5 & 8998 &  & 0.364 & 0.000 & (0.286, 0.445)\\
   $\bb_{3}$ & 24.6 & 10.9 & 8510 &  & 0.356 & 0.000 & (0.283, 0.441)\\
   $\bb_{4}$ & 24.5 & 10.9 & 9236 &  & 0.355 & 0.000 & (0.293, 0.444)\\
   $\bb_{5}$ & 23.1 & 10.2 & 9044 &  & 0.356 & 0.000 & (0.298, 0.454)\\
   $\bb_{6}$ & 20.7 & 9.2 & 7176 &  & 0.359 & 0.000 & (0.294, 0.462)\\
   $\bb_{7}$ & 19.0 & 8.4 & 6089 &  & 0.356 & 0.000 & (0.271, 0.450)\\
   $\bb_{8}$ & 19.2 & 8.5 & 6725 &  & 0.346 & 0.000 & (0.263, 0.436)\\
   $\bb_{9}$ & 19.2 & 8.5 & 6803 &  & 0.339 & 0.000 & (0.261, 0.434)\\
   $\bb_{10}$ & 19.3 & 8.6 & 6925 &  & 0.340 & 0.000 & (0.264, 0.435)\\
   $\bb_{11}$ & 21.1 & 9.3 & 8665 &  & 0.349 & 0.000 & (0.281, 0.434)\\
   $\bb_{12}$ & 21.8 & 9.7 & 9017 &  & 0.363 & 0.000 & (0.298, 0.447)\\
   $\bb_{13}$ & 21.1 & 9.3 & 8683 &  & 0.378 & 0.000 & (0.317, 0.473)\\
   $\bb_{14}$ & 21.8 & 9.6 & 9069 &  & 0.388 & 0.000 & (0.343, 0.494)\\
   $\bb_{15}$ & 22.2 & 9.8 & 10063 &  & 0.396 & 0.000 & (0.355, 0.502)\\
   $\bb_{16}$ & 22.5 & 10.0 & 10210 &  & 0.398 & 0.000 & (0.351, 0.494)\\
   $\bb_{17}$ & 21.3 & 9.4 & 9409 &  & 0.399 & 0.000 & (0.343, 0.459)\\
   $\bb_{18}$ & 20.7 & 9.2 & 9182 &  & 0.393 & 0.000 & (0.338, 0.439)\\
   $\bb_{19}$ & 19.3 & 8.5 & 8468 &  & 0.391 & 0.000 & (0.340, 0.438)\\
   $\bb_{20}$ & 19.4 & 8.6 & 8945 &  & 0.394 & 0.000 & (0.359, 0.444)\\
   $\bb_{21}$ & 20.9 & 9.3 & 10677 &  & 0.399 & 0.000 & (0.373, 0.443)\\
   $\bb_{22}$ & 21.5 & 9.5 & 10221 &  & 0.411 & 0.000 & (0.381, 0.457)\\
   $\bb_{23}$ & 20.3 & 9.0 & 7743 &  & 0.421 & 0.000 & (0.373, 0.467)\\
   $\bb_{24}$ & 20.5 & 9.1 & 7097 &  & 0.431 & 0.000 & (0.365, 0.475)\\
   $\bb_{25}$ & 21.0 & 9.3 & 6188 &  & 0.443 & 0.000 & (0.364, 0.486)\\
   $\bb_{26}$ & 23.2 & 10.3 & 6948 &  & 0.457 & 0.000 & (0.376, 0.499)\\
   $\bb_{27}$ & 25.1 & 11.1 & 7232 &  & 0.471 & 0.000 & (0.390, 0.515)\\
   $\bb_{28}$ & 25.2 & 11.2 & 5984 &  & 0.481 & 0.000 & (0.395, 0.534)\\
   $\bb_{29}$ & 25.6 & 11.3 & 5319 &  & 0.491 & 0.000 & (0.397, 0.552)\\
   $\bb_{30}$ & 26.3 & 11.7 & 4698 &  & 0.506 & 0.000 & (0.393, 0.572)\\
   $\bb_{31}$ & 28.5 & 12.6 & 4759 &  & 0.524 & 0.000 & (0.415, 0.592)\\
   $\bb_{32}$ & 28.9 & 12.8 & 4482 &  & 0.539 & 0.000 & (0.430, 0.616)\\
   $\bb_{33}$ & 30.8 & 13.7 & 4559 &  & 0.555 & 0.000 & (0.449, 0.635)\\
   $\bb_{34}$ & 31.2 & 13.8 & 4048 &  & 0.571 & 0.000 & (0.463, 0.657)\\
   $\bb_{35}$ & 28.1 & 12.5 & 2875 &  & 0.587 & 0.000 & (0.463, 0.689)\\
   $\bb_{36}$ & 26.0 & 11.5 & 2214 &  & 0.605 & 0.000 & (0.469, 0.727)\\
   $\bb_{37}$ & 26.8 & 11.9 & 1987 &  & 0.622 & 0.000 & (0.487, 0.753)\\
   $\bb_{38}$ & 28.5 & 12.6 & 2049 &  & 0.640 & 0.000 & (0.506, 0.772)\\
   $\bb_{39}$ & 33.6 & 14.9 & 2758 &  & 0.656 & 0.000 & (0.540, 0.782)\\
   $\bb_{40}$ & 36.3 & 16.1 & 2916 &  & 0.670 & 0.000 & (0.562, 0.798)\\
  \midrule
   $\mathtt{WBATE}$ &  &  &  &  & 0.447 & 0.000 & (0.423, 0.473)\\
  \midrule
   $\mathtt{LBATE}$ &  &  &  &  & 0.670 &  & (0.561, 0.800)\\
  \bottomrule\bottomrule
\end{tabular}

    \caption{Control-Side Baseline Outcome (SPP Application, ITT Bandwidth Selector).}
    \label{tab:sa-empapp_main_itt0}
\end{table}

\begin{table}[p]
    \centering
\begin{tabular}{lcccccc}
  \toprule\toprule
   & $h_1$ & $h_2$ & $N_{\mathrm{Co}}$ & $N_{\mathrm{Tr}}$ & ITT / ITT.0 (\%) & Fuzzy / ITT.0 (\%) \\
  \midrule
   $\bb_{1}$ & 27.6 & 12.2 & 8921 & 3532 & 83.631 & 151.785\\
   $\bb_{2}$ & 26.0 & 11.5 & 8998 & 3677 & 85.929 & 157.241\\
   $\bb_{3}$ & 24.6 & 10.9 & 8510 & 3678 & 88.773 & 163.114\\
   $\bb_{4}$ & 24.5 & 10.9 & 9236 & 3976 & 87.145 & 162.517\\
   $\bb_{5}$ & 23.1 & 10.2 & 9044 & 3997 & 82.811 & 156.772\\
   $\bb_{6}$ & 20.7 & 9.2 & 7176 & 3525 & 76.957 & 148.838\\
   $\bb_{7}$ & 19.0 & 8.4 & 6089 & 3209 & 75.122 & 148.193\\
   $\bb_{8}$ & 19.2 & 8.5 & 6725 & 3422 & 81.842 & 161.505\\
   $\bb_{9}$ & 19.2 & 8.5 & 6803 & 3516 & 89.783 & 175.477\\
   $\bb_{10}$ & 19.3 & 8.6 & 6925 & 3563 & 95.242 & 181.913\\
   $\bb_{11}$ & 21.1 & 9.3 & 8665 & 4206 & 95.023 & 176.783\\
   $\bb_{12}$ & 21.8 & 9.7 & 9017 & 4456 & 89.931 & 165.864\\
   $\bb_{13}$ & 21.1 & 9.3 & 8683 & 4417 & 85.359 & 154.546\\
   $\bb_{14}$ & 21.8 & 9.6 & 9069 & 4646 & 82.049 & 146.036\\
   $\bb_{15}$ & 22.2 & 9.8 & 10063 & 4977 & 79.337 & 140.579\\
   $\bb_{16}$ & 22.5 & 10.0 & 10210 & 5162 & 79.586 & 139.992\\
   $\bb_{17}$ & 21.3 & 9.4 & 9409 & 4444 & 80.823 & 140.364\\
   $\bb_{18}$ & 20.7 & 9.2 & 9182 & 3850 & 83.501 & 144.816\\
   $\bb_{19}$ & 19.3 & 8.5 & 8468 & 3157 & 85.366 & 148.486\\
   $\bb_{20}$ & 19.4 & 8.6 & 8945 & 2680 & 82.952 & 144.496\\
   $\bb_{21}$ & 20.9 & 9.3 & 10677 & 2442 & 78.739 & 137.393\\
   $\bb_{22}$ & 21.5 & 9.5 & 10221 & 2786 & 74.537 & 129.406\\
   $\bb_{23}$ & 20.3 & 9.0 & 7743 & 2679 & 70.713 & 123.220\\
   $\bb_{24}$ & 20.5 & 9.1 & 7097 & 2856 & 68.327 & 119.007\\
   $\bb_{25}$ & 21.0 & 9.3 & 6188 & 3114 & 65.748 & 113.536\\
   $\bb_{26}$ & 23.2 & 10.3 & 6948 & 3635 & 62.000 & 105.621\\
   $\bb_{27}$ & 25.1 & 11.1 & 7232 & 4098 & 58.005 & 97.800\\
   $\bb_{28}$ & 25.2 & 11.2 & 5984 & 4204 & 56.030 & 93.353\\
   $\bb_{29}$ & 25.6 & 11.3 & 5319 & 4343 & 54.053 & 88.449\\
   $\bb_{30}$ & 26.3 & 11.7 & 4698 & 4560 & 50.756 & 81.532\\
   $\bb_{31}$ & 28.5 & 12.6 & 4759 & 4995 & 47.181 & 74.136\\
   $\bb_{32}$ & 28.9 & 12.8 & 4482 & 4655 & 44.580 & 68.662\\
   $\bb_{33}$ & 30.8 & 13.7 & 4559 & 4766 & 42.138 & 63.743\\
   $\bb_{34}$ & 31.2 & 13.8 & 4048 & 4222 & 39.864 & 59.187\\
   $\bb_{35}$ & 28.1 & 12.5 & 2875 & 2950 & 37.353 & 54.018\\
   $\bb_{36}$ & 26.0 & 11.5 & 2214 & 2243 & 35.513 & 50.042\\
   $\bb_{37}$ & 26.8 & 11.9 & 1987 & 2002 & 33.577 & 46.713\\
   $\bb_{38}$ & 28.5 & 12.6 & 2049 & 2068 & 31.279 & 43.332\\
   $\bb_{39}$ & 33.6 & 14.9 & 2758 & 2787 & 29.159 & 40.427\\
   $\bb_{40}$ & 36.3 & 16.1 & 2916 & 2993 & 27.554 & 38.108\\
  \midrule
   $\mathtt{WBATE}$ &  &  &  &  & 63.008 & 108.897\\
  \midrule
   $\mathtt{LBATE}$ &  &  &  &  & 49.787 & 92.377\\
  \bottomrule\bottomrule
\end{tabular}

    \caption{Treatment Effect Relative to Control-Side Baseline (SPP Application, ITT Bandwidth Selector).}
    \label{tab:sa-empapp_main_percent}
\end{table}

\begin{figure}
    \centering
    \begin{subfigure}[b]{0.70\textwidth}
        \centering
        \includegraphics[width=\linewidth]{figures/empapp_main_itt0_inference_bwitt.png}
        \caption{Point Estimates and Uncertainty Quantification.}
    \end{subfigure}

    \begin{subfigure}[b]{0.70\textwidth}
        \centering
        \includegraphics[width=\linewidth]{figures/empapp_main_itt0_heatmap_bwitt.png}
        \caption{Heatmap of Point Estimates.}
    \end{subfigure}
    \caption{Control-Side Baseline Outcome (SPP Application, ITT Bandwidth Selector).}
    \label{fig:sa-empapp_main_itt0}
\end{figure}

\begin{figure}
    \centering
    \begin{subfigure}[b]{0.70\textwidth}
        \centering
        \includegraphics[width=\linewidth]{figures/empapp_main_percent_itt_pointest_bwitt.png}
        \caption{Intention-to-Treat Effect.}
    \end{subfigure}

    \begin{subfigure}[b]{0.70\textwidth}
        \centering
        \includegraphics[width=\linewidth]{figures/empapp_main_percent_fuzzy_pointest_bwitt.png}
        \caption{Fuzzy Effect.}
    \end{subfigure}
    \caption{Treatment Effect Relative to Control-Side Baseline (SPP Application, ITT Bandwidth Selector).}
    \label{fig:sa-empapp_main_percent}
\end{figure}


\clearpage

\subsection{Fuzzy Bandwidth Selector}

This subsection reports the same empirical quantities using the bandwidth selected for the fuzzy Wald estimand. These results provide a bandwidth-sensitivity check for the main ITT-bandwidth empirical analysis. We provide the corresponding full boundary-grid results for all $40$ evaluation points.
\begin{itemize}
    \item Table~\ref{tab:sa-empapp_main_itt_bwmain} reports the reduced-form, or intention-to-treat (ITT), estimates for college enrollment.
    \item Table~\ref{tab:sa-empapp_main_fs_bwmain} reports the first-stage estimates for treatment take-up.
    \item Table~\ref{tab:sa-empapp_main_fuzzy_bwmain} reports the fuzzy estimates for college enrollment.
    \item Table~\ref{tab:sa-empapp_covbal_itt_bwmain} reports the covariate-balance reduced-form estimates using mother's education as the outcome.
\end{itemize}

The figures in this subsection complement the corresponding confidence intervals and bands plots for the fuzzy bandwidth selector. For each quantity, we report heat maps along the boundary for the point estimates and associated pointwise robust bias-corrected $p$-values. Specifically:
\begin{itemize}
    \item Figure~\ref{fig:sa-heatmaps-itt-main-bwmain} reports the reduced-form, or ITT, estimates for college enrollment.
    \item Figure~\ref{fig:sa-heatmaps-fs-main-bwmain} reports the first-stage estimates for treatment take-up.
    \item Figure~\ref{fig:sa-heatmaps-fuzzy-main-bwmain} reports the fuzzy estimates for college enrollment.
    \item Figure~\ref{fig:sa-heatmaps-covbal-bwmain} reports the covariate-balance reduced-form estimates using mother's education as the outcome.
\end{itemize}
These plots make it easier to inspect how the estimated effects and pointwise significance patterns vary along the boundary.

In addition, we report baseline estimates and treatment effects relative to control-side baselines.
\begin{itemize}
    \item Table~\ref{tab:sa-empapp_main_itt0_bwmain} reports the control-side baseline estimates $\widehat{\mu}_{Y,0}(\bx)$ for college enrollment.
    \item Table~\ref{tab:sa-empapp_main_percent_bwmain} reports ITT and Fuzzy treatment effects relative to control-side baseline estimates $\widehat{\mu}_{Y,0}(\bx)$ for college enrollment.
    \item Figure~\ref{fig:sa-empapp_main_itt0_bwmain} reports baseline estimates and uncertainty quantification.
    \item Figure~\ref{fig:sa-empapp_main_percent_bwmain} reports treatment effects relative to the estimated control-side baseline.
\end{itemize}
The baseline estimates in Table~\ref{tab:sa-empapp_main_itt0_bwmain} help interpret the economic magnitude of the effects, while Table~\ref{tab:sa-empapp_main_percent_bwmain} reports the same treatment effects in percentage terms relative to the local control-side baseline.

\clearpage

\begin{table}[p]
    \centering
\begin{tabular}{lccccccc}
  \toprule\toprule
   & $h_1$ & $h_2$ & $N_{\mathrm{Co}}$ & $N_{\mathrm{Tr}}$ & Estimate & $p$-value & 95\% CI \\
  \midrule
   $\bb_{1}$ & 28.2 & 12.5 & 9704 & 3687 & 0.311 & 0.000 & (0.205, 0.418)\\
   $\bb_{2}$ & 26.0 & 11.5 & 8572 & 3612 & 0.313 & 0.000 & (0.218, 0.429)\\
   $\bb_{3}$ & 24.9 & 11.0 & 8590 & 3720 & 0.316 & 0.000 & (0.220, 0.425)\\
   $\bb_{4}$ & 23.2 & 10.3 & 8397 & 3724 & 0.308 & 0.000 & (0.190, 0.401)\\
   $\bb_{5}$ & 21.8 & 9.6 & 7403 & 3568 & 0.291 & 0.000 & (0.154, 0.371)\\
   $\bb_{6}$ & 20.0 & 8.9 & 6977 & 3427 & 0.273 & 0.000 & (0.122, 0.351)\\
   $\bb_{7}$ & 19.7 & 8.7 & 6633 & 3404 & 0.271 & 0.000 & (0.126, 0.352)\\
   $\bb_{8}$ & 21.3 & 9.4 & 8437 & 4032 & 0.289 & 0.000 & (0.156, 0.362)\\
   $\bb_{9}$ & 21.9 & 9.7 & 8859 & 4218 & 0.306 & 0.000 & (0.178, 0.378)\\
   $\bb_{10}$ & 22.2 & 9.8 & 9828 & 4469 & 0.321 & 0.000 & (0.204, 0.399)\\
   $\bb_{11}$ & 22.3 & 9.9 & 9906 & 4584 & 0.330 & 0.000 & (0.223, 0.415)\\
   $\bb_{12}$ & 22.3 & 9.9 & 9949 & 4710 & 0.327 & 0.000 & (0.224, 0.414)\\
   $\bb_{13}$ & 20.6 & 9.1 & 7995 & 4229 & 0.322 & 0.000 & (0.208, 0.412)\\
   $\bb_{14}$ & 20.7 & 9.2 & 8168 & 4350 & 0.315 & 0.000 & (0.188, 0.390)\\
   $\bb_{15}$ & 21.0 & 9.3 & 8370 & 4493 & 0.309 & 0.000 & (0.182, 0.380)\\
   $\bb_{16}$ & 20.8 & 9.2 & 8278 & 4560 & 0.312 & 0.000 & (0.201, 0.398)\\
   $\bb_{17}$ & 19.7 & 8.7 & 7544 & 4057 & 0.319 & 0.000 & (0.241, 0.428)\\
   $\bb_{18}$ & 19.6 & 8.7 & 7997 & 3623 & 0.326 & 0.000 & (0.262, 0.431)\\
   $\bb_{19}$ & 20.4 & 9.0 & 9592 & 3367 & 0.333 & 0.000 & (0.266, 0.432)\\
   $\bb_{20}$ & 20.0 & 8.9 & 9225 & 2755 & 0.327 & 0.000 & (0.234, 0.430)\\
   $\bb_{21}$ & 20.9 & 9.3 & 10677 & 2442 & 0.314 & 0.000 & (0.159, 0.469)\\
   $\bb_{22}$ & 22.0 & 9.8 & 10830 & 2843 & 0.307 & 0.000 & (0.186, 0.414)\\
   $\bb_{23}$ & 21.3 & 9.4 & 8679 & 2871 & 0.298 & 0.000 & (0.190, 0.401)\\
   $\bb_{24}$ & 21.7 & 9.6 & 8114 & 3140 & 0.294 & 0.000 & (0.201, 0.403)\\
   $\bb_{25}$ & 21.8 & 9.6 & 6858 & 3223 & 0.289 & 0.000 & (0.212, 0.416)\\
   $\bb_{26}$ & 24.3 & 10.8 & 7784 & 3867 & 0.282 & 0.000 & (0.209, 0.401)\\
   $\bb_{27}$ & 27.4 & 12.1 & 8856 & 4564 & 0.271 & 0.000 & (0.190, 0.370)\\
   $\bb_{28}$ & 26.9 & 11.9 & 7294 & 4540 & 0.267 & 0.000 & (0.179, 0.373)\\
   $\bb_{29}$ & 26.7 & 11.8 & 5916 & 4593 & 0.264 & 0.000 & (0.170, 0.377)\\
   $\bb_{30}$ & 27.3 & 12.1 & 5250 & 4759 & 0.256 & 0.000 & (0.155, 0.378)\\
   $\bb_{31}$ & 30.1 & 13.3 & 5758 & 5285 & 0.247 & 0.000 & (0.141, 0.356)\\
   $\bb_{32}$ & 31.8 & 14.1 & 5266 & 5607 & 0.240 & 0.000 & (0.136, 0.359)\\
   $\bb_{33}$ & 32.1 & 14.2 & 4925 & 5218 & 0.234 & 0.000 & (0.122, 0.355)\\
   $\bb_{34}$ & 31.7 & 14.0 & 4323 & 4517 & 0.228 & 0.000 & (0.107, 0.355)\\
   $\bb_{35}$ & 28.8 & 12.7 & 3022 & 3135 & 0.220 & 0.000 & (0.083, 0.368)\\
   $\bb_{36}$ & 27.0 & 12.0 & 2452 & 2467 & 0.215 & 0.000 & (0.072, 0.386)\\
   $\bb_{37}$ & 28.1 & 12.4 & 2318 & 2326 & 0.208 & 0.000 & (0.069, 0.383)\\
   $\bb_{38}$ & 29.7 & 13.1 & 2330 & 2353 & 0.200 & 0.000 & (0.058, 0.370)\\
   $\bb_{39}$ & 35.2 & 15.6 & 2954 & 3036 & 0.191 & 0.000 & (0.061, 0.341)\\
   $\bb_{40}$ & 39.6 & 17.5 & 3538 & 3725 & 0.185 & 0.000 & (0.057, 0.322)\\
  \midrule
   $\mathtt{WBATE}$ &  &  &  &  & 0.281 & 0.000 & (0.246, 0.313)\\
  \midrule
   $\mathtt{LBATE}$ &  &  &  &  & 0.333 &  & (0.264, 0.473)\\
  \bottomrule\bottomrule
\end{tabular}

    \caption{Intention-to-Treat Effects (SPP Application, Fuzzy Bandwidth Selector).}
    \label{tab:sa-empapp_main_itt_bwmain}
\end{table}

\begin{table}[p]
    \centering
\begin{tabular}{lccccccc}
  \toprule\toprule
   & $h_1$ & $h_2$ & $N_{\mathrm{Co}}$ & $N_{\mathrm{Tr}}$ & Estimate & $p$-value & 95\% CI \\
  \midrule
   $\bb_{1}$ & 28.2 & 12.5 & 9704 & 3687 & 0.550 & 0.000 & (0.496, 0.650)\\
   $\bb_{2}$ & 26.0 & 11.5 & 8572 & 3612 & 0.547 & 0.000 & (0.498, 0.650)\\
   $\bb_{3}$ & 24.9 & 11.0 & 8590 & 3720 & 0.543 & 0.000 & (0.495, 0.644)\\
   $\bb_{4}$ & 23.2 & 10.3 & 8397 & 3724 & 0.539 & 0.000 & (0.481, 0.630)\\
   $\bb_{5}$ & 21.8 & 9.6 & 7403 & 3568 & 0.529 & 0.000 & (0.463, 0.613)\\
   $\bb_{6}$ & 20.0 & 8.9 & 6977 & 3427 & 0.516 & 0.000 & (0.440, 0.594)\\
   $\bb_{7}$ & 19.7 & 8.7 & 6633 & 3404 & 0.509 & 0.000 & (0.424, 0.579)\\
   $\bb_{8}$ & 21.3 & 9.4 & 8437 & 4032 & 0.510 & 0.000 & (0.422, 0.569)\\
   $\bb_{9}$ & 21.9 & 9.7 & 8859 & 4218 & 0.514 & 0.000 & (0.419, 0.564)\\
   $\bb_{10}$ & 22.2 & 9.8 & 9828 & 4469 & 0.526 & 0.000 & (0.431, 0.575)\\
   $\bb_{11}$ & 22.3 & 9.9 & 9906 & 4584 & 0.537 & 0.000 & (0.452, 0.594)\\
   $\bb_{12}$ & 22.3 & 9.9 & 9949 & 4710 & 0.542 & 0.000 & (0.465, 0.604)\\
   $\bb_{13}$ & 20.6 & 9.1 & 7995 & 4229 & 0.552 & 0.000 & (0.471, 0.615)\\
   $\bb_{14}$ & 20.7 & 9.2 & 8168 & 4350 & 0.561 & 0.000 & (0.477, 0.616)\\
   $\bb_{15}$ & 21.0 & 9.3 & 8370 & 4493 & 0.563 & 0.000 & (0.478, 0.614)\\
   $\bb_{16}$ & 20.8 & 9.2 & 8278 & 4560 & 0.567 & 0.000 & (0.489, 0.624)\\
   $\bb_{17}$ & 19.7 & 8.7 & 7544 & 4057 & 0.575 & 0.000 & (0.499, 0.638)\\
   $\bb_{18}$ & 19.6 & 8.7 & 7997 & 3623 & 0.575 & 0.000 & (0.497, 0.643)\\
   $\bb_{19}$ & 20.4 & 9.0 & 9592 & 3367 & 0.576 & 0.000 & (0.493, 0.651)\\
   $\bb_{20}$ & 20.0 & 8.9 & 9225 & 2755 & 0.574 & 0.000 & (0.471, 0.670)\\
   $\bb_{21}$ & 20.9 & 9.3 & 10677 & 2442 & 0.573 & 0.000 & (0.428, 0.759)\\
   $\bb_{22}$ & 22.0 & 9.8 & 10830 & 2843 & 0.576 & 0.000 & (0.470, 0.709)\\
   $\bb_{23}$ & 21.3 & 9.4 & 8679 & 2871 & 0.574 & 0.000 & (0.477, 0.689)\\
   $\bb_{24}$ & 21.7 & 9.6 & 8114 & 3140 & 0.575 & 0.000 & (0.478, 0.676)\\
   $\bb_{25}$ & 21.8 & 9.6 & 6858 & 3223 & 0.580 & 0.000 & (0.486, 0.679)\\
   $\bb_{26}$ & 24.3 & 10.8 & 7784 & 3867 & 0.588 & 0.000 & (0.494, 0.672)\\
   $\bb_{27}$ & 27.4 & 12.1 & 8856 & 4564 & 0.597 & 0.000 & (0.498, 0.663)\\
   $\bb_{28}$ & 26.9 & 11.9 & 7294 & 4540 & 0.603 & 0.000 & (0.504, 0.675)\\
   $\bb_{29}$ & 26.7 & 11.8 & 5916 & 4593 & 0.613 & 0.000 & (0.515, 0.691)\\
   $\bb_{30}$ & 27.3 & 12.1 & 5250 & 4759 & 0.624 & 0.000 & (0.526, 0.705)\\
   $\bb_{31}$ & 30.1 & 13.3 & 5758 & 5285 & 0.637 & 0.000 & (0.541, 0.712)\\
   $\bb_{32}$ & 31.8 & 14.1 & 5266 & 5607 & 0.649 & 0.000 & (0.557, 0.726)\\
   $\bb_{33}$ & 32.1 & 14.2 & 4925 & 5218 & 0.660 & 0.000 & (0.568, 0.744)\\
   $\bb_{34}$ & 31.7 & 14.0 & 4323 & 4517 & 0.673 & 0.000 & (0.581, 0.768)\\
   $\bb_{35}$ & 28.8 & 12.7 & 3022 & 3135 & 0.691 & 0.000 & (0.591, 0.805)\\
   $\bb_{36}$ & 27.0 & 12.0 & 2452 & 2467 & 0.707 & 0.000 & (0.606, 0.840)\\
   $\bb_{37}$ & 28.1 & 12.4 & 2318 & 2326 & 0.716 & 0.000 & (0.617, 0.853)\\
   $\bb_{38}$ & 29.7 & 13.1 & 2330 & 2353 & 0.720 & 0.000 & (0.618, 0.856)\\
   $\bb_{39}$ & 35.2 & 15.6 & 2954 & 3036 & 0.719 & 0.000 & (0.626, 0.845)\\
   $\bb_{40}$ & 39.6 & 17.5 & 3538 & 3725 & 0.719 & 0.000 & (0.635, 0.845)\\
  \midrule
   $\mathtt{WBATE}$ &  &  &  &  & 0.593 & 0.000 & (0.567, 0.619)\\
  \midrule
   $\mathtt{LBATE}$ &  &  &  &  & 0.720 &  & (0.636, 0.855)\\
  \bottomrule\bottomrule
\end{tabular}

    \caption{First-Stage Effects (SPP Application, Fuzzy Bandwidth Selector).}
    \label{tab:sa-empapp_main_fs_bwmain}
\end{table}

\begin{table}[p]
    \centering
\begin{tabular}{lccccccc}
  \toprule\toprule
   & $h_1$ & $h_2$ & $N_{\mathrm{Co}}$ & $N_{\mathrm{Tr}}$ & Estimate & $p$-value & 95\% CI \\
  \midrule
   $\bb_{1}$ & 28.2 & 12.5 & 9704 & 3687 & 0.565 & 0.000 & (0.381, 0.708)\\
   $\bb_{2}$ & 26.0 & 11.5 & 8572 & 3612 & 0.572 & 0.000 & (0.401, 0.727)\\
   $\bb_{3}$ & 24.9 & 11.0 & 8590 & 3720 & 0.581 & 0.000 & (0.407, 0.726)\\
   $\bb_{4}$ & 23.2 & 10.3 & 8397 & 3724 & 0.572 & 0.000 & (0.365, 0.700)\\
   $\bb_{5}$ & 21.8 & 9.6 & 7403 & 3568 & 0.551 & 0.000 & (0.308, 0.667)\\
   $\bb_{6}$ & 20.0 & 8.9 & 6977 & 3427 & 0.528 & 0.000 & (0.261, 0.655)\\
   $\bb_{7}$ & 19.7 & 8.7 & 6633 & 3404 & 0.533 & 0.000 & (0.275, 0.679)\\
   $\bb_{8}$ & 21.3 & 9.4 & 8437 & 4032 & 0.567 & 0.000 & (0.338, 0.708)\\
   $\bb_{9}$ & 21.9 & 9.7 & 8859 & 4218 & 0.595 & 0.000 & (0.387, 0.746)\\
   $\bb_{10}$ & 22.2 & 9.8 & 9828 & 4469 & 0.610 & 0.000 & (0.428, 0.771)\\
   $\bb_{11}$ & 22.3 & 9.9 & 9906 & 4584 & 0.613 & 0.000 & (0.446, 0.772)\\
   $\bb_{12}$ & 22.3 & 9.9 & 9949 & 4710 & 0.602 & 0.000 & (0.438, 0.756)\\
   $\bb_{13}$ & 20.6 & 9.1 & 7995 & 4229 & 0.583 & 0.000 & (0.401, 0.741)\\
   $\bb_{14}$ & 20.7 & 9.2 & 8168 & 4350 & 0.562 & 0.000 & (0.363, 0.695)\\
   $\bb_{15}$ & 21.0 & 9.3 & 8370 & 4493 & 0.549 & 0.000 & (0.351, 0.677)\\
   $\bb_{16}$ & 20.8 & 9.2 & 8278 & 4560 & 0.549 & 0.000 & (0.378, 0.698)\\
   $\bb_{17}$ & 19.7 & 8.7 & 7544 & 4057 & 0.556 & 0.000 & (0.441, 0.736)\\
   $\bb_{18}$ & 19.6 & 8.7 & 7997 & 3623 & 0.567 & 0.000 & (0.482, 0.734)\\
   $\bb_{19}$ & 20.4 & 9.0 & 9592 & 3367 & 0.579 & 0.000 & (0.490, 0.731)\\
   $\bb_{20}$ & 20.0 & 8.9 & 9225 & 2755 & 0.569 & 0.000 & (0.446, 0.717)\\
   $\bb_{21}$ & 20.9 & 9.3 & 10677 & 2442 & 0.548 & 0.000 & (0.343, 0.715)\\
   $\bb_{22}$ & 22.0 & 9.8 & 10830 & 2843 & 0.533 & 0.000 & (0.365, 0.652)\\
   $\bb_{23}$ & 21.3 & 9.4 & 8679 & 2871 & 0.519 & 0.000 & (0.364, 0.649)\\
   $\bb_{24}$ & 21.7 & 9.6 & 8114 & 3140 & 0.510 & 0.000 & (0.378, 0.667)\\
   $\bb_{25}$ & 21.8 & 9.6 & 6858 & 3223 & 0.499 & 0.000 & (0.389, 0.689)\\
   $\bb_{26}$ & 24.3 & 10.8 & 7784 & 3867 & 0.479 & 0.000 & (0.380, 0.666)\\
   $\bb_{27}$ & 27.4 & 12.1 & 8856 & 4564 & 0.455 & 0.000 & (0.347, 0.618)\\
   $\bb_{28}$ & 26.9 & 11.9 & 7294 & 4540 & 0.443 & 0.000 & (0.324, 0.613)\\
   $\bb_{29}$ & 26.7 & 11.8 & 5916 & 4593 & 0.430 & 0.000 & (0.301, 0.606)\\
   $\bb_{30}$ & 27.3 & 12.1 & 5250 & 4759 & 0.410 & 0.000 & (0.271, 0.596)\\
   $\bb_{31}$ & 30.1 & 13.3 & 5758 & 5285 & 0.388 & 0.000 & (0.242, 0.551)\\
   $\bb_{32}$ & 31.8 & 14.1 & 5266 & 5607 & 0.370 & 0.000 & (0.228, 0.544)\\
   $\bb_{33}$ & 32.1 & 14.2 & 4925 & 5218 & 0.354 & 0.000 & (0.201, 0.526)\\
   $\bb_{34}$ & 31.7 & 14.0 & 4323 & 4517 & 0.338 & 0.000 & (0.172, 0.512)\\
   $\bb_{35}$ & 28.8 & 12.7 & 3022 & 3135 & 0.319 & 0.000 & (0.133, 0.513)\\
   $\bb_{36}$ & 27.0 & 12.0 & 2452 & 2467 & 0.304 & 0.000 & (0.113, 0.520)\\
   $\bb_{37}$ & 28.1 & 12.4 & 2318 & 2326 & 0.291 & 0.000 & (0.105, 0.510)\\
   $\bb_{38}$ & 29.7 & 13.1 & 2330 & 2353 & 0.278 & 0.000 & (0.089, 0.492)\\
   $\bb_{39}$ & 35.2 & 15.6 & 2954 & 3036 & 0.266 & 0.000 & (0.091, 0.455)\\
   $\bb_{40}$ & 39.6 & 17.5 & 3538 & 3725 & 0.257 & 0.000 & (0.086, 0.426)\\
  \midrule
   $\mathtt{WBATE}$ &  &  &  &  & 0.485 & 0.000 & (0.433, 0.531)\\
  \midrule
   $\mathtt{LBATE}$ &  &  &  &  & 0.613 &  & (0.484, 0.781)\\
  \bottomrule\bottomrule
\end{tabular}

    \caption{Fuzzy Effects (SPP Application, Fuzzy Bandwidth Selector).}
    \label{tab:sa-empapp_main_fuzzy_bwmain}
\end{table}

\begin{table}[p]
    \centering
\begin{tabular}{lccccccc}
  \toprule\toprule
   & $h_1$ & $h_2$ & $N_{\mathrm{Co}}$ & $N_{\mathrm{Tr}}$ & Estimate & $p$-value & 95\% CI \\
  \midrule
   $\bb_{1}$ & 27.2 & 12.1 & 8688 & 3446 & -0.027 & 0.172 & (-0.153, 0.058)\\
   $\bb_{2}$ & 25.8 & 11.5 & 8428 & 3552 & -0.032 & 0.116 & (-0.157, 0.049)\\
   $\bb_{3}$ & 25.4 & 11.3 & 9246 & 3841 & -0.030 & 0.108 & (-0.153, 0.046)\\
   $\bb_{4}$ & 25.4 & 11.3 & 10066 & 4154 & -0.021 & 0.215 & (-0.134, 0.056)\\
   $\bb_{5}$ & 24.4 & 10.8 & 9822 & 4176 & -0.015 & 0.402 & (-0.121, 0.068)\\
   $\bb_{6}$ & 22.7 & 10.1 & 8801 & 3951 & -0.008 & 0.739 & (-0.107, 0.086)\\
   $\bb_{7}$ & 21.9 & 9.7 & 8262 & 3891 & -0.001 & 0.985 & (-0.098, 0.097)\\
   $\bb_{8}$ & 21.4 & 9.5 & 8372 & 3993 & -0.002 & 0.678 & (-0.114, 0.086)\\
   $\bb_{9}$ & 21.6 & 9.6 & 8604 & 4098 & -0.010 & 0.356 & (-0.131, 0.070)\\
   $\bb_{10}$ & 20.9 & 9.3 & 8049 & 3966 & -0.012 & 0.262 & (-0.140, 0.064)\\
   $\bb_{11}$ & 20.7 & 9.2 & 7952 & 3986 & -0.011 & 0.456 & (-0.127, 0.076)\\
   $\bb_{12}$ & 21.5 & 9.5 & 8766 & 4328 & -0.008 & 0.935 & (-0.100, 0.095)\\
   $\bb_{13}$ & 21.0 & 9.3 & 8090 & 4265 & -0.001 & 0.850 & (-0.093, 0.106)\\
   $\bb_{14}$ & 21.0 & 9.3 & 8207 & 4358 & 0.002 & 0.848 & (-0.090, 0.103)\\
   $\bb_{15}$ & 20.4 & 9.0 & 8059 & 4332 & 0.000 & 0.908 & (-0.092, 0.099)\\
   $\bb_{16}$ & 21.6 & 9.6 & 8936 & 4750 & 0.001 & 0.550 & (-0.071, 0.106)\\
   $\bb_{17}$ & 22.5 & 10.0 & 10627 & 4667 & 0.002 & 0.510 & (-0.057, 0.088)\\
   $\bb_{18}$ & 20.1 & 8.9 & 8788 & 3746 & -0.002 & 0.472 & (-0.058, 0.094)\\
   $\bb_{19}$ & 18.7 & 8.3 & 7730 & 3014 & -0.005 & 0.816 & (-0.084, 0.072)\\
   $\bb_{20}$ & 20.1 & 8.9 & 9828 & 2825 & -0.009 & 0.613 & (-0.102, 0.073)\\
   $\bb_{21}$ & 19.9 & 8.8 & 9436 & 2206 & -0.004 & 0.505 & (-0.117, 0.184)\\
   $\bb_{22}$ & 20.8 & 9.2 & 9760 & 2625 & 0.001 & 0.337 & (-0.075, 0.145)\\
   $\bb_{23}$ & 17.8 & 7.9 & 5890 & 2194 & 0.013 & 0.061 & (-0.041, 0.177)\\
   $\bb_{24}$ & 17.7 & 7.9 & 4911 & 2316 & 0.019 & 0.035 & (-0.032, 0.182)\\
   $\bb_{25}$ & 18.5 & 8.2 & 4641 & 2566 & 0.021 & 0.056 & (-0.038, 0.169)\\
   $\bb_{26}$ & 20.2 & 9.0 & 4845 & 3018 & 0.020 & 0.114 & (-0.046, 0.149)\\
   $\bb_{27}$ & 22.0 & 9.8 & 5110 & 3458 & 0.015 & 0.238 & (-0.057, 0.130)\\
   $\bb_{28}$ & 22.5 & 10.0 & 4215 & 3666 & 0.013 & 0.312 & (-0.065, 0.130)\\
   $\bb_{29}$ & 23.0 & 10.2 & 3754 & 3828 & 0.009 & 0.413 & (-0.076, 0.133)\\
   $\bb_{30}$ & 24.4 & 10.8 & 3836 & 3936 & 0.005 & 0.556 & (-0.084, 0.125)\\
   $\bb_{31}$ & 29.3 & 13.0 & 5139 & 5095 & 0.004 & 0.683 & (-0.080, 0.104)\\
   $\bb_{32}$ & 30.0 & 13.3 & 4744 & 4958 & 0.000 & 0.963 & (-0.093, 0.096)\\
   $\bb_{33}$ & 27.7 & 12.3 & 3587 & 3720 & -0.004 & 0.799 & (-0.116, 0.098)\\
   $\bb_{34}$ & 26.1 & 11.6 & 2714 & 2786 & -0.004 & 0.667 & (-0.134, 0.101)\\
   $\bb_{35}$ & 23.1 & 10.3 & 1825 & 1842 & -0.006 & 0.639 & (-0.155, 0.113)\\
   $\bb_{36}$ & 21.2 & 9.4 & 1396 & 1393 & -0.009 & 0.528 & (-0.171, 0.112)\\
   $\bb_{37}$ & 21.9 & 9.7 & 1350 & 1316 & -0.012 & 0.429 & (-0.173, 0.101)\\
   $\bb_{38}$ & 23.1 & 10.2 & 1313 & 1295 & -0.013 & 0.382 & (-0.168, 0.093)\\
   $\bb_{39}$ & 29.5 & 13.1 & 2021 & 2033 & -0.010 & 0.423 & (-0.133, 0.077)\\
   $\bb_{40}$ & 27.8 & 12.3 & 1556 & 1506 & -0.011 & 0.516 & (-0.133, 0.086)\\
  \midrule
   $\mathtt{WBATE}$ &  &  &  &  & -0.004 & 0.951 & (-0.030, 0.028)\\
  \midrule
   $\mathtt{LBATE}$ &  &  &  &  & 0.021 &  & (-0.033, 0.185)\\
  \bottomrule\bottomrule
\end{tabular}

    \caption{Covariate Balance for Mother's Education (SPP Application, Fuzzy Bandwidth Selector).}
    \label{tab:sa-empapp_covbal_itt_bwmain}
\end{table}

\clearpage

\begin{figure}
    \centering
    \begin{subfigure}[b]{0.70\textwidth}
        \centering
        \includegraphics[width=\linewidth]{figures/empapp_main_itt_heatmap_bwmain.png}
        \caption{Point estimates.}
    \end{subfigure}

    \begin{subfigure}[b]{0.70\textwidth}
        \centering
        \includegraphics[width=\linewidth]{figures/empapp_main_itt_heatmap_pval_bwmain.png}
        \caption{Pointwise $p$-values.}
    \end{subfigure}
    \caption{Intention-to-Treat Effects (SPP Application, Fuzzy Bandwidth Selector).}
    \label{fig:sa-heatmaps-itt-main-bwmain}
\end{figure}

\begin{figure}
    \centering
    \begin{subfigure}[b]{0.70\textwidth}
        \centering
        \includegraphics[width=\linewidth]{figures/empapp_main_fs_heatmap_bwmain.png}
        \caption{Point estimates.}
    \end{subfigure}

    \begin{subfigure}[b]{0.70\textwidth}
        \centering
        \includegraphics[width=\linewidth]{figures/empapp_main_fs_heatmap_pval_bwmain.png}
        \caption{Pointwise $p$-values.}
    \end{subfigure}
    \caption{First-Stage Effects (SPP Application, Fuzzy Bandwidth Selector).}
    \label{fig:sa-heatmaps-fs-main-bwmain}
\end{figure}

\begin{figure}
    \centering
    \begin{subfigure}[b]{0.70\textwidth}
        \centering
        \includegraphics[width=\linewidth]{figures/empapp_main_fuzzy_heatmap_bwmain.png}
        \caption{Point estimates.}
    \end{subfigure}

    \begin{subfigure}[b]{0.70\textwidth}
        \centering
        \includegraphics[width=\linewidth]{figures/empapp_main_fuzzy_heatmap_pval_bwmain.png}
        \caption{Pointwise $p$-values.}
    \end{subfigure}
    \caption{Fuzzy Effects (SPP Application, Fuzzy Bandwidth Selector).}
    \label{fig:sa-heatmaps-fuzzy-main-bwmain}
\end{figure}

\begin{figure}
    \centering
    \begin{subfigure}[b]{0.70\textwidth}
        \centering
        \includegraphics[width=\linewidth]{figures/empapp_covbal_itt_heatmap_bwmain.png}
        \caption{Point estimates.}
    \end{subfigure}

    \begin{subfigure}[b]{0.70\textwidth}
        \centering
        \includegraphics[width=\linewidth]{figures/empapp_covbal_itt_heatmap_pval_bwmain.png}
        \caption{Pointwise $p$-values.}
    \end{subfigure}
    \caption{Covariate Balance for Mother's Education (SPP Application, Fuzzy Bandwidth Selector).}
    \label{fig:sa-heatmaps-covbal-bwmain}
\end{figure}

\clearpage

\begin{table}[p]
    \centering
\begin{tabular}{lccccccc}
  \toprule\toprule
   & $h_1$ & $h_2$ & $N_{\mathrm{Co}}$ & $N_{\mathrm{Tr}}$ & Estimate & $p$-value & 95\% CI \\
  \midrule
   $\bb_{1}$ & 28.2 & 12.5 & 9704 &  & 0.371 & 0.000 & (0.299, 0.456)\\
   $\bb_{2}$ & 26.0 & 11.5 & 8572 &  & 0.364 & 0.000 & (0.287, 0.444)\\
   $\bb_{3}$ & 24.9 & 11.0 & 8590 &  & 0.356 & 0.000 & (0.286, 0.439)\\
   $\bb_{4}$ & 23.2 & 10.3 & 8397 &  & 0.357 & 0.000 & (0.290, 0.448)\\
   $\bb_{5}$ & 21.8 & 9.6 & 7403 &  & 0.360 & 0.000 & (0.297, 0.462)\\
   $\bb_{6}$ & 20.0 & 8.9 & 6977 &  & 0.362 & 0.000 & (0.291, 0.465)\\
   $\bb_{7}$ & 19.7 & 8.7 & 6633 &  & 0.355 & 0.000 & (0.277, 0.447)\\
   $\bb_{8}$ & 21.3 & 9.4 & 8437 &  & 0.346 & 0.000 & (0.278, 0.429)\\
   $\bb_{9}$ & 21.9 & 9.7 & 8859 &  & 0.342 & 0.000 & (0.275, 0.421)\\
   $\bb_{10}$ & 22.2 & 9.8 & 9828 &  & 0.343 & 0.000 & (0.277, 0.420)\\
   $\bb_{11}$ & 22.3 & 9.9 & 9906 &  & 0.351 & 0.000 & (0.286, 0.427)\\
   $\bb_{12}$ & 22.3 & 9.9 & 9949 &  & 0.363 & 0.000 & (0.302, 0.444)\\
   $\bb_{13}$ & 20.6 & 9.1 & 7995 &  & 0.379 & 0.000 & (0.316, 0.474)\\
   $\bb_{14}$ & 20.7 & 9.2 & 8168 &  & 0.391 & 0.000 & (0.342, 0.500)\\
   $\bb_{15}$ & 21.0 & 9.3 & 8370 &  & 0.400 & 0.000 & (0.355, 0.510)\\
   $\bb_{16}$ & 20.8 & 9.2 & 8278 &  & 0.403 & 0.000 & (0.346, 0.502)\\
   $\bb_{17}$ & 19.7 & 8.7 & 7544 &  & 0.403 & 0.000 & (0.331, 0.473)\\
   $\bb_{18}$ & 19.6 & 8.7 & 7997 &  & 0.396 & 0.000 & (0.334, 0.443)\\
   $\bb_{19}$ & 20.4 & 9.0 & 9592 &  & 0.391 & 0.000 & (0.345, 0.434)\\
   $\bb_{20}$ & 20.0 & 8.9 & 9225 &  & 0.394 & 0.000 & (0.360, 0.441)\\
   $\bb_{21}$ & 20.9 & 9.3 & 10677 &  & 0.399 & 0.000 & (0.373, 0.442)\\
   $\bb_{22}$ & 22.0 & 9.8 & 10830 &  & 0.411 & 0.000 & (0.382, 0.455)\\
   $\bb_{23}$ & 21.3 & 9.4 & 8679 &  & 0.422 & 0.000 & (0.379, 0.466)\\
   $\bb_{24}$ & 21.7 & 9.6 & 8114 &  & 0.432 & 0.000 & (0.375, 0.475)\\
   $\bb_{25}$ & 21.8 & 9.6 & 6858 &  & 0.445 & 0.000 & (0.371, 0.486)\\
   $\bb_{26}$ & 24.3 & 10.8 & 7784 &  & 0.460 & 0.000 & (0.385, 0.499)\\
   $\bb_{27}$ & 27.4 & 12.1 & 8856 &  & 0.475 & 0.000 & (0.407, 0.517)\\
   $\bb_{28}$ & 26.9 & 11.9 & 7294 &  & 0.485 & 0.000 & (0.408, 0.534)\\
   $\bb_{29}$ & 26.7 & 11.8 & 5916 &  & 0.494 & 0.000 & (0.410, 0.552)\\
   $\bb_{30}$ & 27.3 & 12.1 & 5250 &  & 0.508 & 0.000 & (0.411, 0.572)\\
   $\bb_{31}$ & 30.1 & 13.3 & 5758 &  & 0.525 & 0.000 & (0.434, 0.591)\\
   $\bb_{32}$ & 31.8 & 14.1 & 5266 &  & 0.541 & 0.000 & (0.441, 0.608)\\
   $\bb_{33}$ & 32.1 & 14.2 & 4925 &  & 0.556 & 0.000 & (0.455, 0.630)\\
   $\bb_{34}$ & 31.7 & 14.0 & 4323 &  & 0.571 & 0.000 & (0.466, 0.654)\\
   $\bb_{35}$ & 28.8 & 12.7 & 3022 &  & 0.587 & 0.000 & (0.466, 0.684)\\
   $\bb_{36}$ & 27.0 & 12.0 & 2452 &  & 0.604 & 0.000 & (0.475, 0.720)\\
   $\bb_{37}$ & 28.1 & 12.4 & 2318 &  & 0.622 & 0.000 & (0.492, 0.743)\\
   $\bb_{38}$ & 29.7 & 13.1 & 2330 &  & 0.639 & 0.000 & (0.512, 0.765)\\
   $\bb_{39}$ & 35.2 & 15.6 & 2954 &  & 0.655 & 0.000 & (0.546, 0.772)\\
   $\bb_{40}$ & 39.6 & 17.5 & 3538 &  & 0.667 & 0.000 & (0.570, 0.783)\\
  \midrule
   $\mathtt{WBATE}$ &  &  &  &  & 0.448 & 0.000 & (0.425, 0.474)\\
  \midrule
   $\mathtt{LBATE}$ &  &  &  &  & 0.667 &  & (0.568, 0.784)\\
  \bottomrule\bottomrule
\end{tabular}

    \caption{Control-Side Baseline Outcome (SPP Application, Fuzzy Bandwidth Selector).}
    \label{tab:sa-empapp_main_itt0_bwmain}
\end{table}

\begin{table}[p]
    \centering
\begin{tabular}{lcccccc}
  \toprule\toprule
   & $h_1$ & $h_2$ & $N_{\mathrm{Co}}$ & $N_{\mathrm{Tr}}$ & ITT / ITT.0 (\%) & Fuzzy / ITT.0 (\%) \\
  \midrule
   $\bb_{1}$ & 28.2 & 12.5 & 9704 & 3687 & 83.801 & 152.231\\
   $\bb_{2}$ & 26.0 & 11.5 & 8572 & 3612 & 85.950 & 157.238\\
   $\bb_{3}$ & 24.9 & 11.0 & 8590 & 3720 & 88.686 & 163.182\\
   $\bb_{4}$ & 23.2 & 10.3 & 8397 & 3724 & 86.320 & 160.201\\
   $\bb_{5}$ & 21.8 & 9.6 & 7403 & 3568 & 80.977 & 153.111\\
   $\bb_{6}$ & 20.0 & 8.9 & 6977 & 3427 & 75.410 & 146.062\\
   $\bb_{7}$ & 19.7 & 8.7 & 6633 & 3404 & 76.445 & 150.264\\
   $\bb_{8}$ & 21.3 & 9.4 & 8437 & 4032 & 83.668 & 164.054\\
   $\bb_{9}$ & 21.9 & 9.7 & 8859 & 4218 & 89.489 & 174.020\\
   $\bb_{10}$ & 22.2 & 9.8 & 9828 & 4469 & 93.492 & 177.775\\
   $\bb_{11}$ & 22.3 & 9.9 & 9906 & 4584 & 94.019 & 174.994\\
   $\bb_{12}$ & 22.3 & 9.9 & 9949 & 4710 & 90.043 & 166.056\\
   $\bb_{13}$ & 20.6 & 9.1 & 7995 & 4229 & 85.004 & 153.958\\
   $\bb_{14}$ & 20.7 & 9.2 & 8168 & 4350 & 80.630 & 143.619\\
   $\bb_{15}$ & 21.0 & 9.3 & 8370 & 4493 & 77.426 & 137.434\\
   $\bb_{16}$ & 20.8 & 9.2 & 8278 & 4560 & 77.256 & 136.138\\
   $\bb_{17}$ & 19.7 & 8.7 & 7544 & 4057 & 79.328 & 138.034\\
   $\bb_{18}$ & 19.6 & 8.7 & 7997 & 3623 & 82.308 & 143.110\\
   $\bb_{19}$ & 20.4 & 9.0 & 9592 & 3367 & 85.368 & 148.263\\
   $\bb_{20}$ & 20.0 & 8.9 & 9225 & 2755 & 83.046 & 144.554\\
   $\bb_{21}$ & 20.9 & 9.3 & 10677 & 2442 & 78.747 & 137.406\\
   $\bb_{22}$ & 22.0 & 9.8 & 10830 & 2843 & 74.724 & 129.677\\
   $\bb_{23}$ & 21.3 & 9.4 & 8679 & 2871 & 70.735 & 123.194\\
   $\bb_{24}$ & 21.7 & 9.6 & 8114 & 3140 & 67.914 & 118.071\\
   $\bb_{25}$ & 21.8 & 9.6 & 6858 & 3223 & 65.063 & 112.204\\
   $\bb_{26}$ & 24.3 & 10.8 & 7784 & 3867 & 61.248 & 104.183\\
   $\bb_{27}$ & 27.4 & 12.1 & 8856 & 4564 & 57.067 & 95.654\\
   $\bb_{28}$ & 26.9 & 11.9 & 7294 & 4540 & 55.052 & 91.338\\
   $\bb_{29}$ & 26.7 & 11.8 & 5916 & 4593 & 53.365 & 87.115\\
   $\bb_{30}$ & 27.3 & 12.1 & 5250 & 4759 & 50.399 & 80.804\\
   $\bb_{31}$ & 30.1 & 13.3 & 5758 & 5285 & 47.095 & 73.880\\
   $\bb_{32}$ & 31.8 & 14.1 & 5266 & 5607 & 44.339 & 68.325\\
   $\bb_{33}$ & 32.1 & 14.2 & 4925 & 5218 & 42.053 & 63.694\\
   $\bb_{34}$ & 31.7 & 14.0 & 4323 & 4517 & 39.874 & 59.262\\
   $\bb_{35}$ & 28.8 & 12.7 & 3022 & 3135 & 37.509 & 54.321\\
   $\bb_{36}$ & 27.0 & 12.0 & 2452 & 2467 & 35.590 & 50.375\\
   $\bb_{37}$ & 28.1 & 12.4 & 2318 & 2326 & 33.480 & 46.767\\
   $\bb_{38}$ & 29.7 & 13.1 & 2330 & 2353 & 31.284 & 43.456\\
   $\bb_{39}$ & 35.2 & 15.6 & 2954 & 3036 & 29.205 & 40.612\\
   $\bb_{40}$ & 39.6 & 17.5 & 3538 & 3725 & 27.726 & 38.547\\
  \midrule
   $\mathtt{WBATE}$ &  &  &  &  & 62.655 & 108.222\\
  \midrule
   $\mathtt{LBATE}$ &  &  &  &  & 49.984 & 91.972\\
  \bottomrule\bottomrule
\end{tabular}

    \caption{Treatment Effect Relative to Control-Side Baseline (SPP Application, Fuzzy Bandwidth Selector).}
    \label{tab:sa-empapp_main_percent_bwmain}
\end{table}

\begin{figure}
    \centering
    \begin{subfigure}[b]{0.70\textwidth}
        \centering
        \includegraphics[width=\linewidth]{figures/empapp_main_itt0_inference_bwmain.png}
        \caption{Point Estimates and Uncertainty Quantification.}
    \end{subfigure}

    \begin{subfigure}[b]{0.70\textwidth}
        \centering
        \includegraphics[width=\linewidth]{figures/empapp_main_itt0_heatmap_bwmain.png}
        \caption{Heatmap of Point Estimates.}
    \end{subfigure}
    \caption{Control-Side Baseline Outcome (SPP Application, Fuzzy Bandwidth Selector).}
    \label{fig:sa-empapp_main_itt0_bwmain}
\end{figure}

\begin{figure}
    \centering
    \begin{subfigure}[b]{0.70\textwidth}
        \centering
        \includegraphics[width=\linewidth]{figures/empapp_main_percent_itt_pointest_bwmain.png}
        \caption{Intention-to-Treat Effect.}
    \end{subfigure}

    \begin{subfigure}[b]{0.70\textwidth}
        \centering
        \includegraphics[width=\linewidth]{figures/empapp_main_percent_fuzzy_pointest_bwmain.png}
        \caption{Fuzzy Effect.}
    \end{subfigure}
    \caption{Treatment Effect Relative to Control-Side Baseline (SPP Application, Fuzzy Bandwidth Selector).}
    \label{fig:sa-empapp_main_percent_bwmain}
\end{figure}

\clearpage

%%%%%%%%%%%%%%%%%%%%%%%%%%%%%%%%%%%%%%%%%%%%%%%%%%%%%%%%%%%%%%%%%%%%%%
\section{Simulation Study}
%%%%%%%%%%%%%%%%%%%%%%%%%%%%%%%%%%%%%%%%%%%%%%%%%%%%%%%%%%%%%%%%%%%%%%

The score $\bX_i = (X_{1i}, X_{2i})^\top$ is generated from the distribution $(100\cdot\mathsf{Beta}(3,4)-25,\;100\cdot\mathsf{Beta}(3,4)-25)^\top$ with independent components, which approximately reproduces salient features of the SPP dataset. We employ the same L-shaped treatment assignment boundary as in the empirical application, featuring a kink at $\bx=(0,0)^\top$. To reduce computation burden we consider $21$ evaluation points with the kink corresponding to the evaluation grid point $\bb_{11}$.

The simulation design sets $n=10,000$ and uses $5,000$ Monte Carlo replications. We evaluate performance in terms of point estimation and uncertainty quantification, both pointwise and uniformly over the boundary $\B$. The Monte Carlo design jointly generates the reduced-form outcome and the treatment status, so that the sharp intention-to-treat (ITT), first-stage (FS), and fuzzy analyses can all be performed on the same simulated samples.

\subsection{Data-Generating Process}

For each Monte Carlo replication, draw $\bX_i=(X_{1i},X_{2i})^\top$ independently from the distribution described above and set $T_i=\Indicator(\bX_i\in\A_1)$. We consider DGP polynomial specifications of order $s\in\{1,2,3\}$, using the monomial vector $\br_s(\cdot)$, evaluated at the unscaled score. Thus, for $d=2$, $\br_1(\bx)=(1,x_1,x_2)^\top$ and $\br_2(\bx)=(1,x_1,x_2,x_1^2,x_1x_2,x_2^2)^\top$, with $\br_3(\bx)$ defined analogously by adding all cubic monomials. For each $t\in\{0,1\}$ and $i=1,\ldots,n$, define
\begin{align*}
    \mu_{Y,t}(\bX_i)
    &= \br_s(\bX_i)^\top\bbeta_{Y,t},\\
    \mu_{W,t}(\bX_i)
    &= \Lambda(\eta_{W,t}(\bX_i)),
    \qquad
    \eta_{W,t}(\bX_i) = \br_s(\bX_i)^\top\bbeta_{W,t},
\end{align*}
where $\Lambda(a)=1/(1+\exp(-a))$. Given these conditional mean functions, draw $U_{W,t,i}\thicksim\mathsf{Uniform}(0,1)$ and $\varepsilon_{Y,t,i}\thicksim\mathsf{Normal}(0,\sigma_{Y,t}^2)$ independently across $i$ and $t$ and independently of $\bX_i$. The potential treatment statuses and treatment-indexed potential outcomes are
\begin{align*}
    W_i(t)
    &= \Indicator\{U_{W,t,i}\leq \mu_{W,t}(\bX_i)\},\\
    Y_i(t,r)
    &= \mu_{Y,t}(\bX_i) + \lambda_t\{r-\mu_{W,t}(\bX_i)\}
    + \varepsilon_{Y,t,i},
    \qquad r\in\{0,1\}.
\end{align*}
The reduced-form potential outcome induced by assignment $t$ is $Y_i(t,W_i(t))$. Equivalently,
\begin{align*}
    Y_i(t,W_i(t))
    &=
    \mu_{Y,t}(\bX_i)
    +\lambda_t\{W_i(t)-\mu_{W,t}(\bX_i)\}
    +\varepsilon_{Y,t,i}.
\end{align*}
The observed outcomes and treatment status are
\begin{align*}
    Y_i &= Y_i(0,W_i(0)) (1-T_i) + Y_i(1,W_i(1)) T_i,\\
    W_i &= W_i(0) (1-T_i) + W_i(1) T_i, \qquad  T_i = \Indicator(\bX_i \in \A_1),
\end{align*}
for $i=1,2,\ldots,n$.

In the reduced-form notation used for sharp estimation and inference, $Y_i(t)$ is shorthand for $Y_i(t,W_i(t))$. This centered specification preserves the reduced-form outcome surface because $\E[Y_i(t,W_i(t))|\bX_i=\bx]=\mu_{Y,t}(\bx)$, while allowing the potential outcome and potential treatment status equations to be correlated:
\begin{align*}
    \E[W_i(t)|\bX_i=\bx] = \mu_{W,t}(\bx)
    \qquad\text{and}\qquad
    \Cov(Y_i(t,W_i(t)),W_i(t)|\bX_i=\bx) = \lambda_t\mu_{W,t}(\bx)\{1-\mu_{W,t}(\bx)\}.
\end{align*}
Consequently, for each boundary point $\bb\in\B$,
\begin{align*}
    \tau_W(\bb)
    &=
    \E[W_i(1)-W_i(0)|\bX_i=\bb]
    =
    \mu_{W,1}(\bb)-\mu_{W,0}(\bb),\\
    \tau_Y(\bb)
    &=
    \E[Y_i(1,W_i(1))-Y_i(0,W_i(0))|\bX_i=\bb]
    =
    \mu_{Y,1}(\bb)-\mu_{Y,0}(\bb),\\
    \zeta(\bb)
    &=
    \frac{\tau_Y(\bb)}{\tau_W(\bb)}.
\end{align*}
In each Monte Carlo replication, the sharp analysis applied to $Y_i$ targets the reduced-form or ITT curve $\tau_Y(\bb)$, the same analysis applied to $W_i$ targets the first-stage curve $\tau_W(\bb)$, and the fuzzy analysis targets the Wald curve $\zeta(\bb)$. Thus, the ITT and fuzzy simulation results are different summaries of the same simulated samples rather than results from separate sharp and fuzzy DGPs.

\subsection{Calibration from the SPP data}

The DGP parameters are calibrated once from the SPP data and then held fixed across Monte Carlo replications. We use the following steps separately for $s\in\{1,2,3\}$.

\begin{enumerate}[label=\normalfont Step \arabic*.,leftmargin=*]
    \item Fix the DGP polynomial order $s\in\{1,2,3\}$ and define the side-specific calibration samples $\mathcal I_t=\{i:T_i=t\}$ and $N_t=|\mathcal I_t|$, for $t\in\{0,1\}$.

    \item For each side $t$, estimate the outcome regression coefficients by side-specific least squares:
    \[
        \widehat{\bbeta}_{Y,t}
        =
        \operatorname*{arg\,min}_{\bbeta\in\mathbb R^{\mathfrak p_s}}
        \sum_{i\in\mathcal I_t}\{Y_i-\br_s(\bX_i)^\top\bbeta\}^2.
    \]
    This gives the fitted reduced-form outcome surface $\widehat{\mu}_{Y,t}(\bx)=\br_s(\bx)^\top\widehat{\bbeta}_{Y,t}$.

    \item For each side $t$, estimate the treatment-status regression coefficients by side-specific logit:
    \[
        \widehat{\bbeta}_{W,t}
        =
        \operatorname*{arg\,max}_{\bbeta\in\mathbb R^{\mathfrak p_s}}
        \sum_{i\in\mathcal I_t}
        \left[
            W_i \br_s(\bX_i)^\top\bbeta
            -\log\{1+\exp(\br_s(\bX_i)^\top\bbeta)\}
        \right].
    \]
    This gives the fitted first-stage surface $\widehat{\mu}_{W,t}(\bx)=\Lambda(\br_s(\bx)^\top\widehat{\bbeta}_{W,t})$.

    \item Construct side-specific residuals in the calibration sample:
    \[
        \widehat{\varepsilon}_{Y,i}=Y_i-\widehat{\mu}_{Y,T_i}(\bX_i),
        \qquad
        \widehat{\varepsilon}_{W,i}=W_i-\widehat{\mu}_{W,T_i}(\bX_i).
    \]

    \item Estimate the residual dependence parameter and outcome-shock variance on each side:
\begin{align*}
    \widehat{\lambda}_t
    &=
    \frac{\sum_{i\in\mathcal I_t}\widehat{\varepsilon}_{Y,i}\widehat{\varepsilon}_{W,i}}
    {\sum_{i\in\mathcal I_t}\widehat{\varepsilon}_{W,i}^2},
    \qquad
    \widehat{\sigma}_{Y,t}^2
    =
    \frac{1}{N_t}\sum_{i\in\mathcal I_t}
    \{\widehat{\varepsilon}_{Y,i}-\widehat{\lambda}_t\widehat{\varepsilon}_{W,i}\}^2,
\end{align*}

    \item Set the Monte Carlo DGP parameters equal to the calibration estimates: $\bbeta_{Y,t}=\widehat{\bbeta}_{Y,t}$, $\bbeta_{W,t}=\widehat{\bbeta}_{W,t}$, $\lambda_t=\widehat{\lambda}_t$, and $\sigma_{Y,t}^2=\widehat{\sigma}_{Y,t}^2$ for $t\in\{0,1\}$.

    \item Before running the Monte Carlo replications, compute the true target functions analytically from the calibrated conditional mean surfaces. For each evaluation point $\bb\in\B$,
    \[
        \tau_Y(\bb)
        =
        \br_s(\bb)^\top(\bbeta_{Y,1}-\bbeta_{Y,0}),
        \qquad
        \tau_W(\bb)
        =
        \Lambda(\br_s(\bb)^\top\bbeta_{W,1})
        -
        \Lambda(\br_s(\bb)^\top\bbeta_{W,0}),
        \qquad
        \zeta(\bb)=\frac{\tau_Y(\bb)}{\tau_W(\bb)}.
    \]
\end{enumerate}

This calibration reproduces the empirical residual covariance between the outcome and take-up equations on each side of the boundary, while keeping the marginal reduced-form and first-stage surfaces fixed by the fitted SPP models.

\subsection{Monte Carlo Implementation}

The following steps describe how to simulate from the calibrated model. They are repeated separately for each DGP polynomial order $s\in\{1,2,3\}$.

\begin{enumerate}[label=\normalfont Step \arabic*.,leftmargin=*]
    \item Fix the Monte Carlo inputs: the sample size $n$, the number of replications $M$, the evaluation grid $\B=\{\bb_1,\ldots,\bb_{21}\}$ with $\bb_{11}$ at the kink, the assignment regions $\A_0$ and $\A_1$, and the calibrated parameters
    \[
        \{\bbeta_{Y,t},\bbeta_{W,t},\lambda_t,\sigma_{Y,t}^2:t\in\{0,1\}\}.
    \]

    \item Before entering the Monte Carlo loop, evaluate the true target functions on the grid:
    \[
        \tau_Y(\bb_j),\qquad
        \tau_W(\bb_j),\qquad
        \zeta(\bb_j),
        \qquad j=1,\ldots,21,
    \]
    using the analytic formulas in the calibration step above. These target values are fixed across Monte Carlo replications.

    \item For replication $m=1,\ldots,M$, draw independent score vectors
    \[
        \bX_i=(X_{1i},X_{2i})^\top,
        \qquad
        X_{\ell i}=100B_{\ell i}-25,\quad
        B_{\ell i}\thicksim\mathsf{Beta}(3,4),
        \quad \ell\in\{1,2\},
    \]
    with $B_{1i}$ and $B_{2i}$ independent. Set $T_i=\Indicator(\bX_i\in\A_1)$.

    \item Form the design vectors $\br_s(\bX_i)$ and compute the side-specific conditional means
    \[
        \mu_{Y,t}(\bX_i)=\br_s(\bX_i)^\top\bbeta_{Y,t},
        \qquad
        \mu_{W,t}(\bX_i)=\Lambda(\br_s(\bX_i)^\top\bbeta_{W,t}),
        \qquad t\in\{0,1\}.
    \]

    \item Draw the treatment-status shocks and outcome shocks independently:
    \[
        U_{W,t,i}\thicksim\mathsf{Uniform}(0,1),
        \qquad
        \varepsilon_{Y,t,i}\thicksim\mathsf{Normal}(0,\sigma_{Y,t}^2),
        \qquad t\in\{0,1\}.
    \]
    Then construct the potential treatment statuses and treatment-indexed potential outcomes:
    \begin{align*}
        W_i(t)
        &=
        \Indicator\{U_{W,t,i}\leq \mu_{W,t}(\bX_i)\},\\
        Y_i(t,r)
        &=
        \mu_{Y,t}(\bX_i)
        +\lambda_t\{r-\mu_{W,t}(\bX_i)\}
        +\varepsilon_{Y,t,i},
        \qquad r\in\{0,1\}.
    \end{align*}
    Evaluate the reduced-form potential outcome at the simulated treatment status:
    \begin{align*}
        Y_i(t,W_i(t))
        &=
        \mu_{Y,t}(\bX_i)
        +\lambda_t\{W_i(t)-\mu_{W,t}(\bX_i)\}
        +\varepsilon_{Y,t,i}.
    \end{align*}

    \item Select the observed variables using the assignment indicator:
    \begin{align*}
        W_i
        &=
        W_i(0)(1-T_i)+W_i(1)T_i,\\
        Y_i
        &=
        Y_i(0,W_i(0))(1-T_i)+Y_i(1,W_i(1))T_i.
    \end{align*}
    The simulated dataset for replication $m$ is therefore $\{Y_i,W_i,T_i,\bX_i:i=1,\ldots,n\}$.

    \item \textbf{ITT Analysis}. Apply the sharp estimator to $Y_i$ using the simulated assignment rule and score vectors for the target $\tau_Y(\bb_j)$.

    \item \textbf{FS Analysis}. Apply the sharp estimator to $W_i$ using the simulated assignment rule and score vectors for the target $\tau_W(\bb_j)$.

    \item \textbf{Fuzzy Analysis}. Apply the fuzzy estimator using $Y_i$ as the outcome, $W_i$ as the treatment status, and using the simulated assignment rule and score vectors for the Wald target $\zeta(\bb_j)$.

    \item After all $M$ replications, compute Monte Carlo summaries by comparing the stored estimates and confidence sets with the fixed analytic targets. For each target and each grid point, report bias, standard deviation, root mean squared error, pointwise coverage, and average interval length. For uniform inference, report simultaneous coverage and average band length over the grid $\B$.
    Finally, also recover WBATE and LBATE and report their associated performance.
\end{enumerate}

We use the fuzzy simulations to study the finite-sample behavior of the reduced-form/ITT estimators, the first-stage estimators, and the fuzzy Wald estimator. A complier-average causal interpretation of $\zeta(\bb)$ would require additional assumptions that are not imposed in the simulations. In particular, one would need a local exclusion restriction ruling out a direct effect of assignment on the outcome except through treatment take-up, a monotonicity condition ruling out defiers near the boundary, and a nonzero first-stage effect. Under those conditions, the Wald curve $\zeta(\bb)$ can be interpreted as a local average treatment effect for compliers at the boundary point $\bb$.

\subsection{Results}

The following tables report Monte Carlo summaries separately for the reduced-form/ITT, first-stage, and fuzzy estimators, and separately by DGP polynomial order $s\in\{1,2,3\}$. We report two sets of results: one using the ITT bandwidth selector, matching the main empirical specification, and one using the fuzzy bandwidth selector.

The first set of simulation results uses the ITT bandwidth selector. Specifically:
\begin{itemize}
    \item Table~\ref{tab:sa-simuls-itt-s1} reports the reduced-form (ITT) simulation results for $s=1$.
    \item Table~\ref{tab:sa-simuls-itt-s2} reports the reduced-form (ITT) simulation results for $s=2$.
    \item Table~\ref{tab:sa-simuls-itt-s3} reports the reduced-form (ITT) simulation results for $s=3$.
    \item Table~\ref{tab:sa-simuls-fs-s1} reports the take-up (FS) simulation results for $s=1$.
    \item Table~\ref{tab:sa-simuls-fs-s2} reports the take-up (FS) simulation results for $s=2$.
    \item Table~\ref{tab:sa-simuls-fs-s3} reports the take-up (FS) simulation results for $s=3$.
    \item Table~\ref{tab:sa-simuls-fuzzy-s1} reports the fuzzy simulation results for $s=1$.
    \item Table~\ref{tab:sa-simuls-fuzzy-s2} reports the fuzzy simulation results for $s=2$.
    \item Table~\ref{tab:sa-simuls-fuzzy-s3} reports the fuzzy simulation results for $s=3$.
\end{itemize}

The second set of simulation results uses the fuzzy bandwidth selector. Specifically:
\begin{itemize}
    \item Table~\ref{tab:sa-simuls-itt-s1-bwmain} reports the reduced-form (ITT) simulation results for $s=1$.
    \item Table~\ref{tab:sa-simuls-itt-s2-bwmain} reports the reduced-form (ITT) simulation results for $s=2$.
    \item Table~\ref{tab:sa-simuls-itt-s3-bwmain} reports the reduced-form (ITT) simulation results for $s=3$.
    \item Table~\ref{tab:sa-simuls-fs-s1-bwmain} reports the take-up (FS) simulation results for $s=1$.
    \item Table~\ref{tab:sa-simuls-fs-s2-bwmain} reports the take-up (FS) simulation results for $s=2$.
    \item Table~\ref{tab:sa-simuls-fs-s3-bwmain} reports the take-up (FS) simulation results for $s=3$.
    \item Table~\ref{tab:sa-simuls-fuzzy-s1-bwmain} reports the fuzzy simulation results for $s=1$.
    \item Table~\ref{tab:sa-simuls-fuzzy-s2-bwmain} reports the fuzzy simulation results for $s=2$.
    \item Table~\ref{tab:sa-simuls-fuzzy-s3-bwmain} reports the fuzzy simulation results for $s=3$.
\end{itemize}
The columns report the average bandwidths, average bias, standard deviation, root mean squared error, empirical coverage, and average interval or band length. For the Uniform row, empirical coverage and average length refer to uniform coverage and average uniform band length.

\clearpage

\begin{table}[p]
    \centering
\begin{tabular}{lccccrrrrr}
\toprule\toprule
 & $h_1$ & $h_2$ & $N_{\mathrm{Co}}$ & $N_{\mathrm{Tr}}$ & Bias & SD & RMSE & EC & IL \\
\midrule
$\bb_{1}$ & 13.220 & 13.220 & 732 & 1267 & 0.001 & 0.045 & 0.045 & 0.950 & 0.263 \\
$\bb_{2}$ & 13.525 & 13.525 & 765 & 1343 & 0.001 & 0.044 & 0.044 & 0.949 & 0.255 \\
$\bb_{3}$ & 13.740 & 13.740 & 787 & 1393 & 0.001 & 0.043 & 0.043 & 0.950 & 0.251 \\
$\bb_{4}$ & 13.857 & 13.857 & 808 & 1400 & 0.001 & 0.042 & 0.042 & 0.949 & 0.246 \\
$\bb_{5}$ & 13.904 & 13.904 & 869 & 1329 & 0.001 & 0.041 & 0.041 & 0.953 & 0.234 \\
$\bb_{6}$ & 13.890 & 13.890 & 924 & 1229 & 0.001 & 0.040 & 0.040 & 0.955 & 0.232 \\
$\bb_{7}$ & 13.811 & 13.811 & 959 & 1116 & 0.000 & 0.040 & 0.040 & 0.950 & 0.239 \\
$\bb_{8}$ & 13.639 & 13.639 & 968 & 989 & 0.001 & 0.042 & 0.042 & 0.949 & 0.253 \\
$\bb_{9}$ & 13.993 & 13.993 & 1024 & 928 & 0.000 & 0.046 & 0.046 & 0.948 & 0.271 \\
$\bb_{10}$ & 15.131 & 15.131 & 1147 & 961 & 0.000 & 0.051 & 0.051 & 0.948 & 0.306 \\
$\bb_{11}$ & 14.782 & 14.782 & 1091 & 803 & 0.000 & 0.066 & 0.066 & 0.936 & 0.425 \\
$\bb_{12}$ & 14.679 & 14.679 & 1101 & 949 & 0.000 & 0.048 & 0.048 & 0.941 & 0.288 \\
$\bb_{13}$ & 13.614 & 13.614 & 969 & 965 & 0.000 & 0.043 & 0.043 & 0.949 & 0.257 \\
$\bb_{14}$ & 13.811 & 13.811 & 951 & 1137 & 0.000 & 0.040 & 0.040 & 0.953 & 0.237 \\
$\bb_{15}$ & 13.870 & 13.870 & 889 & 1286 & 0.001 & 0.041 & 0.041 & 0.954 & 0.232 \\
$\bb_{16}$ & 13.837 & 13.837 & 805 & 1397 & 0.001 & 0.043 & 0.043 & 0.949 & 0.246 \\
$\bb_{17}$ & 13.678 & 13.678 & 781 & 1379 & 0.001 & 0.044 & 0.044 & 0.945 & 0.251 \\
$\bb_{18}$ & 13.289 & 13.289 & 739 & 1284 & 0.001 & 0.045 & 0.045 & 0.945 & 0.260 \\
$\bb_{19}$ & 13.011 & 13.011 & 697 & 1195 & 0.001 & 0.047 & 0.047 & 0.947 & 0.269 \\
$\bb_{20}$ & 13.077 & 13.077 & 671 & 1154 & 0.001 & 0.048 & 0.048 & 0.949 & 0.274 \\
$\bb_{21}$ & 13.330 & 13.330 & 651 & 1130 & 0.000 & 0.049 & 0.049 & 0.944 & 0.278 \\
\midrule
$\mathtt{Uniform}$ &  &  &  &  &  &  &  & 0.930 & 0.381 \\
\midrule
$\mathtt{WBATE}$ &  &  &  &  & 0.001 & 0.026 & 0.026 & 0.943 & 0.138 \\
\midrule
$\mathtt{LBATE}$ &  &  &  &  & 0.053 & 0.038 & 0.065 & 0.971 & 0.433 \\
\bottomrule\bottomrule
\end{tabular}

    \caption{Intention-to-Treat Effects (Simulations, $s=1$, ITT Bandwidth Selector).}
    \label{tab:sa-simuls-itt-s1}
\end{table}

\begin{table}[p]
    \centering
\begin{tabular}{lccccrrrrr}
\toprule\toprule
 & $h_1$ & $h_2$ & $N_{\mathrm{Co}}$ & $N_{\mathrm{Tr}}$ & Bias & SD & RMSE & EC & IL \\
\midrule
$\bb_{1}$ & 13.191 & 13.191 & 728 & 1262 & 0.001 & 0.045 & 0.045 & 0.949 & 0.263 \\
$\bb_{2}$ & 13.494 & 13.494 & 762 & 1338 & 0.001 & 0.044 & 0.044 & 0.951 & 0.256 \\
$\bb_{3}$ & 13.717 & 13.717 & 784 & 1390 & 0.002 & 0.043 & 0.043 & 0.951 & 0.251 \\
$\bb_{4}$ & 13.830 & 13.830 & 804 & 1397 & 0.002 & 0.042 & 0.042 & 0.954 & 0.247 \\
$\bb_{5}$ & 13.870 & 13.870 & 864 & 1325 & 0.001 & 0.041 & 0.041 & 0.955 & 0.235 \\
$\bb_{6}$ & 13.846 & 13.846 & 919 & 1223 & 0.000 & 0.040 & 0.040 & 0.950 & 0.233 \\
$\bb_{7}$ & 13.777 & 13.777 & 954 & 1112 & -0.001 & 0.040 & 0.040 & 0.947 & 0.239 \\
$\bb_{8}$ & 13.630 & 13.630 & 966 & 988 & -0.001 & 0.043 & 0.043 & 0.949 & 0.253 \\
$\bb_{9}$ & 14.017 & 14.017 & 1026 & 932 & -0.002 & 0.046 & 0.046 & 0.944 & 0.271 \\
$\bb_{10}$ & 15.156 & 15.156 & 1148 & 965 & -0.003 & 0.051 & 0.051 & 0.942 & 0.305 \\
$\bb_{11}$ & 14.778 & 14.778 & 1090 & 803 & -0.003 & 0.067 & 0.067 & 0.939 & 0.423 \\
$\bb_{12}$ & 14.646 & 14.646 & 1097 & 945 & -0.003 & 0.049 & 0.049 & 0.939 & 0.287 \\
$\bb_{13}$ & 13.581 & 13.581 & 965 & 961 & -0.002 & 0.044 & 0.044 & 0.946 & 0.256 \\
$\bb_{14}$ & 13.768 & 13.768 & 945 & 1131 & -0.001 & 0.041 & 0.041 & 0.947 & 0.236 \\
$\bb_{15}$ & 13.834 & 13.834 & 883 & 1281 & -0.001 & 0.042 & 0.042 & 0.949 & 0.231 \\
$\bb_{16}$ & 13.800 & 13.800 & 801 & 1391 & 0.000 & 0.043 & 0.043 & 0.945 & 0.246 \\
$\bb_{17}$ & 13.629 & 13.629 & 776 & 1370 & 0.000 & 0.044 & 0.044 & 0.945 & 0.251 \\
$\bb_{18}$ & 13.237 & 13.237 & 734 & 1275 & 0.000 & 0.045 & 0.045 & 0.951 & 0.260 \\
$\bb_{19}$ & 12.978 & 12.978 & 693 & 1189 & -0.000 & 0.047 & 0.047 & 0.952 & 0.269 \\
$\bb_{20}$ & 13.065 & 13.065 & 669 & 1153 & -0.001 & 0.047 & 0.047 & 0.947 & 0.273 \\
$\bb_{21}$ & 13.327 & 13.327 & 649 & 1131 & -0.001 & 0.048 & 0.048 & 0.947 & 0.277 \\
\midrule
$\mathtt{Uniform}$ &  &  &  &  &  &  &  & 0.932 & 0.381 \\
\midrule
$\mathtt{WBATE}$ &  &  &  &  & -0.001 & 0.027 & 0.027 & 0.944 & 0.138 \\
\midrule
$\mathtt{LBATE}$ &  &  &  &  & 0.034 & 0.046 & 0.057 & 0.982 & 0.448 \\
\bottomrule\bottomrule
\end{tabular}

    \caption{Intention-to-Treat Effects (Simulations, $s=2$, ITT Bandwidth Selector).}
    \label{tab:sa-simuls-itt-s2}
\end{table}

\begin{table}[p]
    \centering
\begin{tabular}{lccccrrrrr}
\toprule\toprule
 & $h_1$ & $h_2$ & $N_{\mathrm{Co}}$ & $N_{\mathrm{Tr}}$ & Bias & SD & RMSE & EC & IL \\
\midrule
$\bb_{1}$ & 13.216 & 13.216 & 731 & 1266 & 0.002 & 0.045 & 0.045 & 0.948 & 0.262 \\
$\bb_{2}$ & 13.505 & 13.505 & 763 & 1339 & 0.001 & 0.044 & 0.044 & 0.950 & 0.255 \\
$\bb_{3}$ & 13.716 & 13.716 & 784 & 1389 & 0.001 & 0.043 & 0.043 & 0.948 & 0.251 \\
$\bb_{4}$ & 13.816 & 13.816 & 803 & 1393 & 0.000 & 0.043 & 0.043 & 0.946 & 0.247 \\
$\bb_{5}$ & 13.843 & 13.843 & 861 & 1320 & -0.001 & 0.042 & 0.042 & 0.951 & 0.235 \\
$\bb_{6}$ & 13.821 & 13.821 & 915 & 1219 & -0.001 & 0.041 & 0.041 & 0.951 & 0.233 \\
$\bb_{7}$ & 13.740 & 13.740 & 950 & 1107 & -0.002 & 0.041 & 0.041 & 0.951 & 0.239 \\
$\bb_{8}$ & 13.566 & 13.566 & 959 & 980 & -0.002 & 0.043 & 0.043 & 0.950 & 0.253 \\
$\bb_{9}$ & 13.909 & 13.909 & 1014 & 919 & -0.002 & 0.046 & 0.046 & 0.952 & 0.272 \\
$\bb_{10}$ & 15.049 & 15.049 & 1138 & 951 & -0.002 & 0.051 & 0.051 & 0.949 & 0.306 \\
$\bb_{11}$ & 14.775 & 14.775 & 1090 & 803 & -0.001 & 0.064 & 0.064 & 0.944 & 0.423 \\
$\bb_{12}$ & 14.642 & 14.642 & 1097 & 945 & -0.000 & 0.048 & 0.048 & 0.949 & 0.288 \\
$\bb_{13}$ & 13.564 & 13.564 & 963 & 959 & 0.000 & 0.044 & 0.044 & 0.946 & 0.257 \\
$\bb_{14}$ & 13.758 & 13.758 & 944 & 1130 & 0.001 & 0.041 & 0.041 & 0.940 & 0.237 \\
$\bb_{15}$ & 13.839 & 13.839 & 885 & 1281 & 0.003 & 0.042 & 0.042 & 0.949 & 0.232 \\
$\bb_{16}$ & 13.799 & 13.799 & 802 & 1390 & 0.004 & 0.044 & 0.044 & 0.948 & 0.246 \\
$\bb_{17}$ & 13.634 & 13.634 & 777 & 1371 & 0.004 & 0.044 & 0.044 & 0.948 & 0.251 \\
$\bb_{18}$ & 13.260 & 13.260 & 737 & 1279 & 0.004 & 0.045 & 0.045 & 0.949 & 0.260 \\
$\bb_{19}$ & 13.014 & 13.014 & 697 & 1195 & 0.003 & 0.046 & 0.047 & 0.951 & 0.269 \\
$\bb_{20}$ & 13.095 & 13.095 & 672 & 1158 & 0.003 & 0.047 & 0.047 & 0.948 & 0.274 \\
$\bb_{21}$ & 13.341 & 13.341 & 651 & 1133 & 0.003 & 0.048 & 0.048 & 0.948 & 0.277 \\
\midrule
$\mathtt{Uniform}$ &  &  &  &  &  &  &  & 0.932 & 0.381 \\
\midrule
$\mathtt{WBATE}$ &  &  &  &  & 0.001 & 0.026 & 0.026 & 0.941 & 0.138 \\
\midrule
$\mathtt{LBATE}$ &  &  &  &  & 0.035 & 0.044 & 0.056 & 0.980 & 0.444 \\
\bottomrule\bottomrule
\end{tabular}

    \caption{Intention-to-Treat Effects (Simulations, $s=3$, ITT Bandwidth Selector).}
    \label{tab:sa-simuls-itt-s3}
\end{table}

\begin{table}[p]
    \centering
\begin{tabular}{lccccrrrrr}
\toprule\toprule
 & $h_1$ & $h_2$ & $N_{\mathrm{Co}}$ & $N_{\mathrm{Tr}}$ & Bias & SD & RMSE & EC & IL \\
\midrule
$\bb_{1}$ & 13.220 & 13.220 & 732 & 1267 & 0.000 & 0.038 & 0.038 & 0.950 & 0.230 \\
$\bb_{2}$ & 13.525 & 13.525 & 765 & 1343 & 0.000 & 0.037 & 0.037 & 0.947 & 0.223 \\
$\bb_{3}$ & 13.740 & 13.740 & 787 & 1393 & 0.000 & 0.036 & 0.036 & 0.945 & 0.219 \\
$\bb_{4}$ & 13.857 & 13.857 & 808 & 1400 & 0.000 & 0.036 & 0.036 & 0.946 & 0.217 \\
$\bb_{5}$ & 13.904 & 13.904 & 869 & 1329 & 0.000 & 0.036 & 0.036 & 0.944 & 0.219 \\
$\bb_{6}$ & 13.890 & 13.890 & 924 & 1229 & 0.000 & 0.037 & 0.037 & 0.946 & 0.226 \\
$\bb_{7}$ & 13.811 & 13.811 & 959 & 1116 & 0.001 & 0.040 & 0.040 & 0.944 & 0.239 \\
$\bb_{8}$ & 13.639 & 13.639 & 968 & 989 & 0.001 & 0.044 & 0.044 & 0.945 & 0.261 \\
$\bb_{9}$ & 13.993 & 13.993 & 1024 & 928 & 0.001 & 0.048 & 0.048 & 0.944 & 0.289 \\
$\bb_{10}$ & 15.131 & 15.131 & 1147 & 961 & 0.001 & 0.054 & 0.054 & 0.939 & 0.338 \\
$\bb_{11}$ & 14.782 & 14.782 & 1091 & 803 & 0.001 & 0.070 & 0.070 & 0.931 & 0.483 \\
$\bb_{12}$ & 14.679 & 14.679 & 1101 & 949 & 0.000 & 0.051 & 0.051 & 0.942 & 0.313 \\
$\bb_{13}$ & 13.614 & 13.614 & 969 & 965 & 0.000 & 0.044 & 0.044 & 0.947 & 0.266 \\
$\bb_{14}$ & 13.811 & 13.811 & 951 & 1137 & 0.000 & 0.039 & 0.039 & 0.950 & 0.235 \\
$\bb_{15}$ & 13.870 & 13.870 & 889 & 1286 & 0.000 & 0.036 & 0.036 & 0.948 & 0.220 \\
$\bb_{16}$ & 13.837 & 13.837 & 805 & 1397 & 0.000 & 0.036 & 0.036 & 0.951 & 0.215 \\
$\bb_{17}$ & 13.678 & 13.678 & 781 & 1379 & 0.000 & 0.036 & 0.036 & 0.949 & 0.217 \\
$\bb_{18}$ & 13.289 & 13.289 & 739 & 1284 & 0.001 & 0.037 & 0.037 & 0.947 & 0.224 \\
$\bb_{19}$ & 13.011 & 13.011 & 697 & 1195 & 0.001 & 0.038 & 0.038 & 0.947 & 0.232 \\
$\bb_{20}$ & 13.077 & 13.077 & 671 & 1154 & 0.001 & 0.039 & 0.039 & 0.947 & 0.236 \\
$\bb_{21}$ & 13.330 & 13.330 & 651 & 1130 & 0.001 & 0.039 & 0.039 & 0.945 & 0.238 \\
\midrule
$\mathtt{Uniform}$ &  &  &  &  &  &  &  & 0.924 & 0.364 \\
\midrule
$\mathtt{WBATE}$ &  &  &  &  & 0.001 & 0.025 & 0.025 & 0.944 & 0.134 \\
\midrule
$\mathtt{LBATE}$ &  &  &  &  & 0.017 & 0.035 & 0.039 & 0.984 & 0.421 \\
\bottomrule\bottomrule
\end{tabular}

    \caption{First-Stage Effects (Simulations, $s=1$, ITT Bandwidth Selector).}
    \label{tab:sa-simuls-fs-s1}
\end{table}

\begin{table}[p]
    \centering
\begin{tabular}{lccccrrrrr}
\toprule\toprule
 & $h_1$ & $h_2$ & $N_{\mathrm{Co}}$ & $N_{\mathrm{Tr}}$ & Bias & SD & RMSE & EC & IL \\
\midrule
$\bb_{1}$ & 13.191 & 13.191 & 728 & 1262 & 0.002 & 0.039 & 0.039 & 0.948 & 0.231 \\
$\bb_{2}$ & 13.494 & 13.494 & 762 & 1338 & 0.002 & 0.038 & 0.038 & 0.943 & 0.224 \\
$\bb_{3}$ & 13.717 & 13.717 & 784 & 1390 & 0.002 & 0.037 & 0.037 & 0.944 & 0.219 \\
$\bb_{4}$ & 13.830 & 13.830 & 804 & 1397 & 0.002 & 0.036 & 0.037 & 0.948 & 0.218 \\
$\bb_{5}$ & 13.870 & 13.870 & 864 & 1325 & 0.002 & 0.037 & 0.037 & 0.944 & 0.220 \\
$\bb_{6}$ & 13.846 & 13.846 & 919 & 1223 & 0.001 & 0.038 & 0.038 & 0.941 & 0.227 \\
$\bb_{7}$ & 13.777 & 13.777 & 954 & 1112 & 0.001 & 0.040 & 0.040 & 0.942 & 0.239 \\
$\bb_{8}$ & 13.630 & 13.630 & 966 & 988 & 0.000 & 0.043 & 0.043 & 0.936 & 0.259 \\
$\bb_{9}$ & 14.017 & 14.017 & 1026 & 932 & -0.000 & 0.048 & 0.048 & 0.937 & 0.286 \\
$\bb_{10}$ & 15.156 & 15.156 & 1148 & 965 & -0.002 & 0.054 & 0.054 & 0.936 & 0.333 \\
$\bb_{11}$ & 14.778 & 14.778 & 1090 & 803 & -0.002 & 0.069 & 0.069 & 0.932 & 0.476 \\
$\bb_{12}$ & 14.646 & 14.646 & 1097 & 945 & -0.002 & 0.051 & 0.051 & 0.942 & 0.309 \\
$\bb_{13}$ & 13.581 & 13.581 & 965 & 961 & -0.001 & 0.044 & 0.044 & 0.944 & 0.262 \\
$\bb_{14}$ & 13.768 & 13.768 & 945 & 1131 & -0.001 & 0.039 & 0.039 & 0.949 & 0.232 \\
$\bb_{15}$ & 13.834 & 13.834 & 883 & 1281 & -0.001 & 0.036 & 0.036 & 0.945 & 0.218 \\
$\bb_{16}$ & 13.800 & 13.800 & 801 & 1391 & -0.001 & 0.035 & 0.035 & 0.940 & 0.213 \\
$\bb_{17}$ & 13.629 & 13.629 & 776 & 1370 & -0.001 & 0.036 & 0.036 & 0.943 & 0.215 \\
$\bb_{18}$ & 13.237 & 13.237 & 734 & 1275 & -0.001 & 0.037 & 0.037 & 0.941 & 0.222 \\
$\bb_{19}$ & 12.978 & 12.978 & 693 & 1189 & -0.002 & 0.038 & 0.038 & 0.940 & 0.229 \\
$\bb_{20}$ & 13.065 & 13.065 & 669 & 1153 & -0.002 & 0.039 & 0.039 & 0.941 & 0.233 \\
$\bb_{21}$ & 13.327 & 13.327 & 649 & 1131 & -0.002 & 0.039 & 0.039 & 0.939 & 0.235 \\
\midrule
$\mathtt{Uniform}$ &  &  &  &  &  &  &  & 0.917 & 0.361 \\
\midrule
$\mathtt{WBATE}$ &  &  &  &  & -0.000 & 0.024 & 0.024 & 0.938 & 0.133 \\
\midrule
$\mathtt{LBATE}$ &  &  &  &  & 0.021 & 0.038 & 0.043 & 0.978 & 0.427 \\
\bottomrule\bottomrule
\end{tabular}

    \caption{First-Stage Effects (Simulations, $s=2$, ITT Bandwidth Selector).}
    \label{tab:sa-simuls-fs-s2}
\end{table}

\begin{table}[p]
    \centering
\begin{tabular}{lccccrrrrr}
\toprule\toprule
 & $h_1$ & $h_2$ & $N_{\mathrm{Co}}$ & $N_{\mathrm{Tr}}$ & Bias & SD & RMSE & EC & IL \\
\midrule
$\bb_{1}$ & 13.216 & 13.216 & 731 & 1266 & 0.001 & 0.038 & 0.038 & 0.951 & 0.230 \\
$\bb_{2}$ & 13.505 & 13.505 & 763 & 1339 & 0.001 & 0.037 & 0.037 & 0.949 & 0.223 \\
$\bb_{3}$ & 13.716 & 13.716 & 784 & 1389 & 0.001 & 0.036 & 0.036 & 0.950 & 0.219 \\
$\bb_{4}$ & 13.816 & 13.816 & 803 & 1393 & 0.000 & 0.036 & 0.036 & 0.950 & 0.217 \\
$\bb_{5}$ & 13.843 & 13.843 & 861 & 1320 & -0.000 & 0.036 & 0.036 & 0.951 & 0.220 \\
$\bb_{6}$ & 13.821 & 13.821 & 915 & 1219 & -0.000 & 0.037 & 0.037 & 0.949 & 0.226 \\
$\bb_{7}$ & 13.740 & 13.740 & 950 & 1107 & -0.000 & 0.039 & 0.039 & 0.947 & 0.239 \\
$\bb_{8}$ & 13.566 & 13.566 & 959 & 980 & 0.000 & 0.043 & 0.043 & 0.942 & 0.260 \\
$\bb_{9}$ & 13.909 & 13.909 & 1014 & 919 & 0.001 & 0.048 & 0.048 & 0.941 & 0.288 \\
$\bb_{10}$ & 15.049 & 15.049 & 1138 & 951 & 0.002 & 0.053 & 0.053 & 0.940 & 0.336 \\
$\bb_{11}$ & 14.775 & 14.775 & 1090 & 803 & 0.003 & 0.068 & 0.068 & 0.930 & 0.478 \\
$\bb_{12}$ & 14.642 & 14.642 & 1097 & 945 & 0.002 & 0.051 & 0.051 & 0.939 & 0.311 \\
$\bb_{13}$ & 13.564 & 13.564 & 963 & 959 & 0.002 & 0.044 & 0.044 & 0.938 & 0.264 \\
$\bb_{14}$ & 13.758 & 13.758 & 944 & 1130 & 0.002 & 0.039 & 0.039 & 0.943 & 0.235 \\
$\bb_{15}$ & 13.839 & 13.839 & 885 & 1281 & 0.002 & 0.037 & 0.037 & 0.943 & 0.220 \\
$\bb_{16}$ & 13.799 & 13.799 & 802 & 1390 & 0.003 & 0.036 & 0.036 & 0.943 & 0.215 \\
$\bb_{17}$ & 13.634 & 13.634 & 777 & 1371 & 0.002 & 0.036 & 0.036 & 0.945 & 0.217 \\
$\bb_{18}$ & 13.260 & 13.260 & 737 & 1279 & 0.002 & 0.037 & 0.037 & 0.947 & 0.224 \\
$\bb_{19}$ & 13.014 & 13.014 & 697 & 1195 & 0.002 & 0.038 & 0.038 & 0.948 & 0.231 \\
$\bb_{20}$ & 13.095 & 13.095 & 672 & 1158 & 0.002 & 0.038 & 0.038 & 0.947 & 0.235 \\
$\bb_{21}$ & 13.341 & 13.341 & 651 & 1133 & 0.002 & 0.039 & 0.039 & 0.943 & 0.237 \\
\midrule
$\mathtt{Uniform}$ &  &  &  &  &  &  &  & 0.920 & 0.363 \\
\midrule
$\mathtt{WBATE}$ &  &  &  &  & 0.001 & 0.024 & 0.024 & 0.944 & 0.133 \\
\midrule
$\mathtt{LBATE}$ &  &  &  &  & 0.024 & 0.036 & 0.043 & 0.980 & 0.431 \\
\bottomrule\bottomrule
\end{tabular}

    \caption{First-Stage Effects (Simulations, $s=3$, ITT Bandwidth Selector).}
    \label{tab:sa-simuls-fs-s3}
\end{table}

\begin{table}[p]
    \centering
\begin{tabular}{lccccrrrrr}
\toprule\toprule
 & $h_1$ & $h_2$ & $N_{\mathrm{Co}}$ & $N_{\mathrm{Tr}}$ & Bias & SD & RMSE & EC & IL \\
\midrule
$\bb_{1}$ & 13.220 & 13.220 & 732 & 1267 & 0.001 & 0.072 & 0.072 & 0.954 & 0.433 \\
$\bb_{2}$ & 13.525 & 13.525 & 765 & 1343 & 0.001 & 0.069 & 0.069 & 0.958 & 0.419 \\
$\bb_{3}$ & 13.740 & 13.740 & 787 & 1393 & 0.001 & 0.067 & 0.067 & 0.958 & 0.411 \\
$\bb_{4}$ & 13.857 & 13.857 & 808 & 1400 & 0.001 & 0.066 & 0.066 & 0.958 & 0.402 \\
$\bb_{5}$ & 13.904 & 13.904 & 869 & 1329 & 0.001 & 0.064 & 0.064 & 0.957 & 0.373 \\
$\bb_{6}$ & 13.890 & 13.890 & 924 & 1229 & 0.001 & 0.061 & 0.061 & 0.959 & 0.364 \\
$\bb_{7}$ & 13.811 & 13.811 & 959 & 1116 & 0.001 & 0.060 & 0.060 & 0.955 & 0.369 \\
$\bb_{8}$ & 13.639 & 13.639 & 968 & 989 & 0.001 & 0.063 & 0.063 & 0.955 & 0.387 \\
$\bb_{9}$ & 13.993 & 13.993 & 1024 & 928 & 0.001 & 0.067 & 0.067 & 0.958 & 0.411 \\
$\bb_{10}$ & 15.131 & 15.131 & 1147 & 961 & 0.000 & 0.072 & 0.072 & 0.959 & 0.463 \\
$\bb_{11}$ & 14.782 & 14.782 & 1091 & 803 & -0.000 & 0.092 & 0.092 & 0.964 & 0.719 \\
$\bb_{12}$ & 14.679 & 14.679 & 1101 & 949 & 0.000 & 0.068 & 0.068 & 0.951 & 0.425 \\
$\bb_{13}$ & 13.614 & 13.614 & 969 & 965 & 0.000 & 0.061 & 0.061 & 0.954 & 0.378 \\
$\bb_{14}$ & 13.811 & 13.811 & 951 & 1137 & 0.000 & 0.057 & 0.057 & 0.959 & 0.348 \\
$\bb_{15}$ & 13.870 & 13.870 & 889 & 1286 & 0.001 & 0.059 & 0.059 & 0.955 & 0.342 \\
$\bb_{16}$ & 13.837 & 13.837 & 805 & 1397 & 0.001 & 0.062 & 0.062 & 0.954 & 0.369 \\
$\bb_{17}$ & 13.678 & 13.678 & 781 & 1379 & 0.001 & 0.063 & 0.063 & 0.951 & 0.374 \\
$\bb_{18}$ & 13.289 & 13.289 & 739 & 1284 & 0.001 & 0.065 & 0.065 & 0.951 & 0.382 \\
$\bb_{19}$ & 13.011 & 13.011 & 697 & 1195 & 0.001 & 0.066 & 0.066 & 0.952 & 0.390 \\
$\bb_{20}$ & 13.077 & 13.077 & 671 & 1154 & 0.000 & 0.066 & 0.066 & 0.954 & 0.393 \\
$\bb_{21}$ & 13.330 & 13.330 & 651 & 1130 & -0.000 & 0.067 & 0.067 & 0.952 & 0.394 \\
\midrule
$\mathtt{Uniform}$ &  &  &  &  &  &  &  & 0.962 & 0.586 \\
\midrule
$\mathtt{WBATE}$ &  &  &  &  & 0.001 & 0.039 & 0.039 & 0.949 & 0.214 \\
\midrule
$\mathtt{LBATE}$ &  &  &  &  & 0.053 & 0.058 & 0.078 & 0.997 & 0.753 \\
\bottomrule\bottomrule
\end{tabular}

    \caption{Fuzzy Effects (Simulations, $s=1$, ITT Bandwidth Selector).}
    \label{tab:sa-simuls-fuzzy-s1}
\end{table}

\begin{table}[p]
    \centering
\begin{tabular}{lccccrrrrr}
\toprule\toprule
 & $h_1$ & $h_2$ & $N_{\mathrm{Co}}$ & $N_{\mathrm{Tr}}$ & Bias & SD & RMSE & EC & IL \\
\midrule
$\bb_{1}$ & 13.191 & 13.191 & 728 & 1262 & 0.000 & 0.073 & 0.073 & 0.958 & 0.445 \\
$\bb_{2}$ & 13.494 & 13.494 & 762 & 1338 & 0.001 & 0.070 & 0.070 & 0.957 & 0.429 \\
$\bb_{3}$ & 13.717 & 13.717 & 784 & 1390 & 0.001 & 0.068 & 0.068 & 0.959 & 0.417 \\
$\bb_{4}$ & 13.830 & 13.830 & 804 & 1397 & 0.001 & 0.067 & 0.067 & 0.959 & 0.405 \\
$\bb_{5}$ & 13.870 & 13.870 & 864 & 1325 & 0.000 & 0.064 & 0.064 & 0.959 & 0.372 \\
$\bb_{6}$ & 13.846 & 13.846 & 919 & 1223 & -0.001 & 0.061 & 0.061 & 0.954 & 0.359 \\
$\bb_{7}$ & 13.777 & 13.777 & 954 & 1112 & -0.002 & 0.059 & 0.059 & 0.951 & 0.359 \\
$\bb_{8}$ & 13.630 & 13.630 & 966 & 988 & -0.003 & 0.061 & 0.061 & 0.953 & 0.369 \\
$\bb_{9}$ & 14.017 & 14.017 & 1026 & 932 & -0.003 & 0.064 & 0.064 & 0.952 & 0.383 \\
$\bb_{10}$ & 15.156 & 15.156 & 1148 & 965 & -0.003 & 0.068 & 0.068 & 0.949 & 0.421 \\
$\bb_{11}$ & 14.778 & 14.778 & 1090 & 803 & -0.004 & 0.085 & 0.085 & 0.961 & 1.058 \\
$\bb_{12}$ & 14.646 & 14.646 & 1097 & 945 & -0.003 & 0.063 & 0.064 & 0.952 & 0.384 \\
$\bb_{13}$ & 13.581 & 13.581 & 965 & 961 & -0.002 & 0.058 & 0.058 & 0.951 & 0.345 \\
$\bb_{14}$ & 13.768 & 13.768 & 945 & 1131 & -0.001 & 0.054 & 0.054 & 0.953 & 0.321 \\
$\bb_{15}$ & 13.834 & 13.834 & 883 & 1281 & -0.000 & 0.056 & 0.056 & 0.950 & 0.318 \\
$\bb_{16}$ & 13.800 & 13.800 & 801 & 1391 & 0.001 & 0.059 & 0.059 & 0.948 & 0.345 \\
$\bb_{17}$ & 13.629 & 13.629 & 776 & 1370 & 0.001 & 0.060 & 0.060 & 0.951 & 0.351 \\
$\bb_{18}$ & 13.237 & 13.237 & 734 & 1275 & 0.001 & 0.061 & 0.061 & 0.952 & 0.361 \\
$\bb_{19}$ & 12.978 & 12.978 & 693 & 1189 & 0.000 & 0.063 & 0.063 & 0.951 & 0.372 \\
$\bb_{20}$ & 13.065 & 13.065 & 669 & 1153 & -0.001 & 0.064 & 0.064 & 0.951 & 0.376 \\
$\bb_{21}$ & 13.327 & 13.327 & 649 & 1131 & -0.001 & 0.065 & 0.065 & 0.949 & 0.379 \\
\midrule
$\mathtt{Uniform}$ &  &  &  &  &  &  &  & 0.961 & 0.588 \\
\midrule
$\mathtt{WBATE}$ &  &  &  &  & -0.001 & 0.038 & 0.038 & 0.947 & 0.224 \\
\midrule
$\mathtt{LBATE}$ &  &  &  &  & 0.052 & 0.060 & 0.079 & 0.994 & 0.993 \\
\bottomrule\bottomrule
\end{tabular}

    \caption{Fuzzy Effects (Simulations, $s=2$, ITT Bandwidth Selector).}
    \label{tab:sa-simuls-fuzzy-s2}
\end{table}

\begin{table}[p]
    \centering
\begin{tabular}{lccccrrrrr}
\toprule\toprule
 & $h_1$ & $h_2$ & $N_{\mathrm{Co}}$ & $N_{\mathrm{Tr}}$ & Bias & SD & RMSE & EC & IL \\
\midrule
$\bb_{1}$ & 13.216 & 13.216 & 731 & 1266 & 0.001 & 0.071 & 0.071 & 0.958 & 0.428 \\
$\bb_{2}$ & 13.505 & 13.505 & 763 & 1339 & 0.001 & 0.069 & 0.069 & 0.954 & 0.412 \\
$\bb_{3}$ & 13.716 & 13.716 & 784 & 1389 & 0.001 & 0.067 & 0.067 & 0.955 & 0.400 \\
$\bb_{4}$ & 13.816 & 13.816 & 803 & 1393 & 0.000 & 0.066 & 0.066 & 0.952 & 0.388 \\
$\bb_{5}$ & 13.843 & 13.843 & 861 & 1320 & -0.001 & 0.063 & 0.063 & 0.954 & 0.358 \\
$\bb_{6}$ & 13.821 & 13.821 & 915 & 1219 & -0.002 & 0.059 & 0.059 & 0.953 & 0.346 \\
$\bb_{7}$ & 13.740 & 13.740 & 950 & 1107 & -0.003 & 0.058 & 0.058 & 0.954 & 0.349 \\
$\bb_{8}$ & 13.566 & 13.566 & 959 & 980 & -0.003 & 0.060 & 0.060 & 0.960 & 0.363 \\
$\bb_{9}$ & 13.909 & 13.909 & 1014 & 919 & -0.004 & 0.063 & 0.063 & 0.961 & 0.383 \\
$\bb_{10}$ & 15.049 & 15.049 & 1138 & 951 & -0.004 & 0.067 & 0.067 & 0.958 & 0.427 \\
$\bb_{11}$ & 14.775 & 14.775 & 1090 & 803 & -0.005 & 0.084 & 0.084 & 0.964 & 0.636 \\
$\bb_{12}$ & 14.642 & 14.642 & 1097 & 945 & -0.003 & 0.064 & 0.064 & 0.952 & 0.399 \\
$\bb_{13}$ & 13.564 & 13.564 & 963 & 959 & -0.002 & 0.059 & 0.059 & 0.954 & 0.361 \\
$\bb_{14}$ & 13.758 & 13.758 & 944 & 1130 & -0.000 & 0.056 & 0.056 & 0.950 & 0.336 \\
$\bb_{15}$ & 13.839 & 13.839 & 885 & 1281 & 0.002 & 0.059 & 0.059 & 0.951 & 0.332 \\
$\bb_{16}$ & 13.799 & 13.799 & 802 & 1390 & 0.003 & 0.062 & 0.062 & 0.949 & 0.360 \\
$\bb_{17}$ & 13.634 & 13.634 & 777 & 1371 & 0.003 & 0.062 & 0.062 & 0.949 & 0.366 \\
$\bb_{18}$ & 13.260 & 13.260 & 737 & 1279 & 0.004 & 0.064 & 0.064 & 0.956 & 0.376 \\
$\bb_{19}$ & 13.014 & 13.014 & 697 & 1195 & 0.003 & 0.065 & 0.065 & 0.954 & 0.386 \\
$\bb_{20}$ & 13.095 & 13.095 & 672 & 1158 & 0.003 & 0.066 & 0.066 & 0.951 & 0.390 \\
$\bb_{21}$ & 13.341 & 13.341 & 651 & 1133 & 0.002 & 0.066 & 0.066 & 0.954 & 0.392 \\
\midrule
$\mathtt{Uniform}$ &  &  &  &  &  &  &  & 0.961 & 0.561 \\
\midrule
$\mathtt{WBATE}$ &  &  &  &  & -0.000 & 0.037 & 0.037 & 0.950 & 0.203 \\
\midrule
$\mathtt{LBATE}$ &  &  &  &  & 0.050 & 0.059 & 0.078 & 0.995 & 0.687 \\
\bottomrule\bottomrule
\end{tabular}

    \caption{Fuzzy Effects (Simulations, $s=3$, ITT Bandwidth Selector).}
    \label{tab:sa-simuls-fuzzy-s3}
\end{table}

\clearpage

\begin{table}[p]
    \centering
\begin{tabular}{lccccrrrrr}
\toprule\toprule
 & $h_1$ & $h_2$ & $N_{\mathrm{Co}}$ & $N_{\mathrm{Tr}}$ & Bias & SD & RMSE & EC & IL \\
\midrule
$\bb_{1}$ & 13.404 & 13.404 & 746 & 1303 & 0.000 & 0.044 & 0.044 & 0.951 & 0.260 \\
$\bb_{2}$ & 13.823 & 13.823 & 788 & 1402 & 0.000 & 0.042 & 0.042 & 0.951 & 0.251 \\
$\bb_{3}$ & 14.104 & 14.104 & 815 & 1466 & 0.000 & 0.041 & 0.041 & 0.949 & 0.245 \\
$\bb_{4}$ & 14.257 & 14.257 & 849 & 1470 & 0.001 & 0.040 & 0.040 & 0.951 & 0.238 \\
$\bb_{5}$ & 14.318 & 14.318 & 919 & 1394 & 0.001 & 0.039 & 0.039 & 0.951 & 0.227 \\
$\bb_{6}$ & 14.306 & 14.306 & 975 & 1291 & 0.001 & 0.038 & 0.038 & 0.951 & 0.226 \\
$\bb_{7}$ & 14.235 & 14.235 & 1009 & 1176 & 0.000 & 0.038 & 0.038 & 0.949 & 0.233 \\
$\bb_{8}$ & 13.991 & 13.991 & 1009 & 1036 & 0.000 & 0.041 & 0.041 & 0.948 & 0.248 \\
$\bb_{9}$ & 14.133 & 14.133 & 1040 & 946 & 0.000 & 0.045 & 0.045 & 0.948 & 0.270 \\
$\bb_{10}$ & 15.100 & 15.100 & 1144 & 957 & -0.000 & 0.050 & 0.050 & 0.947 & 0.308 \\
$\bb_{11}$ & 14.722 & 14.722 & 1085 & 796 & -0.001 & 0.063 & 0.063 & 0.935 & 0.431 \\
$\bb_{12}$ & 14.698 & 14.698 & 1103 & 951 & -0.000 & 0.047 & 0.047 & 0.941 & 0.288 \\
$\bb_{13}$ & 13.915 & 13.915 & 1004 & 1004 & -0.000 & 0.042 & 0.042 & 0.949 & 0.252 \\
$\bb_{14}$ & 14.252 & 14.252 & 1004 & 1200 & 0.000 & 0.038 & 0.038 & 0.952 & 0.231 \\
$\bb_{15}$ & 14.307 & 14.307 & 941 & 1352 & 0.000 & 0.039 & 0.039 & 0.954 & 0.225 \\
$\bb_{16}$ & 14.254 & 14.254 & 849 & 1469 & 0.001 & 0.041 & 0.041 & 0.948 & 0.237 \\
$\bb_{17}$ & 14.044 & 14.044 & 809 & 1453 & 0.001 & 0.042 & 0.042 & 0.945 & 0.245 \\
$\bb_{18}$ & 13.534 & 13.534 & 758 & 1332 & 0.001 & 0.044 & 0.044 & 0.943 & 0.256 \\
$\bb_{19}$ & 13.137 & 13.137 & 707 & 1218 & 0.001 & 0.046 & 0.046 & 0.947 & 0.267 \\
$\bb_{20}$ & 13.188 & 13.188 & 680 & 1174 & 0.000 & 0.047 & 0.047 & 0.947 & 0.272 \\
$\bb_{21}$ & 13.422 & 13.422 & 657 & 1146 & -0.000 & 0.047 & 0.047 & 0.946 & 0.276 \\
\midrule
$\mathtt{Uniform}$ &  &  &  &  &  &  &  & 0.934 & 0.375 \\
\midrule
$\mathtt{WBATE}$ &  &  &  &  & 0.000 & 0.026 & 0.026 & 0.945 & 0.137 \\
\midrule
$\mathtt{LBATE}$ &  &  &  &  & 0.049 & 0.034 & 0.060 & 0.976 & 0.430 \\
\bottomrule\bottomrule
\end{tabular}

    \caption{Intention-to-Treat Effects (Simulations, $s=1$, Fuzzy Bandwidth Selector).}
    \label{tab:sa-simuls-itt-s1-bwmain}
\end{table}

\begin{table}[p]
    \centering
\begin{tabular}{lccccrrrrr}
\toprule\toprule
 & $h_1$ & $h_2$ & $N_{\mathrm{Co}}$ & $N_{\mathrm{Tr}}$ & Bias & SD & RMSE & EC & IL \\
\midrule
$\bb_{1}$ & 13.386 & 13.386 & 744 & 1300 & 0.001 & 0.044 & 0.044 & 0.949 & 0.260 \\
$\bb_{2}$ & 13.791 & 13.791 & 785 & 1397 & 0.002 & 0.042 & 0.042 & 0.951 & 0.251 \\
$\bb_{3}$ & 14.069 & 14.069 & 811 & 1460 & 0.002 & 0.041 & 0.041 & 0.949 & 0.245 \\
$\bb_{4}$ & 14.218 & 14.218 & 844 & 1464 & 0.002 & 0.040 & 0.040 & 0.953 & 0.239 \\
$\bb_{5}$ & 14.287 & 14.287 & 914 & 1389 & 0.001 & 0.039 & 0.039 & 0.953 & 0.227 \\
$\bb_{6}$ & 14.284 & 14.284 & 972 & 1288 & 0.000 & 0.038 & 0.038 & 0.949 & 0.226 \\
$\bb_{7}$ & 14.225 & 14.225 & 1008 & 1176 & -0.000 & 0.039 & 0.039 & 0.947 & 0.233 \\
$\bb_{8}$ & 13.995 & 13.995 & 1009 & 1037 & -0.001 & 0.041 & 0.041 & 0.951 & 0.248 \\
$\bb_{9}$ & 14.126 & 14.126 & 1038 & 946 & -0.002 & 0.045 & 0.045 & 0.945 & 0.270 \\
$\bb_{10}$ & 15.108 & 15.108 & 1144 & 958 & -0.003 & 0.050 & 0.050 & 0.942 & 0.307 \\
$\bb_{11}$ & 14.748 & 14.748 & 1087 & 799 & -0.005 & 0.063 & 0.064 & 0.936 & 0.429 \\
$\bb_{12}$ & 14.726 & 14.726 & 1106 & 955 & -0.004 & 0.048 & 0.048 & 0.940 & 0.287 \\
$\bb_{13}$ & 13.933 & 13.933 & 1006 & 1007 & -0.003 & 0.043 & 0.043 & 0.946 & 0.251 \\
$\bb_{14}$ & 14.255 & 14.255 & 1003 & 1201 & -0.002 & 0.039 & 0.039 & 0.946 & 0.229 \\
$\bb_{15}$ & 14.312 & 14.312 & 941 & 1353 & -0.001 & 0.039 & 0.039 & 0.947 & 0.223 \\
$\bb_{16}$ & 14.250 & 14.250 & 847 & 1469 & -0.001 & 0.041 & 0.041 & 0.944 & 0.236 \\
$\bb_{17}$ & 14.019 & 14.019 & 806 & 1449 & -0.000 & 0.042 & 0.042 & 0.944 & 0.244 \\
$\bb_{18}$ & 13.498 & 13.498 & 754 & 1326 & -0.000 & 0.044 & 0.044 & 0.949 & 0.255 \\
$\bb_{19}$ & 13.115 & 13.115 & 703 & 1214 & -0.001 & 0.046 & 0.046 & 0.950 & 0.266 \\
$\bb_{20}$ & 13.180 & 13.180 & 678 & 1173 & -0.001 & 0.046 & 0.047 & 0.946 & 0.271 \\
$\bb_{21}$ & 13.415 & 13.415 & 655 & 1146 & -0.002 & 0.047 & 0.047 & 0.946 & 0.275 \\
\midrule
$\mathtt{Uniform}$ &  &  &  &  &  &  &  & 0.934 & 0.374 \\
\midrule
$\mathtt{WBATE}$ &  &  &  &  & -0.001 & 0.026 & 0.026 & 0.947 & 0.136 \\
\midrule
$\mathtt{LBATE}$ &  &  &  &  & 0.030 & 0.041 & 0.051 & 0.984 & 0.445 \\
\bottomrule\bottomrule
\end{tabular}

    \caption{Intention-to-Treat Effects (Simulations, $s=2$, Fuzzy Bandwidth Selector).}
    \label{tab:sa-simuls-itt-s2-bwmain}
\end{table}

\begin{table}[p]
    \centering
\begin{tabular}{lccccrrrrr}
\toprule\toprule
 & $h_1$ & $h_2$ & $N_{\mathrm{Co}}$ & $N_{\mathrm{Tr}}$ & Bias & SD & RMSE & EC & IL \\
\midrule
$\bb_{1}$ & 13.416 & 13.416 & 746 & 1305 & 0.001 & 0.044 & 0.044 & 0.949 & 0.259 \\
$\bb_{2}$ & 13.823 & 13.823 & 788 & 1403 & 0.001 & 0.042 & 0.042 & 0.950 & 0.250 \\
$\bb_{3}$ & 14.093 & 14.093 & 813 & 1465 & 0.000 & 0.041 & 0.041 & 0.947 & 0.245 \\
$\bb_{4}$ & 14.229 & 14.229 & 845 & 1465 & -0.000 & 0.041 & 0.041 & 0.945 & 0.238 \\
$\bb_{5}$ & 14.282 & 14.282 & 913 & 1388 & -0.001 & 0.040 & 0.040 & 0.948 & 0.227 \\
$\bb_{6}$ & 14.281 & 14.281 & 971 & 1288 & -0.002 & 0.038 & 0.039 & 0.950 & 0.226 \\
$\bb_{7}$ & 14.214 & 14.214 & 1006 & 1174 & -0.003 & 0.039 & 0.039 & 0.952 & 0.232 \\
$\bb_{8}$ & 13.955 & 13.955 & 1005 & 1031 & -0.003 & 0.041 & 0.041 & 0.948 & 0.248 \\
$\bb_{9}$ & 14.080 & 14.080 & 1034 & 940 & -0.003 & 0.045 & 0.045 & 0.952 & 0.270 \\
$\bb_{10}$ & 15.038 & 15.038 & 1137 & 949 & -0.002 & 0.049 & 0.049 & 0.949 & 0.308 \\
$\bb_{11}$ & 14.667 & 14.667 & 1079 & 790 & -0.002 & 0.062 & 0.062 & 0.941 & 0.430 \\
$\bb_{12}$ & 14.687 & 14.687 & 1102 & 951 & -0.001 & 0.047 & 0.047 & 0.949 & 0.288 \\
$\bb_{13}$ & 13.902 & 13.902 & 1003 & 1003 & 0.000 & 0.042 & 0.042 & 0.946 & 0.252 \\
$\bb_{14}$ & 14.233 & 14.233 & 1001 & 1197 & 0.001 & 0.039 & 0.039 & 0.941 & 0.230 \\
$\bb_{15}$ & 14.292 & 14.292 & 940 & 1350 & 0.003 & 0.040 & 0.040 & 0.948 & 0.225 \\
$\bb_{16}$ & 14.217 & 14.217 & 845 & 1463 & 0.004 & 0.042 & 0.042 & 0.946 & 0.237 \\
$\bb_{17}$ & 13.995 & 13.995 & 805 & 1443 & 0.004 & 0.042 & 0.043 & 0.945 & 0.245 \\
$\bb_{18}$ & 13.492 & 13.492 & 755 & 1324 & 0.004 & 0.044 & 0.044 & 0.947 & 0.256 \\
$\bb_{19}$ & 13.126 & 13.126 & 705 & 1216 & 0.003 & 0.046 & 0.046 & 0.952 & 0.267 \\
$\bb_{20}$ & 13.183 & 13.183 & 679 & 1173 & 0.003 & 0.046 & 0.046 & 0.947 & 0.272 \\
$\bb_{21}$ & 13.408 & 13.408 & 655 & 1145 & 0.003 & 0.047 & 0.047 & 0.951 & 0.276 \\
\midrule
$\mathtt{Uniform}$ &  &  &  &  &  &  &  & 0.934 & 0.375 \\
\midrule
$\mathtt{WBATE}$ &  &  &  &  & 0.000 & 0.026 & 0.026 & 0.943 & 0.136 \\
\midrule
$\mathtt{LBATE}$ &  &  &  &  & 0.032 & 0.041 & 0.052 & 0.984 & 0.443 \\
\bottomrule\bottomrule
\end{tabular}

    \caption{Intention-to-Treat Effects (Simulations, $s=3$, Fuzzy Bandwidth Selector).}
    \label{tab:sa-simuls-itt-s3-bwmain}
\end{table}

\begin{table}[p]
    \centering
\begin{tabular}{lccccrrrrr}
\toprule\toprule
 & $h_1$ & $h_2$ & $N_{\mathrm{Co}}$ & $N_{\mathrm{Tr}}$ & Bias & SD & RMSE & EC & IL \\
\midrule
$\bb_{1}$ & 13.404 & 13.404 & 746 & 1303 & -0.000 & 0.037 & 0.037 & 0.949 & 0.227 \\
$\bb_{2}$ & 13.823 & 13.823 & 788 & 1402 & -0.000 & 0.035 & 0.035 & 0.948 & 0.218 \\
$\bb_{3}$ & 14.104 & 14.104 & 815 & 1466 & -0.000 & 0.035 & 0.035 & 0.949 & 0.213 \\
$\bb_{4}$ & 14.257 & 14.257 & 849 & 1470 & -0.000 & 0.034 & 0.034 & 0.946 & 0.211 \\
$\bb_{5}$ & 14.318 & 14.318 & 919 & 1394 & -0.000 & 0.035 & 0.035 & 0.945 & 0.214 \\
$\bb_{6}$ & 14.306 & 14.306 & 975 & 1291 & -0.000 & 0.036 & 0.036 & 0.946 & 0.221 \\
$\bb_{7}$ & 14.235 & 14.235 & 1009 & 1176 & 0.000 & 0.038 & 0.038 & 0.945 & 0.234 \\
$\bb_{8}$ & 13.991 & 13.991 & 1009 & 1036 & 0.000 & 0.042 & 0.042 & 0.946 & 0.256 \\
$\bb_{9}$ & 14.133 & 14.133 & 1040 & 946 & 0.000 & 0.047 & 0.047 & 0.941 & 0.287 \\
$\bb_{10}$ & 15.100 & 15.100 & 1144 & 957 & 0.000 & 0.053 & 0.053 & 0.939 & 0.338 \\
$\bb_{11}$ & 14.722 & 14.722 & 1085 & 796 & -0.001 & 0.069 & 0.069 & 0.930 & 0.486 \\
$\bb_{12}$ & 14.698 & 14.698 & 1103 & 951 & 0.000 & 0.050 & 0.050 & 0.944 & 0.313 \\
$\bb_{13}$ & 13.915 & 13.915 & 1004 & 1004 & -0.000 & 0.043 & 0.043 & 0.948 & 0.261 \\
$\bb_{14}$ & 14.252 & 14.252 & 1004 & 1200 & -0.000 & 0.037 & 0.037 & 0.949 & 0.229 \\
$\bb_{15}$ & 14.307 & 14.307 & 941 & 1352 & -0.000 & 0.034 & 0.034 & 0.948 & 0.215 \\
$\bb_{16}$ & 14.254 & 14.254 & 849 & 1469 & -0.000 & 0.034 & 0.034 & 0.948 & 0.209 \\
$\bb_{17}$ & 14.044 & 14.044 & 809 & 1453 & -0.000 & 0.034 & 0.034 & 0.948 & 0.212 \\
$\bb_{18}$ & 13.534 & 13.534 & 758 & 1332 & 0.000 & 0.036 & 0.036 & 0.947 & 0.221 \\
$\bb_{19}$ & 13.137 & 13.137 & 707 & 1218 & 0.000 & 0.037 & 0.037 & 0.947 & 0.230 \\
$\bb_{20}$ & 13.188 & 13.188 & 680 & 1174 & 0.000 & 0.038 & 0.038 & 0.943 & 0.234 \\
$\bb_{21}$ & 13.422 & 13.422 & 657 & 1146 & 0.000 & 0.038 & 0.038 & 0.944 & 0.237 \\
\midrule
$\mathtt{Uniform}$ &  &  &  &  &  &  &  & 0.929 & 0.359 \\
\midrule
$\mathtt{WBATE}$ &  &  &  &  & -0.000 & 0.024 & 0.024 & 0.946 & 0.133 \\
\midrule
$\mathtt{LBATE}$ &  &  &  &  & 0.014 & 0.033 & 0.036 & 0.988 & 0.416 \\
\bottomrule\bottomrule
\end{tabular}

    \caption{First-Stage Effects (Simulations, $s=1$, Fuzzy Bandwidth Selector).}
    \label{tab:sa-simuls-fs-s1-bwmain}
\end{table}

\begin{table}[p]
    \centering
\begin{tabular}{lccccrrrrr}
\toprule\toprule
 & $h_1$ & $h_2$ & $N_{\mathrm{Co}}$ & $N_{\mathrm{Tr}}$ & Bias & SD & RMSE & EC & IL \\
\midrule
$\bb_{1}$ & 13.386 & 13.386 & 744 & 1300 & 0.002 & 0.037 & 0.037 & 0.945 & 0.228 \\
$\bb_{2}$ & 13.791 & 13.791 & 785 & 1397 & 0.002 & 0.036 & 0.036 & 0.942 & 0.219 \\
$\bb_{3}$ & 14.069 & 14.069 & 811 & 1460 & 0.002 & 0.035 & 0.035 & 0.943 & 0.214 \\
$\bb_{4}$ & 14.218 & 14.218 & 844 & 1464 & 0.002 & 0.035 & 0.035 & 0.947 & 0.212 \\
$\bb_{5}$ & 14.287 & 14.287 & 914 & 1389 & 0.002 & 0.035 & 0.035 & 0.945 & 0.214 \\
$\bb_{6}$ & 14.284 & 14.284 & 972 & 1288 & 0.001 & 0.036 & 0.036 & 0.940 & 0.221 \\
$\bb_{7}$ & 14.225 & 14.225 & 1008 & 1176 & 0.001 & 0.038 & 0.038 & 0.941 & 0.233 \\
$\bb_{8}$ & 13.995 & 13.995 & 1009 & 1037 & 0.000 & 0.042 & 0.042 & 0.936 & 0.254 \\
$\bb_{9}$ & 14.126 & 14.126 & 1038 & 946 & -0.001 & 0.047 & 0.047 & 0.938 & 0.285 \\
$\bb_{10}$ & 15.108 & 15.108 & 1144 & 958 & -0.003 & 0.052 & 0.052 & 0.937 & 0.334 \\
$\bb_{11}$ & 14.748 & 14.748 & 1087 & 799 & -0.005 & 0.067 & 0.067 & 0.932 & 0.479 \\
$\bb_{12}$ & 14.726 & 14.726 & 1106 & 955 & -0.003 & 0.049 & 0.049 & 0.942 & 0.308 \\
$\bb_{13}$ & 13.933 & 13.933 & 1006 & 1007 & -0.002 & 0.042 & 0.042 & 0.945 & 0.257 \\
$\bb_{14}$ & 14.255 & 14.255 & 1003 & 1201 & -0.002 & 0.037 & 0.037 & 0.947 & 0.226 \\
$\bb_{15}$ & 14.312 & 14.312 & 941 & 1353 & -0.002 & 0.034 & 0.034 & 0.943 & 0.212 \\
$\bb_{16}$ & 14.250 & 14.250 & 847 & 1469 & -0.002 & 0.033 & 0.033 & 0.943 & 0.206 \\
$\bb_{17}$ & 14.019 & 14.019 & 806 & 1449 & -0.002 & 0.034 & 0.034 & 0.944 & 0.209 \\
$\bb_{18}$ & 13.498 & 13.498 & 754 & 1326 & -0.002 & 0.036 & 0.036 & 0.939 & 0.218 \\
$\bb_{19}$ & 13.115 & 13.115 & 703 & 1214 & -0.002 & 0.037 & 0.037 & 0.939 & 0.227 \\
$\bb_{20}$ & 13.180 & 13.180 & 678 & 1173 & -0.003 & 0.038 & 0.038 & 0.940 & 0.231 \\
$\bb_{21}$ & 13.415 & 13.415 & 655 & 1146 & -0.003 & 0.038 & 0.038 & 0.938 & 0.234 \\
\midrule
$\mathtt{Uniform}$ &  &  &  &  &  &  &  & 0.921 & 0.355 \\
\midrule
$\mathtt{WBATE}$ &  &  &  &  & -0.001 & 0.024 & 0.024 & 0.941 & 0.132 \\
\midrule
$\mathtt{LBATE}$ &  &  &  &  & 0.017 & 0.034 & 0.038 & 0.982 & 0.422 \\
\bottomrule\bottomrule
\end{tabular}

    \caption{First-Stage Effects (Simulations, $s=2$, Fuzzy Bandwidth Selector).}
    \label{tab:sa-simuls-fs-s2-bwmain}
\end{table}

\begin{table}[p]
    \centering
\begin{tabular}{lccccrrrrr}
\toprule\toprule
 & $h_1$ & $h_2$ & $N_{\mathrm{Co}}$ & $N_{\mathrm{Tr}}$ & Bias & SD & RMSE & EC & IL \\
\midrule
$\bb_{1}$ & 13.416 & 13.416 & 746 & 1305 & 0.001 & 0.037 & 0.037 & 0.949 & 0.227 \\
$\bb_{2}$ & 13.823 & 13.823 & 788 & 1403 & 0.000 & 0.035 & 0.035 & 0.947 & 0.219 \\
$\bb_{3}$ & 14.093 & 14.093 & 813 & 1465 & 0.000 & 0.034 & 0.034 & 0.949 & 0.213 \\
$\bb_{4}$ & 14.229 & 14.229 & 845 & 1465 & -0.001 & 0.034 & 0.034 & 0.949 & 0.211 \\
$\bb_{5}$ & 14.282 & 14.282 & 913 & 1388 & -0.001 & 0.034 & 0.034 & 0.951 & 0.214 \\
$\bb_{6}$ & 14.281 & 14.281 & 971 & 1288 & -0.001 & 0.035 & 0.035 & 0.947 & 0.220 \\
$\bb_{7}$ & 14.214 & 14.214 & 1006 & 1174 & -0.001 & 0.037 & 0.037 & 0.945 & 0.233 \\
$\bb_{8}$ & 13.955 & 13.955 & 1005 & 1031 & -0.001 & 0.041 & 0.041 & 0.945 & 0.254 \\
$\bb_{9}$ & 14.080 & 14.080 & 1034 & 940 & -0.000 & 0.047 & 0.047 & 0.940 & 0.286 \\
$\bb_{10}$ & 15.038 & 15.038 & 1137 & 949 & -0.000 & 0.052 & 0.052 & 0.939 & 0.336 \\
$\bb_{11}$ & 14.667 & 14.667 & 1079 & 790 & 0.001 & 0.067 & 0.067 & 0.929 & 0.483 \\
$\bb_{12}$ & 14.687 & 14.687 & 1102 & 951 & 0.002 & 0.050 & 0.050 & 0.939 & 0.310 \\
$\bb_{13}$ & 13.902 & 13.902 & 1003 & 1003 & 0.002 & 0.043 & 0.043 & 0.938 & 0.260 \\
$\bb_{14}$ & 14.233 & 14.233 & 1001 & 1197 & 0.002 & 0.037 & 0.037 & 0.943 & 0.228 \\
$\bb_{15}$ & 14.292 & 14.292 & 940 & 1350 & 0.003 & 0.035 & 0.035 & 0.944 & 0.214 \\
$\bb_{16}$ & 14.217 & 14.217 & 845 & 1463 & 0.003 & 0.034 & 0.034 & 0.943 & 0.209 \\
$\bb_{17}$ & 13.995 & 13.995 & 805 & 1443 & 0.002 & 0.034 & 0.035 & 0.944 & 0.211 \\
$\bb_{18}$ & 13.492 & 13.492 & 755 & 1324 & 0.002 & 0.036 & 0.036 & 0.946 & 0.221 \\
$\bb_{19}$ & 13.126 & 13.126 & 705 & 1216 & 0.002 & 0.037 & 0.037 & 0.946 & 0.229 \\
$\bb_{20}$ & 13.183 & 13.183 & 679 & 1173 & 0.002 & 0.037 & 0.037 & 0.948 & 0.233 \\
$\bb_{21}$ & 13.408 & 13.408 & 655 & 1145 & 0.002 & 0.038 & 0.038 & 0.944 & 0.236 \\
\midrule
$\mathtt{Uniform}$ &  &  &  &  &  &  &  & 0.925 & 0.357 \\
\midrule
$\mathtt{WBATE}$ &  &  &  &  & 0.001 & 0.024 & 0.024 & 0.943 & 0.132 \\
\midrule
$\mathtt{LBATE}$ &  &  &  &  & 0.021 & 0.033 & 0.040 & 0.984 & 0.429 \\
\bottomrule\bottomrule
\end{tabular}

    \caption{First-Stage Effects (Simulations, $s=3$, Fuzzy Bandwidth Selector).}
    \label{tab:sa-simuls-fs-s3-bwmain}
\end{table}

\begin{table}[p]
    \centering
\begin{tabular}{lccccrrrrr}
\toprule\toprule
 & $h_1$ & $h_2$ & $N_{\mathrm{Co}}$ & $N_{\mathrm{Tr}}$ & Bias & SD & RMSE & EC & IL \\
\midrule
$\bb_{1}$ & 13.404 & 13.404 & 746 & 1303 & 0.001 & 0.073 & 0.073 & 0.955 & 0.427 \\
$\bb_{2}$ & 13.823 & 13.823 & 788 & 1402 & 0.001 & 0.069 & 0.069 & 0.958 & 0.411 \\
$\bb_{3}$ & 14.104 & 14.104 & 815 & 1466 & 0.002 & 0.067 & 0.067 & 0.958 & 0.401 \\
$\bb_{4}$ & 14.257 & 14.257 & 849 & 1470 & 0.002 & 0.066 & 0.066 & 0.958 & 0.387 \\
$\bb_{5}$ & 14.318 & 14.318 & 919 & 1394 & 0.002 & 0.064 & 0.064 & 0.959 & 0.360 \\
$\bb_{6}$ & 14.306 & 14.306 & 975 & 1291 & 0.002 & 0.061 & 0.061 & 0.959 & 0.353 \\
$\bb_{7}$ & 14.235 & 14.235 & 1009 & 1176 & 0.001 & 0.060 & 0.060 & 0.955 & 0.359 \\
$\bb_{8}$ & 13.991 & 13.991 & 1009 & 1036 & 0.001 & 0.062 & 0.062 & 0.955 & 0.378 \\
$\bb_{9}$ & 14.133 & 14.133 & 1040 & 946 & 0.001 & 0.067 & 0.067 & 0.958 & 0.407 \\
$\bb_{10}$ & 15.100 & 15.100 & 1144 & 957 & 0.001 & 0.074 & 0.074 & 0.959 & 0.462 \\
$\bb_{11}$ & 14.722 & 14.722 & 1085 & 796 & 0.001 & 0.098 & 0.098 & 0.965 & 0.964 \\
$\bb_{12}$ & 14.698 & 14.698 & 1103 & 951 & 0.000 & 0.070 & 0.070 & 0.953 & 0.424 \\
$\bb_{13}$ & 13.915 & 13.915 & 1004 & 1004 & 0.001 & 0.061 & 0.061 & 0.956 & 0.370 \\
$\bb_{14}$ & 14.252 & 14.252 & 1004 & 1200 & 0.001 & 0.057 & 0.057 & 0.959 & 0.338 \\
$\bb_{15}$ & 14.307 & 14.307 & 941 & 1352 & 0.001 & 0.058 & 0.058 & 0.953 & 0.331 \\
$\bb_{16}$ & 14.254 & 14.254 & 849 & 1469 & 0.001 & 0.061 & 0.061 & 0.956 & 0.355 \\
$\bb_{17}$ & 14.044 & 14.044 & 809 & 1453 & 0.001 & 0.062 & 0.062 & 0.950 & 0.365 \\
$\bb_{18}$ & 13.534 & 13.534 & 758 & 1332 & 0.001 & 0.064 & 0.064 & 0.950 & 0.375 \\
$\bb_{19}$ & 13.137 & 13.137 & 707 & 1218 & 0.001 & 0.066 & 0.066 & 0.952 & 0.387 \\
$\bb_{20}$ & 13.188 & 13.188 & 680 & 1174 & 0.000 & 0.067 & 0.067 & 0.953 & 0.389 \\
$\bb_{21}$ & 13.422 & 13.422 & 657 & 1146 & -0.000 & 0.067 & 0.067 & 0.952 & 0.391 \\
\midrule
$\mathtt{Uniform}$ &  &  &  &  &  &  &  & 0.964 & 0.591 \\
\midrule
$\mathtt{WBATE}$ &  &  &  &  & 0.001 & 0.039 & 0.039 & 0.950 & 0.222 \\
\midrule
$\mathtt{LBATE}$ &  &  &  &  & 0.055 & 0.065 & 0.085 & 0.998 & 0.931 \\
\bottomrule\bottomrule
\end{tabular}

    \caption{Fuzzy Effects (Simulations, $s=1$, Fuzzy Bandwidth Selector).}
    \label{tab:sa-simuls-fuzzy-s1-bwmain}
\end{table}

\begin{table}[p]
    \centering
\begin{tabular}{lccccrrrrr}
\toprule\toprule
 & $h_1$ & $h_2$ & $N_{\mathrm{Co}}$ & $N_{\mathrm{Tr}}$ & Bias & SD & RMSE & EC & IL \\
\midrule
$\bb_{1}$ & 13.386 & 13.386 & 744 & 1300 & 0.001 & 0.074 & 0.074 & 0.959 & 0.439 \\
$\bb_{2}$ & 13.791 & 13.791 & 785 & 1397 & 0.001 & 0.071 & 0.071 & 0.957 & 0.420 \\
$\bb_{3}$ & 14.069 & 14.069 & 811 & 1460 & 0.002 & 0.068 & 0.068 & 0.959 & 0.407 \\
$\bb_{4}$ & 14.218 & 14.218 & 844 & 1464 & 0.002 & 0.067 & 0.067 & 0.960 & 0.390 \\
$\bb_{5}$ & 14.287 & 14.287 & 914 & 1389 & 0.001 & 0.063 & 0.063 & 0.962 & 0.358 \\
$\bb_{6}$ & 14.284 & 14.284 & 972 & 1288 & -0.001 & 0.060 & 0.060 & 0.956 & 0.347 \\
$\bb_{7}$ & 14.225 & 14.225 & 1008 & 1176 & -0.001 & 0.059 & 0.059 & 0.953 & 0.348 \\
$\bb_{8}$ & 13.995 & 13.995 & 1009 & 1037 & -0.002 & 0.061 & 0.061 & 0.952 & 0.360 \\
$\bb_{9}$ & 14.126 & 14.126 & 1038 & 946 & -0.002 & 0.065 & 0.065 & 0.954 & 0.380 \\
$\bb_{10}$ & 15.108 & 15.108 & 1144 & 958 & -0.003 & 0.070 & 0.070 & 0.952 & 0.422 \\
$\bb_{11}$ & 14.748 & 14.748 & 1087 & 799 & -0.003 & 0.089 & 0.089 & 0.963 & 1.656 \\
$\bb_{12}$ & 14.726 & 14.726 & 1106 & 955 & -0.003 & 0.064 & 0.064 & 0.951 & 0.382 \\
$\bb_{13}$ & 13.933 & 13.933 & 1006 & 1007 & -0.002 & 0.057 & 0.057 & 0.950 & 0.337 \\
$\bb_{14}$ & 14.255 & 14.255 & 1003 & 1201 & -0.002 & 0.053 & 0.053 & 0.954 & 0.311 \\
$\bb_{15}$ & 14.312 & 14.312 & 941 & 1353 & -0.001 & 0.055 & 0.055 & 0.952 & 0.306 \\
$\bb_{16}$ & 14.250 & 14.250 & 847 & 1469 & 0.000 & 0.058 & 0.058 & 0.951 & 0.330 \\
$\bb_{17}$ & 14.019 & 14.019 & 806 & 1449 & 0.001 & 0.059 & 0.059 & 0.950 & 0.342 \\
$\bb_{18}$ & 13.498 & 13.498 & 754 & 1326 & 0.001 & 0.061 & 0.061 & 0.950 & 0.355 \\
$\bb_{19}$ & 13.115 & 13.115 & 703 & 1214 & 0.000 & 0.063 & 0.063 & 0.953 & 0.368 \\
$\bb_{20}$ & 13.180 & 13.180 & 678 & 1173 & -0.001 & 0.064 & 0.064 & 0.953 & 0.373 \\
$\bb_{21}$ & 13.415 & 13.415 & 655 & 1146 & -0.001 & 0.065 & 0.065 & 0.950 & 0.377 \\
\midrule
$\mathtt{Uniform}$ &  &  &  &  &  &  &  & 0.961 & 0.619 \\
\midrule
$\mathtt{WBATE}$ &  &  &  &  & -0.001 & 0.038 & 0.038 & 0.943 & 0.249 \\
\midrule
$\mathtt{LBATE}$ &  &  &  &  & 0.054 & 0.064 & 0.084 & 0.995 & 1.461 \\
\bottomrule\bottomrule
\end{tabular}

    \caption{Fuzzy Effects (Simulations, $s=2$, Fuzzy Bandwidth Selector).}
    \label{tab:sa-simuls-fuzzy-s2-bwmain}
\end{table}

\begin{table}[p]
    \centering
\begin{tabular}{lccccrrrrr}
\toprule\toprule
 & $h_1$ & $h_2$ & $N_{\mathrm{Co}}$ & $N_{\mathrm{Tr}}$ & Bias & SD & RMSE & EC & IL \\
\midrule
$\bb_{1}$ & 13.416 & 13.416 & 746 & 1305 & 0.002 & 0.072 & 0.072 & 0.959 & 0.422 \\
$\bb_{2}$ & 13.823 & 13.823 & 788 & 1403 & 0.001 & 0.069 & 0.069 & 0.955 & 0.403 \\
$\bb_{3}$ & 14.093 & 14.093 & 813 & 1465 & 0.001 & 0.067 & 0.067 & 0.956 & 0.390 \\
$\bb_{4}$ & 14.229 & 14.229 & 845 & 1465 & 0.000 & 0.065 & 0.065 & 0.952 & 0.373 \\
$\bb_{5}$ & 14.282 & 14.282 & 913 & 1388 & -0.001 & 0.062 & 0.062 & 0.954 & 0.344 \\
$\bb_{6}$ & 14.281 & 14.281 & 971 & 1288 & -0.003 & 0.059 & 0.059 & 0.954 & 0.335 \\
$\bb_{7}$ & 14.214 & 14.214 & 1006 & 1174 & -0.003 & 0.057 & 0.057 & 0.954 & 0.338 \\
$\bb_{8}$ & 13.955 & 13.955 & 1005 & 1031 & -0.004 & 0.059 & 0.059 & 0.960 & 0.354 \\
$\bb_{9}$ & 14.080 & 14.080 & 1034 & 940 & -0.004 & 0.063 & 0.063 & 0.962 & 0.379 \\
$\bb_{10}$ & 15.038 & 15.038 & 1137 & 949 & -0.004 & 0.069 & 0.069 & 0.960 & 0.428 \\
$\bb_{11}$ & 14.667 & 14.667 & 1079 & 790 & -0.003 & 0.090 & 0.090 & 0.966 & 0.675 \\
$\bb_{12}$ & 14.687 & 14.687 & 1102 & 951 & -0.003 & 0.065 & 0.065 & 0.952 & 0.398 \\
$\bb_{13}$ & 13.902 & 13.902 & 1003 & 1003 & -0.001 & 0.059 & 0.059 & 0.957 & 0.352 \\
$\bb_{14}$ & 14.233 & 14.233 & 1001 & 1197 & -0.000 & 0.055 & 0.055 & 0.952 & 0.325 \\
$\bb_{15}$ & 14.292 & 14.292 & 940 & 1350 & 0.002 & 0.058 & 0.058 & 0.952 & 0.320 \\
$\bb_{16}$ & 14.217 & 14.217 & 845 & 1463 & 0.003 & 0.061 & 0.061 & 0.950 & 0.345 \\
$\bb_{17}$ & 13.995 & 13.995 & 805 & 1443 & 0.004 & 0.062 & 0.062 & 0.948 & 0.357 \\
$\bb_{18}$ & 13.492 & 13.492 & 755 & 1324 & 0.004 & 0.064 & 0.064 & 0.954 & 0.370 \\
$\bb_{19}$ & 13.126 & 13.126 & 705 & 1216 & 0.003 & 0.065 & 0.065 & 0.954 & 0.383 \\
$\bb_{20}$ & 13.183 & 13.183 & 679 & 1173 & 0.003 & 0.066 & 0.066 & 0.951 & 0.387 \\
$\bb_{21}$ & 13.408 & 13.408 & 655 & 1145 & 0.002 & 0.066 & 0.066 & 0.954 & 0.390 \\
\midrule
$\mathtt{Uniform}$ &  &  &  &  &  &  &  & 0.962 & 0.552 \\
\midrule
$\mathtt{WBATE}$ &  &  &  &  & 0.000 & 0.038 & 0.038 & 0.946 & 0.202 \\
\midrule
$\mathtt{LBATE}$ &  &  &  &  & 0.053 & 0.065 & 0.083 & 0.994 & 0.717 \\
\bottomrule\bottomrule
\end{tabular}

    \caption{Fuzzy Effects (Simulations, $s=3$, Fuzzy Bandwidth Selector).}
    \label{tab:sa-simuls-fuzzy-s3-bwmain}
\end{table}

\clearpage


%%%%%%%%%%%%%%%%%%%%%%%%%%%%%%%%%%%%%%%%%%%%%%%%%%%%%%%%%%%%%%%%%%%%%%
%% REFERENCES
%%%%%%%%%%%%%%%%%%%%%%%%%%%%%%%%%%%%%%%%%%%%%%%%%%%%%%%%%%%%%%%%%%%%%%
\clearpage
\bibliography{CTY_2026_JASA--bib}
\bibliographystyle{plainnat}



\end{document}
